% Generated mechanically from frozen input p04/ko/coherent.tex.
% Do not edit this generated file; edit the bound locale source instead.

% OK, start here.
%

\setcounter{chapter}{29}
\chapter{스킴의 코호몰로지}



\phantomsection
\label{coherent-section-phantom}


\section{서론}
\label{coherent-section-introduction}

\noindent
이 장에서는 먼저 준연접층의 코호몰로지에 관한 여러 결과를 증명한다.
기본 참고문헌은 \cite{EGA}이다.
그런 다음 뇌터 상황에서 연접층의 코호몰로지를 자세히 다룬다.
\cite{FAC}를 보라.






\section{준연접층의 {\v C}ech 코호몰로지}
\label{coherent-section-cech-quasi-coherent}

\noindent
$X$를 스킴이라 하자.
$U \subset X$를 아핀 열린집합이라 하자.
$U$의 {\it 표준 열린 덮개}란
$\mathcal{U} : U = \bigcup_{i = 1}^n D(f_i)$ 꼴의 덮개를 말한다.
여기서 $f_1, \ldots, f_n \in \Gamma(U, \mathcal{O}_X)$는 단위 아이디얼을
생성한다. Schemes의 정의 \ref{schemes-definition-standard-covering}를 보라.

\begin{lemma}
\label{coherent-lemma-cech-cohomology-quasi-coherent-trivial}
$X$를 스킴이라 하자.
$\mathcal{F}$를 준연접 $\mathcal{O}_X$-가군이라 하자.
$\mathcal{U} : U = \bigcup_{i = 1}^n D(f_i)$를 $X$의 한 아핀 열린집합의
표준 열린 덮개라 하자.
그러면 $\check{H}^p(\mathcal{U}, \mathcal{F}) = 0$이 모든
$p > 0$에 대해 성립한다.
\end{lemma}

\begin{proof}
$U = \Spec(A)$라 쓰자. 여기서 $A$는 어떤 환이다.
다시 말해 $f_1, \ldots, f_n$은 $A$의 단위 아이디얼을 생성하는
$A$의 원소들이다.
$\mathcal{F}|_U = \widetilde{M}$이라 쓰자. 여기서 어떤 $A$-가군을
$M$으로 나타냈다.
명백히 {\v C}ech 복합체
$\check{\mathcal{C}}^\bullet(\mathcal{U}, \mathcal{F})$는 복합체
$$
\prod\nolimits_{i_0} M_{f_{i_0}} \to
\prod\nolimits_{i_0i_1} M_{f_{i_0}f_{i_1}} \to
\prod\nolimits_{i_0i_1i_2} M_{f_{i_0}f_{i_1}f_{i_2}} \to
\ldots
$$
와 동일시된다. 우리가 보여야 할 것은 확장 복합체
\begin{equation}
\label{coherent-equation-extended}
0 \to
M \to
\prod\nolimits_{i_0} M_{f_{i_0}} \to
\prod\nolimits_{i_0i_1} M_{f_{i_0}f_{i_1}} \to
\prod\nolimits_{i_0i_1i_2} M_{f_{i_0}f_{i_1}f_{i_2}} \to
\ldots
\end{equation}
가 완전하다는 것이다. 그 절단은 Algebra의 보조정리
\ref{algebra-lemma-cover-module}에서 이미 다루었다.
복합체 (\ref{coherent-equation-extended})를 소아이디얼 $\mathfrak p$에서 국소화한 뒤
완전함을 보이면 충분하다. Algebra의 보조정리
\ref{algebra-lemma-characterize-zero-local}을 보라.
사실 우리는 $\mathfrak p$에서 국소화한 확장 복합체가 영사상과
호모토픽임을 보일 것이다.

\medskip\noindent
첨자 $i$ 중 $f_i \not \in \mathfrak p$인 것이 존재한다.
그러한 첨자 하나를 택해 $i_{\text{fix}}$로 고정하자. 그러면
$M_{f_{i_{\text{fix}}}, \mathfrak p} = M_{\mathfrak p}$이다. 마찬가지로
곱 $f_{i_0} \ldots f_{i_p}$에서 국소화한 뒤 다시 $\mathfrak p$에서
국소화할 때, 임의의 $f_{i_j}$는 $i_j = i_{\text{fix}}$이면 생략할 수 있다.
다음 호모토피를 정의하자.
$$
h :
\prod\nolimits_{i_0 \ldots i_{p + 1}}
M_{f_{i_0} \ldots f_{i_{p + 1}}, \mathfrak p}
\longrightarrow
\prod\nolimits_{i_0 \ldots i_p}
M_{f_{i_0} \ldots f_{i_p}, \mathfrak p}
$$
그 규칙은
$$
h(s)_{i_0 \ldots i_p} = s_{i_{\text{fix}} i_0 \ldots i_p}
$$
이다. 이는 Cohomology의 보조정리
\ref{cohomology-lemma-homology-complex}의 증명에 나오는 호모토피와
``쌍대''이다.
다시 말해 $h : \prod_{i_0} M_{f_{i_0}, \mathfrak p} \to M_\mathfrak p$는
인자 $M_{f_{i_{\text{fix}}}, \mathfrak p} = M_{\mathfrak p}$로의 사영이고,
일반적으로 사상 $h$는 인자들
$M_{f_{i_{\text{fix}}} f_{i_1} \ldots f_{i_{p + 1}}, \mathfrak p}
= M_{f_{i_1} \ldots f_{i_{p + 1}}, \mathfrak p}$로의 사영이다. 계산하면
\begin{align*}
(dh + hd)(s)_{i_0 \ldots i_p}
& =
\sum\nolimits_{j = 0}^p
(-1)^j
h(s)_{i_0 \ldots \hat i_j \ldots i_p}
+
d(s)_{i_{\text{fix}} i_0 \ldots i_p}\\
& =
\sum\nolimits_{j = 0}^p
(-1)^j
s_{i_{\text{fix}} i_0 \ldots \hat i_j \ldots i_p}
+
s_{i_0 \ldots i_p}
+
\sum\nolimits_{j = 0}^p
(-1)^{j + 1}
s_{i_{\text{fix}} i_0 \ldots \hat i_j \ldots i_p} \\
& =
s_{i_0 \ldots i_p}
\end{align*}
이다. 이는 원하는 대로 항등사상이 영사상과 호모토픽임을 증명한다.
\end{proof}

\noindent
다음 보조정리는 특히 임의의 아핀 스킴 $X$와 임의의 준연접층
$\mathcal{F}$가 $X$ 위에 주어졌을 때
$$
H^p(X, \mathcal{F}) = 0
$$
가 모든 $p > 0$에 대해 성립함을 말한다.

\begin{lemma}
\label{coherent-lemma-quasi-coherent-affine-cohomology-zero}
\begin{slogan}
Serre 소멸: 아핀 스킴 위의 준연접 가군은 고차 코호몰로지가 소멸한다.
\end{slogan}
$X$를 스킴이라 하자.
$\mathcal{F}$를 준연접 $\mathcal{O}_X$-가군이라 하자.
모든 아핀 열린집합 $U \subset X$에 대해
$H^p(U, \mathcal{F}) = 0$이 모든 $p > 0$에 대해 성립한다.
\end{lemma}

\begin{proof}
Cohomology의 보조정리 \ref{cohomology-lemma-cech-vanish-basis}를 적용한다.
기저 $\mathcal{B}$로 $X$의 아핀 열린집합들을 사용한다. 여기서 고려하는
위상은 $X$의 위상이다.
열린 덮개들의 모음 $\text{Cov}$로 $X$의 아핀 열린집합들의 표준 열린
덮개들을 사용한다. 이제 Cohomology의 보조정리
\ref{cohomology-lemma-cech-vanish-basis}의 조건 (1), (2), (3)을 확인한다.
아핀 안에서 표준 열린집합들의 교집합은 다시 표준 열린집합이다.
따라서 성질 (1)이 성립한다.
이 덮개들은 $\mathcal{B}$의 임의의 원소의 열린 덮개들 가운데 공종계를 이룬다.
Schemes의 보조정리 \ref{schemes-lemma-standard-open}을 보라.
따라서 (2)가 성립한다.
마지막으로 조건 (3)은 보조정리
\ref{coherent-lemma-cech-cohomology-quasi-coherent-trivial}에서 따른다.
\end{proof}

\noindent
다음은 아핀 위의 준연접층 코호몰로지 소멸의 상대적 형태이다.

\begin{lemma}
\label{coherent-lemma-relative-affine-vanishing}
$f : X \to S$를 스킴의 사상이라 하자.
$\mathcal{F}$를 준연접 $\mathcal{O}_X$-가군이라 하자.
$f$가 아핀이면 $R^if_*\mathcal{F} = 0$이 모든 $i > 0$에 대해 성립한다.
\end{lemma}

\begin{proof}
Cohomology의 보조정리
\ref{cohomology-lemma-describe-higher-direct-images}에 따르면 층
$R^if_*\mathcal{F}$는 전층
$V \mapsto H^i(f^{-1}(V), \mathcal{F}|_{f^{-1}(V)})$에 결부된 층이다.
가정에 의해 $V$가 아핀이면 언제나 $f^{-1}(V)$도 아핀이다.
Morphisms의 정의 \ref{morphisms-definition-affine}를 보라.
보조정리 \ref{coherent-lemma-quasi-coherent-affine-cohomology-zero}에 의해
$H^i(f^{-1}(V), \mathcal{F}|_{f^{-1}(V)}) = 0$인데, 이는 $V$가 아핀이면
언제나 성립한다.
$S$에는 아핀 열린집합들로 이루어진 기저가 있으므로 결론이 따른다.
\end{proof}

\begin{lemma}
\label{coherent-lemma-relative-affine-cohomology}
$f : X \to S$를 스킴의 아핀 사상이라 하자.
$\mathcal{F}$를 준연접 $\mathcal{O}_X$-가군이라 하자.
그러면 $H^i(X, \mathcal{F}) = H^i(S, f_*\mathcal{F})$가 모든
$i \geq 0$에 대해 성립한다.
\end{lemma}

\begin{proof}
보조정리 \ref{coherent-lemma-relative-affine-vanishing}과 Leray 스펙트럼 열에서 따른다.
Cohomology의 보조정리 \ref{cohomology-lemma-apply-Leray}를 보라.
\end{proof}

\noindent
다음 두 보조정리는 언제 {\v C}ech 코호몰로지를 사용하여 준연접 가군의
코호몰로지를 계산할 수 있는지를 설명한다.

\begin{lemma}
\label{coherent-lemma-affine-diagonal}
$X$를 스킴이라 하자. 다음은 서로 동치이다.
\begin{enumerate}
\item $X$는 아핀 대각사상 $\Delta : X \to X \times X$를 갖는다.
\item 아핀 열린집합 $U, V \subset X$에 대해 교집합 $U \cap V$는 아핀이다.
\item 열린 덮개 $\mathcal{U} : X = \bigcup_{i \in I} U_i$가 존재하여
$U_{i_0 \ldots i_p}$가 모든 $p \ge 0$과 모든
$i_0, \ldots, i_p \in I$에 대해 아핀 열린집합이다.
\end{enumerate}
특히 $X$가 분리적이면 이들이 성립한다.
\end{lemma}

\begin{proof}
$X$가 아핀 대각사상을 갖는다고 가정하자. $U, V \subset X$를 아핀
열린집합이라 하자. 그러면 $U \cap V = \Delta^{-1}(U \times V)$는 아핀이다.
따라서 (2)가 성립한다. (2)가 (3)을 함의함은 즉시 알 수 있다.
거꾸로 (3)과 같은 $X$의 덮개가 존재하면
$X \times X = \bigcup U_i \times U_{i'}$는 아핀 열린 덮개이고,
$\Delta^{-1}(U_i \times U_{i'}) = U_i \cap U_{i'}$가 아핀임을 알 수 있다.
따라서 Morphisms의 보조정리
\ref{morphisms-lemma-characterize-affine}에 의해 $\Delta$는 아핀 사상이다.
마지막 주장은 Schemes의 보조정리
\ref{schemes-lemma-characterize-separated}에서 따른다.
\end{proof}

\begin{lemma}
\label{coherent-lemma-cech-cohomology-quasi-coherent}
$X$를 스킴이라 하자.
$\mathcal{U} : X = \bigcup_{i \in I} U_i$를 열린 덮개라 하되,
$U_{i_0 \ldots i_p}$가 모든 $p \ge 0$과 모든
$i_0, \ldots, i_p \in I$에 대해 아핀 열린집합이라고 하자.
이 경우 임의의 준연접층 $\mathcal{F}$에 대해
$$
\check{H}^p(\mathcal{U}, \mathcal{F}) = H^p(X, \mathcal{F})
$$
가 $\Gamma(X, \mathcal{O}_X)$-가군으로서 모든 $p$에 대해 성립한다.
\end{lemma}

\begin{proof}
보조정리 \ref{coherent-lemma-quasi-coherent-affine-cohomology-zero}에 비추어 보면 이는
Cohomology의 보조정리
\ref{cohomology-lemma-cech-spectral-sequence-application}의 특수한 경우이다.
\end{proof}






\section{코호몰로지의 소멸}
\label{coherent-section-vanishing}

\noindent
아핀 스킴 위에서는 임의의 준연접층의 고차 코호몰로지 군이
소멸함을 보았다
(보조정리 \ref{coherent-lemma-quasi-coherent-affine-cohomology-zero}).
이 성질은 역으로 아핀 스킴을 특징짓기도 한다.
여기서는 그 두 가지 형태를 제시한다.

\begin{lemma}
\label{coherent-lemma-quasi-compact-h1-zero-covering}
\begin{reference}
\cite{Serre-criterion}, \cite[II, Theorem 5.2.1 (d') and IV (1.7.17)]{EGA}
\end{reference}
\begin{slogan}
아핀성에 대한 세르 판정법.
\end{slogan}
$X$를 스킴이라 하자.
다음을 가정하자.
\begin{enumerate}
\item $X$는 준콤팩트이다.
\item 모든 준연접 아이디얼 층
$\mathcal{I} \subset \mathcal{O}_X$에 대해 $H^1(X, \mathcal{I}) = 0$이다.
\end{enumerate}
그러면 $X$는 아핀이다.
\end{lemma}

\begin{proof}
$x \in X$를 닫힌점이라 하자. $U \subset X$를 $x$의 아핀 열린 근방이라
하자. $U = \Spec(A)$로 쓰고,
$\mathfrak m \subset A$를 $x$에 대응하는 극대 아이디얼이라 하자.
$Z = X \setminus U$ 및 $Z' = Z \cup \{x\}$로 둔다.
Schemes의 보조정리 \ref{schemes-lemma-reduced-closed-subscheme}에 의해
준연접 아이디얼 층
$\mathcal{I}$ 및 각각 $\mathcal{I}'$가 존재하며, 이들은
기약 닫힌 부분스킴 $Z$ 및 각각 $Z'$를 잘라낸다.
다음 짧은 완전열을 생각하자.
$$
0 \to \mathcal{I}' \to \mathcal{I} \to \mathcal{I}/\mathcal{I}' \to 0.
$$
$x$는 $X$의 닫힌점이고 $x \not \in Z$이므로
$\mathcal{I}/\mathcal{I}'$는 $x$에서 지지된다. 실제로
$\mathcal{I}/\mathcal{I'}$를 $U$에 제한한 층은 $A$-가군
$A/\mathfrak m$에 대응한다. 따라서 $\Gamma(X, \mathcal{I}/\mathcal{I'})
= A/\mathfrak m$이다. 가정에 의해 $H^1(X, \mathcal{I}') = 0$이므로,
전역 단면 $f \in \Gamma(X, \mathcal{I})$가 존재하며, 그 상은 원소
$1 \in A/\mathfrak m$, 즉 $\mathcal{I}/\mathcal{I'}$의 단면이다.
분명히 $x \in X_f \subset U$이다. 이는
$X_f = D(f_A)$임을 뜻하며, 여기서 $f_A$는
$f$의 $A = \Gamma(U, \mathcal{O}_X)$ 안에서의 상이다.
특히 $X_f$는 아핀이다.

\medskip\noindent
합집합 $W = \bigcup X_f$를 생각하되, 이 합집합은 모든
$f \in \Gamma(X, \mathcal{O}_X)$ 중 $X_f$가 아핀인 것에 걸친다.
분명히 $W$는 $X$의 열린집합이다.
위 논증에 의해 $X$의 모든 닫힌점은 $W$에 들어간다.
닫힌 부분집합 $X \setminus W$는 $X$ 안에서 준콤팩트이다
(Topology의 보조정리 \ref{topology-lemma-closed-in-quasi-compact} 참조).
따라서 이것이 공집합이 아니면 닫힌점을 갖는다
(Topology의 보조정리 \ref{topology-lemma-quasi-compact-closed-point} 참조).
이는 모든 닫힌점이 $W$에 들어간다는 사실에 모순된다.
따라서 $X = W$이다.

\medskip\noindent
유한 개의 $f_1, \ldots, f_n \in \Gamma(X, \mathcal{O}_X)$를 택하되,
$X = X_{f_1} \cup \ldots \cup X_{f_n}$이고 각
$X_{f_i}$가 아핀이 되게 한다.
위에서 본 바와 같이 이는 가능하다.
Properties의 보조정리 \ref{properties-lemma-characterize-affine}에 의해,
증명을 끝내려면 $f_1, \ldots, f_n$이
$\Gamma(X, \mathcal{O}_X)$의 단위 아이디얼을 생성함을 보이면 충분하다.
다음 짧은 완전열을 생각하자.
$$
\xymatrix{
0 \ar[r] &
\mathcal{F} \ar[r] &
\mathcal{O}_X^{\oplus n} \ar[rr]^{f_1, \ldots, f_n} & &
\mathcal{O}_X \ar[r] &
0
}
$$
$f_1, \ldots, f_n$으로 정의되는 화살표는
열린집합 $X_{f_i}$들이 $X$를 덮으므로 전사이다.
이 전사 사상의 핵을 $\mathcal{F}$라 하자.
$\mathcal{F}$에는 다음과 같은 여과가 있다.
$$
0 = \mathcal{F}_0 \subset \mathcal{F}_1 \subset
\ldots \subset \mathcal{F}_n = \mathcal{F}
$$
각 부분몫 $\mathcal{F}_i/\mathcal{F}_{i - 1}$은
준연접 아이디얼 층과 동형이다.
구체적으로 $\mathcal{F}_i$를 $\mathcal{F}$와 처음 $i$개의
직합 인자로 이루어진 $\mathcal{O}_X^{\oplus n}$ 부분가군의
교집합으로 택할 수 있다.
보조정리의 가정에 의해
$H^1(X, \mathcal{F}_i/\mathcal{F}_{i - 1}) = 0$이 모든 $i$에 대해
성립한다.
이로부터 $H^1(X, \mathcal{F}_2) = 0$인데, 이는
$H^1(X, \mathcal{F}_1)$과
$H^1(X, \mathcal{F}_2/\mathcal{F}_1)$ 사이에 놓이기 때문이다.
이를 계속하면 $H^1(X, \mathcal{F}) = 0$을 얻는다.
따라서 원하는 대로 다음 사상이 전사이다.
$$
\xymatrix{
\bigoplus\nolimits_{i = 1, \ldots, n} \Gamma(X, \mathcal{O}_X)
\ar[rr]^{f_1, \ldots, f_n} & &
\Gamma(X, \mathcal{O}_X)
}
$$
\end{proof}

\noindent
$X$가 뇌터 스킴이면 모든 준연접 아이디얼 층은 자동으로
연접 아이디얼 층이며 유한형 준연접 아이디얼 층이다.
따라서 이 경우에는 앞의 보조정리와 다음 보조정리를 모두 적용할 수 있다.

\begin{lemma}
\label{coherent-lemma-quasi-separated-h1-zero-covering}
\begin{reference}
\cite{Serre-criterion}, \cite[II, Theorem 5.2.1]{EGA}
\end{reference}
\begin{slogan}
아핀성에 대한 세르 판정법.
\end{slogan}
$X$를 스킴이라 하자. 다음을 가정하자.
\begin{enumerate}
\item $X$는 준콤팩트이다.
\item $X$는 준분리이다.
\item $H^1(X, \mathcal{I}) = 0$이다. 여기서 $\mathcal{I}$는
임의의 유한형 준연접 아이디얼 층이다.
\end{enumerate}
그러면 $X$는 아핀이다.
\end{lemma}

\begin{proof}
Properties의 보조정리
\ref{properties-lemma-quasi-coherent-colimit-finite-type}에 의해
모든 준연접 아이디얼 층은 유한형 준연접 아이디얼 층들의
유향 쌍극한이다.
Cohomology의 보조정리
\ref{cohomology-lemma-quasi-separated-cohomology-colimit}에 의해
$X$ 위에서 코호몰로지를 취하는 것은 유향 쌍극한과 가환한다.
따라서 $H^1(X, \mathcal{I}) = 0$이 $X$ 위의 모든
준연접 아이디얼 층에 대해 성립한다. 다시 말해 보조정리
\ref{coherent-lemma-quasi-compact-h1-zero-covering}을 적용할 수 있다.
\end{proof}

\noindent
위의 논증을 이용하면 가역층이 풍부할 충분조건을 얻을 수 있다.
다만 이 조건은 필요조건이 아님을 독자에게 밝혀 둔다.

\begin{lemma}
\label{coherent-lemma-quasi-compact-h1-zero-invertible}
$X$를 스킴이라 하자. $\mathcal{L}$을 가역 $\mathcal{O}_X$-가군이라
하자. 다음을 가정하자.
\begin{enumerate}
\item $X$는 준콤팩트이다.
\item 모든 준연접 아이디얼 층
$\mathcal{I} \subset \mathcal{O}_X$에 대해
$n \geq 1$이 존재하여
$H^1(X, \mathcal{I} \otimes_{\mathcal{O}_X} \mathcal{L}^{\otimes n}) = 0$이
된다.
\end{enumerate}
그러면 $\mathcal{L}$은 풍부하다.
\end{lemma}

\begin{proof}
보조정리 \ref{coherent-lemma-quasi-compact-h1-zero-covering}과
정확히 같은 방식으로 증명한다.
$x \in X$를 닫힌점이라 하자. $U \subset X$를 $x$의 아핀 열린 근방
중 $\mathcal{L}|_U \cong \mathcal{O}_U$인 것으로 택하자.
$U = \Spec(A)$로 쓰고,
$\mathfrak m \subset A$를 $x$에 대응하는 극대 아이디얼이라 하자.
$Z = X \setminus U$ 및 $Z' = Z \cup \{x\}$로 둔다.
Schemes의 보조정리 \ref{schemes-lemma-reduced-closed-subscheme}에 의해
준연접 아이디얼 층
$\mathcal{I}$ 및 각각 $\mathcal{I}'$가 존재하며, 이들은
기약 닫힌 부분스킴 $Z$ 및 각각 $Z'$를 잘라낸다.
다음 짧은 완전열을 생각하자.
$$
0 \to \mathcal{I}' \to \mathcal{I} \to \mathcal{I}/\mathcal{I}' \to 0.
$$
모든 $n \geq 1$에 대해 다음 짧은 완전열을 얻는다.
$$
0 \to \mathcal{I}' \otimes_{\mathcal{O}_X} \mathcal{L}^{\otimes n}
\to \mathcal{I} \otimes_{\mathcal{O}_X} \mathcal{L}^{\otimes n} \to
\mathcal{I}/\mathcal{I}' \otimes_{\mathcal{O}_X} \mathcal{L}^{\otimes n} \to 0.
$$
가정에 의해 $n$을 택하여
$H^1(X, \mathcal{I}' \otimes_{\mathcal{O}_X} \mathcal{L}^{\otimes n}) = 0$이
되게 할 수 있다.
$x$는 $X$의 닫힌점이고 $x \not \in Z$이므로
$\mathcal{I}/\mathcal{I}'$는 $x$에서 지지된다. 실제로
$\mathcal{I}/\mathcal{I'}$를 $U$에 제한한 층은 $A$-가군
$A/\mathfrak m$에 대응한다. $\mathcal{L}$은 $U$ 위에서 자명하므로,
$\mathcal{I}/\mathcal{I}' \otimes_{\mathcal{O}_X} \mathcal{L}^{\otimes n}$을
$U$에 제한한 층도 $A$-가군 $A/\mathfrak m$에 대응한다.
따라서
$\Gamma(X, \mathcal{I}/\mathcal{I'} \otimes_{\mathcal{O}_X}
\mathcal{L}^{\otimes n}) = A/\mathfrak m$이다.
$n$의 선택에 의해 전역 단면
$s \in \Gamma(X, \mathcal{I} \otimes_{\mathcal{O}_X} \mathcal{L}^{\otimes n})$이
존재하며, 이는 원소 $1 \in A/\mathfrak m$에 보내진다.
분명히 $x \in X_s \subset U$인데, 이는 $s$가 $Z$의 점들에서
소멸하기 때문이다.
이는 $X_s = D(f)$임을 뜻하며, 여기서
$f \in A$는 $s$의 $A \cong \Gamma(U, \mathcal{L}^{\otimes n})$
안에서의 상이다. 특히 $X_s$는 아핀이다.

\medskip\noindent
합집합 $W = \bigcup X_s$를 생각하되, 이는 모든
$s \in \Gamma(X, \mathcal{L}^{\otimes n})$와 $n \geq 1$ 중
$X_s$가 아핀인 것에 걸친 합집합이다. 분명히 $W$는 $X$의 열린집합이다.
위 논증에 의해 $X$의 모든 닫힌점은 $W$에 들어간다.
닫힌 부분집합 $X \setminus W$는 $X$ 안에서 준콤팩트이다
(Topology의 보조정리 \ref{topology-lemma-closed-in-quasi-compact} 참조).
따라서 이것이 공집합이 아니면 닫힌점을 갖는다
(Topology의 보조정리 \ref{topology-lemma-quasi-compact-closed-point} 참조).
이는 모든 닫힌점이 $W$에 들어간다는 사실에 모순된다.
따라서 $X = W$이다. Properties의 정의
\ref{properties-definition-ample}에 의해 이는 $\mathcal{L}$이
풍부하다는 뜻이다.
\end{proof}

\noindent
보조정리 \ref{coherent-lemma-quasi-compact-h1-zero-invertible}에는
유한형 아이디얼 층을 사용하는 변형이 있다.
필요하게 되면 여기서 이를 정식화하고 증명할 것이다.

\begin{lemma}
\label{coherent-lemma-criterion-affine-morphism}
$f : X \to Y$를 준콤팩트 사상이라 하고 $X$와 $Y$가 준분리라고 하자.
$R^1f_*\mathcal{I} = 0$이라고 하자. 이 조건은 모든 준연접 아이디얼 층
$\mathcal{I}$, 곧 $X$ 위의 아이디얼 층에 대해 요구한다.
그러면 $f$는 아핀이다.
\end{lemma}

\begin{proof}
$V \subset Y$를 아핀 열린 부분스킴이라 하자.
$U = f^{-1}(V)$가 아핀임을 보이면 된다.
포함 사상 $V \to Y$는 Schemes의 보조정리
\ref{schemes-lemma-quasi-compact-permanence}에 의해 준콤팩트이다.
따라서 밑변환 $U \to X$는 준콤팩트이다. Schemes의 보조정리
\ref{schemes-lemma-quasi-compact-preserved-base-change}를 참조하라.
그러므로 임의의 준연접 아이디얼 층 $\mathcal{I}$가 $U$ 위에 주어지면
$X$ 위의 준연접 아이디얼 층으로 연장된다. Properties의 보조정리
\ref{properties-lemma-extend-trivial}를 참조하라.
$R^1f_*$의 형성은 $Y$ 위에서 국소적이므로
(Cohomology의 절 \ref{cohomology-section-locality}),
보조정리의 가정에 의해 $R^1(U \to V)_*\mathcal{I} = 0$이다.
따라서 르레 스펙트럼 열
(Cohomology의 보조정리 \ref{cohomology-lemma-Leray})에 의해
$H^1(U, \mathcal{I}) = H^1(V, (U \to V)_*\mathcal{I})$이다.
Schemes의 보조정리
\ref{schemes-lemma-push-forward-quasi-coherent}에 의해
$(U \to V)_*\mathcal{I}$는 준연접이므로,
보조정리 \ref{coherent-lemma-quasi-coherent-affine-cohomology-zero}에 의해
$H^1(V, (U \to V)_*\mathcal{I}) = 0$이다.
그러므로 보조정리 \ref{coherent-lemma-quasi-compact-h1-zero-covering}에 의해
$U$는 아핀이다.
\end{proof}













\section{고차 직상의 준연접성}
\label{coherent-section-quasi-coherence}

\noindent
아핀 스킴 위의 준연접 가군은 고차 코호몰로지 군이 영임을 보았다.
(준)분리 준콤팩트 스킴 $X$의 경우에는 이로부터 어느 차수 이후
준연접층의 코호몰로지 군이 소멸함이 따른다. 그러나 $X$가 유한
코호몰로지 차원을 갖는다고 할 수는 없다. 코호몰로지 차원은
{\it 모든} $\mathcal{O}_X$-가군의 코호몰로지 소멸로 정의되기 때문이다.

\begin{lemma}[귀납 원리]
\label{coherent-lemma-induction-principle}
\begin{reference}
\cite[Proposition 3.3.1]{BvdB}
\end{reference}
$X$를 준콤팩트 준분리 스킴이라 하자. $P$를 $X$의 준콤팩트
열린집합들이 갖는 성질이라 하자. 다음을 가정하자.
\begin{enumerate}
\item $P$는 $X$의 모든 아핀 열린집합에 대해 성립한다.
\item $U$가 준콤팩트 열린집합이고 $V$가 아핀 열린집합이며,
$P$가 $U$, $V$, $U \cap V$에 대해 성립하면
$P$는 $U \cup V$에 대해서도 성립한다.
\end{enumerate}
그러면 $P$는 $X$의 모든 준콤팩트 열린집합에 대해 성립하며,
특히 $X$에 대해 성립한다.
\end{lemma}

\begin{proof}
먼저 $P$가 {\it 분리} 준콤팩트 열린집합 $W \subset X$에 대해
성립함을 귀납법으로 보인다. 실제로 그러한 열린집합은
$W = U_1 \cup \ldots \cup U_n$으로 쓸 수 있고, $n$에 대한
귀납법을 쓸 수 있다. 이때 성질 (2)를
$U = U_1 \cup \ldots \cup U_{n - 1}$ 및 $V = U_n$에 적용한다. 이는
$U \cap V = (U_1 \cap U_n) \cup \ldots \cup (U_{n - 1} \cap U_n)$가
$n - 1$개의 아핀 열린 부분스킴의 합집합이기 때문에 허용된다.
여기서는 Schemes의 보조정리
\ref{schemes-lemma-characterize-separated}를 아핀 열린집합
$U_i$와 $U_n$에 적용하였으며, 이들은 $W$의 열린집합들이다.
이제 임의의 준콤팩트 열린집합 $W \subset X$에 대해,
$W$를 덮는 데 필요한 아핀 열린집합의 개수에 관한 귀납법을
앞과 같은 방법으로 적용할 수 있다. 이때 준콤팩트 열린집합
$U_i \cap U_n$은 아핀 스킴 $U_n$의 열린 부분스킴이므로 분리이다.
\end{proof}

\begin{lemma}
\label{coherent-lemma-vanishing-nr-affines}
\begin{slogan}
대각사상이 아핀인 스킴에서는 준연접 가군의 코호몰로지가
덮개에 필요한 아핀 열린집합의 개수보다 큰 차수에서 소멸한다.
\end{slogan}
$X$를 대각사상이 아핀인 준콤팩트 스킴이라 하자
(예를 들어 $X$가 분리이면 그러하다).
$t = t(X)$를 $X$를 덮는 데 필요한 아핀 열린집합의 최소 개수라 하자.
그러면 $H^n(X, \mathcal{F}) = 0$이 모든 $n \geq t$와
모든 준연접층 $\mathcal{F}$에 대해 성립한다.
\end{lemma}

\begin{proof}
첫째 증명.
$t$에 대한 귀납법을 쓴다.
$t = 1$이면 보조정리
\ref{coherent-lemma-quasi-coherent-affine-cohomology-zero}에서 결과가 따른다.
$t > 1$이면 $X = U \cup V$로 쓰되, $V$는 아핀 열린집합이고
$U = U_1 \cup \ldots \cup U_{t - 1}$은 $t - 1$개의 아핀 열린집합의
합집합이게 한다. 이 경우
$U \cap V =  (U_1 \cap V) \cup \ldots (U_{t - 1} \cap V)$도
$t - 1$개의 아핀 열린 부분스킴의 합집합이다.
실제로 대각사상이 아핀이므로 두 아핀 열린집합의 교집합은 아핀이다.
보조정리 \ref{coherent-lemma-affine-diagonal}을 참조하라.
다음 마이어--피토리스 긴 완전열을 적용한다.
$$
0 \to
H^0(X, \mathcal{F}) \to
H^0(U, \mathcal{F}) \oplus H^0(V, \mathcal{F}) \to
H^0(U \cap V, \mathcal{F}) \to
H^1(X, \mathcal{F}) \to \ldots
$$
Cohomology의 보조정리
\ref{cohomology-lemma-mayer-vietoris}를 참조하라.
귀납법에 의해 군 $H^i(U, \mathcal{F})$,
$H^i(V, \mathcal{F})$, $H^i(U \cap V, \mathcal{F})$는
$i \geq t - 1$에서 영이다. 따라서
$H^i(X, \mathcal{F})$는 $i \geq t$에서 곧바로 영이다.

\medskip\noindent
둘째 증명.
$\mathcal{U} : X = \bigcup_{i = 1}^t U_i$를 유한 아핀 열린 덮개라
하자. $X$의 대각사상이 아핀이므로 다중 교집합
$U_{i_0 \ldots i_p}$는 모두 아핀이다.
보조정리 \ref{coherent-lemma-affine-diagonal}을 참조하라.
보조정리 \ref{coherent-lemma-cech-cohomology-quasi-coherent}에 의해 체흐
코호몰로지 군 $\check{H}^p(\mathcal{U}, \mathcal{F})$은
코호몰로지 군과 일치한다.
Cohomology의 보조정리 \ref{cohomology-lemma-alternating-usual}에 의해
체흐 코호몰로지 군은 교대 체흐 복합체
$\check{\mathcal{C}}_{alt}^\bullet(\mathcal{U}, \mathcal{F})$로 계산할 수 있다.
덮개가 $t$개의 원소로 이루어지므로
$\check{\mathcal{C}}_{alt}^p(\mathcal{U}, \mathcal{F}) = 0$이
모든 $p \geq t$에 대해 성립함이 곧바로 보인다. 따라서 결과가 따른다.
\end{proof}

\begin{lemma}
\label{coherent-lemma-affine-diagonal-universal-delta-functor}
$X$를 대각사상이 아핀인 준콤팩트 스킴이라 하자
(예를 들어 $X$가 분리이면 그러하다). 그러면 다음이 성립한다.
\begin{enumerate}
\item 준연접 $\mathcal{O}_X$-가군 $\mathcal{F}$가 주어지면 매입
$\mathcal{F} \to \mathcal{F}'$가 존재한다. 이는 준연접
$\mathcal{O}_X$-가군들 사이의 매입이며,
$H^p(X, \mathcal{F}') = 0$이 모든 $p \geq 1$에 대해 성립한다.
\item $\{H^n(X, -)\}_{n \geq 0}$은 보편 $\delta$-함자로서
$\QCoh(\mathcal{O}_X)$에서 $\textit{Ab}$로 간다.
\end{enumerate}
\end{lemma}

\begin{proof}
$X = \bigcup U_i$를 유한 아핀 열린 덮개라 하자.
$U = \coprod U_i$로 두고 사상을 $j : U \to X$라 하되,
이 사상이 주어진 열린 몰입 $U_i \to X$들을 유도하게 한다.
$U$는 아핀 스킴이고 $X$의 대각사상이 아핀이므로
사상 $j$는 아핀이다. Morphisms의 보조정리
\ref{morphisms-lemma-affine-permanence}를 참조하라.
모든 $\mathcal{O}_X$-가군 $\mathcal{F}$에 대해 표준 사상
$\mathcal{F} \to j_*j^*\mathcal{F}$가 존재한다.
이 사상이 단사임은 줄기에서 확인할 수 있다.
$x \in U_i$이면 다음 분해가 있다.
$$
\mathcal{F}_x \to (j_*j^*\mathcal{F})_x
\to (j^*\mathcal{F})_{x'} = \mathcal{F}_x
$$
여기서 $x' \in U$는 $x$를 $U_i \subset U$의 점으로 본 것이다.
이제 $\mathcal{F}$가 준연접이면 $j^*\mathcal{F}$는 아핀 스킴
$U$ 위에서 준연접이므로 보조정리
\ref{coherent-lemma-quasi-coherent-affine-cohomology-zero}에 의해
고차 코호몰로지가 소멸한다.
그러면 $H^p(X, j_*j^*\mathcal{F}) = 0$이 모든
$p > 0$에 대해 성립한다. 이는 $j$가 아핀이므로 보조정리
\ref{coherent-lemma-relative-affine-cohomology}에서 따른다. 이로써 (1)을 보였다.
마지막으로 사상
$H^p(X, \mathcal{F}) \to H^p(X, j_*j^*\mathcal{F})$가 영임을 알 수 있고,
(2)는 Homology의 보조정리
\ref{homology-lemma-efface-implies-universal}에서 따른다.
\end{proof}

\begin{lemma}
\label{coherent-lemma-vanishing-nr-affines-quasi-separated}
$X$를 준콤팩트 준분리 스킴이라 하자.
$X = U_1 \cup \ldots \cup U_n$을 각 $U_i$가 준콤팩트이고
분리인 열린 덮개라 하자(예를 들어 아핀 열린집합들로 된 덮개).
다음과 같이 둔다.
$$
d = \max\nolimits_{I \subset \{1, \ldots, n\}}
\left(|I| + t(\bigcap\nolimits_{i \in I} U_i) - 1\right)
$$
여기서 $t(U)$는 스킴 $U$를 덮는 데 필요한 아핀 열린집합의
최소 개수이다. 그러면 $H^p(X, \mathcal{F}) = 0$이
모든 $p \geq d$와 모든 준연접층 $\mathcal{F}$에 대해 성립한다.
\end{lemma}

\begin{proof}
$X$는 준분리이고 $U_i$는 준콤팩트이므로 수
$t(\bigcap_{i \in I} U_i)$들은 유한하다.
$n$에 대한 귀납법으로 증명한다.
$n = 1$이면 결과는 보조정리 \ref{coherent-lemma-vanishing-nr-affines}에서 따른다.
$n > 1$이면 $X = U \cup V$로 쓰되,
$U = U_1 \cup \ldots \cup U_{n - 1}$ 및 $V = U_n$으로 둔다.
다음 마이어--피토리스 긴 완전열을 적용한다.
$$
0 \to
H^0(X, \mathcal{F}) \to
H^0(U, \mathcal{F}) \oplus H^0(V, \mathcal{F}) \to
H^0(U \cap V, \mathcal{F}) \to
H^1(X, \mathcal{F}) \to \ldots
$$
Cohomology의 보조정리
\ref{cohomology-lemma-mayer-vietoris}를 참조하라.
$q \geq d$에 대한 증명을 끝내려면
$H^q(V, \mathcal{F})$, $H^q(U, \mathcal{F})$, 그리고
$H^{q - 1}(U \cap V, \mathcal{F})$가 소멸함을 보이면 된다.
$n = 1$인 경우에 의해 $H^q(V, \mathcal{F}) = 0$이
$q \geq t(V) = t(U_n)$에서 성립한다. $t(V) \leq d$이므로 원하는 바가
따른다. 귀납 가정에 의해 다음 범위에서
$H^q(U, \mathcal{F}) = 0$이다.
$$
q \geq \max\nolimits_{I \subset \{1, \ldots, n - 1\}}
\left(|I| + t(\bigcap\nolimits_{i \in I} U_i) - 1\right)
$$
오른쪽 정수는 $d$ 이하이므로 원하는 바가 따른다.
마지막으로 열린집합
$U \cap V = (U_1 \cap U_n) \cup \ldots \cup (U_{n - 1} \cap U_n)$에
귀납 가정을 적용하면 다음 범위에서
$H^q(U \cap V, \mathcal{F}) = 0$이다.
$$
q \geq \max\nolimits_{I \subset \{1, \ldots, n - 1\}}
\left(|I| + t(U_n \cap \bigcap\nolimits_{i \in I} U_i) - 1\right)
$$
오른쪽 정수는 $d$보다 작으므로 보조정리가 따른다.
\end{proof}

\begin{lemma}
\label{coherent-lemma-quasi-coherence-higher-direct-images}
$f : X \to S$를 스킴의 사상이라 하자.
$f$가 준분리이고 준콤팩트이라고 가정하자.
\begin{enumerate}
\item 임의의 준연접 $\mathcal{O}_X$-가군 $\mathcal{F}$에 대해
고차 직상 $R^pf_*\mathcal{F}$는 $S$ 위에서 준연접이다.
\item $S$가 준콤팩트이면 정수 $n = n(X, S, f)$가 존재하여
$R^pf_*\mathcal{F} = 0$이다. 이는 모든 $p \geq n$과 임의의
준연접층 $\mathcal{F}$에 대해 성립하며, 그 층은 $X$ 위에 있다.
\item 실제로 $S$가 준콤팩트이면 $n = n(X, S, f)$를 택하여,
스킴의 모든 사상 $S' \to S$에 대해
$R^p(f')_*\mathcal{F}' = 0$이 되게 할 수 있다. 이는 모든
$p \geq n$과 임의의 준연접층 $\mathcal{F}'$에 대해 성립하며,
그 층은 $X'$ 위에 있다.
여기서 $f' : X' = S' \times_S X \to S'$는 $f$의 밑변환이다.
\end{enumerate}
\end{lemma}

\begin{proof}
먼저 (1)을 증명한다. 보조정리의 가정 아래에서 층
$R^0f_*\mathcal{F} = f_*\mathcal{F}$는 Schemes의 보조정리
\ref{schemes-lemma-push-forward-quasi-coherent}에 의해 준연접이다.
Cohomology의 보조정리
\ref{cohomology-lemma-localize-higher-direct-images}를 사용하면
고차 직상의 형성은 열린 부분스킴으로의 제한과 가환함을 알 수 있다.
준연접성은 $S$ 위에서 국소적이므로 다음 문단에서 다루는 경우로
귀착된다.

\medskip\noindent
$S$가 아핀인 경우의 (1)의 증명. 귀납 원리를 사용한다.
$f$가 준콤팩트 준분리이므로 $X$도 준콤팩트 준분리이다.
$U \subset X$가 준콤팩트 열린집합이고 $a = f|_U$일 때,
$P(U)$를 $R^pa_*\mathcal{F}$가 $S$ 위에서 준연접이라는 성질로
둔다. 이 조건은 모든 준연접 가군 $\mathcal{F}$, 곧 $U$ 위의
가군과 모든 $p \geq 0$에 대해 요구한다. $P(X)$가 (1)이므로, 보조정리
\ref{coherent-lemma-induction-principle}의 조건 (1), (2)가 성립함을 보이면 충분하다.
$U$가 아핀이면 $P(U)$가 성립한다. 실제로
$R^pa_*\mathcal{F} = 0$이 $p \geq 1$에 대해 성립한다
(보조정리 \ref{coherent-lemma-relative-affine-vanishing} 및 Morphisms의 보조정리
\ref{morphisms-lemma-morphism-affines-affine}에 의함).
$p = 0$인 경우에는 첫 문단에서 이미 결과를 확인하였다.
다음으로 $U \subset X$를 준콤팩트 열린집합, $V \subset X$를
아핀 열린집합이라 하고 $P(U)$, $P(V)$, $P(U \cap V)$가 성립한다고
가정하자.
$a = f|_U$, $b = f|_V$, $c = f|_{U \cap V}$,
$g = f|_{U \cup V}$로 두자.
그러면 임의의 준연접 $\mathcal{O}_{U \cup V}$-가군 $\mathcal{F}$에 대해
다음 상대 마이어--피토리스 열이 있다.
$$
0 \to
g_*\mathcal{F} \to
a_*(\mathcal{F}|_U) \oplus b_*(\mathcal{F}|_V) \to
c_*(\mathcal{F}|_{U \cap V}) \to
R^1g_*\mathcal{F} \to \ldots
$$
Cohomology의 보조정리
\ref{cohomology-lemma-relative-mayer-vietoris}를 참조하라.
$P(U)$, $P(V)$, $P(U \cap V)$에 의해
$R^pa_*(\mathcal{F}|_U)$, $R^pb_*(\mathcal{F}|_V)$ 및
$R^pc_*(\mathcal{F}|_{U \cap V})$는 모두 준연접이다.
Schemes의 절 \ref{schemes-section-quasi-coherent}에 있는
준연접층에 관한 결과를 쓰면, 각 층
$R^pg_*\mathcal{F}$는 준연접임이 따른다. 실제로 이 층은 왼쪽에
준연접층 사이 사상의 여핵을, 오른쪽에 준연접층 사이 사상의 핵을 둔
짧은 완전열의 가운데 놓인다.
따라서 $P(U \cup V)$가 성립하며 (1)의 증명이 끝난다.

\medskip\noindent
이제 (3)을 증명하면 특히 (2)도 따른다. 유한 아핀 열린 덮개
$S = \bigcup_{j = 1, \ldots m} S_j$를 택한다. 각 $j$에 대해
유한 아핀 열린 덮개
$f^{-1}(S_j) = \bigcup_{i = 1, \ldots t_j} U_{ji} $를 택한다.
다음 정수를 두자.
$$
d_j = \max\nolimits_{I \subset \{1, \ldots, t_j\}}
\left(|I| + t(\bigcap\nolimits_{i \in I} U_{ji})\right)
$$
이는 보조정리
\ref{coherent-lemma-vanishing-nr-affines-quasi-separated}에서 얻은 정수이다.
$n(X, S, f) = \max d_j$가 조건을 만족한다고 주장한다.

\medskip\noindent
실제로 $S' \to S$를 스킴의 사상이라 하고
$\mathcal{F}'$를 $X' = S' \times_S X$ 위의 준연접층이라 하자.
$R^pf'_*\mathcal{F}' = 0$이 $p \geq n(X, S, f)$에 대해 성립함을 보이고자
한다. 이 문제는 $S'$ 위에서 국소적이므로 $S'$가 아핀이고 어떤
$S_j$ 안으로 사상된다고 가정할 수 있다. 그 지표를 $j$라 하자.
그러면 $X' = S' \times_{S_j} f^{-1}(S_j)$는 아핀 열린집합
$S' \times_{S_j} U_{ji}$, $i = 1, \ldots t_j$로 덮이고, 그 교집합은
$$
\bigcap\nolimits_{i \in I} S' \times_{S_j} U_{ji} =
S' \times_{S_j} \bigcap\nolimits_{i \in I} U_{ji}
$$
이며 밑변환 전과 같은 개수의 아핀 열린집합으로 덮인다.
보조정리
\ref{coherent-lemma-vanishing-nr-affines-quasi-separated}를 적용하면
$H^p(X', \mathcal{F}') = 0$을 얻는다. 증명의 첫 부분에서 이미 각
$R^qf'_*\mathcal{F}'$가 준연접임을 알았으므로, 아핀 스킴
$S'$ 위에서 이 층들의 고차 코호몰로지 군은 소멸한다.
따라서 Cohomology의 보조정리
\ref{cohomology-lemma-apply-Leray}에 의해
$H^0(S', R^pf'_*\mathcal{F}') = H^p(X', \mathcal{F}') = 0$이다.
$R^pf'_*\mathcal{F}'$가 준연접이므로
$R^pf'_*\mathcal{F}' = 0$을 얻는다.
\end{proof}

\begin{lemma}
\label{coherent-lemma-quasi-coherence-higher-direct-images-application}
$f : X \to S$를 스킴의 사상이라 하자.
$f$가 준분리이고 준콤팩트이라고 가정하자.
$S$가 아핀이라고 가정하자.
임의의 준연접 $\mathcal{O}_X$-가군 $\mathcal{F}$에 대해
모든 $q \in \mathbf{Z}$에서
$$
H^q(X, \mathcal{F}) = H^0(S, R^qf_*\mathcal{F})
$$
가 성립한다.
\end{lemma}

\begin{proof}
르레 스펙트럼 열
$E_2^{p, q} = H^p(S, R^qf_*\mathcal{F})$을 생각하자. 이는
$H^{p + q}(X, \mathcal{F})$로 수렴한다.
Cohomology의 보조정리 \ref{cohomology-lemma-Leray}를 참조하라.
보조정리 \ref{coherent-lemma-quasi-coherence-higher-direct-images}에 의해
층 $R^qf_*\mathcal{F}$는 준연접이다.
보조정리 \ref{coherent-lemma-quasi-coherent-affine-cohomology-zero}에 의해
$E_2^{p, q} = 0$이 $p > 0$에서 성립한다.
따라서 스펙트럼 열은 $E_2$에서 퇴화하며 결론을 얻는다.
일반 원리에 대해서는 Cohomology의 보조정리
\ref{cohomology-lemma-apply-Leray} (2)도 참조하라.
\end{proof}







\section{코호몰로지와 밑변환, I}
\label{coherent-section-cohomology-and-base-change}

\noindent
$f : X \to S$를 스킴의 사상이라 하자.
$\mathcal{F}$를 $X$ 위의 준연접층이라 하자.
또한 $g : S' \to S$를 임의의 스킴 사상이라 하자.
밑변환을 $X' = X_{S'} = S' \times_S X$로 쓰며, 이는 $X$의
밑변환이다. 또 $f' : X' \to S'$를 $f$의 밑변환이라 하자.
또 사영을 $g' : X' \to X$로 쓰고
$\mathcal{F}' = (g')^*\mathcal{F}$로 둔다.
이 상황은 다음 그림으로 나타난다.
\begin{equation}
\label{coherent-equation-base-change-diagram}
\vcenter{
\xymatrix{
\mathcal{F}' = (g')^*\mathcal{F} &
X' \ar[r]_{g'} \ar[d]_{f'} &
X \ar[d]^f &
\mathcal{F} \\
Rf'_*\mathcal{F}' &
S' \ar[r]^g &
S &
Rf_*\mathcal{F}
}
}
\end{equation}
염두에 둔 밑변환 성질의 가장 간단한 경우는 다음과 같다.

\begin{lemma}
\label{coherent-lemma-affine-base-change}
$f : X \to S$를 스킴의 사상이라 하자.
$\mathcal{F}$를 준연접 $\mathcal{O}_X$-가군이라 하자.
$f$가 아핀이라고 가정하자.
이 경우 $f_*\mathcal{F} \cong Rf_*\mathcal{F}$는 준연접층이고,
모든 밑변환 그림
(\ref{coherent-equation-base-change-diagram})에 대해
$$
g^*f_*\mathcal{F} = f'_*(g')^*\mathcal{F}.
$$
가 성립한다.
\end{lemma}

\begin{proof}
고차 직상의 소멸은 보조정리
\ref{coherent-lemma-relative-affine-vanishing}이다.
주장은 $S$와 $S'$ 위에서 국소적이다. 따라서
$X = \Spec(A)$, $S = \Spec(R)$,
$S' = \Spec(R')$ 및 $\mathcal{F} = \widetilde{M}$이라고 가정할 수
있으며, 여기서 $A$-가군 $M$을 적절히 택한다.
Schemes의 보조정리 \ref{schemes-lemma-widetilde-pullback}을 사용하여
$\mathcal{F}$의 역상과 직상을 기술한다.
구체적으로 $X' = \Spec(R' \otimes_R A)$이고,
$\mathcal{F}'$는 $(R' \otimes_R A) \otimes_A M$에 대응하는
준연접층이다.
따라서 보조정리는 $R'$-가군으로서의 다음 등식으로 귀착된다.
$$
(R' \otimes_R A) \otimes_A M = R' \otimes_R M
$$
\end{proof}

\noindent
많은 상황에서는 다음과 같은 코호몰로지와 밑변환의 특수한 경우만
알아도 충분하다. 이는 절
\ref{coherent-section-cohomology-and-base-change-derived}의 더 강한 결과에서
곧바로 따르지만, 매우 중요하므로 별도의 증명을 제시할 가치가 있다.

\begin{lemma}[평탄 밑변환]
\label{coherent-lemma-flat-base-change-cohomology}
다음 스킴의 데카르트 그림을 생각하자.
$$
\xymatrix{
X' \ar[d]_{f'} \ar[r]_{g'} & X \ar[d]^f \\
S' \ar[r]^g & S
}
$$
$\mathcal{F}$를 준연접 $\mathcal{O}_X$-가군이라 하고,
그 역상을 $\mathcal{F}' = (g')^*\mathcal{F}$라 하자.
$g$가 평탄이고 $f$가 준콤팩트 준분리라고 가정하자.
임의의 $i \geq 0$에 대해 다음이 성립한다.
\begin{enumerate}
\item Cohomology의 보조정리
\ref{cohomology-lemma-base-change-map-flat-case}의 밑변환 사상
$$
g^*R^if_*\mathcal{F} \longrightarrow R^if'_*\mathcal{F}',
$$
은 동형사상이다.
\item $S = \Spec(A)$이고 $S' = \Spec(B)$이면
$H^i(X, \mathcal{F}) \otimes_A B = H^i(X', \mathcal{F}')$이다.
\end{enumerate}
\end{lemma}

\begin{proof}
(1)에서 Cohomology의 보조정리
\ref{cohomology-lemma-base-change-map-flat-case}을 사용할 수 있는 것은,
Morphisms의 보조정리 \ref{morphisms-lemma-base-change-flat}에 의해
$g'$가 평탄이기 때문이다.
이제 (1)은 (2)에서 따른다. 실제로 (1)은 $S'$ 위에서 국소적이므로
$S$와 $S'$가 아핀이라고 가정할 수 있다.
즉 (2)에서처럼 $S = \Spec(A)$ 및 $S' = \Spec(B)$이다.
그러면 $R^if_*\mathcal{F}$는 준연접이므로
(보조정리 \ref{coherent-lemma-quasi-coherence-higher-direct-images}),
준연접 $\mathcal{O}_S$-가군이며, 여기에 대응하는 $A$-가군은
$H^0(S, R^if_*\mathcal{F}) = H^i(X, \mathcal{F})$이다
(등식은 보조정리
\ref{coherent-lemma-quasi-coherence-higher-direct-images-application}에 의함).
마찬가지로 $R^if'_*\mathcal{F}'$는 준연접
$\mathcal{O}_{S'}$-가군이며, 여기에 대응하는 $B$-가군은
$H^i(X', \mathcal{F}')$이다. $g$에 의한 역상은 가군에서
$- \otimes_A B$에 대응하므로
(Schemes의 보조정리 \ref{schemes-lemma-widetilde-pullback}),
(2)를 증명하면 충분하다.

\medskip\noindent
$A \to B$를 평탄 환 준동형이라 하자.
$X$를 $A$ 위의 준콤팩트 준분리 스킴이라 하자.
$\mathcal{F}$를 준연접 $\mathcal{O}_X$-가군이라 하자.
$X_B = X \times_{\Spec(A)} \Spec(B)$로 두고,
$\mathcal{F}_B$를 $\mathcal{F}$의 역상이라 하자.
다음 사상이 동형사상임을 보이고자 한다.
$$
H^i(X, \mathcal{F}) \otimes_A B \longrightarrow H^i(X_B, \mathcal{F}_B)
$$
이 사상은 보조정리의 명제에 인용된 결과로 주어진다.

\medskip\noindent
$X$가 분리인 경우에는 아핀 열린 덮개
$\mathcal{U} : X = U_1 \cup \ldots \cup U_t$를 택하고
다음을 상기하자.
$$
\check{H}^p(\mathcal{U}, \mathcal{F}) = H^p(X, \mathcal{F}),
$$
보조정리 \ref{coherent-lemma-cech-cohomology-quasi-coherent}를 참조하라.
밑변환으로
$\mathcal{U}_B : X_B = (U_1)_B \cup \ldots \cup (U_t)_B$를 얻으면,
여전히 각 $(U_i)_B$는 아핀이고 $X_B$는 분리이다. 따라서
$$
\check{H}^p(\mathcal{U}_B, \mathcal{F}_B) = H^p(X_B, \mathcal{F}_B).
$$
체흐 복합체들 사이에는 다음 관계가 있다.
$$
\check{\mathcal{C}}^\bullet(\mathcal{U}_B, \mathcal{F}_B) =
\check{\mathcal{C}}^\bullet(\mathcal{U}, \mathcal{F}) \otimes_A B
$$
이는 보조정리 \ref{coherent-lemma-affine-base-change}에서 따른다.
$A \to B$가 평탄이므로 코호몰로지를 취한 뒤에도 같은 명제가 성립한다.

\medskip\noindent
$X$가 준분리인 경우에는 아핀 열린 덮개
$\mathcal{U} : X = U_1 \cup \ldots \cup U_t$를 택한다.
Cohomology의 보조정리
\ref{cohomology-lemma-cech-spectral-sequence}의
체흐-코호몰로지 스펙트럼 열을 사용할 것이다.
이 스펙트럼 열을 피하고 싶은 독자는 보조정리
\ref{coherent-lemma-quasi-coherence-higher-direct-images}의 증명에서처럼
마이어--피토리스 열과 $t$에 대한 귀납법을 사용할 수 있다.
이 스펙트럼 열의 $E_2$-쪽은
$E_2^{p, q} = \check{H}^p(\mathcal{U}, \underline{H}^q(\mathcal{F}))$이고
$H^{p + q}(X, \mathcal{F})$로 수렴한다.
마찬가지로 $E_2$-쪽이
$E_2^{p, q} = \check{H}^p(\mathcal{U}_B, \underline{H}^q(\mathcal{F}_B))$인
스펙트럼 열이 있으며, 이는 $H^{p + q}(X_B, \mathcal{F}_B)$로 수렴한다.
교집합 $U_{i_0 \ldots i_p}$들은 준콤팩트이고 분리이므로,
증명의 둘째 문단의 결과에 의해
$\check{H}^p(\mathcal{U}_B, \underline{H}^q(\mathcal{F}_B)) =
\check{H}^p(\mathcal{U}, \underline{H}^q(\mathcal{F})) \otimes_A B$이다.
$A \to B$의 평탄성을 사용하면
$H^i(X, \mathcal{F}) \otimes_A B \to H^i(X_B, \mathcal{F}_B)$가
모든 $i$에 대해 동형사상임을 얻고 증명이 끝난다.
\end{proof}

\begin{lemma}[유한 국소 자유 밑변환]
\label{coherent-lemma-finite-locally-free-base-change-cohomology}
다음 스킴의 데카르트 그림을 생각하자.
$$
\xymatrix{
Y \ar[d]_{g} \ar[r]_h & X \ar[d]^f \\
\Spec(B) \ar[r] & \Spec(A)
}
$$
$\mathcal{F}$를 준연접 $\mathcal{O}_X$-가군이라 하고,
그 역상을 $\mathcal{G} = h^*\mathcal{F}$라 하자.
$B$가 유한 국소 자유 $A$-가군이면
$H^i(X, \mathcal{F}) \otimes_A B = H^i(Y, \mathcal{G})$이다.
\end{lemma}

\noindent
{\bf 경고}: 이 보조정리와 보조정리
\ref{coherent-lemma-flat-base-change-cohomology}의 차이를 이해하지 못했다면
이 보조정리를 사용하지 말라.

\begin{proof}
$X$가 분리인 경우에는 아핀 열린 덮개
$\mathcal{U} : X = \bigcup_{i \in I} U_i$를 택하고
다음을 상기하자.
$$
\check{H}^p(\mathcal{U}, \mathcal{F}) = H^p(X, \mathcal{F}),
$$
보조정리 \ref{coherent-lemma-cech-cohomology-quasi-coherent}를 참조하라.
$\mathcal{V} : Y = \bigcup_{i \in I} g^{-1}(U_i)$를
$Y$의 대응하는 아핀 열린 덮개라 하자.
열린집합 $V_i = g^{-1}(U_i) = U_i \times_{\Spec(A)} \Spec(B)$는
아핀이고 $Y$는 분리이다. 따라서
$$
\check{H}^p(\mathcal{V}, \mathcal{G}) = H^p(Y, \mathcal{G}).
$$
다음 체흐 복합체의 사상이 동형사상이라고 주장한다.
$$
\check{\mathcal{C}}^\bullet(\mathcal{U}, \mathcal{F}) \otimes_A B
\longrightarrow
\check{\mathcal{C}}^\bullet(\mathcal{V}, \mathcal{G})
$$
실제로 $B$는 $A$-가군으로서 유한 표시이므로,
$B$와 $A$ 위에서 텐서곱을 취하는 것은 곱과 가환한다.
Algebra의 명제 \ref{algebra-proposition-fp-tensor}를 참조하라.
따라서 사상
$\Gamma(U_{i_0 \ldots i_p}, \mathcal{F}) \otimes_A B \to
\Gamma(V_{i_0 \ldots i_p}, \mathcal{G})$ 
이 동형사상임을 보이면 충분하며, 이는 보조정리
\ref{coherent-lemma-affine-base-change}에서 따른다.
$A \to B$가 평탄이므로 코호몰로지를 취한 뒤에도 같은 명제가 성립한다.

\medskip\noindent
일반적인 경우에도 아핀 열린 덮개
$\mathcal{U} : X = \bigcup_{i \in I} U_i$와 이에 대응하는 덮개
$\mathcal{V} : Y = \bigcup_{i \in I} V_i$를 사용하여 정확히 같은
방식으로 논증한다. 여기서 위와 같이 $V_i = g^{-1}(U_i)$이다.
Cohomology의 보조정리
\ref{cohomology-lemma-cech-spectral-sequence}의
체흐-코호몰로지 스펙트럼 열을 사용할 것이다.
이 스펙트럼 열의 $E_2$-쪽은
$E_2^{p, q} = \check{H}^p(\mathcal{U}, \underline{H}^q(\mathcal{F}))$이고
$H^{p + q}(X, \mathcal{F})$로 수렴한다.
마찬가지로 $E_2$-쪽이
$E_2^{p, q} = \check{H}^p(\mathcal{V}, \underline{H}^q(\mathcal{G}))$인
스펙트럼 열이 있으며, 이는 $H^{p + q}(Y, \mathcal{G})$로 수렴한다.
교집합 $U_{i_0 \ldots i_p}$들이 분리이므로 앞 문단의 결과는
다음 동형사상을 준다.
$\Gamma(U_{i_0 \ldots i_p}, \underline{H}^q(\mathcal{F})) \otimes_A B
\to \Gamma(V_{i_0 \ldots i_p}, \underline{H}^q(\mathcal{G}))$.
$- \otimes_A B$가 곱과 가환하고 완전하므로
$\check{H}^p(\mathcal{U}, \underline{H}^q(\mathcal{F})) \otimes_A B
\to \check{H}^p(\mathcal{V}, \underline{H}^q(\mathcal{G}))$가
동형사상임을 얻는다. $A \to B$가 평탄임을 사용하면
$H^i(X, \mathcal{F}) \otimes_A B \to H^i(Y, \mathcal{G})$가
모든 $i$에 대해 동형사상임을 얻고 증명이 끝난다.
\end{proof}








\section{쌍극한과 고차 직상}
\label{coherent-section-colimits}

\noindent
이러한 종류의 일반적인 결과는 Cohomology의 절
\ref{cohomology-section-limits}, Sheaves의 보조정리
\ref{sheaves-lemma-directed-colimits-sections}, 그리고 Modules의
보조정리
\ref{modules-lemma-finite-presentation-quasi-compact-colimit}에서
찾을 수 있다.

\begin{lemma}
\label{coherent-lemma-colimit-cohomology}
$f : X \to S$를 준콤팩트 준분리 스킴 사상이라 하자.
$\mathcal{F} = \colim \mathcal{F}_i$를 $X$ 위의 아벨 층들의
여과 쌍극한이라 하자. 그러면 임의의 $p \geq 0$에 대해
$$
R^pf_*\mathcal{F} = \colim R^pf_*\mathcal{F}_i.
$$
\end{lemma}

\begin{proof}
Cohomology의 보조정리
\ref{cohomology-lemma-higher-direct-image-colimit}를 적용할 것이다.
아핀 열린집합들이 $S$의 위상에 대한 기저를 이루므로,
아핀 열린집합 $U \subset S$에 대해
$H^p(f^{-1}U, \mathcal{F}) = \colim H^p(f^{-1}U, \mathcal{F}_i)$임을
보이면 충분하다.
$f^{-1}U$는 준콤팩트 준분리이므로 Cohomology의 보조정리
\ref{cohomology-lemma-quasi-separated-cohomology-colimit}에서
결론이 따른다.
(이는 $f^{-1}U$의 아핀 열린집합 기저가 그 보조정리의 가정을
만족하기 때문이다.)
\end{proof}









\section{코호몰로지와 밑변환, II}
\label{coherent-section-cohomology-and-base-change-derived}

\noindent
$f : X \to S$를 스킴의 사상이라 하고 $\mathcal{F}$를 준연접
$\mathcal{O}_X$-가군이라 하자. $f$가 준콤팩트 준분리이면
$Rf_*\mathcal{F}$를 $S$ 위의 준연접층들의 복합체로 나타내고자 한다.
이는 층 $R^if_*\mathcal{F}$가 준연접이라는 사실에서 따르며,
여기서 $S$는 준콤팩트이고 대각사상이 아핀이라고 가정한다.
여기서는 $D_\QCoh(S)$가 $D(\QCoh(\mathcal{O}_S))$와 동치임을
사용한다. Derived Categories of Schemes의 명제
\ref{perfect-proposition-quasi-compact-affine-diagonal}을 참조하라.

\medskip\noindent
이 절에서는 좋은 밑변환 성질을 갖는 명시적인 복합체를 만들어 내는
다른 접근법을 사용할 것이다. $f$와 $S$가 분리이거나, 더 일반적으로
대각사상이 아핀인 경우에는 이 구성이 특히 간단하다.
실제로 가장 자주 쓰이는 경우이므로 이를 따로 다룬다.

\begin{lemma}
\label{coherent-lemma-separated-case-relative-cech}
$f : X \to S$를 스킴의 사상이라 하자.
$\mathcal{F}$를 준연접 $\mathcal{O}_X$-가군이라 하자.
$X$가 준콤팩트이고 $X$와 $S$의 대각사상이 아핀이라고 가정하자
(예를 들어 $X$와 $S$가 분리이면 그러하다).
이 경우 $Rf_*\mathcal{F}$를 다음과 같이 계산할 수 있다.
\begin{enumerate}
\item 유한 아핀 열린 덮개
$\mathcal{U} : X = \bigcup_{i = 1, \ldots, n} U_i$를 택한다.
\item $i_0, \ldots, i_p \in \{1, \ldots, n\}$에 대해
$f_{i_0 \ldots i_p} : U_{i_0 \ldots i_p} \to S$를 $f$를 교집합
$U_{i_0 \ldots i_p} = U_{i_0} \cap \ldots \cap U_{i_p}$에 제한한
사상이라 하자.
\item $\mathcal{F}_{i_0 \ldots i_p}$를 $\mathcal{F}$를
$U_{i_0 \ldots i_p}$에 제한한 층으로 둔다.
\item 다음과 같이 둔다.
$$
\check{\mathcal{C}}^p(\mathcal{U}, f, \mathcal{F}) =
\bigoplus\nolimits_{i_0 \ldots i_p}
f_{i_0 \ldots i_p *} \mathcal{F}_{i_0 \ldots i_p}
$$
그리고 미분
$d : \check{\mathcal{C}}^p(\mathcal{U}, f, \mathcal{F})
\to \check{\mathcal{C}}^{p + 1}(\mathcal{U}, f, \mathcal{F})$를
Cohomology의 식 (\ref{cohomology-equation-d-cech})과 같이 정의한다.
\end{enumerate}
그러면 복합체
$\check{\mathcal{C}}^\bullet(\mathcal{U}, f, \mathcal{F})$는
$S$ 위의 준연접층들의 복합체이고, 다음 동형사상을 갖는다.
$$
\check{\mathcal{C}}^\bullet(\mathcal{U}, f, \mathcal{F})
\longrightarrow
Rf_*\mathcal{F}
$$
이는 $D^{+}(S)$ 안의 동형사상이며, 준연접층 $\mathcal{F}$에 대해
함자적이다.
\end{lemma}

\begin{proof}
Cohomology의 보조정리
\ref{cohomology-lemma-covering-resolution}의 분해
$\mathcal{F} \to {\mathfrak C}^\bullet(\mathcal{U}, \mathcal{F})$를
생각하자. 다음 복합체의 등식이 있다.
$\check{\mathcal{C}}^\bullet(\mathcal{U}, f, \mathcal{F}) =
f_*{\mathfrak C}^\bullet(\mathcal{U}, \mathcal{F})$
이는 준연접 $\mathcal{O}_S$-가군의 복합체로서의 등식이다.
사상 $j_{i_0 \ldots i_p} : U_{i_0 \ldots i_p} \to X$와 사상
$f_{i_0 \ldots i_p} : U_{i_0 \ldots i_p} \to S$는
Morphisms의 보조정리 \ref{morphisms-lemma-affine-permanence}와
보조정리 \ref{coherent-lemma-affine-diagonal}에 의해 아핀이다.
따라서 $R^qj_{i_0 \ldots i_p *}\mathcal{F}_{i_0 \ldots i_p}$와
$R^qf_{i_0 \ldots i_p *}\mathcal{F}_{i_0 \ldots i_p}$는 모두
$q > 0$에서 영이다
(보조정리 \ref{coherent-lemma-relative-affine-vanishing}).
$f \circ j_{i_0 \ldots i_p} = f_{i_0 \ldots i_p}$와 Cohomology의
보조정리 \ref{cohomology-lemma-relative-Leray}의 스펙트럼 열을
사용하면
$R^qf_*(j_{i_0 \ldots i_p *}\mathcal{F}_{i_0 \ldots i_p}) = 0$이
$q > 0$에서 성립함을 얻는다.
복합체 ${\mathfrak C}^\bullet(\mathcal{U}, \mathcal{F})$의 항들은 층
$j_{i_0 \ldots i_p *}\mathcal{F}_{i_0 \ldots i_p}$들의 유한 직합이므로,
르레의 비순환성 보조정리
(Derived Categories의 보조정리
\ref{derived-lemma-leray-acyclicity})를 사용하여 원하는 등식을 얻는다.
$$
Rf_* \mathcal{F} = f_*{\mathfrak C}^\bullet(\mathcal{U}, \mathcal{F}) =
\check{\mathcal{C}}^\bullet(\mathcal{U}, f, \mathcal{F})
$$
\end{proof}

\noindent
다음으로 밑변환을 했을 때 어떤 일이 일어나는지 살펴본다.

\begin{lemma}
\label{coherent-lemma-base-change-complex}
그림 (\ref{coherent-equation-base-change-diagram})의 표기를 사용하자.
$f : X \to S$와 $\mathcal{F}$가 보조정리
\ref{coherent-lemma-separated-case-relative-cech}의 가정을 만족한다고 하자.
유한 아핀 열린 덮개
$\mathcal{U} : X = \bigcup U_i$를 택하되, 이는 $X$의 덮개이다.
다음 표준 동형사상이 존재한다.
$$
g^*\check{\mathcal{C}}^\bullet(\mathcal{U}, f, \mathcal{F})
\longrightarrow
Rf'_*\mathcal{F}'
$$
이는 $D^{+}(S')$ 안의 동형사상이다. 더욱이 $S' \to S$가 아핀이면
실제로
$$
g^*\check{\mathcal{C}}^\bullet(\mathcal{U}, f, \mathcal{F})
=
\check{\mathcal{C}}^\bullet(\mathcal{U}', f', \mathcal{F}')
$$
이고, 여기서 $\mathcal{U}' : X' = \bigcup U_i'$이며
$U_i' = (g')^{-1}(U_i) = U_{i, S'}$도 아핀이다.
\end{lemma}

\begin{proof}
사실 $U_i' = (g')^{-1}(U_i) = U_{i, S'}$로 정의할 수 있으며,
이는 $S'$가 $S$ 위에서 아핀인지와 관계없다.
$\mathcal{U}' : X' = \bigcup U_i'$를 $X'$의 유도된 덮개라 하자.
이 경우 다음이 성립한다고 주장한다.
$$
g^*\check{\mathcal{C}}^\bullet(\mathcal{U}, f, \mathcal{F})
=
\check{\mathcal{C}}^\bullet(\mathcal{U}', f', \mathcal{F}')
$$
여기서 $\check{\mathcal{C}}^\bullet(\mathcal{U}', f', \mathcal{F}')$는
보조정리 \ref{coherent-lemma-separated-case-relative-cech}와 정확히 같은
방식으로 정의한다.
이는 $S'$ 위에서 국소적으로 논증하고 아핀 사상의 경우
(보조정리 \ref{coherent-lemma-affine-base-change})를 적용하면 분명하다.
더욱이 보조정리 \ref{coherent-lemma-separated-case-relative-cech}의
증명과 정확히 같은 방식으로 다음 동형사상을 얻는다.
$$
\check{\mathcal{C}}^\bullet(\mathcal{U}', f', \mathcal{F}')
\longrightarrow
Rf'_*\mathcal{F}'
$$
이는 $D^{+}(S')$ 안의 동형사상이다. 실제로 사상
$U_i' \to X'$와 $U_i' \to S'$는 여전히 아핀이다
(아핀 사상의 밑변환이기 때문이다).
세부사항은 생략한다.
\end{proof}

\noindent
위 보조정리는 복합체
$$
\mathcal{K}^\bullet = \check{\mathcal{C}}^\bullet(\mathcal{U}, f, \mathcal{F})
$$
가 $S$ 위의 준연접층들로 이루어진 아래로 유계인 복합체이며,
$f : X \to S$의 고차 직상을 {\it 보편적으로} 계산한다는 뜻이다.
이는 이 특정 복합체에 관한 성질이며, 일반적으로
$\check{\mathcal{C}}^\bullet(\mathcal{U}, f, \mathcal{F})$를
준동형인 복합체로 바꾸면 보존되지 않는다!
다시 말해 유도 범주에서 의미를 갖는 명제가 아니다.
일반적으로 역상 $g^*\mathcal{K}^\bullet$은
유도 역상 $Lg^*\mathcal{K}^\bullet$, 곧
$\mathcal{K}^\bullet$의 유도 역상과 {\it 같지 않기} 때문이다!

\medskip\noindent
다음은 이 명제를 증명할 수 있는 더 일반적인 경우이다.
$S$가 분리라는 조건은 대부분의 응용에서 무해하다.
보통 전체 유도상의 어떤 국소적 성질을 증명하기 위해 이 결과를
사용하기 때문이다. 증명은 훨씬 더 복잡하고 초덮개를 사용한다.
이는 초덮개를 때때로 어떻게 활용할 수 있는지 보여 주는 좋은 예이다.

\begin{lemma}
\label{coherent-lemma-hypercoverings}
$f : X \to S$를 스킴의 사상이라 하자.
$\mathcal{F}$를 $X$ 위의 준연접층이라 하자.
$f$가 준콤팩트 준분리이고 $S$가 준콤팩트 분리라고 가정하자.
다음 성질을 갖는 아래로 유계인 복합체 $\mathcal{K}^\bullet$가
존재하며, 이는 준연접 $\mathcal{O}_S$-가군들의 복합체이다. 모든 사상
$g : S' \to S$에 대해 복합체 $g^*\mathcal{K}^\bullet$는
그림 (\ref{coherent-equation-base-change-diagram})의 표기를 사용할 때
$Rf'_*\mathcal{F}'$의 대표이다.
\end{lemma}

\begin{proof}
($f$도 분리인 경우에는 보조정리
\ref{coherent-lemma-base-change-complex}을 참조하라.)
가정에 의해 특히 $X$는 스킴으로서 준콤팩트 준분리이다.
$\mathcal{B}$를 $X$의 아핀 열린집합들의 집합이라 하자.
Hypercoverings의 보조정리
\ref{hypercovering-lemma-quasi-separated-quasi-compact-hypercovering}에
의해 초덮개 $K = (I, \{U_i\})$를 찾을 수 있다. 여기서 각
$I_n$은 유한이고 각 $U_i$는 $X$의 아핀 열린집합이다.
Hypercoverings의 보조정리
\ref{hypercovering-lemma-cech-spectral-sequence}에 의해
$E_2$-쪽이
$$
E_2^{p, q} = \check{H}^p(K, \underline{H}^q(\mathcal{F}))
$$
인 스펙트럼 열이 있고, 이는 $H^{p + q}(X, \mathcal{F})$로 수렴한다.
$\check{H}^p(K, \underline{H}^q(\mathcal{F}))$는 다음 복합체의
$p$번째 코호몰로지 군임에 유의하라.
$$
\prod\nolimits_{i \in I_0} H^q(U_i, \mathcal{F})
\to
\prod\nolimits_{i \in I_1} H^q(U_i, \mathcal{F})
\to
\prod\nolimits_{i \in I_2} H^q(U_i, \mathcal{F})
\to \ldots
$$
각 $U_i$가 아핀이므로 $q = 0$이 아니면 이것은 영이고,
나머지 경우에는 다음 복합체를 얻는다.
$$
\prod\nolimits_{i \in I_0} \mathcal{F}(U_i)
\to
\prod\nolimits_{i \in I_1} \mathcal{F}(U_i)
\to
\prod\nolimits_{i \in I_2} \mathcal{F}(U_i)
\to \ldots
$$
따라서 $R\Gamma(X, \mathcal{F})$는 이 복합체로 계산된다.

\medskip\noindent
임의의 $n$과 $i \in I_n$에 대해
$f_i : U_i \to S$를 $f$를 $U_i$에 제한한 사상이라 하자.
$S$는 분리이고 $U_i$는 아핀이므로 이 사상은 아핀이다.
다음 준연접층의 복합체를 생각하자.
$$
\mathcal{K}^\bullet = (
\prod\nolimits_{i \in I_0} f_{i, *}\mathcal{F}|_{U_i}
\to
\prod\nolimits_{i \in I_1} f_{i, *}\mathcal{F}|_{U_i}
\to
\prod\nolimits_{i \in I_2} f_{i, *}\mathcal{F}|_{U_i}
\to \ldots )
$$
이는 $S$ 위의 복합체이다. Hypercoverings의 보조정리
\ref{hypercovering-lemma-cech-spectral-sequence}에서와 같이
사상 $\mathcal{K}^\bullet \to Rf_*\mathcal{F}$를 얻는데, 이는
$D(\mathcal{O}_S)$ 안의 사상이고 $\mathcal{F}$의 단사 분해를 택하여
구성한다
(세부사항은 생략한다). 임의의 아핀 스킴 $V$와 사상
$g : V \to S$를 생각하자. 그러면 밑변환 $X_V$는 밑변환으로
얻은 초덮개 $K_V = (I, \{U_{i, V}\})$를 갖는다. 더욱이
$g^*f_{i, *}\mathcal{F} = f_{i, V, *}(g')^*\mathcal{F}|_{U_{i, V}}$이다.
따라서 위 논증에 의해 $\Gamma(V, g^*\mathcal{K}^\bullet)$는
$R\Gamma(X_V, (g')^*\mathcal{F})$를 계산한다.
복합체의 등식은 $S'$ 위에서 자리스키 국소적으로 증명하면
충분하므로, 이로써 보조정리의 증명이 끝난다.
\end{proof}

\noindent
다음 보조정리는 평탄 밑변환의 한 변형이다.

\begin{lemma}
\label{coherent-lemma-base-change-tensor-with-flat}
다음 스킴의 데카르트 그림을 생각하자.
$$
\xymatrix{
X' \ar[d]_{f'} \ar[r]_{g'} & X \ar[d]^f \\
S' \ar[r]^g & S
}
$$
$\mathcal{F}$를 준연접 $\mathcal{O}_X$-가군이라 하자.
$\mathcal{G}$를 준연접 $\mathcal{O}_{S'}$-가군이라 하고,
이 가군이 $S$ 위에서 평탄이라고 하자.
$f$가 준콤팩트 준분리라고 가정하자.
임의의 $i \geq 0$에 대해 다음 동일시가 있다.
$$
\mathcal{G} \otimes_{\mathcal{O}_{S'}} g^*R^if_*\mathcal{F} =
R^if'_*\left((f')^*\mathcal{G}
\otimes_{\mathcal{O}_{X'}} (g')^*\mathcal{F}\right)
$$
\end{lemma}

\begin{proof}
왼쪽에서 오른쪽으로 가는 사상을 구성하자. 먼저 밑변환 사상
$Lg^*Rf_*\mathcal{F} \to Rf'_*L(g')^*\mathcal{F}$가 있다.
또 수반 사상 $\mathcal{G} \to Rf'_*L(f')^*\mathcal{G}$가 있다.
상대 컵곱을 사용하면 다음을 얻는다.
\begin{align*}
\mathcal{G}
\otimes_{\mathcal{O}_{S'}}^\mathbf{L}
Lg^*Rf_*\mathcal{F}
& \to
Rf'_*L(f')^*\mathcal{G}
\otimes_{\mathcal{O}_{S'}}^\mathbf{L}
Rf'_*L(g')^*\mathcal{F} \\
& \to
Rf'_*\left(L(f')^*\mathcal{G}
\otimes_{\mathcal{O}_{X'}}^\mathbf{L}
L(g')^*\mathcal{F}\right) \\
& \to
Rf'_*\left((f')^*\mathcal{G}
\otimes_{\mathcal{O}_{X'}} (g')^*\mathcal{F}\right)
\end{align*}
가운데 화살표에는 상대 컵곱을 사용하였다. Cohomology의 주석
\ref{cohomology-remark-cup-product}을 참조하라.
합성의 시원은 다음과 같다.
$$
\mathcal{G} \otimes_{\mathcal{O}_{S'}}^\mathbf{L} Lg^*Rf_*\mathcal{F} =
\mathcal{G} \otimes_{g^{-1}\mathcal{O}_S}^\mathbf{L} g^{-1}Rf_*\mathcal{F}
$$
이는 Cohomology의 보조정리
\ref{cohomology-lemma-variant-derived-pullback}에서 따른다.
$\mathcal{G}$는 $g^{-1}\mathcal{O}_S$-가군의 층으로서 평탄이고
$g^{-1}$은 완전 함자이므로, 이 복합체의 $i$번째 코호몰로지 층은
$\mathcal{G} \otimes_{g^{-1}\mathcal{O}_S} g^{-1}R^if_*\mathcal{F} =
\mathcal{G} \otimes_{\mathcal{O}_{S'}} g^*R^if_*\mathcal{F}$이다.
이렇게 하여 보조정리의 등식에서 왼쪽에서 오른쪽으로 가는
전역 사상들을 얻는다. 이 사상이 동형사상임을 보이기 위해
$S'$ 위에서 국소적으로 논증할 수 있다. 따라서 $S$와 $S'$가
아핀 스킴이라고 가정한다.

\medskip\noindent
$S$와 $S'$가 아핀인 경우의 증명. $S = \Spec(A)$ 및
$S' = \Spec(B)$라 하고, $\mathcal{G}$가 $B$-가군 $N$에 대응하며
이 가군은 $A$-평탄이라고 하자. $S$가 아핀이므로 $X$는 준콤팩트
준분리이다. 초덮개 논증으로 증명을 끝낼 것이다.
$X$가 분리이거나 대각사상이 아핀이면 체흐 덮개를 사용할 수 있다.
$\mathcal{B}$를 $X$의 아핀 열린집합들의 집합이라 하자.
Hypercoverings의 보조정리
\ref{hypercovering-lemma-quasi-separated-quasi-compact-hypercovering}에
의해 초덮개 $K = (I, \{U_i\})$를 찾을 수 있으며, 이는 $X$의
초덮개이고,
각 $I_n$은 유한이고 각 $U_i$는 $X$의 아핀 열린집합이다.
Hypercoverings의 보조정리
\ref{hypercovering-lemma-cech-spectral-sequence}에 의해
$E_2$-쪽이 다음과 같은 스펙트럼 열이 있다.
$$
E_2^{p, q} = \check{H}^p(K, \underline{H}^q(\mathcal{F}))
$$
이는 $H^{p + q}(X, \mathcal{F})$로 수렴한다.
각 $U_i$가 아핀이고 $\mathcal{F}$가 준연접이므로,
$\underline{H}^q(\mathcal{F})$의 $U_i$에서의 값은 영이며,
이는 $q > 0$일 때 성립한다.
따라서 스펙트럼 열은 퇴화하고, 코호몰로지 가군
$H^q(X, \mathcal{F})$는 다음 복합체로 계산된다.
$$
\prod\nolimits_{i \in I_0} \mathcal{F}(U_i)
\to
\prod\nolimits_{i \in I_1} \mathcal{F}(U_i)
\to
\prod\nolimits_{i \in I_2} \mathcal{F}(U_i)
\to \ldots
$$
다음으로 초덮개를 $S'$로 밑변환하면
$X' = S' \times_S X$의 초덮개가 됨에 유의하라.
스킴 $S' \times_S U_i$도 아핀이고 다음 등식이 성립한다.
$$
\left((f')^*\mathcal{G} \otimes_{\mathcal{O}_{S'}} (g')^*\mathcal{F}\right)
(S' \times_S U_i) =
N \otimes_A \mathcal{F}(U_i)
$$
따라서 코호몰로지 가군
$H^q(X', (f')^*\mathcal{G} \otimes_{\mathcal{O}_{S'}} (g')^*\mathcal{F})$는
다음 복합체로 계산된다.
$$
N \otimes_A \left(
\prod\nolimits_{i \in I_0} \mathcal{F}(U_i)
\to
\prod\nolimits_{i \in I_1} \mathcal{F}(U_i)
\to
\prod\nolimits_{i \in I_2} \mathcal{F}(U_i)
\to \ldots
\right)
$$
$N$은 $A$ 위에서 평탄이므로 다음 등식을 얻는다.
$$
H^q(X', (f')^*\mathcal{G} \otimes_{\mathcal{O}_{S'}} (g')^*\mathcal{F})
=
N \otimes_A H^q(X, \mathcal{F})
$$
이는 증명해야 할 명제를 대수로 옮긴 것이므로 증명이 끝난다.
\end{proof}







\section{사영공간의 코호몰로지}
\label{coherent-section-cohomology-projective-space}

\noindent
이 절에서는 $\mathbf{P}^n_S$ 위 구조층의 꼬임들의 코호몰로지를
스킴 $S$ 위에서 계산한다.
$\mathbf{P}^n_S$가 Constructions, Definition
\ref{constructions-definition-projective-space}에서 섬유곱
$
\mathbf{P}^n_S = S \times_{\Spec(\mathbf{Z})} \mathbf{P}^n_{\mathbf{Z}}
$
으로 정의되었음을 상기하자. 또한 Constructions, Lemma
\ref{constructions-lemma-projective-space-bundle}에서 이것이
$$
\mathbf{P}^n_S = \underline{\text{Proj}}_S(\mathcal{O}_S[T_0, \ldots, T_n])
$$
와 같음을 보였다. 특히 사영공간은 사영다발의 한 특수한 경우이다.
$S = \Spec(R)$가 아핀이라면
$$
\mathbf{P}^n_S = \mathbf{P}^n_R = \text{Proj}(R[T_0, \ldots, T_n]).
$$
이 모든 동일시는 서로 양립하며, 꼬인 구조층
$\mathcal{O}_{\mathbf{P}^n_S}(d)$의 구성과도 양립한다.

\medskip\noindent
결과를 진술하기 전에 몇 가지 기호가 필요하다.
$R$을 환이라 하자.
$R[T_0, \ldots, T_n]$은 등급 $R$-대수이고, 각 $T_i$는
차수 $1$의 동차원소임을 상기하자.
그 성분 $(R[T_0, \ldots, T_n])_d$를 차수 $d$ 성분이라 나타내자.
이는 유한 자유 $R$-가군이며, 그 계수는 $\binom{n + d}{d}$이다
($d \geq 0$일 때). 그 밖의 경우에는 영이다.
이는 단항식 $T_0^{e_0} \ldots T_n^{e_n}$들 중
$\sum e_i = d$를 만족하는 것들로 이루어진 기저를 갖는다.
또 다음 기호를 사용한다. $R[\frac{1}{T_0}, \ldots, \frac{1}{T_n}]$은
$\mathbf{Z}$-등급 환이며, 각 $\frac{1}{T_i}$의 차수를 $-1$로 둔다.
특히 아래 진술에 등장하는 $\mathbf{Z}$-등급
$R[\frac{1}{T_0}, \ldots, \frac{1}{T_n}]$-가군
$$
\frac{1}{T_0 \ldots T_n} R[\frac{1}{T_0}, \ldots, \frac{1}{T_n}]
$$
은 차수 $\geq -n$에서 영이고, 생성원 $\frac{1}{T_0 \ldots T_n}$
위에서 차수 $-n - 1$에 자유이며, 계수가
$(-1)^n\binom{n + d}{d}$인 자유 가군이다($d \leq -n - 1$일 때).

\begin{lemma}
\label{coherent-lemma-cohomology-projective-space-over-ring}
\begin{reference}
\cite[III Proposition 2.1.12]{EGA}
\end{reference}
$R$을 환이라 하고 $n \geq 0$을 정수라 하자. 그러면 $R$-가군으로서
$$
H^q(\mathbf{P}^n, \mathcal{O}_{\mathbf{P}^n_R}(d)) =
\left\{
\begin{matrix}
(R[T_0, \ldots, T_n])_d & \text{if} & q = 0,\ d \geq 0 \\
\left(\frac{1}{T_0 \ldots T_n} R[\frac{1}{T_0}, \ldots, \frac{1}{T_n}]\right)_d
& \text{if} & q = n,\ d < 0 \\
0 & \text{if} & q, n, d\text{ not as above}
\end{matrix}
\right.
$$
가 성립한다.
\end{lemma}

\begin{proof}
표준 아핀 열린 덮개
$$
\mathcal{U} : \mathbf{P}^n_R = \bigcup\nolimits_{i = 0}^n D_{+}(T_i)
$$
와 체흐 복합체를 사용하여 코호몰로지를 계산하겠다.
Lemma \ref{coherent-lemma-cech-cohomology-quasi-coherent}에 따라 이렇게 해도 된다.
실제로 유한 개의 아핀 $D_{+}(T_i)$들의 임의의 교집합도 표준 아핀
열린집합이다(Constructions, Section \ref{constructions-section-proj} 참조).
실제로 Cohomology, Lemmas \ref{cohomology-lemma-ordered-alternating}와
\ref{cohomology-lemma-alternating-usual}에 따라 교대 체흐 복합체 또는
순서화된 체흐 복합체를 사용할 수 있다.

\medskip\noindent
$\{0, \ldots, n\}$에는 통상적인 순서를 사용한다. 따라서 우리가 살펴볼
복합체의 항들은
$$
\check{\mathcal{C}}_{ord}^p(\mathcal{U}, \mathcal{O}_{\mathbf{P}_R}(d))
=
\bigoplus\nolimits_{i_0 < \ldots < i_p}
(R[T_0, \ldots, T_n, \frac{1}{T_{i_0} \ldots T_{i_p}}])_d
$$
이다. 또한 사상들은 통상적인 공식
$$
d(s)_{i_0 \ldots i_{p + 1}} =
\sum\nolimits_{j = 0}^{p + 1} (-1)^j s_{i_0 \ldots \hat i_j \ldots i_{p + 1}}
$$
으로 주어진다. Cohomology, Section
\ref{cohomology-section-alternating-cech}을 보라.
이 복합체의 각 항에는 자연스러운 $\mathbf{Z}^{n + 1}$-등급이 있다.
즉 단항식 $T_0^{e_0} \ldots T_n^{e_n}$의 가중치를
$(e_0, \ldots, e_n) \in \mathbf{Z}^{n + 1}$로 정하여 이를 얻는다.
위의 미분이 이 등급을 보존함은 명백하다. 공식으로 쓰면
$$
\check{\mathcal{C}}_{ord}^\bullet(\mathcal{U}, \mathcal{O}_{\mathbf{P}_R}(d))
=
\bigoplus\nolimits_{\vec{e} \in \mathbf{Z}^{n + 1}}
\check{\mathcal{C}}^\bullet(\vec{e})
$$
이다. 여기서 우변의 모든 직합항이 나타나는 것은 아니다(아래 참조).
그러므로 이 복합체의 코호몰로지 가군을 계산하려면 등급 성분들의
코호몰로지를 계산한 뒤 마지막에 직합을 취하면 충분하다.

\medskip\noindent
$\vec{e} = (e_0, \ldots, e_n) \in \mathbf{Z}^{n + 1}$를 고정하자.
이 가중치가 위 복합체에 나타나려면 $e_0 + \ldots + e_n = d$라고
가정해야 한다(그렇지 않으면 물론 구조층의 다른 꼬임에 나타난다).
이를 가정하고
$$
NEG(\vec{e}) = \{i \in \{0, \ldots, n\} \mid e_i < 0\}.
$$
로 두자. 이 기호 아래에서 위 체흐 복합체의 가중치 $\vec{e}$ 직합항
$\check{\mathcal{C}}^\bullet(\vec{e})$의 항들은
$$
\check{\mathcal{C}}^p(\vec{e})
=
\bigoplus\nolimits_{i_0 < \ldots < i_p,
\ NEG(\vec{e}) \subset \{i_0, \ldots, i_p\}}
R \cdot T_0^{e_0} \ldots T_n^{e_n}
$$
이다. 다시 말해 $i_0 < \ldots < i_p$ 중에서 $NEG(\vec{e})$가
$\{i_0 \ldots i_p\}$에 포함되지 않는 것에 대응하는 항들은 영이다.
복합체 $\check{\mathcal{C}}^\bullet(\vec{e})$의 미분은 여전히 위와
완전히 같은 공식으로 주어진다.

\medskip\noindent
$NEG(\vec{e}) = \{0, \ldots, n\}$, 즉 모든 지수 $e_i$가 음수라고 하자.
이 경우 복합체 $\check{\mathcal{C}}^\bullet(\vec{e})$에는 한 항만 있으며,
그 항은 $\check{\mathcal{C}}^n(\vec{e}) =
R \cdot \frac{1}{T_0^{-e_0} \ldots T_n^{-e_n}}$이다. 따라서 이 경우
$$
H^q(\check{\mathcal{C}}^\bullet(\vec{e})) =
\left\{
\begin{matrix}
R \cdot \frac{1}{T_0^{-e_0} \ldots T_n^{-e_n}} & \text{if} & q = n \\
0 & \text{if} & \text{else}
\end{matrix}
\right.
$$
이다. 이 항들을 모두 직합하면 보조정리에서 말한 차수 $n$의 값
$$
\left(\frac{1}{T_0 \ldots T_n} R[\frac{1}{T_0}, \ldots, \frac{1}{T_n}]\right)_d
$$
를 얻음은 명백하다. 또한 이 항들은 다른 차수의 코호몰로지에는 기여하지
않는다(이 역시 보조정리의 진술과 일치한다).

\medskip\noindent
$NEG(\vec{e}) = \emptyset$이라 하자. 이 경우 복합체
$\check{\mathcal{C}}^\bullet(\vec{e})$에는 직합항 $R$이
각 $i_0 < \ldots < i_p$에 대응하여 있다. 복합체
$\check{\mathcal{C}}^\bullet(\vec{e})$를 다른 복합체와 비교하자.
구체적으로, $V_i = \Spec(R)$로 모든 $i$에 대해 둔 아핀 열린 덮개
$$
\mathcal{V} : \Spec(R) = \bigcup\nolimits_{i \in \{0, \ldots, n\}} V_i
$$
를 생각하자. 교대 체흐 복합체
$$
\check{\mathcal{C}}_{ord}^\bullet(\mathcal{V}, \mathcal{O}_{\Spec(R)})
$$
를 생각한다. 위와 같은 논리에 따라 이는 $\Spec(R)$ 위 구조층의
코호몰로지를 계산한다. 따라서
$H^p(
\check{\mathcal{C}}_{ord}^\bullet(\mathcal{V}, \mathcal{O}_{\Spec(R)})
) = R$은 $p = 0$일 때 성립하고, $0$은 $p > 0$일 때의 값이다.
이 사실들에 대해서는 Lemma
\ref{coherent-lemma-cech-cohomology-quasi-coherent-trivial}과 그 증명을 보라.
또한 $\check{\mathcal{C}}_{ord}^\bullet(\mathcal{V}, \mathcal{O}_{\Spec(R)})$도
직합항 $R$을 각 $i_0 < \ldots < i_p$마다 가지며,
$\check{\mathcal{C}}^\bullet(\vec{e})$와 정확히 같은 미분을 갖는다.
다시 말해 두 복합체는 서로 동형인 복합체이므로 같은 코호몰로지를 갖는다.
따라서 $NEG(\vec{e}) = \emptyset$인 경우
$$
H^q(\check{\mathcal{C}}^\bullet(\vec{e})) =
\left\{
\begin{matrix}
R \cdot T_0^{e_0} \ldots T_n^{e_n} & \text{if} & q = 0 \\
0 & \text{if} & \text{else}
\end{matrix}
\right.
$$
이다. 이 항들을 모두 직합하면 보조정리에서 말한 차수 $0$의 값
$$
(R[T_0, \ldots, T_n])_d
$$
를 얻음은 명백하다. 또한 이 항들은 다른 차수의 코호몰로지에는 기여하지
않는다(이 역시 보조정리의 진술과 일치한다).

\medskip\noindent
보조정리의 증명을 끝내려면 복합체
$\check{\mathcal{C}}^\bullet(\vec{e})$가 비순환임을 보여야 한다.
여기서 $NEG(\vec{e})$는 공집합도 아니고 $\{0, \ldots, n\}$도 아니다.
지표 $i_{\text{fix}} \not \in NEG(\vec{e})$를 하나 택하자
(그러한 지표가 존재한다). $i_0 < \ldots < i_p$에 대해 다음 규칙으로
주어지는 사상
$$
h :
\check{\mathcal{C}}^{p + 1}(\vec{e})
\to
\check{\mathcal{C}}^p(\vec{e})
$$
을 생각하자.
$$
h(s)_{i_0 \ldots i_p} =
\left\{
\begin{matrix}
0 & \text{if} & p \not \in \{0, \ldots, n - 1\} \\
0 & \text{if} & i_{\text{fix}} \in \{i_0, \ldots, i_p\} \\
s_{i_{\text{fix}} i_0 \ldots i_p} & \text{if} & i_{\text{fix}} < i_0 \\
(-1)^a s_{i_0 \ldots i_{a - 1} i_{\text{fix}} i_a \ldots i_p} &
\text{if} & i_{a - 1} < i_{\text{fix}} < i_a \\
(-1)^{p + 1} s_{i_0 \ldots i_p i_{\text{fix}}} &
\text{if} & i_p < i_{\text{fix}}
\end{matrix}
\right.
$$
Lemma \ref{coherent-lemma-cech-cohomology-quasi-coherent-trivial}의 증명과 비교하라.
다음 동치가 성립하므로 이 정의는 타당하다.
$$
NEG(\vec{e}) \subset \{i_0, \ldots, i_p\}
\Leftrightarrow
NEG(\vec{e}) \subset \{i_{\text{fix}}, i_0, \ldots, i_p\}
$$
Lemma \ref{coherent-lemma-cech-cohomology-quasi-coherent-trivial}의 증명에서와
완전히 같은 조합론적 계산\footnote{
예를 들어 $i_0 < \ldots < i_p$가
$i_{\text{fix}} \not \in \{i_0, \ldots, i_p\}$를 만족하고,
$i_{a - 1} < i_{\text{fix}} < i_a$가 어떤 $1 \leq a \leq p$에 대해
성립한다고 하자. 그러면
\begin{align*}
&
(dh + hd)(s)_{i_0 \ldots i_p} \\
& =
\sum\nolimits_{j = 0}^p
(-1)^j
h(s)_{i_0 \ldots \hat i_j \ldots i_p}
+
(-1)^a d(s)_{i_0 \ldots i_{a - 1} i_{\text{fix}} i_a \ldots i_p}\\
& =
\sum\nolimits_{j = 0}^{a - 1}
(-1)^{j + a - 1}
s_{i_0 \ldots \hat i_j \ldots i_{a - 1} i_{\text{fix}} i_a \ldots i_p} +
\sum\nolimits_{j = a}^p
(-1)^{j + a}
s_{i_0 \ldots i_{a - 1} i_{\text{fix}} i_a \ldots \hat i_j \ldots i_p}
+ \\
&
\sum\nolimits_{j = 0}^{a - 1}
(-1)^{a + j}
s_{i_0 \ldots \hat i_j \ldots i_{a - 1} i_{\text{fix}} i_a \ldots i_p} +
(-1)^{2a} s_{i_0 \ldots i_p} +
\sum\nolimits_{j = a}^p
(-1)^{a + j + 1}
s_{i_0 \ldots i_{a - 1} i_{\text{fix}} i_a \ldots \hat i_j \ldots i_p}
\\
& =
s_{i_0 \ldots i_p}
\end{align*}
이며, 이것이 원하는 바이다. 다른 경우들도 마찬가지이다.}
으로부터
$$
(hd + dh)(s)_{i_0 \ldots i_p}
=
s_{i_0 \ldots i_p}
$$
를 얻는다. 따라서 복합체 $\check{\mathcal{C}}^\bullet(\vec{e})$의
항등사상은 영 사상과 호모토픽이고, 이에 따라 이 복합체는 비순환이다.
\end{proof}

\noindent
다음 보조정리에서는 자유 $R$-가군들의 쌍대쌍
$$
R[T_0, \ldots, T_n]
\times
\frac{1}{T_0 \ldots T_n} R[\frac{1}{T_0}, \ldots, \frac{1}{T_n}]
\longrightarrow
R
$$
을 사용할 것이다. 이는 다음 규칙으로 정의된다.
$$
(f, g)
\longmapsto
\text{coefficient of }
\frac{1}{T_0 \ldots T_n}
\text{ in }fg.
$$
다시 말해 기저원소 $T_0^{e_0} \ldots T_n^{e_n}$와 기저원소
$T_0^{d_0} \ldots T_n^{d_n}$의 쌍대값은 $1$인데, 이는
$e_i + d_i = -1$이 모든 $i$에 대해 성립할 때 그리고 그때에만
그러하며, 그 밖의 모든 경우에는 영이다. 이 쌍대쌍을 사용하면
다음 동일시를 얻는다.
$$
\left(\frac{1}{T_0 \ldots T_n} R[\frac{1}{T_0}, \ldots, \frac{1}{T_n}]\right)_d
=
\Hom_R((R[T_0, \ldots, T_n])_{-n - 1 - d}, R)
$$
따라서 Lemma \ref{coherent-lemma-cohomology-projective-space-over-ring}의 결과를
다음과 같이 다시 쓸 수 있다.
\begin{equation}
\label{coherent-equation-identify}
H^q(\mathbf{P}^n, \mathcal{O}_{\mathbf{P}^n_R}(d)) =
\left\{
\begin{matrix}
(R[T_0, \ldots, T_n])_d & \text{if} & q = 0 \\
0 & \text{if} & q \not = 0, n \\
\Hom_R((R[T_0, \ldots, T_n])_{-n - 1 - d}, R)
& \text{if} & q = n
\end{matrix}
\right.
\end{equation}

\begin{lemma}
\label{coherent-lemma-identify-functorially}
Equation (\ref{coherent-equation-identify})의 동일시들은 환 준동형
$R \to R'$에 관한 밑변환과 양립한다. 또한 임의의 동차원소
$f \in R[T_0, \ldots, T_n]$가 차수 $m$이면 $f$를 곱하는 사상
$$
\mathcal{O}_{\mathbf{P}^n_R}(d)
\longrightarrow
\mathcal{O}_{\mathbf{P}^n_R}(d + m)
$$
은 Equation (\ref{coherent-equation-identify})의 동일시를 통해 코호몰로지군 위에
다음 사상을 유도한다. 이것은 $f$의 곱셈이며 $H^0$에 작용한다.
또한 $f$를 곱하는 사상의 반변 사상
$$
(R[T_0, \ldots, T_n])_{-n - 1 - (d + m)}
\longrightarrow
(R[T_0, \ldots, T_n])_{-n - 1 - d}
$$
으로 $H^n$에 작용한다.
\end{lemma}

\begin{proof}
$R \to R'$이 환 준동형이라고 하자.
$\mathcal{U}$를 $\mathbf{P}^n_R$의 표준 아핀 열린 덮개라 하고,
$\mathcal{U}'$를 $\mathbf{P}^n_{R'}$의 표준 아핀 열린 덮개라 하자.
$\mathcal{U}'$는 덮개 $\mathcal{U}$의 당김이며, 이는 표준 사상
$\mathbf{P}^n_{R'} \to \mathbf{P}^n_R$ 아래에서의 당김이다.
따라서 체흐 복합체의 사상
$$
\gamma :
\check{\mathcal{C}}_{ord}^\bullet(\mathcal{U},
\mathcal{O}_{\mathbf{P}_R}(d))
\longrightarrow
\check{\mathcal{C}}_{ord}^\bullet(\mathcal{U}',
\mathcal{O}_{\mathbf{P}_{R'}}(d))
$$
이 있으며, Cohomology, Lemma
\ref{cohomology-lemma-functoriality-cech}에 따라 이는 코호몰로지 위의
사상과 양립한다. Lemma
\ref{coherent-lemma-cohomology-projective-space-over-ring}의 증명에서 수행한
계산으로부터, 이 체흐 복합체의 사상이 해당 코호몰로지군의 동일시와
양립함은 명백하다. (즉 $R$ 위 체흐 복합체의 기저원소들은 단순히
$R'$ 위 체흐 복합체의 대응하는 기저원소들로 간다.) 이로써 보조정리의
첫 번째 진술을 얻는다.

\medskip\noindent
이제 환 $R$을 고정하고 두 동차다항식
$f, g \in R[T_0, \ldots, T_n]$을 생각하되, 둘의 차수는 같은 $m$이라 하자.
코호몰로지는 가법 함자이므로
$f + g$를 곱하여 유도되는 사상은 $f$를 곱하여 유도되는 사상과
$g$를 곱하여 유도되는 사상의 합과 같다. 또한 코호몰로지가 함자이므로,
곱 $fg$를 곱하는 것에 대해서도 비슷한 결과가 성립한다. 여기서
$f, g$는 모두 동차원소이지만 두 차수가 같을 필요는 없다. 그러므로 보조정리의
두 번째 진술을 확인하려면 $f = x \in R$인 경우와 $f = T_i$인 경우만
증명하면 충분하다. 원소 $x \in R$를 곱하는 경우에는 눈앞의 각
코호몰로지군이나 복합체가 $R$-가군 또는 $R$-가군들의 복합체 구조를
가지므로 결과가 성립한다. 마지막으로 $T_i$를 곱하는 것을
$\mathcal{O}_{\mathbf{P}^n_R}$-선형 사상
$$
\mathcal{O}_{\mathbf{P}^n_R}(d)
\longrightarrow
\mathcal{O}_{\mathbf{P}^n_R}(d + 1)
$$
으로 생각하자. $H^0$에 관한 진술은 명백하다. $H^n$에 관한 진술을
위해서는 $T_i$를 곱하는 것을 체흐 복합체 위의 사상
$$
\check{\mathcal{C}}_{ord}^\bullet(\mathcal{U},
\mathcal{O}_{\mathbf{P}_R}(d))
\longrightarrow
\check{\mathcal{C}}_{ord}^\bullet(\mathcal{U},
\mathcal{O}_{\mathbf{P}_R}(d + 1))
$$
으로 생각한다. Lemma
\ref{coherent-lemma-cohomology-projective-space-over-ring}의 증명에서 도입한
기호를 사용하겠다. 다음 분해들에 관하여 $T_i$를 곱하는 효과를 살펴본다.
$$
\check{\mathcal{C}}_{ord}^\bullet(\mathcal{U}, \mathcal{O}_{\mathbf{P}_R}(d))
=
\bigoplus\nolimits_{\vec{e} \in \mathbf{Z}^{n + 1}, \ \sum e_i = d}
\check{\mathcal{C}}^\bullet(\vec{e})
$$
그리고
$$
\check{\mathcal{C}}_{ord}^\bullet(\mathcal{U},
\mathcal{O}_{\mathbf{P}_R}(d + 1))
=
\bigoplus\nolimits_{\vec{e} \in \mathbf{Z}^{n + 1}, \ \sum e_i = d + 1}
\check{\mathcal{C}}^\bullet(\vec{e})
$$
직합복합체 $\check{\mathcal{C}}^\bullet(\vec{e})$가 직합복합체
$\check{\mathcal{C}}^\bullet(\vec{e} + \vec{b}_i)$로 감은 명백하다.
여기서 $\vec{b}_i = (0, \ldots, 0, 1, 0, \ldots, 0))$는 $i$번째
기저벡터이다. 다시 말해 $H^n$에서 $\vec{e}$에 대응하고
$e_i < 0$이며 $\sum e_i = d$인 직합항을
$H^n$에서 $\vec{e} + \vec{b}_i$에 대응하는 직합항으로 보낸다
(이는 $e_i + b_i \geq 0$이면 영이다).
이것은 정확히 $T_i$를 곱하는 사상
$$
(R[T_0, \ldots, T_n])_{-n - 1 - (d + 1)}
\longrightarrow
(R[T_0, \ldots, T_n])_{-n - 1 - d}
$$
의 반변 사상의 작용에 대응함을 쉽게 알 수 있다. 이로써 보조정리가
증명되었다.
\end{proof}

\noindent
상대적 형태를 진술하기 전에 몇 가지 기호가 필요하다. 즉
$\mathcal{O}_S[T_0, \ldots, T_n]$은 등급 $\mathcal{O}_S$-가군이며,
각 $T_i$는 차수 $1$의 동차원소임을 상기하자. 그 성분
$(\mathcal{O}_S[T_0, \ldots, T_n])_d$를 차수 $d$ 성분이라 나타내자.
이는 계수가 $\binom{n + d}{d}$인 유한 국소 자유층이며 $S$ 위에 놓인다.

\begin{lemma}
\label{coherent-lemma-cohomology-projective-space-over-base}
$S$를 스킴이라 하고 $n \geq 0$을 정수라 하자. 구조 사상
$$
f : \mathbf{P}^n_S \longrightarrow S.
$$
을 생각하자. 그러면
$$
R^qf_*(\mathcal{O}_{\mathbf{P}^n_S}(d)) =
\left\{
\begin{matrix}
(\mathcal{O}_S[T_0, \ldots, T_n])_d & \text{if} & q = 0 \\
0 & \text{if} & q \not = 0, n \\
\SheafHom_{\mathcal{O}_S}(
(\mathcal{O}_S[T_0, \ldots, T_n])_{- n - 1 - d}, \mathcal{O}_S)
& \text{if} & q = n
\end{matrix}
\right.
$$
가 성립한다.
\end{lemma}

\begin{proof}
생략한다. 힌트: (\ref{coherent-equation-identify})의 동일시들이 Lemma
\ref{coherent-lemma-identify-functorially}에 따라 아핀 밑변환과 양립하므로
결과가 따라온다.
\end{proof}

\noindent
다음으로 유한 국소 자유층에 대응하는 사영다발의 형태를 진술한다.
$S$를 스킴이라 하자. $\mathcal{E}$를 유한 국소 자유
$\mathcal{O}_S$-가군이라 하고 그 계수를 $n + 1$이라 하자. Modules, Section
\ref{modules-section-locally-free}을 보라. 이 경우
$\text{Sym}(\mathcal{E})$를 등급 $\mathcal{O}_S$-가군으로 보며,
$\mathcal{E}$는 차수 $1$인 등급 부분이다. 그리고
$\text{Sym}^d(\mathcal{E})$는 차수 $d$ 성분이다. 이는 계수가
$\binom{n + d}{d}$인 유한 국소 자유층이며 $S$ 위에 놓인다. 우리의 정규화가
$$
\pi :
\mathbf{P}(\mathcal{E})
=
\underline{\text{Proj}}_S(\text{Sym}(\mathcal{E}))
\longrightarrow
S
$$
라는 것과, 자연스러운 사상
$\text{Sym}^d(\mathcal{E}) \to \pi_*\mathcal{O}_{\mathbf{P}(\mathcal{E})}(d)$
들이 있음을 상기하자.

\begin{lemma}
\label{coherent-lemma-cohomology-projective-bundle}
$S$를 스킴이라 하고 $n \geq 1$이라 하자. $\mathcal{E}$를 유한 국소 자유
$\mathcal{O}_S$-가군이라 하고 그 계수를 $n + 1$이라 하자. 구조 사상
$$
\pi : \mathbf{P}(\mathcal{E}) \longrightarrow S.
$$
을 생각하자. 그러면
$$
R^q\pi_*(\mathcal{O}_{\mathbf{P}(\mathcal{E})}(d)) =
\left\{
\begin{matrix}
\text{Sym}^d(\mathcal{E}) & \text{if} & q = 0 \\
0 & \text{if} & q \not = 0, n \\
\SheafHom_{\mathcal{O}_S}(
\text{Sym}^{- n - 1 - d}(\mathcal{E})
\otimes_{\mathcal{O}_S}
\wedge^{n + 1}\mathcal{E},
\mathcal{O}_S)
& \text{if} & q = n
\end{matrix}
\right.
$$
가 성립한다. 이 동일시들은 밑변환 및 국소 자유층 사이의 동형들과
양립한다.
\end{lemma}

\begin{proof}
표준 사상
$$
\pi^*\mathcal{E} \longrightarrow \mathcal{O}_{\mathbf{P}(\mathcal{E})}(1)
$$
을 생각하고, 이를 $1$만큼 아래로 꼬아서
$$
\pi^*(\mathcal{E})(-1) \longrightarrow \mathcal{O}_{\mathbf{P}(\mathcal{E})}
$$
을 얻는다. 이는 계수 $n + 1$인 국소 자유층에서 구조층으로 가는
전사 사상이다. 따라서 대응하는 코쥘 복합체는 완전하다(More on
Algebra, Lemma \ref{more-algebra-lemma-homotopy-koszul-abstract}).
다시 말해 완전 복합체
$$
0 \to
\pi^*(\wedge^{n + 1}\mathcal{E})(-n - 1) \to
\ldots \to
\pi^*(\wedge^i\mathcal{E})(-i) \to
\ldots \to
\pi^*\mathcal{E}(-1) \to
\mathcal{O}_{\mathbf{P}(\mathcal{E})} \to 0
$$
가 있다. 항 $\pi^*(\wedge^i\mathcal{E})(-i)$를 차수 $-i$에 놓인 것으로
생각하겠다. Derived Categories, Lemma
\ref{derived-lemma-two-ss-complex-functor}의 첫 번째 스펙트럼열을 사용하여
이 비순환 복합체의 고차 직접상들을 계산하겠다. 즉 영으로 수렴하는,
다음 항들을 갖는 스펙트럼열이 있다.
$$
E_1^{p, q} = R^q\pi_*\left(\pi^*(\wedge^{-p}\mathcal{E})(p)\right)
$$
사영 공식(Cohomology, Lemma
\ref{cohomology-lemma-projection-formula})에 따라
$$
E_1^{p, q} = \wedge^{-p} \mathcal{E} \otimes_{\mathcal{O}_S}
R^q\pi_*\left(\mathcal{O}_{\mathbf{P}(\mathcal{E})}(p)\right).
$$
$S$ 위에서 국소적으로 층 $\mathcal{E}$는 자명하다. 즉
$\mathcal{O}_S^{\oplus n + 1}$과 동형이므로, $S$ 위에서 국소적으로
사상 $\mathbf{P}(\mathcal{E}) \to S$를 $\mathbf{P}^n_S \to S$와
동일시할 수 있다. 따라서 $S$ 위에서 국소적으로 Lemmas
\ref{coherent-lemma-cohomology-projective-space-over-ring},
\ref{coherent-lemma-identify-functorially}, 또는
\ref{coherent-lemma-cohomology-projective-space-over-base}의 결과를 사용할 수 있다.
따라서 $E_1^{p, q} = 0$은 $(p, q)$가 $(0, 0)$ 또는
$(-n - 1, n)$이 아닌 한 성립한다. 영이 아닌 항들은
\begin{align*}
E_1^{0, 0} & = \pi_*\mathcal{O}_{\mathbf{P}(\mathcal{E})} = \mathcal{O}_S \\
E_1^{-n - 1, n} & =
R^n\pi_*\left(\pi^*(\wedge^{n + 1}\mathcal{E})(-n - 1)\right) =
\wedge^{n + 1}\mathcal{E} \otimes_{\mathcal{O}_S}
R^n\pi_*\left(\mathcal{O}_{\mathbf{P}(\mathcal{E})}(-n - 1)\right)
\end{align*}
이다. 그러므로 이 스펙트럼열에는 영이 아닌 미분이 하나만 있을 수 있으며,
그것은 사상
$d_{n + 1}^{-n - 1, n} : E_{n + 1}^{-n - 1, n} \to E_{n + 1}^{0, 0}$
이다. 스펙트럼열이 $0$ 층으로 수렴하므로 이는 동형이어야 한다.
따라서 $E_1^{p, q} = E_{n + 1}^{p, q}$이고 표준 동형
$$
\wedge^{n + 1}\mathcal{E} \otimes_{\mathcal{O}_S}
R^n\pi_*\left(\mathcal{O}_{\mathbf{P}(\mathcal{E})}(-n - 1)\right) =
R^n\pi_*\left(\pi^*(\wedge^{n + 1}\mathcal{E})(-n - 1)\right)
\xrightarrow{d_{n + 1}^{-n - 1, n}}
\mathcal{O}_S
$$
을 얻는다. $\wedge^{n + 1}\mathcal{E}$가 가역층이므로, 이는
$R^n\pi_*\mathcal{O}_{\mathbf{P}(\mathcal{E})}(-n - 1)$도 가역이며
$\wedge^{n + 1}\mathcal{E}$의 역과 표준적으로 동형임을 뜻한다.
다시 말해 보조정리에서 $d = - n - 1$인 경우를 증명하였다.

\medskip\noindent
$S$ 위에서 국소적으로 작업하면 Lemmas
\ref{coherent-lemma-cohomology-projective-space-over-ring},
\ref{coherent-lemma-identify-functorially}, 또는
\ref{coherent-lemma-cohomology-projective-space-over-base}의 코호몰로지 계산으로부터
보조정리의 소멸 진술들이 즉시 따라온다. 또한 $R^0\pi_*$에 관한 결과도
명백하다. 실제로 표준 사상
$\text{Sym}^d(\mathcal{E}) \to \pi_* \mathcal{O}_{\mathbf{P}(\mathcal{E})}(d)$
들이 모든 $d$에 대해 있다. 이제
$R^n\pi_*\mathcal{O}_{\mathbf{P}(\mathcal{E})}(d)$에 대한 서술이
$d < -n - 1$일 때 옳음을 보이면 된다. 이를 위해 사상
$$
\pi^*(\text{Sym}^{-d - n - 1}(\mathcal{E}))
\otimes_{\mathcal{O}_{\mathbf{P}(\mathcal{E})}}
\mathcal{O}_{\mathbf{P}(\mathcal{E})}(d)
\longrightarrow
\mathcal{O}_{\mathbf{P}(\mathcal{E})}(-n - 1)
$$
을 생각하자. $R^n\pi_*$를 적용하고 사영 공식(위 참조)을 사용하면 사상
$$
\text{Sym}^{-d - n - 1}(\mathcal{E})
\otimes_{\mathcal{O}_S}
R^n\pi_*(\mathcal{O}_{\mathbf{P}(\mathcal{E})}(d))
\longrightarrow
R^n\pi_*\mathcal{O}_{\mathbf{P}(\mathcal{E})}(-n - 1) =
(\wedge^{n + 1}\mathcal{E})^{\otimes -1}
$$
을 얻는다(마지막 등식은 위에서 보였다). 다시 Lemmas
\ref{coherent-lemma-cohomology-projective-space-over-ring},
\ref{coherent-lemma-identify-functorially}, 또는
\ref{coherent-lemma-cohomology-projective-space-over-base}의 국소 계산에 따라,
이 사상은 $R^n\pi_*(\mathcal{O}_{\mathbf{P}(\mathcal{E})}(d))$와
$\text{Sym}^{-d - n - 1}(\mathcal{E}) \otimes \wedge^{n + 1}(\mathcal{E})$
사이에 원하는 완전 쌍대쌍을 유도한다.
\end{proof}

\section{국소 뇌터 스킴 위의 연접층}
\label{coherent-section-coherent-sheaves}

\noindent
임의의 환 달린 공간 위의 연접 가군 개념을 Modules, Section
\ref{modules-section-coherent}에서 정의하였다. 비뇌터 스킴 위의 연접층도
생각할 수 있지만, 여기서는 연접층을 다룰 때마다 바탕 스킴이 국소 뇌터라고
가정한다. 국소 뇌터 스킴 위 연접층의 다음 특성화를 얻는다.

\begin{lemma}
\label{coherent-lemma-coherent-Noetherian}
$X$를 국소 뇌터 스킴이라 하고 $\mathcal{F}$를 $\mathcal{O}_X$-가군이라
하자. 다음 조건들은 서로 동치이다.
\begin{enumerate}
\item $\mathcal{F}$는 연접이다.
\item $\mathcal{F}$는 준연접 유한형 $\mathcal{O}_X$-가군이다.
\item $\mathcal{F}$는 유한 표시 $\mathcal{O}_X$-가군이다.
\item 임의의 아핀 열린집합 $\Spec(A) = U \subset X$에 대해
$\mathcal{F}|_U = \widetilde M$이고, 여기서 $M$은 유한 $A$-가군이다.
\item 아핀 열린 덮개 $X = \bigcup U_i$, $U_i = \Spec(A_i)$가 존재하여
각각에 대해 $\mathcal{F}|_{U_i} = \widetilde M_i$이고, 여기서 $M_i$는
유한 $A_i$-가군이다.
\end{enumerate}
특히 $\mathcal{O}_X$는 연접이고, 모든 가역 $\mathcal{O}_X$-가군은
연접이며, 더 일반적으로 모든 유한 국소 자유 $\mathcal{O}_X$-가군은
연접이다.
\end{lemma}

\begin{proof}
(1) $\Rightarrow$ (2) 및 (1) $\Rightarrow$ (3)은 일반적으로 성립한다.
Modules, Lemma \ref{modules-lemma-coherent-finite-presentation}을 보라.
$\mathcal{F}$가 유한 표시이면 Modules, Lemma
\ref{modules-lemma-finite-presentation-quasi-coherent}에 따라
$\mathcal{F}$는 준연접이다. 따라서 (3) $\Rightarrow$ (2)도 성립한다.

\medskip\noindent
$\mathcal{F}$가 준연접 유한형 $\mathcal{O}_X$-가군이라고 하자.
Properties, Lemma \ref{properties-lemma-finite-type-module}에 따라 임의의
아핀 열린집합 $\Spec(A) = U \subset X$ 위에서
$\mathcal{F}|_U = \widetilde M$이고, 여기서 $M$은 유한 $A$-가군이다.
$A$가 뇌터이므로 $M$에는 유한 분해
$$
A^{\oplus m} \to A^{\oplus n} \to M \to 0.
$$
가 있다. 따라서 Properties, Lemma
\ref{properties-lemma-finite-presentation-module}에 의해 $\mathcal{F}$는
유한 표시이다. 다시 말해 (2) $\Rightarrow$ (3)이다.

\medskip\noindent
Modules, Lemma \ref{modules-lemma-coherent-structure-sheaf}에 따라 (3)이
(1)을 함의함을 보이려면 $\mathcal{O}_X$가 연접임을 보이면 충분하다.
따라서 다음을 보여야 한다. 임의의 열린집합 $U \subset X$와 유한한
절단들의 모음 $f_i \in \mathcal{O}_X(U)$, $i = 1, \ldots, n$이 주어졌을 때
사상 $\bigoplus_{i = 1, \ldots, n} \mathcal{O}_U \to \mathcal{O}_U$의
핵은 유한형이다. 유한형이라는 것은 국소 성질이므로 임의의 $x \in U$의
근방에서 이를 확인하면 충분하다. 따라서 $U = \Spec(A)$가 아핀이라고
가정해도 된다. 이 경우 $f_1, \ldots, f_n \in A$는 $A$의 원소들이다.
$A$가 뇌터이므로(Properties, Lemma
\ref{properties-lemma-locally-Noetherian} 참조), 그 핵 $K$, 즉 사상
$\bigoplus_{i = 1, \ldots, n} A \to A$의 핵은 유한 $A$-가군이다.
예를 들어 Algebra, Lemma \ref{algebra-lemma-Noetherian-basic}을 보라.
함자\ $\widetilde{ }$\ 는 완전하므로(Schemes, Lemma
\ref{schemes-lemma-spec-sheaves} 참조) 완전열
$$
\widetilde K \to
\bigoplus\nolimits_{i = 1, \ldots, n} \mathcal{O}_U \to
\mathcal{O}_U
$$
을 얻는다. Properties, Lemma
\ref{properties-lemma-finite-type-module}를 다시 적용하면
$\widetilde K$가 유한형임을 알 수 있다. 따라서 (1), (2), (3)은 모두
동치이다.

\medskip\noindent
Properties, Lemma \ref{properties-lemma-finite-type-module}에 의해 (2)는
(4)를 함의한다. (4)가 (5)를 함의함은 자명하다. Schemes, Section
\ref{schemes-section-quasi-coherent}의 논의에 의해 (5)는 $\mathcal{F}$가
준연접임을 함의하며, (5)가 $\mathcal{F}$의 유한형성을 함의함도 명백하다.
따라서 (5)는 (2)를 함의하고 증명이 끝난다.
\end{proof}

\begin{lemma}
\label{coherent-lemma-coherent-abelian-Noetherian}
$X$를 국소 뇌터 스킴이라 하자. 연접 $\mathcal{O}_X$-가군들의 범주는
아벨 범주이다. 더 정확히 말하면 연접 $\mathcal{O}_X$-가군들의 사상의
핵과 여핵은 연접이다. 연접층들의 모든 확대는 연접이다.
\end{lemma}

\begin{proof}
이는 Modules, Lemma \ref{modules-lemma-coherent-abelian}을 특수한 경우에
다시 진술한 것이다.
\end{proof}

\noindent
다음 보조정리는 일반 환 달린 공간의 연접 $\mathcal{O}_X$-가군 범주에서
공간 $X$에 대해 항상 성립하지는 않는다.

\begin{lemma}
\label{coherent-lemma-coherent-Noetherian-quasi-coherent-sub-quotient}
$X$를 국소 뇌터 스킴이라 하자. $\mathcal{F}$를 연접
$\mathcal{O}_X$-가군이라 하자. $\mathcal{F}$의 모든 준연접 부분가군은
연접이고, $\mathcal{F}$의 모든 준연접 몫가군은 연접이다.
\end{lemma}

\begin{proof}
$X$가 아핀이라고 가정해도 되며, $X = \Spec(A)$라 쓰자.
Properties, Lemma \ref{properties-lemma-locally-Noetherian}에 따라
$A$는 뇌터이다. Lemma \ref{coherent-lemma-coherent-Noetherian}을 적용하면 이는
대수의 명제가 된다. 그 대수적 대응물은 유한 $A$-가군의 몫가군이나
부분가군이 유한 $A$-가군이라는 사실이다. 예를 들어 Algebra, Lemma
\ref{algebra-lemma-Noetherian-basic}을 보라.
\end{proof}

\begin{lemma}
\label{coherent-lemma-tensor-hom-coherent}
$X$를 국소 뇌터 스킴이라 하자. $\mathcal{F}$, $\mathcal{G}$를 연접
$\mathcal{O}_X$-가군들이라 하자. $\mathcal{O}_X$-가군들
$\mathcal{F} \otimes_{\mathcal{O}_X} \mathcal{G}$와
$\SheafHom_{\mathcal{O}_X}(\mathcal{F}, \mathcal{G})$는 연접이다.
\end{lemma}

\begin{proof}
Modules, Lemma
\ref{modules-lemma-internal-hom-locally-kernel-direct-sum}에서
$\SheafHom_{\mathcal{O}_X}(\mathcal{F}, \mathcal{G})$가 연접임을 보였다.
텐서곱에 관한 결과는 Modules, Lemma
\ref{modules-lemma-tensor-product-permanence}이다.
\end{proof}

\begin{lemma}
\label{coherent-lemma-local-isomorphism}
$X$를 국소 뇌터 스킴이라 하자. $\mathcal{F}$, $\mathcal{G}$를 연접
$\mathcal{O}_X$-가군들이라 하고, $\varphi : \mathcal{G} \to \mathcal{F}$를
$\mathcal{O}_X$-가군들의 준동형이라 하자. $x \in X$라 하자.
\begin{enumerate}
\item $\mathcal{F}_x = 0$이면 열린 근방 $U \subset X$가 $x$ 주위에 존재하여
$\mathcal{F}|_U = 0$이다.
\item $\varphi_x : \mathcal{G}_x \to \mathcal{F}_x$가 단사이면
열린 근방 $U \subset X$가 $x$ 주위에 존재하여 $\varphi|_U$는 단사이다.
\item $\varphi_x : \mathcal{G}_x \to \mathcal{F}_x$가 전사이면
열린 근방 $U \subset X$가 $x$ 주위에 존재하여 $\varphi|_U$는 전사이다.
\item $\varphi_x : \mathcal{G}_x \to \mathcal{F}_x$가 전단사이면
열린 근방 $U \subset X$가 $x$ 주위에 존재하여 $\varphi|_U$는 동형이다.
\end{enumerate}
\end{lemma}

\begin{proof}
Modules, Lemmas
\ref{modules-lemma-finite-type-surjective-on-stalk},
\ref{modules-lemma-finite-type-stalk-zero}, 그리고
\ref{modules-lemma-finite-type-to-coherent-injective-on-stalk}을 보라.
\end{proof}

\begin{lemma}
\label{coherent-lemma-map-stalks-local-map}
$X$를 국소 뇌터 스킴이라 하자. $\mathcal{F}$, $\mathcal{G}$를 연접
$\mathcal{O}_X$-가군들이라 하자. $x \in X$라 하자.
$\psi : \mathcal{G}_x \to \mathcal{F}_x$가 $\mathcal{O}_{X, x}$-가군들의
사상이라고 하자. 그러면 열린 근방 $U \subset X$가 $x$ 주위에 존재하고 사상
$\varphi : \mathcal{G}|_U \to \mathcal{F}|_U$가 존재하여
$\varphi_x = \psi$이다.
\end{lemma}

\begin{proof}
Lemma \ref{coherent-lemma-coherent-Noetherian}을 고려하면 이는 Modules, Lemma
\ref{modules-lemma-stalk-internal-hom}을 다시 쓴 것이다.
\end{proof}

\begin{lemma}
\label{coherent-lemma-coherent-support-closed}
$X$를 국소 뇌터 스킴이라 하고 $\mathcal{F}$를 연접
$\mathcal{O}_X$-가군이라 하자. 그러면 $\text{Supp}(\mathcal{F})$는
닫혀 있고, $\mathcal{F}$는 $\mathcal{F}$의 스킴론적 지지 위의 연접층에서
온다. Morphisms, Definition
\ref{morphisms-definition-scheme-theoretic-support}를 보라.
\end{lemma}

\begin{proof}
$i : Z \to X$를 $\mathcal{F}$의 스킴론적 지지라 하고,
$\mathcal{G}$를 $Z$ 위의 유한형 준연접층으로서
$i_*\mathcal{G} \cong \mathcal{F}$인 것이라 하자.
$Z = \text{Supp}(\mathcal{F})$이므로 지지는 닫혀 있다.
Morphisms, Lemmas \ref{morphisms-lemma-immersion-locally-finite-type}와
\ref{morphisms-lemma-finite-type-noetherian}에 따라 스킴 $Z$는 국소
뇌터이다. 마지막으로 Lemma \ref{coherent-lemma-coherent-Noetherian}에 따라
$\mathcal{G}$는 연접 $\mathcal{O}_Z$-가군이다.
\end{proof}

\begin{lemma}
\label{coherent-lemma-i-star-equivalence}
$i : Z \to X$를 국소 뇌터 스킴들의 닫힌 몰입이라 하자.
$\mathcal{I} \subset \mathcal{O}_X$를 $Z$를 잘라내는 준연접 아이디얼층이라
하자. 함자 $i_*$는 연접 $\mathcal{O}_X$-가군들 중 $\mathcal{I}$에 의해
소멸되는 것들의 범주와 연접 $\mathcal{O}_Z$-가군들의 범주 사이의
동치를 유도한다.
\end{lemma}

\begin{proof}
Morphisms, Lemma \ref{morphisms-lemma-i-star-equivalence}에 따라 이 함자는
충실충만이다. $\mathcal{F}$를 연접 $\mathcal{O}_X$-가군으로서
$\mathcal{I}$에 의해 소멸되는 것이라 하자. Morphisms, Lemma
\ref{morphisms-lemma-i-star-equivalence}에 따라
$\mathcal{F} = i_*\mathcal{G}$로 쓸 수 있고, 여기서 $\mathcal{G}$는
$Z$ 위의 어떤 준연접층이다.
Modules, Lemma \ref{modules-lemma-i-star-reflects-finite-type}에 따라
$\mathcal{G}$는 유한형이다. 따라서 Lemma
\ref{coherent-lemma-coherent-Noetherian}에 따라 $\mathcal{G}$는 연접이다.
그러므로 이 함자는 원하는 대로 본질적 전사이기도 하다.
\end{proof}

\begin{lemma}
\label{coherent-lemma-finite-pushforward-coherent}
$f : X \to Y$를 스킴들의 사상이라 하고 $\mathcal{F}$를 준연접
$\mathcal{O}_X$-가군이라 하자. $f$가 유한이고 $Y$가 국소 뇌터라고
가정하자. 그러면 $R^pf_*\mathcal{F} = 0$이며 이는 $p > 0$에 대해
성립하고, $f_*\mathcal{F}$는 $\mathcal{F}$가 연접이면 연접이다.
\end{lemma}

\begin{proof}
Lemma \ref{coherent-lemma-relative-affine-vanishing} 및 유한 사상이 정의상
아핀이라는 사실에 의해 고차 직접상들은 소멸한다. 또한 가정들로부터
$X$도 국소 뇌터이므로(Morphisms, Lemma
\ref{morphisms-lemma-finite-type-noetherian} 참조) 진술은 의미가 있다.
$\Spec(A) = V \subset Y$를 아핀 열린부분집합이라 하자. Morphisms,
Definition \ref{morphisms-definition-integral}에 따라
$f^{-1}(V) = \Spec(B)$이고 $A \to B$는 유한이다. Lemma
\ref{coherent-lemma-coherent-Noetherian}을 적용하면 보조정리의 진술은 다음 대수적
사실이 된다. $M$이 유한 $B$-가군이면 $M$은 $A$-가군으로 보아도
유한이다. Algebra, Lemma
\ref{algebra-lemma-finite-module-over-finite-extension}을 보라.
\end{proof}

\noindent
보조정리의 상황에서는 고차 직접상들이 소멸하므로 그것들 역시 연접이다.
국소 뇌터 스킴들 사이의 고유 사상에 대해서도 항상 그러함을
(Proposition \ref{coherent-proposition-proper-pushforward-coherent}) 보일 것이다.

\begin{lemma}
\label{coherent-lemma-coherent-support-dimension-0}
$X$를 국소 뇌터 스킴이라 하자. $\mathcal{F}$를
$\dim(\text{Supp}(\mathcal{F})) \leq 0$인 연접층이라 하자.
그러면 $\mathcal{F}$는 대역 절단들로 생성되고,
$H^i(X, \mathcal{F}) = 0$이며 이는 $i > 0$에 대해 성립한다.
\end{lemma}

\begin{proof}
Lemma \ref{coherent-lemma-coherent-support-closed}에 따라
$\mathcal{F} = i_*\mathcal{G}$이고, 여기서 $i : Z \to X$는
$\mathcal{F}$의 스킴론적 지지의 포함이며 $\mathcal{G}$는 연접
$\mathcal{O}_Z$-가군이다. $Z$의 차원이 $0$이므로 $Z$는 아핀들의
서로소 합집합이다(Properties, Lemma
\ref{properties-lemma-locally-Noetherian-dimension-0}). 따라서
$\mathcal{G}$는 대역적으로 생성되고 $\mathcal{G}$의 고차 코호몰로지군들은
영이다(Lemma \ref{coherent-lemma-quasi-coherent-affine-cohomology-zero}).
따라서 $\mathcal{F} = i_*\mathcal{G}$는 대역적으로 생성된다.
$\mathcal{F}$와 $\mathcal{G}$의 코호몰로지들은 일치하므로
(닫힌 몰입은 아핀이므로 Lemma
\ref{coherent-lemma-relative-affine-cohomology}가 적용된다), $\mathcal{F}$의 고차
코호몰로지군들은 영이다.
\end{proof}

\begin{lemma}
\label{coherent-lemma-pushforward-coherent-on-open}
$X$를 스킴이라 하고 $j : U \to X$를 열린집합의 포함이라 하자.
$T \subset X$를 $U$에 포함되는 닫힌 부분집합이라 하자.
$\mathcal{F}$가 연접 $\mathcal{O}_U$-가군이고
$\text{Supp}(\mathcal{F}) \subset T$이면 $j_*\mathcal{F}$는 연접
$\mathcal{O}_X$-가군이다.
\end{lemma}

\begin{proof}
열린 덮개 $X = U \cup (X \setminus T)$를 생각하자.
그러면 $j_*\mathcal{F}|_U = \mathcal{F}$는 연접이고,
$j_*\mathcal{F}|_{X \setminus T} = 0$도 연접이다. 따라서
$j_*\mathcal{F}$는 연접이다.
\end{proof}

\section{뇌터 스킴 위의 연접층}
\label{coherent-section-coherent-quasi-compact}

\noindent
이 절에서는 뇌터 스킴 위 연접층의 몇 가지 성질을 언급한다.

\begin{lemma}
\label{coherent-lemma-acc-coherent}
$X$를 뇌터 스킴이라 하고 $\mathcal{F}$를 연접 $\mathcal{O}_X$-가군이라
하자. $\mathcal{F}$의 준연접 부분가군들에 대해 오름사슬 조건이
성립한다. 다시 말해 준연접 부분가군들의 임의의 열
$$
\mathcal{F}_1 \subset \mathcal{F}_2 \subset \ldots \subset \mathcal{F}
$$
이 주어지면 $\mathcal{F}_n = \mathcal{F}_{n + 1} = \ldots $이고,
이는 어떤 $n \geq 0$에 대해 성립한다.
\end{lemma}

\begin{proof}
유한 아핀 열린 덮개를 택하자. 덮개의 각 원소 위에서 Algebra, Lemma
\ref{algebra-lemma-Noetherian-basic}에 따라 안정화를 얻는다.
따라서 보조정리가 성립한다.
\end{proof}

\begin{lemma}
\label{coherent-lemma-power-ideal-kills-sheaf}
$X$를 뇌터 스킴이라 하고 $\mathcal{F}$를 $X$ 위의 연접층이라 하자.
$\mathcal{I} \subset \mathcal{O}_X$를 닫힌 부분스킴 $Z \subset X$에
대응하는 준연접 아이디얼층이라 하자. 그러면 어떤 $n \geq 0$에 대해
$\mathcal{I}^n\mathcal{F} = 0$일 필요충분조건은
$\text{Supp}(\mathcal{F}) \subset Z$가 집합론적으로 성립하는 것이다.
\end{lemma}

\begin{proof}
$X$가 뇌터 환들의 스펙트럼들로 이루어진 유한 덮개를 가지므로,
이는 Algebra, Lemma
\ref{algebra-lemma-Noetherian-power-ideal-kills-module}에서 즉시 따라온다.
\end{proof}

\begin{lemma}[아르틴--리스]
\label{coherent-lemma-Artin-Rees}
$X$를 뇌터 스킴이라 하고 $\mathcal{F}$를 $X$ 위의 연접층이라 하자.
$\mathcal{G} \subset \mathcal{F}$를 준연접 부분층이라 하고,
$\mathcal{I} \subset \mathcal{O}_X$를 준연접 아이디얼층이라 하자.
그러면 어떤 $c \geq 0$이 존재하여 모든 $n \geq c$에 대해
$$
\mathcal{I}^{n - c}(\mathcal{I}^c\mathcal{F} \cap \mathcal{G})
=
\mathcal{I}^n\mathcal{F} \cap \mathcal{G}
$$
이다.
\end{lemma}

\begin{proof}
$X$가 뇌터 환들의 스펙트럼들로 이루어진 유한 덮개를 가지므로,
이는 Algebra, Lemma \ref{algebra-lemma-Artin-Rees}에서 즉시 따라온다.
\end{proof}

\begin{lemma}
\label{coherent-lemma-directed-colimit-coherent}
$X$를 뇌터 스킴이라 하자. 모든 준연접 $\mathcal{O}_X$-가군은 그 연접
부분가군들의 여과 쌍극한이다.
\end{lemma}

\begin{proof}
이는 Properties, Lemma
\ref{properties-lemma-quasi-coherent-colimit-finite-type}을 다시 쓴 것이다.
실제로 유한형 준연접 $\mathcal{O}_X$-가군은 Lemma
\ref{coherent-lemma-coherent-Noetherian}에 따라 연접이다.
\end{proof}

\begin{lemma}
\label{coherent-lemma-homs-over-open}
$X$를 뇌터 스킴이라 하자. $\mathcal{F}$를 준연접
$\mathcal{O}_X$-가군이라 하고, $\mathcal{G}$를 연접
$\mathcal{O}_X$-가군이라 하자. $\mathcal{I} \subset \mathcal{O}_X$를
준연접 아이디얼층이라 하자. 이에 대응하는 닫힌 부분스킴을
$Z \subset X$로 나타내고 $U = X \setminus Z$로 두자. 표준 동형
$$
\colim_n \Hom_{\mathcal{O}_X}(\mathcal{I}^n\mathcal{G}, \mathcal{F})
\longrightarrow
\Hom_{\mathcal{O}_U}(\mathcal{G}|_U, \mathcal{F}|_U).
$$
이 있다. 특히 동형
$$
\colim_n \Hom_{\mathcal{O}_X}(
\mathcal{I}^n, \mathcal{F})
\longrightarrow
\Gamma(U, \mathcal{F}).
$$
을 얻는다.
\end{lemma}

\begin{proof}
먼저 두 번째 사상이 동형임을 증명한다. 이는 Properties, Lemma
\ref{properties-lemma-sections-over-quasi-compact-open}에 따라 단사이다.
$\mathcal{F}$는 그 연접 부분가군들의 합집합이므로(Properties, Lemma
\ref{properties-lemma-quasi-coherent-colimit-finite-type} 및 Lemma
\ref{coherent-lemma-coherent-Noetherian} 참조), 전사성을 증명할 때
$\mathcal{F}$가 연접이라고 가정해도 되고 그렇게 하겠다.
$\mathcal{F}_n$을 $\mathcal{F}$의 준연접 부분층 중
$\mathcal{I}^n$에 의해 소멸되는 절단들로 이루어진 것이라 하자.
Properties, Lemma
\ref{properties-lemma-sections-over-quasi-compact-open}을 보라.
$\mathcal{F}_1 \subset \mathcal{F}_2 \subset \ldots$이므로 Lemma
\ref{coherent-lemma-acc-coherent}에 따라
$\mathcal{F}_n = \mathcal{F}_{n + 1} = \ldots $이고, 이는 어떤
$n \geq 0$에 대해 성립한다. $\mathcal{H} = \mathcal{F}_n$으로 두되
이때의 $n$을 사용하자. 아르틴--리스
(Lemma \ref{coherent-lemma-Artin-Rees})에 따라 어떤 $c \geq 0$이 존재하여
$\mathcal{I}^m\mathcal{F} \cap \mathcal{H}
\subset \mathcal{I}^{m - c}\mathcal{H}$이다. $m = n + c$를 택하면
$\mathcal{I}^m\mathcal{F} \cap \mathcal{H} \subset \mathcal{I}^n\mathcal{H}
= 0$을 얻는다. 따라서 $\mathcal{F}' = \mathcal{I}^m\mathcal{F}$로 두면
$\mathcal{F}' \cap \mathcal{F}_n = 0$이고
$\mathcal{F}'|_U = \mathcal{F}|_U$이다. 특히 절단들의 부분층
$(\mathcal{F}')_N$, 즉 $\mathcal{I}^N$에 의해 소멸되는 것들은 모든
$N \geq 0$에 대해
영이다. 따라서 Properties, Lemma
\ref{properties-lemma-sections-over-quasi-compact-open}에 의해 다음 가환
도표의 위쪽 수평 화살표가 전단사임을 알 수 있다.
$$
\xymatrix{
\colim_n \Hom_{\mathcal{O}_X}(
\mathcal{I}^n, \mathcal{F}')
\ar[r] \ar[d] &
\Gamma(U, \mathcal{F}') \ar[d] \\
\colim_n \Hom_{\mathcal{O}_X}(
\mathcal{I}^n, \mathcal{F})
\ar[r] &
\Gamma(U, \mathcal{F})
}
$$
오른쪽 수직 화살표도 전단사이므로, 아래쪽 수평 화살표가 원하는 대로
전사임을 얻는다.

\medskip\noindent
다음으로 보조정리의 첫 번째 화살표가 전단사임을 증명한다.
Lemma \ref{coherent-lemma-coherent-Noetherian}에 따라 층 $\mathcal{G}$는 유한
표시이고, 따라서 층
$\mathcal{H} = \SheafHom_{\mathcal{O}_X}(\mathcal{G}, \mathcal{F})$는
준연접이다. Schemes, Section
\ref{schemes-section-quasi-coherent}을 보라. 정의에 따라
$$
\mathcal{H}(U)
=
\Hom_{\mathcal{O}_U}(\mathcal{G}|_U, \mathcal{F}|_U)
$$
이다. 보조정리의 첫 번째 화살표의 우변에서 $\psi$를 하나 택하자.
즉 $\psi \in \mathcal{H}(U)$이다. 방금 증명한 결과를 $\mathcal{H}$에
적용하면 어떤 $n \geq 0$과 사상
$\varphi : \mathcal{I}^n \to \mathcal{H}$가 존재하여 $\psi$를 복원하며,
그 제한은 $U$ 위에서 취한다. Modules, Lemma
\ref{modules-lemma-internal-hom}에 따라
$\varphi$는 사상
$$
\varphi :
\mathcal{I}^n \otimes_{\mathcal{O}_X} \mathcal{G}
\longrightarrow
\mathcal{F}.
$$
에 대응한다. 이는 거의 원하는 것이지만 화살표의 원천은
$\mathcal{I}^n$과 $\mathcal{G}$의 곱이 아니라 텐서곱이다.
$n$을 늘리는 대가로 이 차이가 무관해짐을 보이겠다. 짧은 완전열
$$
0 \to \mathcal{K} \to
\mathcal{I}^n \otimes_{\mathcal{O}_X} \mathcal{G} \to
\mathcal{I}^n\mathcal{G} \to 0
$$
을 생각하자. 여기서 $\mathcal{K}$는 핵으로 정의된다.
$\mathcal{I}^n\mathcal{K} = 0$임에 유의하라(증명은 생략한다).
아르틴--리스를 다시 적용하면 충분히 큰 어떤 $m$에 대해
$$
\mathcal{K}
\cap
\mathcal{I}^m(\mathcal{I}^n \otimes_{\mathcal{O}_X} \mathcal{G})
=
0
$$
임을 알 수 있다. 다시 말해
$$
\mathcal{I}^m(\mathcal{I}^n \otimes_{\mathcal{O}_X} \mathcal{G})
\longrightarrow
\mathcal{I}^{n + m}\mathcal{G}
$$
은 동형이다. $\varphi'$를 이 부분가군에 대한 $\varphi$의 제한으로 두고,
이를 사상 $\mathcal{I}^{m + n}\mathcal{G} \to \mathcal{F}$로 보자.
그러면 $\varphi'$는 보조정리의 첫 번째 화살표의 좌변의 원소를 주며,
그 화살표 아래에서 $\psi$로 간다. 다시 말해 이 화살표의 전사성을
증명하였다. 단사성의 증명은 생략한다.
\end{proof}

\begin{lemma}
\label{coherent-lemma-extend-coherent}
$X$를 국소 뇌터 스킴이라 하자. $\mathcal{F}$, $\mathcal{G}$를 연접
$\mathcal{O}_X$-가군들이라 하자. $U \subset X$를 열린집합이라 하고,
$\varphi : \mathcal{F}|_U \to \mathcal{G}|_U$를
$\mathcal{O}_U$-가군 사상이라 하자. 그러면 연접 부분가군
$\mathcal{F}' \subset \mathcal{F}$가 존재하며, 이는 $\mathcal{F}$와
$U$ 위에서 일치한다. 또한 $\varphi$는 사상
$\varphi' : \mathcal{F}' \to \mathcal{G}$로 연장된다.
\end{lemma}

\begin{proof}
$\mathcal{I} \subset \mathcal{O}_X$를 $X \setminus U$ 위의 축소 유도
스킴 구조를 잘라내는 연접 아이디얼층이라 하자. $X$가 뇌터이면 Lemma
\ref{coherent-lemma-homs-over-open}에 따라
$\mathcal{F}' = \mathcal{I}^n\mathcal{F}$로 택할 수 있고, 여기서 $n$은
어떤 정수이다. 일반적인 경우는
초른의 보조정리를 사용하여 이 경우로부터 따라온다.

\medskip\noindent
삼중항 $(U', \mathcal{F}', \varphi')$들의 집합을 생각하자. 여기서
$U \subset U' \subset X$는 열린집합이고,
$\mathcal{F}' \subset \mathcal{F}|_{U'}$는 $\mathcal{F}$와 일치하되
$U$ 위에서 일치하는 연접 부분층이며,
$\varphi' : \mathcal{F}' \to \mathcal{G}|_{U'}$는 $\varphi$로 제한되되
그 제한은 $U$ 위에서 취한다. 다음 조건이 성립할 때 그리고 그때에만
$(U'', \mathcal{F}'', \varphi'') \geq (U', \mathcal{F}', \varphi')$라 하자.
$U'' \supset U'$, $\mathcal{F}''|_{U'} = \mathcal{F}'$, 그리고
$\varphi''|_{U'} = \varphi'$.
삼중항들의 전순서화된 모음 $(U_i, \mathcal{F}_i, \varphi_i)$가 주어지면
$\mathcal{F}_i$들을 붙여 부분층 $\mathcal{F}'$을 얻는다. 이는
$\mathcal{F}$의 부분층이며 $U' = \bigcup U_i$ 위에 놓인다. 또한
$\varphi$를 사상
$\varphi' : \mathcal{F}' \to \mathcal{G}|_{U'}$로 연장할 수 있음은
명백하다. 따라서 삼중항들의 모든 전순서화된 부분집합은 상계를 갖는다.
마지막으로 $(U', \mathcal{F}', \varphi')$가 임의의 삼중항이지만
$U' \not = X$라고 하자. 아핀 열린집합 $W \subset X$를 택하되,
이는 $U'$에 포함되지 않도록 할 수 있다. 첫 문단의 결과에 따라 부분층
$\mathcal{F}'|_{W \cap U'}$와 제한 $\varphi'|_{W \cap U'}$를 어떤 부분층
$\mathcal{F}'' \subset \mathcal{F}|_W$와 사상
$\varphi'' : \mathcal{F}'' \to \mathcal{G}|_W$로 연장할 수 있다.
$(\mathcal{F}', \varphi')$와 $(\mathcal{F}'', \varphi'')$의 일치가
$W \cap U'$ 위에서 성립한다는 것은 정확히 이것을 삼중항
$(U' \cup W, \mathcal{F}''', \varphi''')$로 연장할 수 있다는 뜻이다.
따라서 모든 극대 삼중항 $(U', \mathcal{F}', \varphi')$은
(초른의 보조정리에 의해 존재하며) $U' = X$를 만족한다. 이로써 증명이
끝난다.
\end{proof}

\section{깊이}
\label{coherent-section-depth}

\noindent
이 절에서는 국소 뇌터 스킴 위 연접 가군의 깊이와 성질 $(S_k)$에 대해
조금 논한다. 이 개념은 이미 국소 뇌터 스킴에 대해 Properties, Section
\ref{properties-section-Rk}에서 논하였다.

\begin{definition}
\label{coherent-definition-depth}
$X$를 국소 뇌터 스킴이라 하고 $\mathcal{F}$를 연접
$\mathcal{O}_X$-가군이라 하자. $k \geq 0$을 정수라 하자.
\begin{enumerate}
\item $\mathcal{F}$가 {\it 깊이 $k$를 점 $x$에서 갖는다}고 함은,
그 점이 $X$에 속하고
$\text{depth}_{\mathcal{O}_{X, x}}(\mathcal{F}_x) = k$일 때이다.
\item $X$가 {\it 깊이 $k$를 점 $x$에서 갖는다}고 함은, 이 점이
$X$에 속하고 $\text{depth}(\mathcal{O}_{X, x}) = k$일 때이다.
\item $\mathcal{F}$가 성질 {\it $(S_k)$}를 갖는다고 함은 모든
$x \in X$에 대해
$$
\text{depth}_{\mathcal{O}_{X, x}}(\mathcal{F}_x)
\geq \min(k, \dim(\text{Supp}(\mathcal{F}_x)))
$$
일 때이다.
\item $X$가 성질 {\it $(S_k)$}를 갖는다고 함은 $\mathcal{O}_X$가
성질 $(S_k)$를 가질 때이다.
\end{enumerate}
\end{definition}

\noindent
모든 연접층은 조건 $(S_0)$를 만족한다. 조건 $(S_1)$은 매장된 연관점을
갖지 않는 것과 동치이다. Divisors, Lemma
\ref{divisors-lemma-S1-no-embedded}를 보라.

\begin{lemma}
\label{coherent-lemma-hom-into-depth}
$X$를 국소 뇌터 스킴이라 하자. $\mathcal{F}$, $\mathcal{G}$를 연접
$\mathcal{O}_X$-가군들이라 하고 $x \in X$라 하자.
\begin{enumerate}
\item $\mathcal{G}_x$의 깊이가 $\geq 1$이면
$\SheafHom_{\mathcal{O}_X}(\mathcal{F}, \mathcal{G})_x$의 깊이도
$\geq 1$이다.
\item $\mathcal{G}_x$의 깊이가 $\geq 2$이면
$\Hom_{\mathcal{O}_X}(\mathcal{F}, \mathcal{G})_x$의 깊이도
$\geq 2$이다.
\end{enumerate}
\end{lemma}

\begin{proof}
$\SheafHom_{\mathcal{O}_X}(\mathcal{F}, \mathcal{G})$는 Lemma
\ref{coherent-lemma-tensor-hom-coherent}에 따라 연접 $\mathcal{O}_X$-가군이다.
연접 가군은 유한 표시이므로(Lemma \ref{coherent-lemma-coherent-Noetherian}),
싹을 취하는 것은 $\SheafHom$ 및 $\Hom$을 취하는 것과 가환한다.
Modules, Lemma \ref{modules-lemma-stalk-internal-hom}을 보라.
따라서 국소환 위 유한 가군의 경우로 환원되며, 이는 More on Algebra,
Lemma \ref{more-algebra-lemma-hom-into-depth}이다.
\end{proof}

\begin{lemma}
\label{coherent-lemma-hom-into-S2}
$X$를 국소 뇌터 스킴이라 하자. $\mathcal{F}$, $\mathcal{G}$를 연접
$\mathcal{O}_X$-가군들이라 하자.
\begin{enumerate}
\item $\mathcal{G}$가 성질 $(S_1)$을 가지면
$\SheafHom_{\mathcal{O}_X}(\mathcal{F}, \mathcal{G})$도 성질 $(S_1)$을 갖는다.
\item $\mathcal{G}$가 성질 $(S_2)$를 가지면
$\SheafHom_{\mathcal{O}_X}(\mathcal{F}, \mathcal{G})$도 성질 $(S_2)$를 갖는다.
\end{enumerate}
\end{lemma}

\begin{proof}
Lemma \ref{coherent-lemma-hom-into-depth}와 정의들에서 즉시 따라온다.
\end{proof}

\noindent
Properties, Lemma \ref{properties-lemma-scheme-CM-iff-all-Sk}에서 국소 뇌터
스킴이 코언--매콜리일 필요충분조건은 $(S_k)$가 모든 $k$에 대해
성립하는 것임을 보았다. 따라서 다음 정의를 도입하는 것은 자연스럽다.
이는 모든 싹이 코언--매콜리 가군이라는 조건과 동치이다.

\begin{definition}
\label{coherent-definition-Cohen-Macaulay}
$X$를 국소 뇌터 스킴이라 하고 $\mathcal{F}$를 연접
$\mathcal{O}_X$-가군이라 하자. $\mathcal{F}$가 {\it 코언--매콜리}라고
함은 $(S_k)$가 모든 $k \geq 0$에 대해 성립할 때이다.
\end{definition}

\begin{lemma}
\label{coherent-lemma-Cohen-Macaulay-over-regular}
$X$를 정칙 스킴이라 하고 $\mathcal{F}$를 연접 $\mathcal{O}_X$-가군이라
하자. 다음은 서로 동치이다.
\begin{enumerate}
\item $\mathcal{F}$는 코언--매콜리이고 $\text{Supp}(\mathcal{F}) = X$이다.
\item $\mathcal{F}$는 계수 $> 0$인 유한 국소 자유 가군이다.
\end{enumerate}
\end{lemma}

\begin{proof}
$x \in X$라 하자. (2)가 성립하면 $\mathcal{F}_x$는 자유
$\mathcal{O}_{X, x}$-가군이고 그 계수는 $> 0$이다. 따라서
$\text{depth}(\mathcal{F}_x) = \dim(\mathcal{O}_{X, x})$이다. 실제로 정칙
국소환은 코언--매콜리이다(Algebra, Lemma
\ref{algebra-lemma-regular-ring-CM}). 반대로 (1)이 성립하면
$\mathcal{F}_x$는 $\mathcal{O}_{X, x}$ 위의 극대 코언--매콜리 가군이다
(Algebra, Definition \ref{algebra-definition-maximal-CM}). 따라서 Algebra,
Lemma \ref{algebra-lemma-regular-mcm-free}에 따라 $\mathcal{F}_x$는
자유이다.
\end{proof}

\section{연접층의 데비사주}
\label{coherent-section-devissage}

\noindent
$X$를 뇌터 스킴이라 하자. 정수적 닫힌 부분스킴
$i : Z \to X$를 생각하자. $i_*\mathcal{G}$ 꼴의 연접층을 생각하는
것이 흔히 편리하며, 여기서 $\mathcal{G}$는 $Z$ 위의 연접층이다.
특히 $\mathcal{G}$가 꼬임 없는 계수 $1$의 층인 경우에 이러한 층들에
관심이 있다. 예를 들어 $\mathcal{G}$는 $Z$ 위의 영이 아닌 아이디얼층일
수 있고, 더 구체적으로는 $\mathcal{G} = \mathcal{O}_Z$일 수도 있다.

\medskip\noindent
이 절 전체에서 연접층은 유한형 준연접층과 같은 것이며, 연접층의
준연접 부분몫이 연접이라는 사실을 사용한다. 절
\ref{coherent-section-coherent-sheaves}를 보라. 연접층의 지지집합은 닫혀 있다.
Modules, Lemma \ref{modules-lemma-support-finite-type-closed}를 보라.

\begin{lemma}
\label{coherent-lemma-prepare-filter-support}
$X$를 뇌터 스킴이라 하자. $\mathcal{F}$를 $X$ 위의 연접층이라 하자.
$\text{Supp}(\mathcal{F}) = Z \cup Z'$이고 $Z$, $Z'$가 닫혀 있다고
가정하자.
그러면 연접층들의 짧은 완전열
$$
0 \to \mathcal{G}' \to \mathcal{F} \to \mathcal{G} \to 0
$$
이 존재하여 $\text{Supp}(\mathcal{G}') \subset Z'$이고
$\text{Supp}(\mathcal{G}) \subset Z$이다.
\end{lemma}

\begin{proof}
$\mathcal{I} \subset \mathcal{O}_X$를 $Z$ 위의 유도된 축소 닫힌
부분스킴 구조를 정의하는 아이디얼층이라 하자. Schemes, Lemma
\ref{schemes-lemma-reduced-closed-subscheme}를 보라.
부분층 $\mathcal{G}'_n = \mathcal{I}^n\mathcal{F}$와 몫
$\mathcal{G}_n = \mathcal{F}/\mathcal{I}^n\mathcal{F}$를 생각하자.
각 $n$에 대해 짧은 완전열
$$
0 \to \mathcal{G}'_n \to \mathcal{F} \to \mathcal{G}_n \to 0
$$
이 있다. 모든 점 $x$ 가운데 $Z' \setminus Z$에 속하는 것에 대해
$\mathcal{I}_x = \mathcal{O}_{X, x}$이고, 따라서
$\mathcal{G}_{n, x} = 0$이다. 그러므로
$\text{Supp}(\mathcal{G}_n) \subset Z$임을 알 수 있다.
$X \setminus Z'$는 뇌터 스킴임을 주목하자. 따라서 Lemma
\ref{coherent-lemma-power-ideal-kills-sheaf}에 의해 어떤 $n$이 존재하여
$\mathcal{G}'_n|_{X \setminus Z'} =
\mathcal{I}^n\mathcal{F}|_{X \setminus Z'} = 0$이다.
그러한 $n$에 대해
$\text{Supp}(\mathcal{G}'_n) \subset Z'$임을 알 수 있다.
따라서 $\mathcal{G}' = \mathcal{G}'_n$ 및
$\mathcal{G} = \mathcal{G}_n$으로 놓으면 된다.
\end{proof}

\begin{lemma}
\label{coherent-lemma-prepare-filter-irreducible}
$X$를 뇌터 스킴이라 하자. $i : Z \to X$를 정수적 닫힌 부분스킴이라
하고, $\xi \in Z$를 그 일반점이라 하자. $\mathcal{F}$를 $X$ 위의
연접층이라 하자. $\mathcal{F}_\xi$가 $\mathfrak m_\xi$에 의해
소멸된다고 가정하자. 그러면 정수 $r \geq 0$, 연접 아이디얼층
$\mathcal{I} \subset \mathcal{O}_Z$, 그리고 연접층들의 단사 사상
$$
i_*\left(\mathcal{I}^{\oplus r}\right) \to \mathcal{F}
$$
이 존재하며, 이 사상은 $\xi$의 한 근방에서 동형이다.
\end{lemma}

\begin{proof}
$\mathcal{J} \subset \mathcal{O}_X$를 $Z$의 아이디얼층이라 하자.
$\mathcal{F}' \subset \mathcal{F}$를 $\mathcal{F}$의 국소 단면들 중
$\mathcal{J}$에 의해 소멸되는 것들로 이루어진 부분층이라 하자. 이는
Properties, Lemma
\ref{properties-lemma-sections-annihilated-by-ideal}에 의해
준연접층이다. 또한 $\mathcal{F}'_\xi = \mathcal{F}_\xi$이다.
실제로 $\mathcal{J}_\xi = \mathfrak m_\xi$이고 Properties, Lemma
\ref{properties-lemma-sections-annihilated-by-ideal}의 (3)을 적용한다.
Lemma \ref{coherent-lemma-local-isomorphism}에 의해
$\mathcal{F}' \to \mathcal{F}$는 $\xi$의 한 근방에서 동형을 유도한다.
따라서 $\mathcal{F}$를 $\mathcal{F}'$로 바꾸어
$\mathcal{F}$가 $\mathcal{J}$에 의해 소멸된다고 가정해도 된다.

\medskip\noindent
$\mathcal{J}\mathcal{F} = 0$이라 가정하자. Lemma
\ref{coherent-lemma-i-star-equivalence}에 의해
$\mathcal{F} = i_*\mathcal{G}$로 쓸 수 있으며, 여기서
$\mathcal{G}$는 $Z$ 위의 어떤 연접층이다. 사상
$\mathcal{I}^{\oplus r} \to \mathcal{G}$를 찾아서 일반점
$\xi$의 한 근방에서 동형이 되게 할 수 있다고 하자. 이 일반점은
$Z$의 일반점이다.
그러면 $i_*$를 적용하여(이는 왼쪽 완전이다) 보조정리의 결론을 얻는다.
따라서 $X = Z$인 경우로 환원되었다.

\medskip\noindent
$Z = X$가 일반점 $\xi$를 갖는 정수적 뇌터 스킴이라고 하자.
이 경우 $\mathcal{O}_{X, \xi} = \kappa(\xi)$는 $X$의 함수체임을
주목하자. $\mathcal{F}_\xi$는 유한 $\mathcal{O}_\xi$-가군이므로
$r = \dim_{\kappa(\xi)} \mathcal{F}_\xi$는 유한하다.
따라서 층 $\mathcal{O}_X^{\oplus r}$와 $\mathcal{F}$의 $\xi$에서의
싹들은 서로 동형이다. Lemma \ref{coherent-lemma-map-stalks-local-map}에 의해
영이 아닌 열린집합 $U \subset X$와 사상
$\psi : \mathcal{O}_X^{\oplus r}|_U \to \mathcal{F}|_U$가 존재하여
$\xi$에서 동형이다. 따라서 Lemma \ref{coherent-lemma-local-isomorphism}에 의해
이는 $\xi$의 한 근방에서 동형이다. Schemes, Lemma
\ref{schemes-lemma-reduced-closed-subscheme}에 의해 준연접 아이디얼층
$\mathcal{I} \subset \mathcal{O}_X$가 존재하며, 이에 연관된 닫힌
부분스킴 $Y \subset X$는 $U$의 여집합이다.
Lemma \ref{coherent-lemma-homs-over-open}에 의해 어떤 $n \geq 0$와 사상
$\mathcal{I}^n(\mathcal{O}_X^{\oplus r}) \to \mathcal{F}$가 존재하여
우리의 $\psi$를 $U$ 위에서 복원한다.
$\mathcal{I}^n(\mathcal{O}_X^{\oplus r}) = (\mathcal{I}^n)^{\oplus r}$이므로
보조정리에서와 같은 사상을 얻는다. $X$가 정수적이고 이 사상이
$X$의 일반점에서 단사이므로, 이 사상은 단사이다
(쉬운 증명은 생략한다).
\end{proof}

\begin{lemma}
\label{coherent-lemma-coherent-filter}
$X$를 뇌터 스킴이라 하고 $\mathcal{F}$를 $X$ 위의 연접층이라 하자.
연접 부분층들에 의한 여과
$$
0 = \mathcal{F}_0 \subset \mathcal{F}_1 \subset
\ldots \subset \mathcal{F}_m = \mathcal{F}
$$
가 존재하여, 각 $j = 1, \ldots, m$에 대해 정수적 닫힌 부분스킴
$Z_j \subset X$와 영이 아닌 연접 아이디얼층
$\mathcal{I}_j \subset \mathcal{O}_{Z_j}$가 존재하고
$$
\mathcal{F}_j/\mathcal{F}_{j - 1}
\cong (Z_j \to X)_* \mathcal{I}_j
$$
이다.
\end{lemma}

\begin{proof}
다음 모음을 생각하자.
$$
\mathcal{T} =
\left\{
\begin{matrix}
Z \subset X
\text{ closed such that there exists a coherent sheaf }
\mathcal{F} \\
\text{ with }
\text{Supp}(\mathcal{F}) = Z
\text{ for which the lemma is wrong}
\end{matrix}
\right\}
$$
우리는 $\mathcal{T}$가 공집합임을 보이려 한다. 그렇지 않다면 $X$가
뇌터이므로 최소 원소 $Z \in \mathcal{T}$를 택할 수 있다. 이는
연접층 $\mathcal{F}$가 $X$ 위에 존재하여 지지집합이 $Z$이고
보조정리의 결론이 성립하지 않는다는 뜻이다. 지지집합이 공집합인 유일한 층은
영층이므로 분명히 $Z \not = \emptyset$이다. 영층에 대해서는
보조정리가 성립한다($m = 0$으로 두면 된다).

\medskip\noindent
$Z$가 기약이 아니면 $Z = Z_1 \cup Z_2$로 쓸 수 있으며, 여기서
$Z_1, Z_2$는 닫혀 있고 $Z$보다 진정으로 작다. 그러면 Lemma
\ref{coherent-lemma-prepare-filter-support}를 적용하여 연접층들의 짧은 완전열
$$
0 \to
\mathcal{G}_1 \to
\mathcal{F} \to
\mathcal{G}_2 \to 0
$$
을 얻으며, 여기서 $\text{Supp}(\mathcal{G}_i) \subset Z_i$이다.
$Z$의 최소성에 의해 각 $\mathcal{G}_i$는 보조정리의 명제에 나오는
여과를 갖는다. $\mathcal{F}$ 위에 유도된 여과를 생각하면 모순에
이른다. 따라서 $Z$는 기약이다.

\medskip\noindent
$Z$가 기약이라고 하자. $\mathcal{J}$를 $Z$ 위의 유도된 축소 닫힌
부분스킴 구조를 잘라내는 아이디얼층이라 하자. Schemes, Lemma
\ref{schemes-lemma-reduced-closed-subscheme}를 보라.
Lemma \ref{coherent-lemma-power-ideal-kills-sheaf}에 의해 어떤 $n \geq 0$가
존재하여 $\mathcal{J}^n\mathcal{F} = 0$이다. 따라서 여과
$$
0 = \mathcal{J}^n\mathcal{F} \subset \mathcal{J}^{n - 1}\mathcal{F}
\subset \ldots \subset \mathcal{J}\mathcal{F} \subset \mathcal{F}
$$
를 얻으며, 그 연속하는 각 부분몫은 $\mathcal{J}$에 의해 소멸된다.
각 부분몫이 보조정리의 명제에 나오는 여과를 가지면
$\mathcal{F}$도 그러한 여과를 갖는다. 다시 말해
$\mathcal{J}$가 $\mathcal{F}$를 소멸시킨다고 가정해도 된다.

\medskip\noindent
$Z$가 기약이고 $\mathcal{J}\mathcal{F} = 0$인 경우 Lemma
\ref{coherent-lemma-prepare-filter-irreducible}를 적용할 수 있다.
이로부터 짧은 완전열
$$
0 \to
i_*(\mathcal{I}^{\oplus r}) \to
\mathcal{F} \to
\mathcal{Q} \to 0
$$
을 얻으며, 여기서 $\mathcal{Q}$는 몫으로 정의된다. 방금 인용한
보조정리에 의해 $\mathcal{Q}$는 $\xi$의 한 근방에서 영이므로,
$\mathcal{Q}$의 지지집합은 $Z$보다 진정으로 작다. 따라서
$\mathcal{Q}$는 $Z$의 최소성에 의해 원하는 종류의 여과를 갖는다.
그러면 분명히 $\mathcal{F}$도 그러한 여과를 가지므로, 마지막 모순을
얻는다.
\end{proof}

\begin{lemma}
\label{coherent-lemma-property-initial}
$X$를 뇌터 스킴이라 하고, $\mathcal{P}$를 $X$ 위 연접층들의 한
성질이라 하자. 다음을 가정하자.
\begin{enumerate}
\item 연접층들의 임의의 짧은 완전열
$$
0 \to \mathcal{F}_1 \to \mathcal{F} \to \mathcal{F}_2 \to 0
$$
에 대해, $\mathcal{F}_i$, $i = 1, 2$가 성질 $\mathcal{P}$를 가지면
$\mathcal{F}$도 이 성질을 갖는다.
\item 모든 정수적 닫힌 부분스킴 $Z \subset X$와 모든 준연접
아이디얼층 $\mathcal{I} \subset \mathcal{O}_Z$에 대해
$\mathcal{P}$가 $i_*\mathcal{I}$에 대해 성립한다.
\end{enumerate}
그러면 성질 $\mathcal{P}$는 $X$ 위의 모든 연접층에 대해 성립한다.
\end{lemma}

\begin{proof}
먼저 연접층 $\mathcal{F}$가 연접 부분층들에 의한 여과
$$
0 = \mathcal{F}_0 \subset \mathcal{F}_1 \subset
\ldots \subset \mathcal{F}_m = \mathcal{F}
$$
를 갖고 각 $\mathcal{F}_i/\mathcal{F}_{i - 1}$이 성질
$\mathcal{P}$를 가지면, $\mathcal{F}$도 이 성질을 가짐을 주목하자.
이는 $\mathcal{P}$의 성질 (1)에서 따른다. 한편 Lemma
\ref{coherent-lemma-coherent-filter}에 의해 임의의 $\mathcal{F}$를
(2)에서와 같은 연속 부분몫들을 갖도록 여과할 수 있다.
따라서 보조정리가 따른다.
\end{proof}

\begin{lemma}
\label{coherent-lemma-property-irreducible}
$X$를 뇌터 스킴이라 하자. $Z_0 \subset X$를 일반점 $\xi$를 갖는
기약 닫힌 부분집합이라 하자. $\mathcal{P}$를 $X$ 위 연접층들 가운데
지지집합이 $Z_0$에 포함되는 것들의 성질이라 하고, 다음을 가정하자.
\begin{enumerate}
\item 연접층들의 임의의 짧은 완전열에 대해 그 세 층 중 둘이 성질
$\mathcal{P}$를 가지면 나머지 하나도 이 성질을 갖는다.
\item 모든 정수적 닫힌 부분스킴 $Z \subset Z_0 \subset X$,
$Z \not = Z_0$와 모든 준연접 아이디얼층
$\mathcal{I} \subset \mathcal{O}_Z$에 대해
$\mathcal{P}$가 $(Z \to X)_*\mathcal{I}$에 대해 성립한다.
\item 다음을 만족하는 어떤 연접층 $\mathcal{G}$가 $X$ 위에 존재한다.
\begin{enumerate}
\item $\text{Supp}(\mathcal{G}) = Z_0$,
\item $\mathcal{G}_\xi$는 $\mathfrak m_\xi$에 의해 소멸되고,
\item $\dim_{\kappa(\xi)} \mathcal{G}_\xi = 1$이며,
\item 성질 $\mathcal{P}$가 $\mathcal{G}$에 대해 성립한다.
\end{enumerate}
\end{enumerate}
그러면 성질 $\mathcal{P}$는 모든 연접층 $\mathcal{F}$ 가운데
$X$ 위에서 지지집합이 $Z_0$에 포함되는 것에 대해 성립한다.
\end{lemma}

\begin{proof}
먼저 연접층 $\mathcal{F}$의 지지집합이 $Z_0$에 포함되고 이 층이 연접
부분층들에 의한 여과
$$
0 = \mathcal{F}_0 \subset \mathcal{F}_1 \subset
\ldots \subset \mathcal{F}_m = \mathcal{F}
$$
를 갖고 각 $\mathcal{F}_i/\mathcal{F}_{i - 1}$이 성질
$\mathcal{P}$를 가지면, $\mathcal{F}$도 이 성질을 가짐을 주목하자.
또는 $\mathcal{F}$가 성질 $\mathcal{P}$를 갖고
$\mathcal{F}_i/\mathcal{F}_{i - 1}$ 중 하나를 제외한 모두가 성질
$\mathcal{P}$를 가지면, 마지막 하나도 이 성질을 갖는다.
이는 가정 (1)에서 따른다.

\medskip\noindent
첫 응용으로, 지지집합이 $Z_0$에 진정으로 포함되는 모든 연접층이
성질 $\mathcal{P}$를 가짐을 얻는다. 실제로 그러한 층은 여과를 가지며
(Lemma \ref{coherent-lemma-coherent-filter}를 보라), 그 부분몫들은 (2)에 의해
성질 $\mathcal{P}$를 갖는다.

\medskip\noindent
$\mathcal{G}$를 (3)에서와 같이 택하자. Lemma
\ref{coherent-lemma-prepare-filter-irreducible}에 의해 아이디얼층
$\mathcal{I}$가 $Z_0$ 위에 존재하고 정수 $r \geq 1$도 존재한다.
또한 짧은 완전열
$$
0 \to
\left((Z_0 \to X)_*\mathcal{I}\right)^{\oplus r} \to
\mathcal{G} \to
\mathcal{Q} \to 0
$$
이 존재하며, $\mathcal{Q}$의 지지집합은 $Z_0$에 진정으로 포함된다.
(3)(c)에 의해 $r = 1$임을 알 수 있다. $\mathcal{Q}$도 성질
$\mathcal{P}$를 가지므로 $(Z_0 \to X)_*\mathcal{I}$가 성질
$\mathcal{P}$를 갖는다고 결론 내린다.

\medskip\noindent
다음으로 $\mathcal{I}' \not = 0$을 $Z_0$ 위의 또 다른 준연접
아이디얼층이라 하자. 교집합
$\mathcal{I}'' = \mathcal{I}' \cap \mathcal{I}$을 생각하면 두 짧은 완전열
$$
0 \to
(Z_0 \to X)_*\mathcal{I}'' \to
(Z_0 \to X)_*\mathcal{I} \to
\mathcal{Q} \to 0
$$
및
$$
0 \to
(Z_0 \to X)_*\mathcal{I}'' \to
(Z_0 \to X)_*\mathcal{I}' \to
\mathcal{Q}' \to 0.
$$
을 얻는다. 연접층 $\mathcal{Q}$와 $\mathcal{Q}'$의 지지집합들은
$Z_0$에 진정으로 포함됨을 주목하자. 따라서 $\mathcal{Q}$와
$\mathcal{Q}'$는 성질 $\mathcal{P}$를 갖는다(위 참조).
그러므로 (1)을 사용하여 $(Z_0 \to X)_*\mathcal{I}''$와
$(Z_0 \to X)_*\mathcal{I}'$도 모두 $\mathcal{P}$를 갖는다고
결론 내린다.

\medskip\noindent
마지막으로 임의의 연접층 $\mathcal{F}$가 $X$ 위에 있고 그
지지집합이 $Z_0$에 포함되면 이 층은 여과를 가지며(Lemma
\ref{coherent-lemma-coherent-filter}를
다시 보라), 그 모든 부분몫은 방금 보인 바에 의해 성질
$\mathcal{P}$를 가짐을 주목하면 증명이 끝난다.
\end{proof}

\begin{lemma}
\label{coherent-lemma-property}
$X$를 뇌터 스킴이라 하자. $\mathcal{P}$를 $X$ 위 연접층들의 한
성질이라 하고, 다음을 가정하자.
\begin{enumerate}
\item 연접층들의 임의의 짧은 완전열에 대해 그 세 층 중 둘이 성질
$\mathcal{P}$를 가지면 나머지 하나도 이 성질을 갖는다.
\item 모든 정수적 닫힌 부분스킴 $Z \subset X$와 그 일반점 $\xi$에
대해 다음을 만족하는 어떤 연접층 $\mathcal{G}$가 존재한다.
\begin{enumerate}
\item $\text{Supp}(\mathcal{G}) = Z$,
\item $\mathcal{G}_\xi$는 $\mathfrak m_\xi$에 의해 소멸되고,
\item $\dim_{\kappa(\xi)} \mathcal{G}_\xi = 1$이며,
\item 성질 $\mathcal{P}$가 $\mathcal{G}$에 대해 성립한다.
\end{enumerate}
\end{enumerate}
그러면 성질 $\mathcal{P}$는 $X$ 위의 모든 연접층에 대해 성립한다.
\end{lemma}

\begin{proof}
Lemma \ref{coherent-lemma-property-initial}에 따라, 모든 정수적 닫힌 부분스킴
$Z \subset X$와 모든 준연접 아이디얼층
$\mathcal{I} \subset \mathcal{O}_Z$에 대해
$\mathcal{P}$가 $(Z \to X)_*\mathcal{I}$에 대해 성립함을 보이면
충분하다. 그렇지 않다고 하자. $X$가 뇌터이므로 최소 정수적 닫힌
부분스킴 $Z_0 \subset X$가 존재하여, 성질 $\mathcal{P}$는
$(Z_0 \to X)_*\mathcal{I}_0$에 대해 성립하지 않는다. 여기서
$\mathcal{I}_0 \subset \mathcal{O}_{Z_0}$는 어떤 준연접
아이디얼층이다. 그러나 성질 $\mathcal{P}$는
$(Z \to X)_*\mathcal{I}$에 대해 모든 정수적 닫힌 부분스킴
$Z \subset Z_0$, $Z \not = Z_0$와 준연접 아이디얼층
$\mathcal{I} \subset \mathcal{O}_Z$마다 성립한다. (2)에 의해
$\mathcal{G}$가 $Z_0$에 대해 존재하므로, Lemma
\ref{coherent-lemma-property-irreducible}에 따르면 이러한 일은 일어날 수 없다.
\end{proof}

\begin{lemma}
\label{coherent-lemma-property-irreducible-higher-rank-cohomological}
$X$를 뇌터 스킴이라 하자. $Z_0 \subset X$를 일반점 $\xi$를 갖는
기약 닫힌 부분집합이라 하자. $\mathcal{P}$를 $X$ 위 연접층들의 한
성질이라 하고, 다음을 가정하자.
\begin{enumerate}
\item 연접층들의 임의의 짧은 완전열
$$
0 \to \mathcal{F}_1 \to \mathcal{F} \to \mathcal{F}_2 \to 0
$$
에 대해, $\mathcal{F}_i$, $i = 1, 2$가 성질 $\mathcal{P}$를 가지면
$\mathcal{F}$도 이 성질을 갖는다.
\item $\mathcal{P}$가 $\mathcal{F}^{\oplus r}$에 대해 어떤
$r \geq 1$에 대하여 성립하면, 이는 $\mathcal{F}$에 대해서도 성립한다.
\item 모든 정수적 닫힌 부분스킴 $Z \subset Z_0 \subset X$,
$Z \not = Z_0$와 모든 준연접 아이디얼층
$\mathcal{I} \subset \mathcal{O}_Z$에 대해
$\mathcal{P}$가 $(Z \to X)_*\mathcal{I}$에 대해 성립한다.
\item 다음을 만족하는 어떤 연접층 $\mathcal{G}$가 존재한다.
\begin{enumerate}
\item $\text{Supp}(\mathcal{G}) = Z_0$,
\item $\mathcal{G}_\xi$는 $\mathfrak m_\xi$에 의해 소멸되며,
\item 모든 준연접 아이디얼층
$\mathcal{J} \subset \mathcal{O}_X$ 가운데
$\mathcal{J}_\xi = \mathcal{O}_{X, \xi}$를 만족하는 것에 대해,
준연접 부분층 $\mathcal{G}' \subset \mathcal{J}\mathcal{G}$가
존재하여 $\mathcal{G}'_\xi = \mathcal{G}_\xi$이고 성질
$\mathcal{P}$가 $\mathcal{G}'$에 대해 성립한다.
\end{enumerate}
\end{enumerate}
그러면 성질 $\mathcal{P}$는 모든 연접층 $\mathcal{F}$ 가운데
$X$ 위에서 지지집합이 $Z_0$에 포함되는 것에 대해 성립한다.
\end{lemma}

\begin{proof}
연접층 $\mathcal{F}$가 연접 부분층들에 의한 여과
$$
0 = \mathcal{F}_0 \subset \mathcal{F}_1 \subset
\ldots \subset \mathcal{F}_m = \mathcal{F}
$$
를 갖고 각 $\mathcal{F}_i/\mathcal{F}_{i - 1}$이 성질
$\mathcal{P}$를 가지면, $\mathcal{F}$도 이 성질을 가짐을 주목하자.
이는 가정 (1)에서 따른다.

\medskip\noindent
첫 응용으로, 지지집합이 $Z_0$에 진정으로 포함되는 모든 연접층이
성질 $\mathcal{P}$를 가짐을 얻는다. 실제로 그러한 층은 여과를 가지며
(Lemma \ref{coherent-lemma-coherent-filter}를 보라), 그 부분몫들은 (3)에 의해
성질 $\mathcal{P}$를 갖는다.

\medskip\noindent
$i : Z_0 \to X$로 닫힌 몰입을 나타내자. $\mathcal{G}$를 (4)에서와
같은 연접층이라 하자. Lemma
\ref{coherent-lemma-prepare-filter-irreducible}에 의해 아이디얼층
$\mathcal{I}$가 $Z_0$ 위에 존재하며, 또한 짧은 완전열
$$
0 \to
i_*\mathcal{I}^{\oplus r} \to
\mathcal{G} \to
\mathcal{Q} \to 0
$$
이 존재하며, $\mathcal{Q}$의 지지집합은 $Z_0$에 진정으로 포함된다.
특히 $r > 0$이고 $\mathcal{I}$는 영이 아닌데, 이는
$\mathcal{G}$의 지지집합이 $Z_0$와 같기 때문이다.
$\mathcal{I}' \subset \mathcal{I}$를 $Z_0$ 위의 임의의 영이 아닌
준연접 아이디얼층으로 택하되 $\mathcal{I}$에 포함되게 하자.
그러면 짧은 완전열
$$
0 \to
i_*(\mathcal{I}')^{\oplus r} \to
\mathcal{G} \to
\mathcal{Q}' \to 0
$$
도 얻으며, $\mathcal{Q}'$의 지지집합은 $Z_0$에 진정으로 포함된다.
$\mathcal{J} \subset \mathcal{O}_X$를 $\mathcal{Q}'$의 지지집합을
잘라내는 준연접 아이디얼층이라 하자(예를 들어 $\mathcal{Q}'$의
지지집합 위의 유도된 축소 닫힌 부분스킴 구조에 대응하는 아이디얼).
그러면 $\mathcal{J}_\xi = \mathcal{O}_{X, \xi}$이다.
Lemma \ref{coherent-lemma-power-ideal-kills-sheaf}에 의해
$\mathcal{J}^n\mathcal{Q}' = 0$인 어떤 $n$이 존재한다. 따라서
$\mathcal{J}^n\mathcal{G} \subset i_*(\mathcal{I}')^{\oplus r}$이다.
보조정리의 가정 (4)(c)에 의해 준연접 부분층
$\mathcal{G}' \subset \mathcal{J}^n\mathcal{G}$가 존재하여
$\mathcal{G}'_\xi = \mathcal{G}_\xi$이고 성질 $\mathcal{P}$를 갖는다.
따라서 짧은 완전열
$$
0 \to \mathcal{G}' \to
i_*(\mathcal{I}')^{\oplus r} \to
\mathcal{Q}'' \to 0
$$
을 얻으며, 여기서 $\mathcal{Q}''$의 지지집합은 $Z_0$에 진정으로
포함된다. 그러므로 처음의 관찰과 보조정리의 성질 (1)에 의해
$i_*(\mathcal{I}')^{\oplus r}$이 $\mathcal{P}$를 만족한다고
결론 내린다. 따라서 (2)에 의해 $i_*\mathcal{I}'$도
$\mathcal{P}$를 만족한다. 마지막으로 임의의 준연접 아이디얼층
$\mathcal{I}'' \subset \mathcal{O}_{Z_0}$에 대해
$\mathcal{I}' = \mathcal{I}'' \cap \mathcal{I}$로 놓으면 짧은 완전열
$$
0 \to
i_*(\mathcal{I}') \to
i_*(\mathcal{I}'') \to
\mathcal{Q}''' \to 0
$$
을 얻으며, $\mathcal{Q}'''$의 지지집합은 $Z_0$에 진정으로 포함된다.
따라서 성질 $\mathcal{P}$는 $i_*\mathcal{I}''$에 대해서도 성립한다.

\medskip\noindent
마지막으로 임의의 연접층 $\mathcal{F}$가 $X$ 위에 있고 그
지지집합이 $Z_0$에 포함되면 이 층은 여과를 가지며(Lemma
\ref{coherent-lemma-coherent-filter}를
다시 보라), 그 모든 부분몫은 방금 보인 바에 의해 성질
$\mathcal{P}$를 가짐을 주목하면 증명이 끝난다.
\end{proof}

\begin{lemma}
\label{coherent-lemma-property-higher-rank-cohomological}
$X$를 뇌터 스킴이라 하자. $\mathcal{P}$를 $X$ 위 연접층들의 한
성질이라 하고, 다음을 가정하자.
\begin{enumerate}
\item 연접층들의 임의의 짧은 완전열
$$
0 \to \mathcal{F}_1 \to \mathcal{F} \to \mathcal{F}_2 \to 0
$$
에 대해, $\mathcal{F}_i$, $i = 1, 2$가 성질 $\mathcal{P}$를 가지면
$\mathcal{F}$도 이 성질을 갖는다.
\item $\mathcal{P}$가 $\mathcal{F}^{\oplus r}$에 대해 어떤
$r \geq 1$에 대하여 성립하면, 이는 $\mathcal{F}$에 대해서도 성립한다.
\item 모든 정수적 닫힌 부분스킴 $Z \subset X$와 그 일반점 $\xi$에
대해 다음을 만족하는 어떤 연접층 $\mathcal{G}$가 존재한다.
\begin{enumerate}
\item $\text{Supp}(\mathcal{G}) = Z$,
\item $\mathcal{G}_\xi$는 $\mathfrak m_\xi$에 의해 소멸되며,
\item 모든 준연접 아이디얼층
$\mathcal{J} \subset \mathcal{O}_X$ 가운데
$\mathcal{J}_\xi = \mathcal{O}_{X, \xi}$를 만족하는 것에 대해,
준연접 부분층 $\mathcal{G}' \subset \mathcal{J}\mathcal{G}$가
존재하여 $\mathcal{G}'_\xi = \mathcal{G}_\xi$이고 성질
$\mathcal{P}$가 $\mathcal{G}'$에 대해 성립한다.
\end{enumerate}
\end{enumerate}
그러면 성질 $\mathcal{P}$는 $X$ 위의 모든 연접층에 대해 성립한다.
\end{lemma}

\begin{proof}
Lemma \ref{coherent-lemma-property-irreducible-higher-rank-cohomological}에서
Lemma \ref{coherent-lemma-property}가 Lemma
\ref{coherent-lemma-property-irreducible}에서 따라오는 것과 정확히 같은
방식으로 따라온다.
\end{proof}

\section{유한 사상과 아핀 스킴}
\label{coherent-section-finite-affine}

\noindent
이 절에서는 앞 절들의 결과를 사용하여 뇌터 아핀 스킴의 유한 사상에
의한 상이 아핀임을 보인다. 나중에 이 결과가 더 일반적으로 성립함을
볼 것이다(Limits, Lemma \ref{limits-lemma-affine} 및 Proposition
\ref{limits-proposition-affine} 참조).

\begin{lemma}
\label{coherent-lemma-finite-morphism-Noetherian}
$f : Y \to X$를 스킴의 사상이라 하자. $f$가 유한이고 전사이며
$X$가 국소 뇌터라고 가정하자. $Z \subset X$를 일반점 $\xi$를 갖는
정수적 닫힌 부분스킴이라 하자. 그러면 연접층 $\mathcal{F}$가
$Y$ 위에 존재하여 $f_*\mathcal{F}$의 지지집합은 $Z$와 같고
$(f_*\mathcal{F})_\xi$는 $\mathfrak m_\xi$에 의해 소멸된다.
\end{lemma}

\begin{proof}
Morphisms, Lemma \ref{morphisms-lemma-finite-type-noetherian}에 의해
$Y$는 국소 뇌터임을 주목하자. $f$가 전사이므로 올 $Y_\xi$는
공집합이 아니다. $\xi' \in Y$를 택하여 $\xi$로 가게 하고
$Z' = \overline{\{\xi'\}}$로 놓자. $Z' \subset Y$를 축소 닫힌
부분스킴으로 생각할 수 있다. Schemes, Lemma
\ref{schemes-lemma-reduced-closed-subscheme}를 보라. 따라서 층
$\mathcal{F} = (Z' \to Y)_*\mathcal{O}_{Z'}$는 $Y$ 위의 연접층이다
(Lemma \ref{coherent-lemma-finite-pushforward-coherent} 참조).
다음 가환 그림을 보자.
$$
\xymatrix{
Z' \ar[r]_{i'} \ar[d]_{f'} &
Y \ar[d]^f \\
Z \ar[r]^i &
X
}
$$
$f_*\mathcal{F} = i_*f'_*\mathcal{O}_{Z'}$임을 알 수 있다. 따라서
$f_*\mathcal{F}$의 $\xi$에서의 싹은
$f'_*\mathcal{O}_{Z'}$의 $\xi$에서의 싹이다. $Z'$는 일반점
$\xi'$를 갖는 정수적 스킴이므로, $\xi'$는 $Z'$의 점 가운데
$\xi$ 위에 놓이는 유일한 점임을 주목하자. Algebra, Lemmas
\ref{algebra-lemma-finite-is-integral} 및
\ref{algebra-lemma-integral-no-inclusion}을 보라. 따라서
$f'_*\mathcal{O}_{Z'}$의 $\xi$에서의 싹은
$\mathcal{O}_{Z', \xi'} = \kappa(\xi')$와 같다. 특히
$f_*\mathcal{F}$의 $\xi$에서의 싹은 영이 아니다. 이 사실과
$f_*\mathcal{F}$가 $i_*f'_*(\text{something})$ 꼴이라는 사실을
결합하면 보조정리가 따른다.
\end{proof}

\begin{lemma}
\label{coherent-lemma-affine-morphism-projection-ideal}
$f : Y \to X$를 스킴의 사상이라 하자. $\mathcal{F}$를 $Y$ 위의
준연접층이라 하고, $\mathcal{I}$를 $X$ 위의 준연접 아이디얼층이라
하자. 사상 $f$가 아핀이면
$\mathcal{I}f_*\mathcal{F} = f_*(f^{-1}\mathcal{I}\mathcal{F})$이다.
\end{lemma}

\begin{proof}
이 표기법의 뜻은 다음과 같다. $f^{-1}$은 완전 함자이므로
$f^{-1}\mathcal{I}$는 $f^{-1}\mathcal{O}_X$의 아이디얼층이다.
사상 $f^\sharp : f^{-1}\mathcal{O}_X \to \mathcal{O}_Y$를 통해
이는 $\mathcal{F}$에 작용한다. 이제
$f^{-1}\mathcal{I}\mathcal{F}$는 $as$ 꼴의 국소 단면들의 합들로
생성되는 부분층이며, 여기서 $a$는 $f^{-1}\mathcal{I}$의 국소 단면이고
$s$는 $\mathcal{F}$의 국소 단면이다. 이는 준연접
$\mathcal{O}_Y$-부분가군으로서 $\mathcal{F}$에 포함되는데, 자연 사상
$f^*\mathcal{I} \otimes_{\mathcal{O}_Y} \mathcal{F} \to \mathcal{F}$의
상이기도 하기 때문이다.

\medskip\noindent
이렇게 표기법을 설명하면 증명은 바로 된다. 문제는 국소적이므로
$X$가 아핀이라고 가정해도 된다. $f$가 아핀이므로 $Y$도 아핀이다.
따라서
$Y = \Spec(B)$, $X = \Spec(A)$, $\mathcal{F} = \widetilde{M}$,
그리고 $\mathcal{I} = \widetilde{I}$로 쓸 수 있다. 이 경우
보조정리의 주장은
$$
I(M_A) = ((IB)M)_A
$$
라는 명제로 환원된다. 여기서 $M_A$는 $A$-가군을 나타내며, 이는
$B$-가군 $M$에 딸린 것이다.
\end{proof}

\begin{lemma}
\label{coherent-lemma-image-affine-finite-morphism-affine-Noetherian}
$f : Y \to X$를 스킴의 사상이라 하고 다음을 가정하자.
\begin{enumerate}
\item $f$는 유한이고,
\item $f$는 전사이며,
\item $Y$는 아핀이고,
\item $X$는 뇌터이다.
\end{enumerate}
그러면 $X$는 아핀이다.
\end{lemma}

\begin{proof}
보조정리의 가정 아래에서 임의의 연접 $\mathcal{O}_X$-가군
$\mathcal{F}$에 대해 $H^1(X, \mathcal{F}) = 0$임을 보이겠다.
특히 $H^1(X, \mathcal{I}) = 0$이고, 이는 $\mathcal{O}_X$의 모든
준연접 아이디얼층에 대해 성립한다. 그러면 Lemma
\ref{coherent-lemma-quasi-compact-h1-zero-covering} 또는 Lemma
\ref{coherent-lemma-quasi-separated-h1-zero-covering}에 의해 $X$가 아핀임이
따른다.

\medskip\noindent
$\mathcal{P}$를 연접층 $\mathcal{F}$에 대해 정의하되, 이 층은
$X$ 위에 있고 규칙은 다음과 같다.
$$
\mathcal{P}(\mathcal{F}) \Leftrightarrow H^1(X, \mathcal{F}) = 0.
$$
Lemma \ref{coherent-lemma-property-higher-rank-cohomological}를 적용하겠다.
따라서 그 보조정리의 (1), (2), (3)이 $\mathcal{P}$에 대해 성립함을
확인해야 한다. 성질 (1)은 층들의 짧은 완전열에 딸린 긴 코호몰로지
완전열에서 따른다. $H^1(X, -)$가 가법 함자이므로 성질 (2)가 따른다.
(3)을 보이기 위해 $Z \subset X$를 일반점 $\xi$를 갖는 정수적 닫힌
부분스킴이라 하자. $\mathcal{F}$를 $Y$ 위의 연접층으로 택하여
$f_*\mathcal{F}$의 지지집합이 $Z$와 같고
$(f_*\mathcal{F})_\xi$가 $\mathfrak m_\xi$에 의해 소멸되게 하자.
Lemma \ref{coherent-lemma-finite-morphism-Noetherian}를 보라.
$\mathcal{G} = f_*\mathcal{F}$로 택하면 된다고 주장한다.
Lemma \ref{coherent-lemma-property-higher-rank-cohomological}의 (3)(c)만
확인하면 된다. 따라서 $\mathcal{J} \subset \mathcal{O}_X$를
$\mathcal{J}_\xi = \mathcal{O}_{X, \xi}$를 만족하는 준연접
아이디얼층이라 하자. 유한 사상은 아핀이므로 Lemma
\ref{coherent-lemma-affine-morphism-projection-ideal}에 의해
$\mathcal{J}\mathcal{G} = f_*(f^{-1}\mathcal{J}\mathcal{F})$이다.
또한 Lemma \ref{coherent-lemma-affine-morphism-projection-ideal}의 증명에서
지적했듯이, 층 $f^{-1}\mathcal{J}\mathcal{F}$는 준연접
$\mathcal{O}_Y$-가군이다. $Y$가 아핀이므로
$H^1(Y, f^{-1}\mathcal{J}\mathcal{F}) = 0$이다. Lemma
\ref{coherent-lemma-quasi-coherent-affine-cohomology-zero}를 보라.
$f$는 유한이고 따라서 아핀이므로
$$
H^1(X, \mathcal{J}\mathcal{G}) =
H^1(X, f_*(f^{-1}\mathcal{J}\mathcal{F})) =
H^1(Y, f^{-1}\mathcal{J}\mathcal{F}) = 0
$$
이다. 이는 Lemma \ref{coherent-lemma-relative-affine-cohomology}에 따른다.
따라서 준연접 부분층
$\mathcal{G}' = \mathcal{J}\mathcal{G}$는 $\mathcal{P}$를 만족한다.
이로써 Lemma \ref{coherent-lemma-property-higher-rank-cohomological}의 성질
(3)(c)가 원하는 대로 확인된다.
\end{proof}

\section{Proj 위의 연접층, I}
\label{coherent-section-coherent-proj}

\noindent
이 절에서는 $\text{Proj}(A)$ 위의 연접층을 논의한다. 여기서
$A$는 $A_1$에 의해 $A_0$ 위에서 생성되는 뇌터 등급환이다.
다음 절에서는 $A$가 차수 $1$인 원소들로 생성되지 않는 경우를
논의한다. 먼저 뇌터 환 위의 사영공간에 관한 결과들을 하나로 모아
서술한다.

\begin{lemma}
\label{coherent-lemma-coherent-projective}
$R$을 뇌터 환이라 하고 $n \geq 0$을 정수라 하자.
모든 연접층 $\mathcal{F}$에 대해, 그것이 $\mathbf{P}^n_R$ 위에 있으면 다음이
성립한다.
\begin{enumerate}
\item 어떤 $r \geq 0$와 $d_1, \ldots, d_r \in \mathbf{Z}$ 및 전사
$$
\bigoplus\nolimits_{j = 1, \ldots, r}
\mathcal{O}_{\mathbf{P}^n_R}(d_j)
\longrightarrow
\mathcal{F}.
$$
가 존재한다.
\item $H^i(\mathbf{P}^n_R, \mathcal{F}) = 0$이며, 예외는
$0 \leq i \leq n$인 경우뿐이다.
\item 임의의 $i$에 대해 코호몰로지 군
$H^i(\mathbf{P}^n_R, \mathcal{F})$는 유한 $R$-가군이다.
\item $i > 0$이면 $H^i(\mathbf{P}^n_R, \mathcal{F}(d)) = 0$은
충분히 큰 모든 $d$에 대해 성립한다.
\item 임의의 $k \in \mathbf{Z}$에 대해 등급
$R[T_0, \ldots, T_n]$-가군
$$
\bigoplus\nolimits_{d \geq k} H^0(\mathbf{P}^n_R, \mathcal{F}(d))
$$
은 유한 $R[T_0, \ldots, T_n]$-가군이다.
\end{enumerate}
\end{lemma}

\begin{proof}
$\mathcal{O}_{\mathbf{P}^n_R}(1)$이 스킴
$\mathbf{P}^n_R$ 위의 풍부한 가역층임을 사용하겠다. 이는 정의에서
직접 따른다. 실제로 $\mathbf{P}^n_R$는 표준 아핀 열린집합
$D_{+}(T_i)$들로 덮인다. 따라서 Properties, Proposition
\ref{properties-proposition-characterize-ample}에 의해 모든 유한형
준연접 $\mathcal{O}_{\mathbf{P}^n_R}$-가군은
$\mathcal{O}_{\mathbf{P}^n_R}(1)$의 텐서 거듭제곱들의 유한 직합의
몫이다. 한편 $R$ 위의 사영공간에서는 Lemma
\ref{coherent-lemma-coherent-Noetherian}에 의해 연접층과 유한형 준연접층이
같다. 따라서 (1)을 얻는다.

\medskip\noindent
사영 $n$-공간 $\mathbf{P}^n_R$는 $n + 1$개의 아핀, 즉 표준
열린집합 $D_{+}(T_i)$, $i = 0, \ldots, n$으로 덮인다. Constructions,
Lemma \ref{constructions-lemma-standard-covering-projective-space}를
보라. 따라서 임의의 준연접층 $\mathcal{F}$에 대해 $\mathbf{P}^n_R$
위에서 $H^i(\mathbf{P}^n_R, \mathcal{F}) = 0$이며, 이는
$i \geq n + 1$일 때 성립한다. Lemma
\ref{coherent-lemma-vanishing-nr-affines}를 보라. 그러므로 (2)가 성립한다.

\medskip\noindent
$\mathbf{P}^n_R$ 위의 모든 연접층에 대해 $i$에 대한 내림 귀납법으로
(3)과 (4)를 동시에 증명하자. (2)에 의해 $i \geq n + 1$이면 결과가
분명히 성립한다. $i + 1$에 대한 결과를 알고 있고 $i$에 대한 결과를
보이려 한다고 하자. ($i = 0$이면 (4)는 공허하다.)
$\mathcal{F}$를 $\mathbf{P}^n_R$ 위의 연접층이라 하자. (1)에서와
같은 전사를 택하고 그 핵을 $\mathcal{G}$로 나타내면 짧은 완전열
$$
0 \to \mathcal{G} \to
\bigoplus\nolimits_{j = 1, \ldots, r}
\mathcal{O}_{\mathbf{P}^n_R}(d_j)
\to
\mathcal{F} \to 0
$$
을 얻는다. Lemma \ref{coherent-lemma-coherent-abelian-Noetherian}에 의해
$\mathcal{G}$는 연접이다. 긴 코호몰로지 완전열은 완전열
$$
H^i(\mathbf{P}^n_R, \bigoplus\nolimits_{j = 1, \ldots, r}
\mathcal{O}_{\mathbf{P}^n_R}(d_j))
\to
H^i(\mathbf{P}^n_R, \mathcal{F})
\to
H^{i + 1}(\mathbf{P}^n_R, \mathcal{G}).
$$
을 준다. 귀납 가정에 의해 오른쪽 $R$-가군은 유한이고, Lemma
\ref{coherent-lemma-cohomology-projective-space-over-ring}에 의해 왼쪽
$R$-가군도 유한이다. $R$이 뇌터이므로
$H^i(\mathbf{P}^n_R, \mathcal{F})$는 유한 $R$-가군이다.
이로써 주장 (3)의 귀납 단계를 증명했다.
$\mathcal{O}_{\mathbf{P}^n_R}(d)$는 가역이므로
$\mathbf{P}^n_R$ 위에서 꼬임을 취하는 것은 완전 함자이다
(가역층과 텐서곱하여 얻기 때문이다. Constructions, Definition
\ref{constructions-definition-twist}를 보라). 이는 모든
$d \in \mathbf{Z}$에 대해 수열
$$
0 \to \mathcal{G}(d) \to
\bigoplus\nolimits_{j = 1, \ldots, r}
\mathcal{O}_{\mathbf{P}^n_R}(d_j + d)
\to
\mathcal{F}(d) \to 0
$$
이 짧은 완전열이라는 뜻이다. 그 결과로 얻는 코호몰로지 수열은
$$
H^i(\mathbf{P}^n_R, \bigoplus\nolimits_{j = 1, \ldots, r}
\mathcal{O}_{\mathbf{P}^n_R}(d_j + d))
\to
H^i(\mathbf{P}^n_R, \mathcal{F}(d))
\to
H^{i + 1}(\mathbf{P}^n_R, \mathcal{G}(d)).
$$
이다. 귀납 가정에 의해 오른쪽 가군은 $d \gg 0$일 때 영이고,
Lemma \ref{coherent-lemma-cohomology-projective-space-over-ring}의 계산에
의해 왼쪽 가군은 $d \geq -\min\{d_j\}$이고 $i \geq 1$인 순간부터
영이다. 따라서 주장 (4)의 귀납 단계가 성립한다. 이로써 (3)과 (4)의
증명이 끝난다.

\medskip\noindent
(5)를 증명하기 위해, 모든 충분히 큰 $d$에 대해 사상
$$
H^0(\mathbf{P}^n_R, \bigoplus\nolimits_{j = 1, \ldots, r}
\mathcal{O}_{\mathbf{P}^n_R}(d_j + d))
\to
H^0(\mathbf{P}^n_R, \mathcal{F}(d))
$$
이 전사임을 주목하자. 이는 방금 증명한
$H^1(\mathbf{P}^n_R, \mathcal{G}(d))$의 소멸에 따른다.
다시 말해 가군
$$
M_k
=
\bigoplus\nolimits_{d \geq k} H^0(\mathbf{P}^n_R, \mathcal{F}(d))
$$
은 충분히 큰 $k$에 대해 대응하는 가군
$$
N_k
=
\bigoplus\nolimits_{d \geq k} H^0(\mathbf{P}^n_R,
\bigoplus\nolimits_{j = 1, \ldots, r}
\mathcal{O}_{\mathbf{P}^n_R}(d_j + d)
)
$$
의 몫이다. $k$가 충분히 작으면(예를 들어 $k < -d_j$가 모든
$j$에 대해 성립하면)
$$
N_k = \bigoplus\nolimits_{j = 1, \ldots, r}
R[T_0, \ldots, T_n](d_j)
$$
이다. 이는 절 \ref{coherent-section-cohomology-projective-space}의 계산에
따른다. 특히 이는 유한 생성이다. 이제 $k \in \mathbf{Z}$가
임의라고 하자. $k_{-} \ll k \ll k_{+}$를 택하고 그림
$$
\xymatrix{
N_{k_{-}} & N_{k_{+}} \ar[d] \ar[l] \\
M_k & M_{k_{+}} \ar[l]
}
$$
을 생각하자. 여기서 수직 화살표는 위의 전사이고 수평 화살표들은
명백한 포함 사상들이다. 위에서 말한 바에 의해 $N_{k_{-}}$는 유한
생성 $R[T_0, \ldots, T_n]$-가군이다. 따라서 $N_{k_{+}}$도 유한 생성
$R[T_0, \ldots, T_n]$-가군이다. 실제로 이는 유한 생성 가군의
부분가군이고 환 $R[T_0, \ldots, T_n]$은 뇌터이다. 수직 화살표가
전사이므로 $M_{k_{+}}$는 유한 생성
$R[T_0, \ldots, T_n]$-가군이다. 몫 $M_k/M_{k_{+}}$는 유한
$R$-가군이다. 이는 유한개의 유한 $R$-가군인
$H^0(\mathbf{P}^n_R, \mathcal{F}(d))$ ($k \leq d < k_{+}$)의
직합이기 때문이다. 여기서
$i = 0$인 경우의 (3)을 사용함을 주목하자. 따라서
$M_k/M_{k_{+}}$는 더구나 유한 $R[T_0, \ldots, T_n]$-가군이다.
다시 말해 $M_k$를 두 유한 $R[T_0, \ldots, T_n]$-가군 사이에
끼워 넣었으므로 결론을 얻는다.
\end{proof}

\begin{lemma}
\label{coherent-lemma-coherent-on-proj}
$A$를 등급환이라 하고, $A_0$는 뇌터이며 $A$는 $A_1$의 유한개의
원소로 $A_0$ 위에서 생성된다고 하자.
$X = \text{Proj}(A)$로 놓자. 그러면 $X$는 뇌터 스킴이다.
$\mathcal{F}$를 연접 $\mathcal{O}_X$-가군이라 하자.
\begin{enumerate}
\item 어떤 $r \geq 0$와 $d_1, \ldots, d_r \in \mathbf{Z}$ 및 전사
$$
\bigoplus\nolimits_{j = 1, \ldots, r} \mathcal{O}_X(d_j)
\longrightarrow \mathcal{F}.
$$
가 존재한다.
\item 임의의 $p$에 대해 코호몰로지 군 $H^p(X, \mathcal{F})$는
유한 $A_0$-가군이다.
\item $p > 0$이면 $H^p(X, \mathcal{F}(d)) = 0$은 충분히 큰 모든
$d$에 대해 성립한다.
\item 임의의 $k \in \mathbf{Z}$에 대해 등급 $A$-가군
$$
\bigoplus\nolimits_{d \geq k} H^0(X, \mathcal{F}(d))
$$
은 유한 $A$-가군이다.
\end{enumerate}
\end{lemma}

\begin{proof}
가정에 의해 등급 $A_0$-대수의 전사
$$
A_0[T_0, \ldots, T_n] \longrightarrow A
$$
가 존재하며, 여기서 $\deg(T_j) = 1$이고
$j = 0, \ldots, n$이다. Constructions, Lemma
\ref{constructions-lemma-surjective-graded-rings-generated-degree-1-map-proj}
에 의해 이는 닫힌 몰입 $i : X \to \mathbf{P}^n_{A_0}$를 정의하며
$i^*\mathcal{O}_{\mathbf{P}^n_{A_0}}(1) = \mathcal{O}_X(1)$이다.
특히 $X$는 뇌터 스킴 $\mathbf{P}^n_{A_0}$의 닫힌 부분스킴이므로
뇌터이다. $\mathcal{F}$에 대한 보조정리의 결과들은 연접층
$i_*\mathcal{F}$에 대한 Lemma \ref{coherent-lemma-coherent-projective}의
대응하는 결과들에서 따른다고 주장한다. 여기서 Lemma
\ref{coherent-lemma-i-star-equivalence}를 사용하며 이 층은
$\mathbf{P}^n_{A_0}$ 위에 있다. 예를 들어 이 보조정리에 의해 전사
$$
\bigoplus\nolimits_{j = 1, \ldots, r} \mathcal{O}_{\mathbf{P}^n_{A_0}}(d_j)
\longrightarrow i_*\mathcal{F}.
$$
가 존재한다. 수반성에 의해 이는 사상
$\bigoplus_{j = 1, \ldots, r} \mathcal{O}_X(d_j) \longrightarrow \mathcal{F}$
에 대응하며, 이 사상 역시 전사이다. 코호몰로지에 관한 명제들은
Lemma \ref{coherent-lemma-relative-affine-cohomology}에 의해
$H^p(X, \mathcal{F}(d)) = H^p(\mathbf{P}^n_{A_0}, i_*\mathcal{F}(d))$
인 사실에서 따른다.
\end{proof}

\begin{lemma}
\label{coherent-lemma-recover-tail-graded-module}
$A$를 등급환이라 하고, $A_0$는 뇌터이며 $A$는 $A_1$의 유한개의
원소로 $A_0$ 위에서 생성된다고 하자. $M$을 유한 등급
$A$-가군이라 하자. $X = \text{Proj}(A)$로 놓고, $\widetilde{M}$을
준연접 $\mathcal{O}_X$-가군이라 하자. 이 가군은 $X$ 위에 있고
$M$에 딸린다.
Constructions, Lemma \ref{constructions-lemma-apply-modules}의 사상들
$$
M_n \longrightarrow \Gamma(X, \widetilde{M}(n))
$$
은 모든 충분히 큰 $n$에 대해 동형이다.
\end{lemma}

\begin{proof}
$M$은 유한 $A$-가군이므로 $\widetilde{M}$은 유한형
$\mathcal{O}_X$-가군, 즉 연접 $\mathcal{O}_X$-가군이다.
$N = \bigoplus_{n \in \mathbf{Z}} \Gamma(X, \widetilde{M}(n))$로 놓자.
사상 $M \to N$이 등급 $A$-가군의 사상으로서 모든 충분히 큰 차수에서
동형임을 보여야 한다. Properties, Lemma
\ref{properties-lemma-proj-quasi-coherent}에 의해 표준 동형
$\widetilde{N} \to \widetilde{M}$이 있고, 유도된 사상
$N_n \to N_n = \Gamma(X, \widetilde{M}(n))$은 항등 사상이다.
따라서 사상 $\widetilde{M} \to \widetilde{N} \to \widetilde{M}$이
있으며, 모든 $n$에 대해 그림
$$
\xymatrix{
M_n \ar[d] \ar[r] & N_n \ar[d] \ar@{=}[rd] \\
\Gamma(X, \widetilde{M}(n)) \ar[r] &
\Gamma(X, \widetilde{N}(n)) \ar[r]^{\cong} &
\Gamma(X, \widetilde{M}(n))
}
$$
은 가환한다. 이는 합성
$$
M_n \to \Gamma(X, \widetilde{M}(n))
\to \Gamma(X, \widetilde{N}(n)) \to \Gamma(X, \widetilde{M}(n))
$$
이 표준 사상 $M_n \to \Gamma(X, \widetilde{M}(n))$과 같다는 뜻이다.
분명히 이는 합성
$\widetilde{M} \to \widetilde{N} \to \widetilde{M}$이 항등임을
함의한다. 따라서 $\widetilde{M} \to \widetilde{N}$은 동형이다.
$K = \Ker(M \to N)$ 및 $Q = \Coker(M \to N)$로 놓자. 함자
$M \mapsto \widetilde{M}$은 완전임을 상기하자. Constructions, Lemma
\ref{constructions-lemma-proj-sheaves}를 보라. 따라서
$\widetilde{K} = 0$이고 $\widetilde{Q} = 0$이다.
$A$는 뇌터 환이고 $M$은 유한 생성 $A$-가군이며, $N$은 등급
$A$-가군으로서
$N' = \bigoplus_{n \geq 0} N_n$은 Lemma
\ref{coherent-lemma-coherent-on-proj}의 마지막 부분에 의해 유한 생성이다.
따라서 $K' = \bigoplus_{n \geq 0} K_n$과
$Q' = \bigoplus_{n \geq 0} Q_n$은 유한 $A$-가군이다.
$\widetilde{Q} = \widetilde{Q'}$이고
$\widetilde{K} = \widetilde{K'}$임을 관찰하자. 그러므로 증명을
끝내려면 유한 $A$-가군 $K$가 $\widetilde{K} = 0$이고 영이 아닌
동차 성분 $K_d$를 $d \geq 0$에 대해 유한개만 가짐을
보이면 충분하다. 이를 위해 $x_1, \ldots, x_r \in K$를 동차 생성원으로
택하고 그 차수를 $d_1, \ldots, d_r$라 하자.
$f_1, \ldots, f_n \in A_1$을 $A$를 $A_0$ 위에서 생성하는 원소들이라
하자. 각 $i$와 $j$에 대해 어떤 $n_{ij} \geq 0$가 존재하여
$f_i^{n_{ij}} x_j = 0$이고 이는 $K_{d_j + n_{ij}}$ 안의 등식이다. 그렇지 않으면
$x_i/f_i^{d_i} \in K_{(f_i)}$가 영이 아니어서 $\widetilde{K}$가
영이 아니게 된다. 이제 $K_d$는 $d > \max_j(d_j + \sum_i n_{ij})$이면
영이다. 실제로 $K_d$의 모든 원소는 항들의 합이며, 각 항은
$f_i$들의 단항식과 $x_j$ 중 하나의 곱이고 전체 차수는 $d$이다.
\end{proof}

\noindent
$A$를 등급환이라 하고, $A_0$는 뇌터이며 $A$는 $A_1$의 유한개의
원소로 $A_0$ 위에서 생성된다고 하자.
$A_+ = \bigoplus_{n > 0} A_n$이 무관 아이디얼임을 상기하자.
$M$을 등급 $A$-가군이라 하자. $M$이 $A_+$-거듭제곱 꼬임 가군이라는
것은, 모든 $x \in M$에 대해 어떤 $n \geq 1$이 존재하여
$(A_+)^n x = 0$이라는 뜻이다. More on Algebra, Definition
\ref{more-algebra-definition-f-power-torsion}을 보라. $M$이 유한
생성이면 이는 $M_n = 0$이 $n \gg 0$에 대해 성립하는 것과 동치이다.
때로는 $A_+$-거듭제곱 꼬임 가군을 꼬임 가군이라고 부른다.
때로는 등급 $A$-가군 $M$에 대해, $x \in M$이고
$(A_+)^n x = 0$, $n > 0$이면 $x = 0$일 때 꼬임이 없다고 한다.
$\text{Mod}_A$로 등급 $A$-가군들의 범주를,
$\text{Mod}^{fg}_A$로 유한 생성인 것들의 충만한 부분범주를,
$\text{Mod}^{fg}_{A, torsion}$으로 가군 $M$ 중 $M_n = 0$이
$n \gg 0$에 대해 성립하는 것들의 충만한 부분범주를 나타낸다.

\begin{proposition}
\label{coherent-proposition-coherent-modules-on-proj}
$A$를 등급환이라 하고, $A_0$는 뇌터이며 $A$는 $A_1$의 유한개의
원소로 $A_0$ 위에서 생성된다고 하자.
$X = \text{Proj}(A)$로 놓자. 함자 $M \mapsto \widetilde M$은 동치
$$
\text{Mod}^{fg}_A/\text{Mod}^{fg}_{A, torsion}
\longrightarrow
\textit{Coh}(\mathcal{O}_X)
$$
를 유도하며, 그 준역은
$\mathcal{F} \longmapsto \bigoplus_{n \geq 0} \Gamma(X, \mathcal{F}(n))$
으로 주어진다.
\end{proposition}

\begin{proof}
부분범주 $\text{Mod}^{fg}_{A, torsion}$은
$\text{Mod}^{fg}_A$의 세르 부분범주이다. Homology, Definition
\ref{homology-definition-serre-subcategory}를 보라. 이는 위에서
제시한 대상들의 기술로부터 분명하지만, More on Algebra, Lemma
\ref{more-algebra-lemma-I-power-torsion}에서도 따른다. 따라서 화살표
왼쪽의 몫 범주는 Homology, Lemma
\ref{homology-lemma-serre-subcategory-is-kernel}에서 정의된다.
명제의 함자를 정의하려면 함자 $M \mapsto \widetilde M$이 꼬임
가군들을 $0$으로 보냄을 보이면 충분하다. 이는 분명하다. 실제로 임의의
동차 원소 $f \in A_+$에 대해 가군 $M_f$는 영이고, 따라서 값
$M_{(f)}$도 영이다. 이것은 $\widetilde M$의 $D_+(f)$에서의 값이다.

\medskip\noindent
Lemma \ref{coherent-lemma-coherent-on-proj}에 의해 제안된 준역은 잘 정의된다.
실제로 이 보조정리는
$\mathcal{F} \longmapsto \bigoplus_{n \geq 0} \Gamma(X, \mathcal{F}(n))$
이 함자 $\textit{Coh}(\mathcal{O}_X) \to \text{Mod}^{fg}_A$임을 보인다.
이를 몫 함자
$\text{Mod}^{fg}_A \to \text{Mod}^{fg}_A/\text{Mod}^{fg}_{A, torsion}$
와 합성할 수 있다.

\medskip\noindent
Lemma \ref{coherent-lemma-recover-tail-graded-module}에 의해 왼쪽에서 오른쪽을
거쳐 다시 왼쪽으로 가는 합성은 항등 함자와 동형이다.

\medskip\noindent
마지막으로 $\mathcal{F}$를 연접 $\mathcal{O}_X$-가군이라 하자.
$M = \bigoplus_{n \in \mathbf{Z}} \Gamma(X, \mathcal{F}(n))$로 놓고 이를
등급 $A$-가군으로 보자. 그러면 우리의 함자는 $\mathcal{F}$를
$M_{\geq 0} = \bigoplus_{n \geq 0} M_n$으로 보낸다.
Properties, Lemma \ref{properties-lemma-proj-quasi-coherent}에 의해
표준 사상 $\widetilde M \to \mathcal{F}$는 동형이다. 포함 사상
$M_{\geq 0} \to M$은 동형
$\widetilde{M_{\geq 0}} \to \widetilde M$을 정의하므로, 오른쪽에서
왼쪽을 거쳐 다시 오른쪽으로 가는 합성도 항등 함자와 동형이라고
결론 내린다.
\end{proof}

\section{Proj 위의 연접층, II}
\label{coherent-section-coherent-proj-general}

\noindent
이 절에서는 $\text{Proj}(A)$ 위의 연접층을 논의한다. 여기서 $A$는
뇌터 등급환이다. 결과 대부분은 앞 절에서 $A$가 차수 $1$인 원소들로
생성될 때 얻은 대응하는 결과로부터 약간의 기교를 써서 도출한다.

\begin{lemma}
\label{coherent-lemma-coherent-on-proj-general}
$A$를 뇌터 등급환이라 하자. $X = \text{Proj}(A)$로 놓자. 그러면
$X$는 뇌터 스킴이다. $\mathcal{F}$를 연접 $\mathcal{O}_X$-가군이라 하자.
\begin{enumerate}
\item 어떤 $r \geq 0$와 $d_1, \ldots, d_r \in \mathbf{Z}$ 및 전사
$$
\bigoplus\nolimits_{j = 1, \ldots, r} \mathcal{O}_X(d_j)
\longrightarrow \mathcal{F}.
$$
가 존재한다.
\item 임의의 $p$에 대해 코호몰로지 군 $H^p(X, \mathcal{F})$는 유한
$A_0$-가군이다.
\item $p > 0$이면 $H^p(X, \mathcal{F}(d)) = 0$은 충분히 큰 모든
$d$에 대해 성립한다.
\item 임의의 $k \in \mathbf{Z}$에 대해 등급 $A$-가군
$$
\bigoplus\nolimits_{d \geq k} H^0(X, \mathcal{F}(d))
$$
은 유한 $A$-가군이다.
\end{enumerate}
\end{lemma}

\begin{proof}
이를 Lemma \ref{coherent-lemma-coherent-on-proj}로 환원하여 증명하겠다.
Algebra, Lemmas \ref{algebra-lemma-graded-Noetherian}와
\ref{algebra-lemma-S-plus-generated}에 의해 환 $A_0$는 뇌터이고,
$A$는 $A_0$ 위에서 양의 차수를 갖는 유한개의 동차 원소
$f_1, \ldots, f_r$로 생성된다.
$d$를 충분히 나누어지는 정수라 하자. Algebra, Section
\ref{algebra-section-graded}의 표기법으로 $A' = A^{(d)}$라 놓는다.
그러면 $A'$은 $A'_0 = A_0$ 위에서 차수 $1$인 원소들로 생성된다.
Algebra, Lemma \ref{algebra-lemma-uple-generated-degree-1}를 보라.
따라서 Lemma \ref{coherent-lemma-coherent-on-proj}를
$X' = \text{Proj}(A')$에 적용할 수 있다.

\medskip\noindent
Constructions, Lemma \ref{constructions-lemma-d-uple}에 의해 스킴의
동형 $i : X \to X'$와 동형들
$\mathcal{O}_X(nd) \to i^*\mathcal{O}_{X'}(n)$이 존재하며, 이들은
사상 $A' \to A$ 및 사상들
$A_n \to H^0(X, \mathcal{O}_X(n))$과
$A'_n \to H^0(X', \mathcal{O}_{X'}(n))$과 양립한다.
따라서 Lemma \ref{coherent-lemma-coherent-on-proj}는 $X$가 뇌터이고 (1)과
(2)가 성립함을 함의한다. (3)과 (4)를 보이기 위해, 고정된 임의의
$k$, $p$, $q$에 대해 위 양립성으로부터
$$
\bigoplus\nolimits_{dn + q \geq k} H^p(X, \mathcal{F}(dn + q)) =
\bigoplus\nolimits_{dn + q \geq k} H^p(X', (i_*\mathcal{F}(q))(n)
$$
임을 사용할 수 있다. $p > 0$이면 Lemma
\ref{coherent-lemma-coherent-on-proj}에 의해 $k$가 $q$에 의존하여 충분히 크면
대해 우변이 소멸한다. 정수의 $d$를 법으로 하는 합동류는 유한개뿐이므로
$\mathcal{F}$에 대해 $X$ 위에서 (3)이 성립한다.
$p = 0$이면 Lemma \ref{coherent-lemma-coherent-on-proj}에 의해 우변은 유한
$A'$-가군이다. 합동류의 유한성을 다시 사용하면
$\bigoplus_{n \geq k} H^0(X, \mathcal{F}(n))$도 유한 $A'$-가군임을
얻는다. $A'$-가군 구조는 $A$-가군 구조에서 오므로(위에서 말한
양립성에 의해), 이것은 유한 $A$-가군이기도 하다.
\end{proof}

\begin{lemma}
\label{coherent-lemma-recover-tail-graded-module-general}
$A$를 뇌터 등급환이라 하고, $d$를 $A$가 $A_0$ 위에서 갖는
생성원들의 최소공배수라 하자. $M$을 유한 등급 $A$-가군이라 하자.
$X = \text{Proj}(A)$로 놓고, $\widetilde{M}$을 준연접
$\mathcal{O}_X$-가군이라 하자. 이 가군은 $X$ 위에 있고 $M$에 딸린다.
$k \in \mathbf{Z}$라 하자.
\begin{enumerate}
\item $N' = \bigoplus_{n \geq k} H^0(X, \widetilde{M(n)})$은
유한 $A$-가군이다.
\item $N = \bigoplus_{n \geq k} H^0(X, \widetilde{M}(n))$은
유한 $A$-가군이다.
\item 표준 사상 $N \to N'$이 있다.
\item $k$가 충분히 작으면 표준 사상 $M \to N'$이 있다.
\item 사상 $M_n \to N'_n$은 $n \gg 0$에 대해 동형이다.
\item $N_n \to N'_n$은 $d | n$일 때 동형이다.
\end{enumerate}
\end{lemma}

\begin{proof}
(3)의 사상 $N \to N'$은 Constructions, Equation
(\ref{constructions-equation-multiply-more-generally})에서
대역 절단을 취하여 얻는다.

\medskip\noindent
Constructions, Equation
(\ref{constructions-equation-global-sections-more-generally})에 의해
등급 $A$-가군의 사상
$M \to \bigoplus_{n \in \mathbf{Z}} H^0(X, \widetilde{M(n)})$이 있다.
$M$의 생성원들이 차수 $\geq k$에 놓이면 그 상은 부분가군
$N' \subset \bigoplus_{n \in \mathbf{Z}} H^0(X, \widetilde{M(n)})$에
포함되며, (4)의 사상을 얻는다.

\medskip\noindent
Algebra, Lemmas \ref{algebra-lemma-graded-Noetherian}와
\ref{algebra-lemma-S-plus-generated}에 의해 환 $A_0$는 뇌터이고,
$A$는 $A_0$ 위에서 양의 차수를 갖는 유한개의 동차 원소
$f_1, \ldots, f_r$로 생성된다.
$d = \text{lcm}(\deg(f_i))$로 놓는다. 그러면 예를 들어
Constructions, Lemma \ref{constructions-lemma-where-invertible}에
의해 (6)이 성립한다.

\medskip\noindent
$M$은 유한 $A$-가군이므로 $\widetilde{M}$은 유한형
$\mathcal{O}_X$-가군, 즉 연접 $\mathcal{O}_X$-가군이다.
따라서 (2)는 Lemma \ref{coherent-lemma-coherent-on-proj-general}에서 따른다.

\medskip\noindent
한 가지 기교를 써서 (2)로부터 (1)을 도출하겠다.
$q \in \{0, \ldots, d - 1\}$에 대해
$$
{}^qN = \bigoplus\nolimits_{n + q \geq k} H^0(X, \widetilde{M(q)}(n))
$$
로 쓴다. (2)에 의해 이들은 유한 $A$-가군이다. 뇌터 환 $A$는
$A^{(d)} = \bigoplus_{n \geq 0} A_{dn}$ 위에서 유한이다. 실제로
이 환은 $f_i$들로 $A^{(d)}$ 위에서 생성되고 $f_i^d \in A^{(d)}$이다.
따라서 ${}^qN$은 유한 $A^{(d)}$-가군이다. 더욱이 $A^{(d)}$는
뇌터이다(Algebra, Lemma
\ref{algebra-lemma-dehomogenize-finite-type}에서 따른다).
따라서 $A^{(d)}$-부분가군
${}^qN^{(d)} = \bigoplus_{n \in \mathbf{Z}} {}^qN_{dn}$은
$A^{(d)}$ 위의 유한 가군이다. 동형들
$\widetilde{M(dn + q)} = \widetilde{M(q)}(dn)$을 사용하면
$$
N' =
\bigoplus\nolimits_{q \in \{0, \ldots, d - 1\}}
\bigoplus\nolimits_{dn + q \geq k}
H^0(X, \widetilde{M(q)}(dn)) =
\bigoplus\nolimits_{q \in \{0, \ldots, d - 1\}} {}^qN^{(d)}
$$
으로 쓸 수 있다. 따라서 $N'$은 $A^{(d)}$ 위에서 유한이고, 더구나
$A$ 위에서도 유한이다. 그러므로 (1)이 참이다.

\medskip\noindent
$K$를 유한 $A$-가군이라 하고 $\widetilde{K} = 0$이라 하자.
$K_n = 0$이 $d|n$이고 $n \gg 0$일 때 성립한다고 주장한다.
위와 같이 논증하면 $K^{(d)}$는 유한 $A^{(d)}$-가군이다.
$x_1, \ldots, x_m \in K$를 $K^{(d)}$의 동차 생성원으로서
$A^{(d)}$ 위에서 택하고, 이들이 차수 $d_1, \ldots, d_m$에 놓이며
$d | d_j$라고 하자. 각 $i$와 $j$에 대해 어떤
$n_{ij} \geq 0$가 존재하여 $f_i^{n_{ij}} x_j = 0$이고, 이는
$K_{d_j + n_{ij}}$ 안의 등식이다. 그렇지 않으면
$x_j/f_i^{d_i/\deg(f_i)} \in K_{(f_i)}$가 영이 아니어서
$\widetilde{K}$가 영이 아니게 된다. 여기서
$\deg(f_i) | d | d_j$임을 모든 $i, j$에 대해 사용한다. 따라서
$K_n$은 $n$이 $d | n$ 및
$n > \max_j (d_j + \sum_i n_{ij} \deg(f_i))$을 만족할 때 영이다.
실제로 $K_n$의 모든 원소는 항들의 합이며, 각 항은 $f_i$들의
단항식과 $x_j$ 중 하나의 곱이고 전체 차수는 $n$이다.

\medskip\noindent
증명을 끝내려면 $M \to N'$이 모든 충분히 큰 차수에서 동형임을
보여야 한다. 사상 $N \to N'$은 동형
$\widetilde{N} \to \widetilde{N'}$을 유도한다. 실제로 아핀 열린집합
$D_+(f_i) = D_+(f_i^d)$ 위에서 대응하는 가군들은 (6)에 의해 동형이다:
$N_{(f_i)} \cong N_{(f_i^d)} \cong N'_{(f_i^d)} \cong N'_{(f_i)}$.
Properties, Lemma \ref{properties-lemma-proj-quasi-coherent}에 의해
표준 동형 $\widetilde{N} \to \widetilde{M}$이 있다.
합성 $\widetilde{N} \to \widetilde{M} \to \widetilde{N'}$은 위의
동형이다(증명은 생략한다. 힌트: 표준 아핀 열린집합에서 확인하라).
따라서 사상 $M \to N'$은 동형
$\widetilde{M} \to \widetilde{N'}$을 유도한다.
$K = \Ker(M \to N')$ 및 $Q = \Coker(M \to N')$로 놓자.
함자 $M \mapsto \widetilde{M}$은 완전임을 상기하자. Constructions,
Lemma \ref{constructions-lemma-proj-sheaves}를 보라. 따라서
$\widetilde{K} = 0$이고 $\widetilde{Q} = 0$이다.
앞 문단의 결과에 의해 $K_n = 0$이고 $Q_n = 0$이며, 이는
$d | n$이고 $n \gg 0$일 때 성립한다. 여기서 마침내 $N'$을
$N$보다 사용한 이점이 드러난다. 함자 $M \leadsto N'$은 이동과
양립한다(구성에서 곧바로 알 수 있다). 따라서 전체 논증을
$M$ 대신 $M(q)$에 대해 반복하면 $K_n = 0$이고 $Q_n = 0$이며,
이는 $n \equiv q \bmod d$이고 $n \gg 0$일 때 성립한다.
$n$을 법으로 하는 합동류는 유한개뿐이므로 증명이 끝난다.
\end{proof}

\noindent
$A$를 뇌터 등급환이라 하자.
$A_+ = \bigoplus_{n > 0} A_n$이 무관 아이디얼임을 상기하자.
Algebra, Lemmas \ref{algebra-lemma-graded-Noetherian}와
\ref{algebra-lemma-S-plus-generated}에 의해 환 $A_0$는 뇌터이고,
$A$는 $A_0$ 위에서 양의 차수를 갖는 유한개의 동차 원소
$f_1, \ldots, f_r$로 생성된다.
$d = \text{lcm}(\deg(f_i))$로 놓자. $M$을 등급 $A$-가군이라 하자.
이 상황에서 동차 원소 $x \in M$가
{\it 무관하다}\footnote{이는 표준적인 표기법이 아니다.}고 하는 것은
$$
(A_+ x)_{nd} = 0\text{ for all }n \gg 0
$$
일 때이다. $x \in M$가 동차이고 무관하며 $f \in A$가 동차이면,
$fx$도 무관하다. 따라서 무관한 원소들의 집합은 등급 부분가군
$M_{irrelevant} \subset M$을 생성한다. $M$이 {\it 무관하다}고
하는 것은 $M$의 모든 동차 원소가 무관할 때, 즉
$M_{irrelevant} = M$일 때이다. $M$이 유한 생성이면 이는
$M_{nd} = 0$이 $n \gg 0$에 대해 성립하는 것과 동치이다.
$\text{Mod}_A$로 등급 $A$-가군들의 범주를,
$\text{Mod}^{fg}_A$로 유한 생성인 것들의 충만한 부분범주를,
$\text{Mod}^{fg}_{A, irrelevant}$로 무관 가군들의 충만한
부분범주를 나타낸다.

\begin{proposition}
\label{coherent-proposition-coherent-modules-on-proj-general}
$A$를 뇌터 등급환이라 하자. $X = \text{Proj}(A)$로 놓자. 함자
$M \mapsto \widetilde M$은 동치
$$
\text{Mod}^{fg}_A/\text{Mod}^{fg}_{A, irrelevant}
\longrightarrow
\textit{Coh}(\mathcal{O}_X)
$$
를 유도하며, 그 준역은
$\mathcal{F} \longmapsto \bigoplus_{n \geq 0} \Gamma(X, \mathcal{F}(n))$
으로 주어진다.
\end{proposition}

\begin{proof}
먼저 $A$가 차수 $1$에서 생성되는 경우의 증명을 읽기를 독자에게
권한다. Proposition \ref{coherent-proposition-coherent-modules-on-proj}를 보라.
$f_1, \ldots, f_r \in A$를 양의 차수의 동차 원소들로서 $A$를
$A_0$ 위에서 생성한다고 하자. $d$를 차수들 $d_i$의 최소공배수라
하자. 이들은 $f_i$의 차수이다. $M$을 유한 $A$-가군이라 하자.
$\widetilde{M}$이 영일 필요충분조건이 $M$이 무관 등급
$A$-가군인 것임을(명제의 서술 앞에서 정의한 의미로) 보이자.
실제로 $x \in M$를 동차 원소라 하자.
$k \in \mathbf{Z}$를 충분히 작게 택하고,
Lemma \ref{coherent-lemma-recover-tail-graded-module-general}에서와 같이
$N \to N'$ 및 $M \to N'$을 택한다.
또한 $l$을 충분히 크게 택하여 $M_n \to N_n$이 $n \geq l$에
대해 동형이 되게 할 수 있다. $\widetilde{M}$이 영이면 $N = 0$이다.
따라서 임의의 동차 원소 $f \in A_+$에 대해
$\deg(f) + \deg(x) = nd$이고 $nd > l$이면 $fx$는 영이다. 실제로
$N_{nd} \to N'_{nd}$와 $M_{nd} \to N'_{nd}$가 동형이기 때문이다.
따라서 $x$는 무관하다. 역으로 $M$이 무관하다고 하자. 그러면
$M_{nd}$는 $n \gg 0$에 대해 영이다(명제 앞의 논의를 보라).
이는 분명히
$M_{(f_i)} = M_{(f_i^{d/\deg(f_i)})} = 0$을 함의하며, 따라서 구성에
의해 $\widetilde{M} = 0$이다.

\medskip\noindent
따라서 부분범주 $\text{Mod}^{fg}_{A, irrelevant}$는
$\text{Mod}^{fg}_A$의 세르 부분범주이며, 완전 함자
$M \mapsto \widetilde M$의 핵이다. Homology, Lemma
\ref{homology-lemma-kernel-exact-functor}와 Constructions, Lemma
\ref{constructions-lemma-proj-sheaves}를 보라.
그러므로 화살표 왼쪽의 몫 범주는 Homology, Lemma
\ref{homology-lemma-serre-subcategory-is-kernel}에서 정의된다.
명제의 함자를 정의하려면 함자 $M \mapsto \widetilde M$이 무관
가군들을 $0$으로 보냄을 보이면 충분하며, 이는 위에서 보였다.

\medskip\noindent
Lemma \ref{coherent-lemma-coherent-on-proj-general}에 의해 제안된 준역은
잘 정의된다. 실제로 이 보조정리는
$\mathcal{F} \longmapsto \bigoplus_{n \geq 0} \Gamma(X, \mathcal{F}(n))$
이 함자 $\textit{Coh}(\mathcal{O}_X) \to \text{Mod}^{fg}_A$임을 보인다.
이를 몫 함자
$\text{Mod}^{fg}_A \to \text{Mod}^{fg}_A/\text{Mod}^{fg}_{A, irrelevant}$
와 합성할 수 있다.

\medskip\noindent
Lemma \ref{coherent-lemma-recover-tail-graded-module-general}에 의해 왼쪽에서
오른쪽을 거쳐 다시 왼쪽으로 가는 합성은 항등 함자와 동형이다.
실제로 $M$을 유한 등급 $A$-가군이라 하고, $k \in \mathbf{Z}$를
충분히 작게 택하며, Lemma
\ref{coherent-lemma-recover-tail-graded-module-general}에서와 같이
$N \to N'$ 및 $M \to N'$을 택하자. 그러면 $M \to N'$의 핵과
여핵은 유한개의 차수에서만 영이 아니므로 무관하다. 또한 사상
$N \to N'$의 핵과 여핵은 $d$로 나누어지는 모든 충분히 큰
차수에서 영이므로 이들도 무관 가군이다. 따라서 $M \to N'$과
$N \to N'$은 모두 몫 범주에서 동형이며, 이것이 원한 바이다.

\medskip\noindent
마지막으로 $\mathcal{F}$를 연접 $\mathcal{O}_X$-가군이라 하자.
$M = \bigoplus_{n \in \mathbf{Z}} \Gamma(X, \mathcal{F}(n))$로 놓고
이를 등급 $A$-가군으로 보자. 그러면 우리의 함자는 $\mathcal{F}$를
$M_{\geq 0} = \bigoplus_{n \geq 0} M_n$으로 보낸다.
Properties, Lemma \ref{properties-lemma-proj-quasi-coherent}에 의해
표준 사상 $\widetilde M \to \mathcal{F}$는 동형이다. 포함 사상
$M_{\geq 0} \to M$은 동형
$\widetilde{M_{\geq 0}} \to \widetilde M$을 정의하므로, 오른쪽에서
왼쪽을 거쳐 다시 오른쪽으로 가는 합성도 항등 함자와 동형이라고
결론 내린다.
\end{proof}

\section{사영 사상에 따른 고차 직상}
\label{coherent-section-projective-pushforward}

\noindent
먼저 밑이 아핀인 경우의 결과를 서술하고 증명한 뒤, 사영 사상에 관한
몇 가지 결과를 도출한다.

\begin{lemma}
\label{coherent-lemma-coherent-proper-ample}
$R$을 뇌터 환이라 하자. $X \to \Spec(R)$을 고유 사상이라 하자.
$\mathcal{L}$을 $X$ 위의 풍부한 가역층이라 하자. $\mathcal{F}$를
연접 $\mathcal{O}_X$-가군이라 하자.
\begin{enumerate}
\item 등급환
$A = \bigoplus_{d \geq 0} H^0(X, \mathcal{L}^{\otimes d})$는
유한 생성 $R$-대수이다.
\item 어떤 $r \geq 0$와 $d_1, \ldots, d_r \in \mathbf{Z}$ 및 전사
$$
\bigoplus\nolimits_{j = 1, \ldots, r} \mathcal{L}^{\otimes d_j}
\longrightarrow \mathcal{F}.
$$
가 존재한다.
\item 임의의 $p$에 대해 코호몰로지 군 $H^p(X, \mathcal{F})$는 유한
$R$-가군이다.
\item $p > 0$이면
$H^p(X, \mathcal{F} \otimes_{\mathcal{O}_X} \mathcal{L}^{\otimes d}) = 0$은
충분히 큰 모든 $d$에 대해 성립한다.
\item 임의의 $k \in \mathbf{Z}$에 대해 등급 $A$-가군
$$
\bigoplus\nolimits_{d \geq k}
H^0(X, \mathcal{F} \otimes_{\mathcal{O}_X} \mathcal{L}^{\otimes d})
$$
은 유한 $A$-가군이다.
\end{enumerate}
\end{lemma}

\begin{proof}
Morphisms, Lemma
\ref{morphisms-lemma-finite-type-over-affine-ample-very-ample}에 의해
어떤 $d > 0$와 몰입 $i : X \to \mathbf{P}^n_R$가 존재하여
$\mathcal{L}^{\otimes d} \cong i^*\mathcal{O}_{\mathbf{P}^n_R}(1)$이다.
$X$가 $R$ 위에서 고유이므로 사상 $i$는 닫힌 몰입이다
(Morphisms, Lemma \ref{morphisms-lemma-image-proper-scheme-closed}).
따라서 $H^i(X, \mathcal{G}) = H^i(\mathbf{P}^n_R, i_*\mathcal{G})$이고,
이는 임의의 준연접층 $\mathcal{G}$에 대해 $X$ 위에서 성립한다
(Lemma \ref{coherent-lemma-relative-affine-cohomology} 및 닫힌 몰입은
아핀이라는 사실에 의한다. Morphisms, Lemma
\ref{morphisms-lemma-closed-immersion-affine}를 보라).
더욱이 $\mathcal{G}$가 연접이면 $i_*\mathcal{G}$도 연접이다
(Lemma \ref{coherent-lemma-i-star-equivalence}). 이하에서는 이 사실들을
별도의 언급 없이 사용한다.

\medskip\noindent
(1)의 증명. $S = R[T_0, \ldots, T_n]$로 놓아서
$\mathbf{P}^n_R = \text{Proj}(S)$가 되게 하자.
$A$는 $S$-대수임을 관찰하자(다만 환 사상 $S \to A$는 등급환의
준동형이 아니다. $S_n$이 $A_{dn}$으로 가기 때문이다).
사영 공식(Cohomology, Lemma
\ref{cohomology-lemma-projection-formula})에 의해
$$
i_*(\mathcal{L}^{\otimes nd + q}) =
i_*(\mathcal{L}^{\otimes q})
\otimes_{\mathcal{O}_{\mathbf{P}^n_R}}
\mathcal{O}_{\mathbf{P}^n_R}(n)
$$
이고 이는 모든 $n \in \mathbf{Z}$에 대해 성립한다. 따라서 Lemma
\ref{coherent-lemma-coherent-projective}에 의해
$\bigoplus_{n \geq 0} A_{nd + q}$는 유한 등급 $S$-가군이다.
한편
$A = \bigoplus_{q \in \{0, \ldots, d - 1} \bigoplus_{n \geq 0} A_{nd + q}$
이므로 $A$는 $S$-대수로서 유한이다. 따라서 (1)이 참이다.

\medskip\noindent
(2)의 증명. 이는 Properties, Proposition
\ref{properties-proposition-characterize-ample}에서 따른다.

\medskip\noindent
(3)의 증명. Lemma \ref{coherent-lemma-coherent-projective}를 적용하고
$H^p(X, \mathcal{F}) = H^p(\mathbf{P}^n_R, i_*\mathcal{F})$를 사용한다.

\medskip\noindent
(4)의 증명. $p > 0$을 고정하자. 사영 공식에 의해
$$
i_*(\mathcal{F} \otimes_{\mathcal{O}_X} \mathcal{L}^{\otimes nd + q}) =
i_*(\mathcal{F} \otimes_{\mathcal{O}_X} \mathcal{L}^{\otimes q})
\otimes_{\mathcal{O}_{\mathbf{P}^n_R}}
\mathcal{O}_{\mathbf{P}^n_R}(n)
$$
이고 이는 모든 $n \in \mathbf{Z}$에 대해 성립한다. Lemma
\ref{coherent-lemma-coherent-projective}에 의해
$H^p(X, \mathcal{F} \otimes \mathcal{L}^{nd + q}) = 0$은
$n \gg 0$에 대해 성립한다. 정수의 $d$를 법으로 하는 합동류는
유한개뿐이므로 이것으로 (4)가 증명된다.

\medskip\noindent
(5)의 증명. 정수 $k$를 고정하자.
$M = \bigoplus_{n \geq k} H^0(X, \mathcal{F} \otimes \mathcal{L}^{\otimes n})$로
놓는다. 위와 같이 논증하면
$\bigoplus_{nd + q \geq k} M_{nd + q}$가 유한 등급 $S$-가군임을 얻는다.
한편
$M = \bigoplus_{q \in \{0, \ldots, d - 1\}}
\bigoplus_{nd + q \geq k} M_{nd + q}$
이므로 $M$은 유한 $S$-가군이다. $S$-가군 구조는 환 사상
$S \to A$를 통해 분해되므로, $M$은 유한 $A$-가군이라고 결론 내린다.
\end{proof}

\begin{lemma}
\label{coherent-lemma-kill-by-twisting}
$f : X \to S$를 스킴의 사상이라 하자. $\mathcal{F}$를 준연접
$\mathcal{O}_X$-가군이라 하자. $\mathcal{L}$을 $X$ 위의 가역층이라 하자.
다음을 가정한다.
\begin{enumerate}
\item $S$는 뇌터이다.
\item $f$는 고유이다.
\item $\mathcal{F}$는 연접이다.
\item $\mathcal{L}$은 $X/S$ 위에서 상대적으로 풍부하다.
\end{enumerate}
그러면 어떤 $n_0$가 존재하여 모든 $n \geq n_0$에 대해
$$
R^pf_*\left(\mathcal{F} \otimes_{\mathcal{O}_X} \mathcal{L}^{\otimes n}\right)
=
0
$$
이고 이는 모든 $p > 0$에 대해 성립한다.
\end{lemma}

\begin{proof}
유한 아핀 열린 덮개 $S = \bigcup V_j$를 택하고
$X_j = f^{-1}(V_j)$로 놓는다. 유한개의 각 계
$(X_j \to V_j, \mathcal{L}|_{X_j}, \mathcal{F}|_{V_j})$에 대해 문제를
풀면 결과가 따름은 분명하다. 따라서 $S$가 아핀이라고 가정할 수 있다.
이 경우
$R^pf_*(\mathcal{F} \otimes \mathcal{L}^{\otimes n})$의 소멸은
$H^p(X, \mathcal{F} \otimes \mathcal{L}^{\otimes n})$의 소멸과 동치이다.
Lemma \ref{coherent-lemma-quasi-coherence-higher-direct-images-application}를 보라.
따라서 필요한 소멸은 Lemma \ref{coherent-lemma-coherent-proper-ample}에서
따른다(이 보조정리를 적용할 수 있는 것은 Morphisms, Lemma
\ref{morphisms-lemma-finite-type-over-affine-ample-very-ample}에 의해
$\mathcal{L}$이 $X$ 위에서 풍부하기 때문이다).
\end{proof}

\begin{lemma}
\label{coherent-lemma-locally-projective-pushforward}
$S$를 국소 뇌터 스킴이라 하자. $f : X \to S$를 국소 사영
사상이라 하자. $\mathcal{F}$를 연접 $\mathcal{O}_X$-가군이라 하자.
그러면 $R^if_*\mathcal{F}$는 연접 $\mathcal{O}_S$-가군이고, 이는
모든 $i \geq 0$에 대해 성립한다.
\end{lemma}

\begin{proof}
먼저 국소 사영 사상은 고유이고(Morphisms, Lemma
\ref{morphisms-lemma-locally-projective-proper}), 따라서 유한형임을
언급해 둔다. 특히 $X$는 국소 뇌터이고(Morphisms, Lemma
\ref{morphisms-lemma-finite-type-noetherian}), 따라서 이 명제는
의미가 있다. 더욱이 Lemma
\ref{coherent-lemma-quasi-coherence-higher-direct-images}에 의해 층들
$R^pf_*\mathcal{F}$는 준연접이다.

\medskip\noindent
이제 명제는 $S$ 위에서 국소적이다(예를 들어 Cohomology, Lemma
\ref{cohomology-lemma-localize-higher-direct-images}에 의해).
따라서 $S = \Spec(R)$가 뇌터 환의 스펙트럼이고, $X$가
$\mathbf{P}^n_R$의 닫힌 부분스킴이라고 가정할 수 있다. 여기서
$n$은 어떤 정수이다.
Morphisms, Lemma
\ref{morphisms-lemma-characterize-locally-projective}를 보라.
이 경우 층들 $R^pf_*\mathcal{F}$는 $R$-가군들
$H^p(X, \mathcal{F})$에 딸린 준연접층이다. Lemma
\ref{coherent-lemma-quasi-coherence-higher-direct-images-application}를 보라.
따라서 $R$-가군들 $H^p(X, \mathcal{F})$가 유한 $R$-가군임을
보이면 충분하다(Lemma \ref{coherent-lemma-coherent-Noetherian}).
이는 Lemma \ref{coherent-lemma-coherent-proper-ample}에서 따른다
($\mathcal{O}_{\mathbf{P}^n_R}(1)$을 $X$에 제한한 층이 $X$ 위에서
풍부하기 때문이다).
\end{proof}

\section{풍부한 가역층과 코호몰로지}
\label{coherent-section-ample-cohomology}

\noindent
아핀 밑 위의 고유 스킴에서, 꼬임 뒤의 코호몰로지 소멸로 풍부성을
판정하는 기준을 제시한다.

\begin{lemma}
\label{coherent-lemma-vanshing-gives-ample}
\begin{reference}
\cite[III Proposition 2.6.1]{EGA}
\end{reference}
$R$을 뇌터 환이라 하자. $f : X \to \Spec(R)$을 고유 사상이라 하자.
$\mathcal{L}$을 가역 $\mathcal{O}_X$-가군이라 하자. 다음은 서로 동치이다.
\begin{enumerate}
\item $\mathcal{L}$은 $X$ 위에서 풍부하다(이는 다른 많은 조건과
동치이다. Properties, Proposition
\ref{properties-proposition-characterize-ample} 및 Morphisms, Lemma
\ref{morphisms-lemma-finite-type-over-affine-ample-very-ample}를 보라).
\item 모든 연접 $\mathcal{O}_X$-가군 $\mathcal{F}$에 대해 어떤
$n_0 \geq 0$가 존재하여
$H^p(X, \mathcal{F} \otimes \mathcal{L}^{\otimes n}) = 0$이고, 이는
모든 $n \geq n_0$ 및 $p > 0$에 대해 성립한다.
\item 모든 준연접 아이디얼 층
$\mathcal{I} \subset \mathcal{O}_X$에 대해 어떤 $n \geq 1$이 존재하여
$H^1(X, \mathcal{I} \otimes \mathcal{L}^{\otimes n}) = 0$이다.
\end{enumerate}
\end{lemma}

\begin{proof}
함의 (1) $\Rightarrow$ (2)는 Lemma
\ref{coherent-lemma-coherent-proper-ample}에서 따른다.
함의 (2) $\Rightarrow$ (3)은 자명하다.
함의 (3) $\Rightarrow$ (1)은 Lemma
\ref{coherent-lemma-quasi-compact-h1-zero-invertible}이다.
\end{proof}

\begin{lemma}
\label{coherent-lemma-surjective-finite-morphism-ample}
$R$을 뇌터 환이라 하자. $f : Y \to X$를 $R$ 위에서 고유인 스킴들의
사상이라 하자. $\mathcal{L}$을 가역 $\mathcal{O}_X$-가군이라 하자.
$f$가 유한이고 전사라고 가정하자. 그러면 $\mathcal{L}$이 풍부일
필요충분조건은 $f^*\mathcal{L}$이 풍부인 것이다.
\end{lemma}

\begin{proof}
풍부한 가역층을 준아핀 사상으로 당겨 얻은 층은 풍부하다. Morphisms,
Lemma \ref{morphisms-lemma-pullback-ample-tensor-relatively-ample}를
보라. 유한 사상은 정의에 의해 아핀이므로 이것으로 한쪽 함의가 증명된다.

\medskip\noindent
$f^*\mathcal{L}$이 풍부하다고 가정하자. 다음 성질을 $P$라 하자.
이는 연접 $\mathcal{O}_X$-가군 $\mathcal{F}$에 관한 성질이다:
어떤 $n_0$가 존재하여
$H^p(X, \mathcal{F} \otimes \mathcal{L}^{\otimes n}) = 0$이고, 이는
모든 $n \geq n_0$ 및 $p > 0$에 대해 성립한다.
$P$가 임의의 연접 $\mathcal{O}_X$-가군 $\mathcal{F}$에 대해
성립함을 보이겠다. 그러면 Lemma
\ref{coherent-lemma-vanshing-gives-ample}에 의해 $\mathcal{L}$이 풍부이다.
Lemma \ref{coherent-lemma-property-higher-rank-cohomological}를 적용하겠다.
따라서 $P$에 대해 그 보조정리의 (1), (2), (3)을 확인해야 한다.
성질 (1)은 층의 짧은 완전열에 딸린 긴 코호몰로지 완전열 및 가역층과의
텐서곱이 완전 함자라는 사실에서 따른다. 성질 (2)는
$H^p(X, -)$가 가법 함자이므로 성립한다. (3)을 보이기 위해
$Z \subset X$를 일반점 $\xi$를 갖는 정수 닫힌 부분스킴이라 하자.
$\mathcal{F}$를 $Y$ 위의 연접층으로서 $f_*\mathcal{F}$의 지지가
$Z$와 같고 $(f_*\mathcal{F})_\xi$가 $\mathfrak m_\xi$에 의해
소멸되는 층이라 하자. Lemma
\ref{coherent-lemma-finite-morphism-Noetherian}를 보라.
$\mathcal{G} = f_*\mathcal{F}$로 택하면 된다고 주장한다.
Lemma \ref{coherent-lemma-property-higher-rank-cohomological}의 (3)(c)만
확인하면 된다. 따라서 $\mathcal{J} \subset \mathcal{O}_X$가
$\mathcal{J}_\xi = \mathcal{O}_{X, \xi}$를 만족하는 준연접 아이디얼
층이라고 가정하자. 유한 사상은 아핀이므로 Lemma
\ref{coherent-lemma-affine-morphism-projection-ideal}에 의해
$\mathcal{J}\mathcal{G} = f_*(f^{-1}\mathcal{J}\mathcal{F})$이다.
또한 Lemma \ref{coherent-lemma-affine-morphism-projection-ideal}의 증명에서
지적했듯이 층 $f^{-1}\mathcal{J}\mathcal{F}$는 연접
$\mathcal{O}_Y$-가군이다. $\mathcal{L}$이 풍부이므로 Lemma
\ref{coherent-lemma-vanshing-gives-ample}에서 어떤 $n_0$가 존재하여
$$
H^p(Y, f^{-1}\mathcal{J}\mathcal{F}
\otimes_{\mathcal{O}_Y} f^*\mathcal{L}^{\otimes n}) = 0,
$$
이고 이는 $n \geq n_0$ 및 $p > 0$에 대해 성립한다. $f$는
유한이므로 아핀이고, 따라서
\begin{align*}
H^p(X, \mathcal{J}\mathcal{G} \otimes_{\mathcal{O}_X}
\mathcal{L}^{\otimes n})
& =
H^p(X, f_*(f^{-1}\mathcal{J}\mathcal{F}) \otimes_{\mathcal{O}_X}
\mathcal{L}^{\otimes n}) \\
& =
H^p(X, f_*(f^{-1}\mathcal{J}\mathcal{F} \otimes_{\mathcal{O}_Y}
f^*\mathcal{L}^{\otimes n})) \\
& =
H^p(Y, f^{-1}\mathcal{J}\mathcal{F} \otimes_{\mathcal{O}_Y}
f^*\mathcal{L}^{\otimes n}) = 0
\end{align*}
이다. 여기서 사영 공식(Cohomology, Lemma
\ref{cohomology-lemma-projection-formula})과 Lemma
\ref{coherent-lemma-relative-affine-cohomology}를 사용했다.
따라서 준연접 부분층 $\mathcal{G}' = \mathcal{J}\mathcal{G}$는
$P$를 만족한다. 이로써 Lemma
\ref{coherent-lemma-property-higher-rank-cohomological}의 성질 (3)(c)가
원하는 대로 확인된다.
\end{proof}

\noindent
코호몰로지는 함자적이다. 특히 환 달린 공간 $X$, 가역
$\mathcal{O}_X$-가군 $\mathcal{L}$, 절단
$s \in \Gamma(X, \mathcal{L})$가 주어지면 사상들
$$
H^p(X, \mathcal{F})
\longrightarrow
H^p(X, \mathcal{F} \otimes_{\mathcal{O}_X} \mathcal{L}), \quad
\xi \longmapsto s\xi
$$
을 얻는다. 이는 사상
$\mathcal{F} \to \mathcal{F} \otimes_{\mathcal{O}_X} \mathcal{L}$에
의해 유도되며, 그 사상은 $s$를 곱하는 사상이다. 등급환으로
$\Gamma_*(X, \mathcal{L}) =
\bigoplus_{n \geq 0} \Gamma(X, \mathcal{L}^{\otimes n})$로 놓는다.
Modules, Definition \ref{modules-definition-gamma-star}를 보라.
$\mathcal{O}_X$-가군의 층 $\mathcal{F}$와 정수 $p \geq 0$가
주어지면
$$
H^p_*(X, \mathcal{L}, \mathcal{F}) =
\bigoplus\nolimits_{n \in \mathbf{Z}} H^p(X,
\mathcal{F} \otimes_{\mathcal{O}_X} \mathcal{L}^{\otimes n})
$$
로 놓는다. 이는 위에서 정의한 곱셈에 의해 등급
$\Gamma_*(X, \mathcal{L})$-가군이다. 주의:
$H^p_*(X, \mathcal{L}, \mathcal{F})$라는 표기법은 표준적이지 않다.

\begin{lemma}
\label{coherent-lemma-invert-s-cohomology}
$X$를 스킴이라 하자. $\mathcal{L}$을 $X$ 위의 가역층이라 하자.
$s \in \Gamma(X, \mathcal{L})$라 하자. $\mathcal{F}$를 준연접
$\mathcal{O}_X$-가군이라 하자. $X$가 준콤팩트이고 준분리이면 표준 사상
$$
H^p_*(X, \mathcal{L}, \mathcal{F})_{(s)}
\longrightarrow
H^p(X_s, \mathcal{F})
$$
은 동형이다. 이 사상은 $\xi/s^n$을 $s^{-n}\xi$로 보낸다.
\end{lemma}

\begin{proof}
$p = 0$인 경우 이는 Properties, Lemma
\ref{properties-lemma-invert-s-sections}이다.
귀납 원리(Lemma \ref{coherent-lemma-induction-principle})를 사용하여 명제를
증명하겠다. 준콤팩트 열린집합 $U \subset X$에 대해 다음 성질을
$P(U)$라 하자: 모든 $p \geq 0$에 대해 사상
$$
H^p_*(U, \mathcal{L}, \mathcal{F})_{(s)}
\longrightarrow
H^p(U_s, \mathcal{F})
$$
은 동형이다.

\medskip\noindent
$U$가 아핀이면 위에 표시한 화살표의 양변은 $p > 0$에 대해 영이다.
이는 Lemma \ref{coherent-lemma-quasi-coherent-affine-cohomology-zero}와
Properties, Lemma \ref{properties-lemma-affine-cap-s-open}에
의하며, 따라서 명제가 참이다. $P$가 $U$, $V$, $U \cap V$에 대해
참이면 Mayer--Vietoris 수열(Cohomology, Lemma
\ref{cohomology-lemma-mayer-vietoris})을 사용하여 긴 완전열의 사상
$$
\xymatrix{
H^{p - 1}_*(U \cap V, \mathcal{L}, \mathcal{F})_{(s)} \ar[r] \ar[d] &
H^p_*(U \cup V, \mathcal{L}, \mathcal{F})_{(s)} \ar[r] \ar[d] &
H^p_*(U, \mathcal{L}, \mathcal{F})_{(s)}
\oplus
H^p_*(V, \mathcal{L}, \mathcal{F})_{(s)} \ar[d] \\
H^{p - 1}(U_s \cap V_s, \mathcal{F}) \ar[r]&
H^p(U_s \cup V_s, \mathcal{F}) \ar[r] &
H^p(U_s, \mathcal{F})
\oplus
H^p(V_s, \mathcal{F})
}
$$
을 얻는다(일부만 표시했다). $U_s \cap V_s = (U \cap V)_s$이고
$U_s \cup V_s = (U \cup V)_s$임을 관찰하자. 따라서 왼쪽과 오른쪽의
수직 사상은 동형이다(그림에 표시하지 않은 오른쪽의 사상 하나와
왼쪽의 사상 하나도 동형이다). 5-보조정리(Homology, Lemma
\ref{homology-lemma-five-lemma})에 의해 $P(U \cup V)$가
성립한다고 결론 내린다. 이것으로 증명이 끝난다.
\end{proof}

\begin{lemma}
\label{coherent-lemma-section-affine-open-kills-classes}
$X$를 스킴이라 하자. $\mathcal{L}$을 가역 $\mathcal{O}_X$-가군이라
하자. $s \in \Gamma(X, \mathcal{L})$를 절단이라 하자. 다음을 가정한다.
\begin{enumerate}
\item $X$는 준콤팩트이고 준분리이다.
\item $X_s$는 아핀이다.
\end{enumerate}
그러면 모든 준연접 $\mathcal{O}_X$-가군 $\mathcal{F}$, 모든
$p > 0$, 모든 $\xi \in H^p(X, \mathcal{F})$에 대해 어떤
$n \geq 0$가 존재하여 $s^n\xi = 0$이고, 이는
$H^p(X, \mathcal{F} \otimes_{\mathcal{O}_X} \mathcal{L}^{\otimes n})$
안의 등식이다.
\end{lemma}

\begin{proof}
$H^p(X_s, \mathcal{G})$가 모든 준연접 가군 $\mathcal{G}$에 대해
영임을 상기하자. 이는 Lemma
\ref{coherent-lemma-quasi-coherent-affine-cohomology-zero}에 따른다.
따라서 보조정리는 Lemma \ref{coherent-lemma-invert-s-cohomology}에서 따른다.
\end{proof}

\noindent
다음 보조정리의 더 일반적인 형태는 Limits, Lemma
\ref{limits-lemma-ample-on-reduction}를 보라.

\begin{lemma}
\label{coherent-lemma-ample-on-reduction}
$i : Z \to X$를 뇌터 스킴들의 닫힌 몰입으로서 밑 위상공간의
위상동형을 유도한다고 하자. $\mathcal{L}$을 $X$ 위의 가역층이라 하자.
그러면 $i^*\mathcal{L}$이 $Z$ 위에서 풍부일 필요충분조건은
$\mathcal{L}$이 $X$ 위에서 풍부인 것이다.
\end{lemma}

\begin{proof}
$\mathcal{L}$이 풍부이면 예를 들어 Morphisms, Lemma
\ref{morphisms-lemma-pullback-ample-tensor-relatively-ample}에 의해
$i^*\mathcal{L}$이 풍부이다. $i^*\mathcal{L}$이 풍부하다고 가정하자.
$\mathcal{L}$이 $X$ 위에서 풍부임을 보여야 한다.
$\mathcal{I} \subset \mathcal{O}_X$를 닫힌 부분스킴 $Z$를 잘라내는
연접 아이디얼 층이라 하자. 집합론적으로 $i(Z) = X$이므로 Lemma
\ref{coherent-lemma-power-ideal-kills-sheaf}에 의해 $\mathcal{I}^n = 0$이고,
이를 만족하는 어떤 $n$이 존재한다. 닫힌 부분스킴들의 열
$$
X = Z_n \supset Z_{n - 1} \supset Z_{n - 2} \supset \ldots \supset Z_1 = Z
$$
을 생각하자. 이들은
$0 = \mathcal{I}^n \subset \mathcal{I}^{n - 1} \subset \ldots \subset
\mathcal{I}$에 의해 잘려 나온다. 그러면 각 닫힌 몰입
$Z_{i - 1} \to Z_{i}$는 제곱이 영인 연접 아이디얼 층으로 정의된다.
이렇게 하여 $\mathcal{I}^2 = 0$인 경우로 환원한다.

\medskip\noindent
준연접 $\mathcal{O}_X$-가군들의 짧은 완전열
$$
0 \to \mathcal{I} \to \mathcal{O}_X \to i_*\mathcal{O}_Z \to 0
$$
을 생각하자. $\mathcal{L}^{\otimes n}$과 텐서곱하면 짧은 완전열들
\begin{equation}
\label{coherent-equation-ses}
0 \to \mathcal{I} \otimes_{\mathcal{O}_X} \mathcal{L}^{\otimes n}
\to \mathcal{L}^{\otimes n} \to i_*i^*\mathcal{L}^{\otimes n} \to 0
\end{equation}
을 얻는다. $\mathcal{I}^2 = 0$이므로 Morphisms, Lemma
\ref{morphisms-lemma-i-star-equivalence}를 사용하여
$\mathcal{I}$를 준연접 $\mathcal{O}_Z$-가군으로 생각할 수 있고,
명백히 표기법을 남용하면
$\mathcal{I} \otimes_{\mathcal{O}_X} \mathcal{L}^{\otimes n} =
\mathcal{I} \otimes_{\mathcal{O}_Z} i^*\mathcal{L}^{\otimes n}$이다.
더욱이 이 층의 $Z$ 위 코호몰로지는 표준적으로 이 층의 $X$ 위
코호몰로지와 같다($i$가 위상동형이므로).

\medskip\noindent
$x \in X$를 점이라 하고, 이에 대응하는 점을 $z \in Z$로 나타내자.
$i^*\mathcal{L}$이 풍부이므로 어떤 $n$과 절단
$s \in \Gamma(Z, i^*\mathcal{L}^{\otimes n})$가 존재하여
$z \in Z_s$이고 $Z_s$는 아핀이다. $s$를
$\mathcal{L}^{\otimes n}$의 절단으로 $X$ 위에서 올리는 데 대한 장애물은
경계류
$$
\xi = \partial s \in
H^1(X, \mathcal{I} \otimes_{\mathcal{O}_X} \mathcal{L}^{\otimes n}) =
H^1(Z, \mathcal{I} \otimes_{\mathcal{O}_Z} i^*\mathcal{L}^{\otimes n})
$$
이다. 이는 층의 짧은 완전열 (\ref{coherent-equation-ses})에서 나온다.
$s$를 $s^{e + 1}$로 바꾸면 $\xi$는
$\partial(s^{e + 1}) = (e + 1) s^e \xi$로 바뀌며, 이는
$H^1(Z, \mathcal{I} \otimes_{\mathcal{O}_Z} i^*\mathcal{L}^{\otimes (e + 1)n})$
안의 등식이다. 실제로
$$
0 \to
\bigoplus\nolimits_{m \geq 0}
\mathcal{I} \otimes_{\mathcal{O}_X} \mathcal{L}^{\otimes m} \to
\bigoplus\nolimits_{m \geq 0}
\mathcal{L}^{\otimes m} \to
\bigoplus\nolimits_{m \geq 0}
i_*i^*\mathcal{L}^{\otimes m} \to 0
$$
의 경계 사상은 Cohomology, Lemma
\ref{cohomology-lemma-boundary-derivation-over-cup-product}에 의해
미분이다. Lemma \ref{coherent-lemma-section-affine-open-kills-classes}에 의해
$s^e \xi$는 충분히 큰 $e$에 대해 영이다. 따라서 $s$를 거듭제곱으로
바꾼 뒤에는 $s$가 어떤 절단
$s' \in \Gamma(X, \mathcal{L}^{\otimes n})$의 상이라고 가정할 수 있다.
그러면 $X_{s'}$은 열린 부분스킴이고
$Z_s \to X_{s'}$은 뇌터 스킴들의 전사 닫힌 몰입이며 $Z_s$는
아핀이다. 따라서 Lemma
\ref{coherent-lemma-image-affine-finite-morphism-affine-Noetherian}에 의해
$X_s$는 아핀이고, $\mathcal{L}$이 풍부라고 결론 내린다.
\end{proof}

\noindent
다음 보조정리의 더 일반적인 형태는 Limits, Lemma
\ref{limits-lemma-thickening-quasi-affine}를 보라.

\begin{lemma}
\label{coherent-lemma-thickening-quasi-affine}
$i : Z \to X$를 뇌터 스킴들의 닫힌 몰입으로서 밑 위상공간의
위상동형을 유도한다고 하자. 그러면 $X$가 준아핀일 필요충분조건은
$Z$가 준아핀인 것이다.
\end{lemma}

\begin{proof}
스킴이 준아핀일 필요충분조건은 구조층이 풍부인 것임을 상기하자.
Properties, Lemma \ref{properties-lemma-quasi-affine-O-ample}을 보라.
따라서 $Z$가 준아핀이면 $\mathcal{O}_Z$가 풍부이고, Lemma
\ref{coherent-lemma-ample-on-reduction}에 의해 $\mathcal{O}_X$가 풍부이며,
그러므로 $X$는 준아핀이다. 그 역도 방금 제시한 논증을 거꾸로
읽으면 증명되며, 초등적인 방법으로도 알 수 있다.
\end{proof}

\begin{lemma}
\label{coherent-lemma-affine-in-presence-ample}
$X$를 스킴이라 하자. $\mathcal{L}$을 풍부한 가역
$\mathcal{O}_X$-가군이라 하자. $n_0$를 정수라 하자.
$H^p(X, \mathcal{L}^{\otimes -n}) = 0$이 $n \geq n_0$ 및
$p > 0$에 대해 성립하면 $X$는 아핀이다.
\end{lemma}

\begin{proof}
$H^p(X, \mathcal{F}) = 0$이라고 주장한다. 이는 모든 준연접
$\mathcal{O}_X$-가군과 $p > 0$에 대해 성립한다. Properties, Definition
\ref{properties-definition-ample}에 의해 $X$는 준콤팩트이므로,
이 주장은 Lemma \ref{coherent-lemma-quasi-compact-h1-zero-covering}에 의해
증명을 끝낸다. Properties, Lemma
\ref{properties-lemma-ample-separated}에 의해 스킴 $X$는 분리이다.
$X$가 $e + 1$개의 아핀 열린집합으로 덮인다고 하자. 그러면 Lemma
\ref{coherent-lemma-vanishing-nr-affines}에 의해 $H^p(X, \mathcal{F}) = 0$이고,
이는 $p > e$일 때 성립한다. 따라서 $p$에 대한 내림 귀납법으로
주장을 증명할 수 있다. $\mathcal{F}$를 유한형 준연접 가군들의
필터 여극한으로 쓰고(Properties, Lemma
\ref{properties-lemma-quasi-coherent-colimit-finite-type}),
Cohomology, Lemma
\ref{cohomology-lemma-quasi-separated-cohomology-colimit}를 사용하면
$\mathcal{F}$가 유한형이라고 가정할 수 있다.
그러면 어떤 $n > n_0$를 택하여
$\mathcal{F} \otimes_{\mathcal{O}_X} \mathcal{L}^{\otimes n}$이
대역 생성되게 할 수 있다. Properties, Proposition
\ref{properties-proposition-characterize-ample}을 보라.
이는 어떤 집합 $I$에 대해 짧은 완전열
$$
0 \to \mathcal{F}' \to
\bigoplus\nolimits_{i \in I} \mathcal{L}^{\otimes -n}
\to \mathcal{F} \to 0
$$
이 있다는 뜻이다(실제로 $I$를 유한으로 택할 수 있다). 귀납 가정에
의해 $H^{p + 1}(X, \mathcal{F}') = 0$이고, 가정 및 이미 사용한
코호몰로지와 여극한의 교환 가능성을 함께 사용하면
$H^p(X, \bigoplus_{i \in I} \mathcal{L}^{\otimes -n}) = 0$이다.
긴 코호몰로지 완전열에서 원하는 대로
$H^p(X, \mathcal{F}) = 0$이라고 결론 내린다.
\end{proof}

\begin{lemma}
\label{coherent-lemma-affine-if-quasi-affine}
$X$를 준아핀 스킴이라 하자.
$H^p(X, \mathcal{O}_X) = 0$이 $p > 0$에 대해 성립하면
$X$는 아핀이다.
\end{lemma}

\begin{proof}
Properties, Lemma \ref{properties-lemma-quasi-affine-O-ample}에 의해
$\mathcal{O}_X$가 풍부이므로, 이는 Lemma
\ref{coherent-lemma-affine-in-presence-ample}에서 따른다.
\end{proof}

\section{Chow 보조정리}
\label{coherent-section-chows-lemma}

\noindent
이 절에서는 뇌터인 경우의 Chow 보조정리(Lemma
\ref{coherent-lemma-chow-Noetherian})를 증명한다. 비뇌터인 경우의 몇 가지
변형은 Limits, Section \ref{limits-section-chows-lemma}에서 증명한다.

\begin{lemma}
\label{coherent-lemma-chow-Noetherian}
\begin{reference}
\cite[II Theorem 5.6.1(a)]{EGA}
\end{reference}
$S$를 뇌터 스킴이라 하자. $f : X \to S$를 분리 유한형 사상이라 하자.
그러면 어떤 $n \geq 0$와 그림
$$
\xymatrix{
X \ar[rd] & X' \ar[d] \ar[l]^\pi \ar[r] & \mathbf{P}^n_S \ar[dl] \\
& S &
}
$$
이 존재한다. 여기서 $X' \to \mathbf{P}^n_S$는 몰입이고,
$\pi : X' \to X$는 고유 전사이다. 더욱이 어떤 조밀 열린
부분스킴 $U \subset X$가 존재하여 $\pi^{-1}(U) \to U$가 동형이
되도록 정할 수 있다.
\end{lemma}

\begin{proof}
이하 증명에서 만나는 모든 스킴은 뇌터 스킴 $S$ 위에서 유한형이고,
따라서 뇌터이다(Morphisms, Lemma
\ref{morphisms-lemma-finite-type-noetherian}를 보라).
그 사이의 모든 사상은 자동으로 준콤팩트이고, 국소 유한형이며,
준분리이다. Morphisms, Lemma
\ref{morphisms-lemma-permanence-finite-type}와 Properties, Lemmas
\ref{properties-lemma-locally-Noetherian-quasi-separated} 및
\ref{properties-lemma-morphism-Noetherian-schemes-quasi-compact}를 보라.

\medskip\noindent
스킴 $X$는 기약 성분을 유한개만 갖는다(Properties, Lemma
\ref{properties-lemma-Noetherian-irreducible-components}).
$X = X_1 \cup \ldots \cup X_r$를 $X$의 기약 성분 분해라 하자.
$\eta_i \in X_i$를 일반점이라 하자. 모든 점 $x \in X$에 대해 아핀
열린집합 $U_x \subset X$가 존재하여 $x$와 각 일반점 $\eta_i$를
포함한다. Properties, Lemma
\ref{properties-lemma-point-and-maximal-points-affine}를 보라.
$X$는 준콤팩트이므로 유한 아핀 열린 덮개
$X = U_1 \cup \ldots \cup U_m$을 찾아서 각 $U_i$가
$\eta_1, \ldots, \eta_r$를 포함하게 할 수 있다.
특히 열린집합
$U = U_1 \cap \ldots \cap U_m \subset X$는 조밀 열린집합이다.
이 사실과 $U_i$들이 $X$를 덮는 아핀 열린집합이라는 사실만을
이하에서 사용한다.

\medskip\noindent
$X^* \subset X$를 $U \to X$의 스킴론적 폐포라 하자. Morphisms,
Definition \ref{morphisms-definition-scheme-theoretic-image}를 보라.
$U_i^* = X^* \cap U_i$로 놓자. $U_i^*$는 $U_i$의 닫힌 부분스킴임을
주목하자. 따라서 $U_i^*$는 아핀이다. $U$가 $X$에서 조밀하므로
사상 $X^* \to X$는 전사 닫힌 몰입이다. 이는 $U$ 위에서 동형이다.
따라서 $X$를 $X^*$로, $U_i$를 $U_i^*$로 바꾸고 $U$가 $X$에서
스킴론적으로 조밀하다고 가정할 수 있다. Morphisms, Definition
\ref{morphisms-definition-scheme-theoretically-dense}를 보라.

\medskip\noindent
Morphisms, Lemma
\ref{morphisms-lemma-quasi-projective-finite-type-over-S}에 의해 몰입
$j_i : U_i \to \mathbf{P}_S^{n_i}$를 각 $i$에 대해 찾을 수 있다.
Morphisms, Lemma \ref{morphisms-lemma-quasi-compact-immersion}에 의해
닫힌 부분스킴 $Z_i \subset \mathbf{P}_S^{n_i}$를 찾아서
$j_i : U_i \to Z_i$가 스킴론적으로 조밀한 열린 몰입이 되게 할 수 있다.
$Z_i \to S$는 고유임을 주목하자. Morphisms, Lemma
\ref{morphisms-lemma-locally-projective-proper}를 보라. 사상
$$
j = (j_1|_U, \ldots, j_m|_U) : U \longrightarrow
\mathbf{P}_S^{n_1} \times_S \ldots \times_S \mathbf{P}_S^{n_m}.
$$
을 생각하자. 위에서 인용한 보조정리에 의해 닫힌 부분스킴 $Z$를
$\mathbf{P}_S^{n_1} \times_S \ldots \times_S \mathbf{P}_S^{n_m}$ 안에서
찾아서 $j : U \to Z$가 열린 몰입이고 $U$가
$Z$에서 스킴론적으로 조밀하게 할 수 있다. 사상 $Z \to S$는 고유이다.
$i$번째 사영
$$
\text{pr}_i|_Z : Z \longrightarrow \mathbf{P}^{n_i}_S.
$$
을 생각하자. 이 사상은 $Z_i$를 통해 분해된다(Morphisms, Lemma
\ref{morphisms-lemma-factor-factor}를 보라). 유도된 사상을
$p_i : Z \to Z_i$로 나타내자. 이는 고유 사상이다. 예를 들어
Morphisms, Lemma
\ref{morphisms-lemma-image-proper-scheme-closed}를 보라. 이제
$U \subset U_i \subset Z_i$는 스킴론적으로 조밀한 열린 몰입들이다.
더욱이 $Z$를 “대각” 사상
$U \to Z_1 \times_S \ldots \times_S Z_m$의 스킴론적 상으로 생각할
수 있다.

\medskip\noindent
$V_i = p_i^{-1}(U_i)$로 놓자.
$p_i|_{V_i} : V_i \to U_i$는 고유임을 주목하자.
$X' = V_1 \cup \ldots \cup V_m$로 놓자. 구성에 의해 $X'$은 스킴
$\mathbf{P}^{n_1}_S \times_S \ldots \times_S \mathbf{P}^{n_m}_S$로의
몰입을 갖는다. 따라서 Segre 매장(Constructions, Lemma
\ref{constructions-lemma-segre-embedding}를 보라)에 의해 $X'$은
$S$ 위의 사영공간으로의 몰입을 갖는다.

\medskip\noindent
사상들 $p_i|_{V_i}: V_i \to U_i$가 사상 $X' \to X$로 접합된다고
주장한다. 실제로 $p_i|_U$가 $U$에서 $U$로 가는 항등 사상임은
분명하다. 구성에 의해 $U \subset X'$은 스킴론적으로 조밀하므로,
열린 부분스킴 $V_i \cap V_j$에서도 스킴론적으로 조밀하다. 따라서
$p_i|_{V_i \cap V_j} = p_j|_{V_i \cap V_j}$이다. 이들은 분리된
$S$-스킴 $X$로 가는 사상이다. Morphisms, Lemma
\ref{morphisms-lemma-equality-of-morphisms}를 보라. 그 결과로 얻는
사상을 $\pi : X' \to X$로 나타낸다.

\medskip\noindent
$\pi^{-1}(U_i) = V_i$라고 주장한다.
$\pi|_{V_i} = p_i|_{V_i}$이므로
$V_i \subset \pi^{-1}(U_i)$이다. 그림
$$
\xymatrix{
V_i \ar[r] \ar[rd]_{p_i|_{V_i}} & \pi^{-1}(U_i) \ar[d] \\
& U_i
}
$$
을 생각하자. $V_i \to U_i$는 고유이므로 수평 화살표의 상은 닫혀 있다.
Morphisms, Lemma
\ref{morphisms-lemma-image-proper-scheme-closed}를 보라.
$V_i \subset \pi^{-1}(U_i)$는 $U$를 포함하므로 스킴론적으로
조밀하다. 따라서 주장한 대로 $V_i = \pi^{-1}(U_i)$이다.

\medskip\noindent
이로써 $\pi^{-1}(U_i) \to U_i$가 고유 사상
$p_i|_{V_i} : V_i \to U_i$와 동일시됨을 알 수 있다. 따라서 $X$는
유한 아핀 덮개 $X = \bigcup U_i$를 가지며, 각 덮개 원소에 제한한
$\pi$는 고유이다. 그러므로 Morphisms, Lemma
\ref{morphisms-lemma-proper-local-on-the-base}에 의해 $\pi$는 고유이다.

\medskip\noindent
마지막으로 $\pi^{-1}(U) = U$임을 보여야 한다. 이를 위해 위와
마찬가지로 그림
$$
\xymatrix{
U \ar[r] \ar[rd] & \pi^{-1}(U) \ar[d] \\
& U
}
$$
을 사용하여 논증한다. 여기서 $\text{id}_U : U \to U$가 고유이고
$U$가 $\pi^{-1}(U)$에서 스킴론적으로 조밀함을 사용한다.
\end{proof}

\begin{remark}
\label{coherent-remark-chow-Noetherian}
Chow Lemma \ref{coherent-lemma-chow-Noetherian}의 상황에서 다음이 성립한다.
\begin{enumerate}
\item 사상 $\pi$는 실제로 H-사영이다(따라서 사영이다. Morphisms,
Lemma \ref{morphisms-lemma-H-projective}를 보라). 실제로 사상
$X' \to \mathbf{P}^n_S \times_S X = \mathbf{P}^n_X$는 닫힌 몰입이다
($\pi$가 고유라는 사실을 사용하라. Morphisms, Lemma
\ref{morphisms-lemma-image-proper-scheme-closed}를 보라).
\item $\pi^{-1}(U)$가 $X'$에서 스킴론적으로 조밀하다고 가정할 수 있다.
실제로 $X'$을 $\pi^{-1}(U)$의 스킴론적 폐포로 바꾸기만 하면 된다.
이 경우 $U$를 $X'$의 스킴론적으로 조밀한 열린 부분스킴으로 생각할
수 있다. Morphisms, Section
\ref{morphisms-section-scheme-theoretic-image}를 보라.
\item $X$가 축약이면 $X'$도 축약으로 택할 수 있다. 이는 (2)에서
분명하다.
\end{enumerate}
\end{remark}

\section{연접층의 고차 직상}
\label{coherent-section-proper-pushforward}

\noindent
이 절에서는 고유 사상에 의한 연접층의 고차 직상이 연접이라는
기본 사실을 증명한다.

\begin{proposition}
\label{coherent-proposition-proper-pushforward-coherent}
\begin{reference}
\cite[III Theorem 3.2.1]{EGA}
\end{reference}
$S$를 국소 뇌터 스킴이라 하자. $f : X \to S$를 고유 사상이라 하자.
$\mathcal{F}$를 연접 $\mathcal{O}_X$-가군이라 하자.
그러면 $R^if_*\mathcal{F}$는 연접 $\mathcal{O}_S$-가군이고, 이는
모든 $i \geq 0$에 대해 성립한다.
\end{proposition}

\begin{proof}
문제는 $S$ 위에서 국소적이므로 $S$가 뇌터 스킴이라고 가정할 수 있다.
고유 사상은 유한형이므로 이 경우 $X$도 뇌터 스킴이다.
연접층의 성질 $\mathcal{P}$를 $X$ 위에서 다음 규칙으로 정의하자:
$$
\mathcal{P}(\mathcal{F}) \Leftrightarrow
R^pf_*\mathcal{F}\text{ is coherent for all }p \geq 0
$$
Lemma \ref{coherent-lemma-property}의 결과를 사용하여 $\mathcal{P}$가
$X$ 위의 모든 연접층에 대해 성립함을 증명하겠다.

\medskip\noindent
$X$ 위의 연접층들의 짧은 완전열
$$
0 \to \mathcal{F}_1 \to \mathcal{F}_2 \to \mathcal{F}_3 \to 0
$$
을 생각하자. 고차 직상의 긴 완전열
$$
R^{p - 1}f_*\mathcal{F}_3 \to
R^pf_*\mathcal{F}_1 \to
R^pf_*\mathcal{F}_2 \to
R^pf_*\mathcal{F}_3 \to
R^{p + 1}f_*\mathcal{F}_1
$$
을 생각하자. 층들 $\mathcal{F}_i$ 가운데 셋 중 둘이 성질
$\mathcal{P}$를 가지면, 셋째 층의 고차 직상은 이 완전 복합체에서
두 연접층 사이에 끼인다. 따라서 Lemmas
\ref{coherent-lemma-coherent-abelian-Noetherian}와
\ref{coherent-lemma-coherent-Noetherian-quasi-coherent-sub-quotient}에 의해
이 고차 직상들도 연접이다. 그러므로 성질 $\mathcal{P}$는
셋째 층에 대해서도 성립한다.

\medskip\noindent
$Z \subset X$를 정수 닫힌 부분스킴이라 하자. 연접층 $\mathcal{F}$를
$X$ 위에서 찾아야 한다. 이 층의 지지는 $Z$에 포함되고, 그 줄기는
일반점 $\xi$, 즉 $Z$의 일반점에서 $1$차원 $\kappa(\xi)$-벡터공간이며,
성질 $\mathcal{P}$가 $\mathcal{F}$에 대해 성립해야 한다.
$g = f|_Z : Z \to S$로 $f$의 제한을 나타내자. 연접층 $\mathcal{G}$를
$Z$ 위에서 찾아서 다음을 만족시킬 수 있다고 하자:
(a) $\mathcal{G}_\xi$는 $1$차원 $\kappa(\xi)$-벡터공간이다.
(b) $R^pg_*\mathcal{G} = 0$은 $p > 0$에 대해 성립한다.
(c) $g_*\mathcal{G}$는 연접이다. 그러면
$\mathcal{F} = (Z \to X)_*\mathcal{G}$를 생각할 수 있다.
$Z \to X$가 닫힌 몰입이므로 $(Z \to X)_*\mathcal{G}$는 $X$ 위에서
연접이고, $R^p(Z \to X)_*\mathcal{G} = 0$은 $p > 0$에 대해 성립한다
(Lemma \ref{coherent-lemma-finite-pushforward-coherent}).
따라서 상대 Leray 스펙트럼 수열(Cohomology, Lemma
\ref{cohomology-lemma-relative-Leray})에 의해
$R^pf_*\mathcal{F} = R^pg_*\mathcal{G} = 0$은 $p > 0$에 대해
성립하고, $f_*\mathcal{F} = g_*\mathcal{G}$는 연접이다.
마지막으로
$\mathcal{F}_\xi = ((Z \to X)_*\mathcal{G})_\xi = \mathcal{G}_\xi$이므로
$\xi$에서의 줄기에 관한 조건이 확인된다. 따라서 모든 것은 (a), (b),
(c)를 만족하는 연접층 $\mathcal{G}$를 $Z$ 위에서 찾는 데 달려 있다.

\medskip\noindent
Chow's Lemma \ref{coherent-lemma-chow-Noetherian}를 사상 $Z \to S$에 적용할
수 있다. 따라서 Chow 보조정리의 서술과 같은 그림
$$
\xymatrix{
Z \ar[rd]_g & Z' \ar[d]^-{g'} \ar[l]^\pi \ar[r]_i & \mathbf{P}^m_S \ar[dl] \\
& S &
}
$$
을 얻는다. 또한 $U \subset Z$를 $\pi^{-1}(U) \to U$가 동형이 되는
조밀 열린 부분스킴이라 하자. Remark
\ref{coherent-remark-chow-Noetherian}의 논의에 의해
$i' = (i, \pi) : Z' \to \mathbf{P}^m_Z$는 닫힌 몰입이다. 따라서
$$
\mathcal{L} = i^*\mathcal{O}_{\mathbf{P}^m_S}(1) \cong
(i')^*\mathcal{O}_{\mathbf{P}^m_Z}(1)
$$
은 $g'$-상대적으로 풍부하고 $\pi$-상대적으로 풍부하다(예를 들어
Morphisms, Lemma
\ref{morphisms-lemma-characterize-ample-on-finite-type}에 의해).
따라서 Lemma \ref{coherent-lemma-kill-by-twisting}에 의해 어떤
$n \geq 0$가 존재하여 $R^p\pi_*\mathcal{L}^{\otimes n} = 0$이 모든
$p > 0$에 대해 성립하고, $R^p(g')_*\mathcal{L}^{\otimes n} = 0$도
모든 $p > 0$에 대해 성립한다.
$\mathcal{G} = \pi_*\mathcal{L}^{\otimes n}$로 놓자.
$\pi_*\mathcal{L}^{\otimes n}|_U$가 가역층이므로 성질 (a)가 성립한다
($\pi^{-1}(U) \to U$가 동형이기 때문이다). 상대 Leray 스펙트럼
수열(Cohomology, Lemma \ref{cohomology-lemma-relative-Leray})에 의해
$$
E_2^{p, q} = R^pg_* R^q\pi_*\mathcal{L}^{\otimes n}
\Rightarrow
R^{p + q}(g')_*\mathcal{L}^{\otimes n}
$$
이므로 성질 (b)와 (c)가 성립한다. 실제로 $n$의 선택에 의해
$E_2^{p, q}$에서 영이 아닌 항은 $q = 0$인 항뿐이고,
$R^{p + q}(g')_*\mathcal{L}^{\otimes n}$에서 영이 아닌 항은
$p = q = 0$인 항뿐이다. 이는 $R^pg_*\mathcal{G} = 0$이
$p > 0$에 대해 성립하고
$g_*\mathcal{G} = (g')_*\mathcal{L}^{\otimes n}$임을 함의한다.
마지막으로 앞의 Lemma
\ref{coherent-lemma-locally-projective-pushforward}를 적용하면
$g_*\mathcal{G} = (g')_*\mathcal{L}^{\otimes n}$이 원하는 대로
연접임을 알 수 있다.
\end{proof}

\begin{lemma}
\label{coherent-lemma-proper-over-affine-cohomology-finite}
$S = \Spec(A)$이고 $A$가 뇌터 환이라고 하자.
$f : X \to S$를 고유 사상이라 하자. $\mathcal{F}$를 연접
$\mathcal{O}_X$-가군이라 하자. 그러면 $H^i(X, \mathcal{F})$는
유한 $A$-가군이고, 이는 모든 $i \geq 0$에 대해 성립한다.
\end{lemma}

\begin{proof}
이는 Proposition \ref{coherent-proposition-proper-pushforward-coherent}의
아핀인 경우일 뿐이다. 실제로 Lemmas
\ref{coherent-lemma-quasi-coherence-higher-direct-images}와
\ref{coherent-lemma-quasi-coherence-higher-direct-images-application}에 의해
$R^if_*\mathcal{F}$는 $A$-가군 $H^i(X, \mathcal{F})$에 딸린
준연접층이고, Lemma \ref{coherent-lemma-coherent-Noetherian}에 의해 이것이
연접층일 필요충분조건은 $H^i(X, \mathcal{F})$가 유한형
$A$-가군인 것이다.
\end{proof}

\begin{lemma}
\label{coherent-lemma-graded-finiteness}
$A$를 뇌터 환이라 하자. $B$를 유한 생성 등급 $A$-대수라 하자.
$f : X \to \Spec(A)$를 고유 사상이라 하자.
$\mathcal{B} = f^*\widetilde B$로 놓자. $\mathcal{F}$를 유한형
준연접 등급 $\mathcal{B}$-가군이라 하자.
\begin{enumerate}
\item 모든 $p \geq 0$에 대해 등급 $B$-가군
$H^p(X, \mathcal{F})$는 유한 $B$-가군이다.
\item $\mathcal{L}$이 풍부한 가역 $\mathcal{O}_X$-가군이면, 어떤
정수 $d_0$가 존재하여
$H^p(X, \mathcal{F} \otimes \mathcal{L}^{\otimes d}) = 0$이고 이는
모든 $p > 0$ 및 $d \geq d_0$에 대해 성립한다.
\end{enumerate}
\end{lemma}

\begin{proof}
이를 증명하기 위해 올곱 그림
$$
\xymatrix{
X' = \Spec(B) \times_{\Spec(A)} X
\ar[r]_-\pi \ar[d]_{f'} &
X \ar[d]^f \\
\Spec(B) \ar[r] &
\Spec(A)
}
$$
을 생각하자. $f'$는 고유 사상임을 주목하자. Morphisms, Lemma
\ref{morphisms-lemma-base-change-proper}를 보라. 또한 $B$는 유한
생성 $A$-대수이므로 뇌터이다(Algebra, Lemma
\ref{algebra-lemma-Noetherian-permanence}). 이는 $X'$이 뇌터
스킴임을 함의한다(Morphisms, Lemma
\ref{morphisms-lemma-finite-type-noetherian}).
$X'$은 Constructions, Lemma
\ref{constructions-lemma-spec-properties}에 의해 준연접
$\mathcal{O}_X$-대수 $\mathcal{B}$의 상대 스펙트럼임을 주목하자.
$\mathcal{F}$는 준연접 $\mathcal{B}$-가군이므로 유일한 준연접
$\mathcal{O}_{X'}$-가군 $\mathcal{F}'$가 존재하여
$\pi_*\mathcal{F}' = \mathcal{F}$이다. Morphisms, Lemma
\ref{morphisms-lemma-affine-equivalence-modules}를 보라.
$\mathcal{F}$가 $\mathcal{B}$-가군으로서 유한형이므로
$\mathcal{F}'$는 유한형 $\mathcal{O}_{X'}$-가군이다(세부사항은
생략한다). 다시 말해 $\mathcal{F}'$는 연접
$\mathcal{O}_{X'}$-가군이다(Lemma \ref{coherent-lemma-coherent-Noetherian}).
사상 $\pi : X' \to X$는 아핀이므로
$$
H^p(X, \mathcal{F}) = H^p(X', \mathcal{F}')
$$
이다. 이는 Lemma \ref{coherent-lemma-relative-affine-cohomology}에 따른다.
따라서 (1)은 Lemma
\ref{coherent-lemma-proper-over-affine-cohomology-finite}에서 따른다.
(2)에서와 같은 $\mathcal{L}$이 주어지면
$\mathcal{L}' = \pi^*\mathcal{L}$로 놓는다. $\mathcal{L}'$은
$X'$ 위에서 풍부임을 주목하자. 이는 Morphisms, Lemma
\ref{morphisms-lemma-pullback-ample-tensor-relatively-ample}에 따른다.
사영 공식(Cohomology, Lemma
\ref{cohomology-lemma-projection-formula})에 의해
$\pi_*(\mathcal{F}' \otimes \mathcal{L}') = \mathcal{F} \otimes \mathcal{L}$이다.
따라서 위와 같은 논증과 Lemma \ref{coherent-lemma-kill-by-twisting}에 의해
(2)가 따른다.
\end{proof}

\section{형식 함수 정리}
\label{coherent-section-theorem-formal-functions}

\noindent
이 절에서는 $\{I^n\mathcal{F}\}_{n \geq 0}$ 꼴 또는
$\{\mathcal{F}/I^n\mathcal{F}\}_{n \geq 0}$ 꼴인 층들의 열의
코호몰로지가 $n$에 따라 어떻게 변하는지 연구한다.

\medskip\noindent
여기와 아래에서는 다음 표기를 사용한다. 스킴의 사상 $f : X \to Y$,
준연접층 $\mathcal{F}$를 $X$ 위에서, 그리고 준연접 아이디얼 층
$\mathcal{I} \subset \mathcal{O}_Y$가 주어졌을 때,
$\mathcal{I}^n\mathcal{F}$는 $f^{-1}(\mathcal{I}^n)$의 국소 단면과
$\mathcal{F}$의 국소 단면의 곱들로 생성되는 준연접 부분층을 뜻한다.
식으로 쓰면
$$
\mathcal{I}^n\mathcal{F}
=
\Im\left(
f^*(\mathcal{I}^n) \otimes_{\mathcal{O}_X} \mathcal{F}
\longrightarrow
\mathcal{F}
\right).
$$
다음 자연스러운 사상들이 있음을 주목하자.
$$
f^{-1}(\mathcal{I}^n) \otimes_{f^{-1}\mathcal{O}_Y} \mathcal{I}^m\mathcal{F}
\longrightarrow
f^*(\mathcal{I}^n) \otimes_{\mathcal{O}_X} \mathcal{I}^m\mathcal{F}
\longrightarrow
\mathcal{I}^{n + m}\mathcal{F}
$$
따라서 $\mathcal{I}^n$의 한 단면은 고차 직상의 함자성에 의해
$R^pf_*(\mathcal{I}^m\mathcal{F}) \to
R^pf_*(\mathcal{I}^{n + m}\mathcal{F})$라는 사상을 준다. 국소화한 뒤
층화하면 다음 $\mathcal{O}_Y$-가군 사상들을 얻는다.
$$
\mathcal{I}^n \otimes_{\mathcal{O}_Y} R^pf_*(\mathcal{I}^m\mathcal{F})
\longrightarrow
R^pf_*(\mathcal{I}^{n + m}\mathcal{F}).
$$
다시 말해 $\bigoplus_{n \geq 0} R^pf_*(\mathcal{I}^n\mathcal{F})$는
등급 $\bigoplus_{n \geq 0} \mathcal{I}^n$-가군이다.

\medskip\noindent
$Y = \Spec(A)$이고 $\mathcal{I} = \widetilde{I}$이면
$\mathcal{I}^n\mathcal{F}$를 간단히 $I^n\mathcal{F}$로 쓴다. 위에서
도입한 사상들은 $M = \bigoplus H^p(X, I^n\mathcal{F})$에 등급
$S = \bigoplus I^n$-가군의 구조를 준다. $f$가 고유이고 $A$가
뇌터이며 $\mathcal{F}$가 연접이면, 이것은 유한형 가군임이 드러난다.

\begin{lemma}
\label{coherent-lemma-cohomology-powers-ideal-times-F}
\begin{reference}
\cite[III Cor 3.3.2]{EGA}
\end{reference}
$A$를 뇌터 환이라 하자. $I \subset A$를 아이디얼이라 하고
$B = \bigoplus_{n \geq 0} I^n$으로 놓자.
$f : X \to \Spec(A)$를 고유 사상이라 하고, 연접층 $\mathcal{F}$를
$X$ 위에서 택하자. 그러면 모든 $p \geq 0$에 대해 등급 $B$-가군
$\bigoplus_{n \geq 0} H^p(X, I^n\mathcal{F})$는 유한 $B$-가군이다.
\end{lemma}

\begin{proof}
$\mathcal{B} = f^*\widetilde{B}$로 놓자. $\mathcal{B}$에서
$\bigoplus I^n\mathcal{O}_X$로 가는 전사가 있으므로
$\bigoplus I^n\mathcal{F}$는 유한형 등급 $\mathcal{B}$-가군임이
분명하다. 따라서 결과는 Lemma \ref{coherent-lemma-graded-finiteness}의 (1)에서
따른다.
\end{proof}

\begin{lemma}
\label{coherent-lemma-cohomology-powers-ideal-times-sheaf}
스킴의 사상 $f : X \to Y$, 준연접층 $\mathcal{F}$를 $X$ 위에서, 그리고
준연접 아이디얼 층 $\mathcal{I} \subset \mathcal{O}_Y$가 주어졌다고
하자. $Y$가 국소 뇌터이고, $f$가 고유이며, $\mathcal{F}$가 연접이라고
가정하자. 그러면
$$
\mathcal{M} =
\bigoplus\nolimits_{n \geq 0} R^pf_*(\mathcal{I}^n\mathcal{F})
$$
은 준연접이고 유한형인 등급
$\mathcal{A} = \bigoplus_{n \geq 0} \mathcal{I}^n$-가군이다.
\end{lemma}

\begin{proof}
명제는 $Y$ 위에서 국소적이므로 $Y$가 아핀인 경우로 환원된다.
아핀인 경우 결과는 Lemma
\ref{coherent-lemma-cohomology-powers-ideal-times-F}에서 따른다. 세부사항은
생략한다.
\end{proof}

\begin{lemma}
\label{coherent-lemma-cohomology-powers-ideal-application}
$A$를 뇌터 환이라 하자. $I \subset A$를 아이디얼이라 하자.
$f : X \to \Spec(A)$를 고유 사상이라 하고, $\mathcal{F}$를 $X$ 위의
연접층이라 하자. 그러면 모든 $p \geq 0$에 대해 어떤 정수
$c \geq 0$가 존재하여 다음이 성립한다.
\begin{enumerate}
\item 곱셈 사상
$I^{n - c} \otimes H^p(X, I^c\mathcal{F}) \to H^p(X, I^n\mathcal{F})$
은 모든 $n \geq c$에 대해 전사이다.
\item $H^p(X, I^{n + m}\mathcal{F}) \to H^p(X, I^n\mathcal{F})$의 상은
부분가군 $I^{m - e} H^p(X, I^n\mathcal{F})$에 포함된다. 여기서
$e = \max(0, c - n)$이고 $n + m \geq c$, $n, m \geq 0$이다.
\item 다음 등식이 성립한다.
$$
\Ker(H^p(X, I^n\mathcal{F}) \to H^p(X, \mathcal{F})) =
\Ker(H^p(X, I^n\mathcal{F}) \to H^p(X, I^{n - c}\mathcal{F}))
$$
이는 $n \geq c$에 대해 성립한다.
\item 사상
$I^nH^p(X, \mathcal{F}) \to H^p(X, I^{n - c}\mathcal{F})$들이
$n \geq c$에 대해 존재하여,
합성
$$
H^p(X, I^n\mathcal{F}) \to 
I^{n - c}H^p(X, \mathcal{F}) \to
H^p(X, I^{n - 2c}\mathcal{F})
$$
과
$$
I^nH^p(X, \mathcal{F}) \to
H^p(X, I^{n - c}\mathcal{F}) \to
I^{n - 2c}H^p(X, \mathcal{F})
$$
은 $n \geq 2c$에 대해 표준 사상이다.
\item 역계 $(H^p(X, I^n\mathcal{F}))$와
$(I^nH^p(X, \mathcal{F}))$는 pro-동형이다.
\end{enumerate}
\end{lemma}

\begin{proof}
$M_n = H^p(X, I^n\mathcal{F})$로 $n \geq 1$에 대해 쓰고,
$M_0 = H^p(X, \mathcal{F})$로 쓰자. 그러면 사상
$\ldots \to M_3 \to M_2 \to M_1 \to M_0$가 있다.
$B = \bigoplus_{n \geq 0} I^n$으로 놓으면
$M = \bigoplus_{n \geq 0} M_n$은 유한 등급 $B$-가군이다. Lemma
\ref{coherent-lemma-cohomology-powers-ideal-times-F}를 보라. 곱
$B_n \otimes M_m \to M_{m + n}$, $a \otimes m \mapsto a \cdot m$은
우리 역계의 사상들과 양립한다. 즉 그림
$$
\xymatrix{
B_n \otimes_A M_m \ar[r] \ar[d] & M_{n + m} \ar[d] \\
B_n \otimes_A M_{m'} \ar[r] & M_{n + m'}
}
$$
은 $n, m' \geq 0$ 및 $m \geq m'$에 대해 가환한다.

\medskip\noindent
(1)의 증명. $d_1, \ldots, d_t \geq 0$ 및 $x_i \in M_{d_i}$를 택하여
$M$이 $x_1, \ldots, x_t$로 $B$ 위에서 생성되게 하자. 임의의
$c \geq \max\{d_i\}$에 대해 $B_{n - c} \cdot M_c = M_n$임을
$n \geq c$에 대해 얻으므로 (1)이 성립한다.

\medskip\noindent
(2)의 증명. $c$를 (1)의 증명에서와 같이 택하고
$n + m \geq c$라 하자. 그러면
$M_{n + m} = B_{n + m - c} \cdot M_c$이다. $c > n$이면
$M_c \to M_n$과 위에서 말한 곱과 전이 사상의 양립성을 사용하여
$M_{n + m} \to M_n$의 상이 $I^{n + m - c}M_n$에 포함됨을 얻는다.
$c \leq n$이면 $M_{n + m} = B_m \cdot B_{n - c} \cdot M_c =
B_m \cdot M_n$으로 써서 그 상이 $I^m M_n$에 포함됨을 알 수 있다.
이로써 (2)가 증명되었다.

\medskip\noindent
$K_n \subset M_n$을 사상 $M_n \to M_0$의 핵이라 하자. 위에서 말한
곱과 전이 사상의 양립성에 의해
$K = \bigoplus K_n \subset M$은 등급 $B$-부분가군이다. $B$는
뇌터이고 $M$은 유한 생성 등급 $B$-가군이므로 $K$는 유한 생성 등급
$B$-가군이다. $d'_1, \ldots, d'_{t'} \geq 0$ 및
$y_i \in K_{d'_i}$를 택하여 $K$가 $y_1, \ldots, y_{t'}$로
$B$ 위에서 생성되게 하자.
$c = \max(d'_i, d'_j)$로 놓자. $y_i \in \Ker(M_{d'_i} \to M_0)$이므로
$B_n \cdot y_i \subset \Ker(M_{n + d'_i} \to M_n)$임을 알 수 있다.
이로부터 $K_n = \Ker(M_n \to M_{n - c})$가 $n \geq c$에 대해
성립함을 알 수 있다. 이로써 (3)이 증명되었다.

\medskip\noindent
다음 가환하는 실선 그림을 생각하자.
$$
\xymatrix{
I^n \otimes_A M_0 \ar[r] \ar[d] &
I^nM_0 \ar[r] \ar@{..>}[d] &
M_0 \ar[d] \\
M_n \ar[r] &
M_{n - c} \ar[r] &
M_0
}
$$
위의 논의에 의해 전사 $I^n \otimes_A M_0 \to I^nM_0$의 핵은
$K_n$으로 사상된다. 따라서 점선 화살표를 얻고 합성
$I^nM_0 \to M_{n - c} \to M_0$는 표준 사상이다. 그러면 합성
$I^nM_0 \to M_{n - c} \to I^{n - 2c}M_0$도 $n \geq 2c$에 대해
분명히 표준 사상이다. 합성
$M_n \to I^{n - c}M_0 \to M_{n - 2c}$를 생각하자. 첫째 사상은
$a \cdot m$ 꼴의 원소를 보낸다. 여기서 $a \in I^{n - c}$이고
$m \in M_c$이다. 그 상은 $a m'$인데, 여기서 $m'$는 $m$의
$M_0$에서의 상이다. 그러면 둘째 사상은 이것을 $a \cdot m'$로
$M_{n - 2c}$ 안에서 보낸다. 따라서
(4)가 성립한다.

\medskip\noindent
(5)는 (4)와 pro-대상의 사상의 정의에서 즉시 따른다.
\end{proof}

\noindent
Lemmas \ref{coherent-lemma-cohomology-powers-ideal-times-F}와
\ref{coherent-lemma-cohomology-powers-ideal-application}의 상황에서 역계
$$
\mathcal{F}/I\mathcal{F} \leftarrow
\mathcal{F}/I^2\mathcal{F} \leftarrow
\mathcal{F}/I^3\mathcal{F} \leftarrow \ldots
$$
를 생각하자. 코호몰로지 군들에 어떤 일이 일어나는지 알고자 한다.
첫 번째 결과는 다음과 같다.

\begin{lemma}
\label{coherent-lemma-ML-cohomology-powers-ideal}
$A$를 뇌터 환이라 하자. $I \subset A$를 아이디얼이라 하자.
$f : X \to \Spec(A)$를 고유 사상이라 하자. 연접층 $\mathcal{F}$를
$X$ 위에서 택하자. $p \geq 0$를 고정하자. 다음을 만족하는
$c \geq 0$가 존재한다.
\begin{enumerate}
\item 모든 $n \geq c$에 대해
$$
\Ker(H^p(X, \mathcal{F}) \to H^p(X, \mathcal{F}/I^n\mathcal{F})) \subset
I^{n - c}H^p(X, \mathcal{F}).
$$
\item 역계
$$
\left(H^p(X, \mathcal{F}/I^n\mathcal{F})\right)_{n \in \mathbf{N}}
$$
는 Mittag-Leffler 조건을 만족한다(Homology, Definition
\ref{homology-definition-Mittag-Leffler}을 보라).
\item 모든 $k \geq n + c$에 대해
$$
\Im(H^p(X, \mathcal{F}/I^k\mathcal{F})
\to H^p(X, \mathcal{F}/I^n\mathcal{F}))
=
\Im(H^p(X, \mathcal{F})
\to H^p(X, \mathcal{F}/I^n\mathcal{F}))
$$
이다.
\end{enumerate}
\end{lemma}

\begin{proof}
$c = \max\{c_p, c_{p + 1}\}$로 놓자. 여기서 $c_p, c_{p + 1}$는
Lemma \ref{coherent-lemma-cohomology-powers-ideal-application}을
$H^p$와 $H^{p + 1}$에 적용하여 얻은 정수들이다.

\medskip\noindent
(1)을 증명하자. 짧은 완전열
$$
0 \to I^n\mathcal{F} \to \mathcal{F} \to \mathcal{F}/I^n\mathcal{F} \to 0
$$
을 생각하자. 긴 코호몰로지 완전열에서
$$
\Ker(
H^p(X, \mathcal{F}) \to H^p(X, \mathcal{F}/I^n\mathcal{F})
)
=
\Im(
H^p(X, I^n\mathcal{F}) \to H^p(X, \mathcal{F})
)
$$
임을 알 수 있다. 따라서 Lemma
\ref{coherent-lemma-cohomology-powers-ideal-application}의 (2)에 의해 이것은
$I^{n - c}H^p(X, \mathcal{F})$에 포함되며, 이는 $n \geq c$일 때
성립한다.

\medskip\noindent
(3)이 Mittag-Leffler 계의 정의에 의해 (2)를 함의함을 주목하자.

\medskip\noindent
(3)을 증명하자. $n$을 고정하고 다음 가환 그림을 생각하자.
$$
\xymatrix{
0 \ar[r] &
I^n\mathcal{F} \ar[r] &
\mathcal{F} \ar[r] &
\mathcal{F}/I^n\mathcal{F} \ar[r] &
0 \\
0 \ar[r] &
I^{n + m}\mathcal{F} \ar[r] \ar[u] &
\mathcal{F} \ar[r] \ar[u] &
\mathcal{F}/I^{n + m}\mathcal{F} \ar[r] \ar[u] &
0
}
$$
이로부터 완전한 행들을 갖는 다음 가환 그림을 얻는다.
$$
\xymatrix{
H^p(X, \mathcal{F}) \ar[r] &
H^p(X, \mathcal{F}/I^n\mathcal{F}) \ar[r]_\delta &
H^{p + 1}(X, I^n\mathcal{F}) \ar[r] &
H^{p + 1}(X, \mathcal{F}) \\
H^p(X, \mathcal{F}) \ar[r] \ar[u]^1 &
H^p(X, \mathcal{F}/I^{n + m}\mathcal{F}) \ar[r] \ar[u]^\gamma &
H^{p + 1}(X, I^{n + m}\mathcal{F}) \ar[u]^\alpha \ar[r]^-\beta &
H^{p + 1}(X, \mathcal{F}) \ar[u]_1
}
$$
Lemma \ref{coherent-lemma-cohomology-powers-ideal-application}의 (4)에 의해
$\beta$의 핵은 $\alpha$의 핵과 같으며, 이는 $m \geq c$일 때
성립한다. 그림 추적을 하면
$\gamma$의 상이 $\delta$의 핵에 포함됨을 얻으며, 이는 (3)을 증명한다
($k = n + m$로 놓으면 된다).
\end{proof}

\begin{theorem}[형식 함수 정리]
\label{coherent-theorem-formal-functions}
$A$를 뇌터 환이라 하자. $I \subset A$를 아이디얼이라 하자.
$f : X \to \Spec(A)$를 고유 사상이라 하자. 연접층 $\mathcal{F}$를
$X$ 위에서 택하자. $p \geq 0$를 고정하자. 사상들의 계
$$
H^p(X, \mathcal{F})/I^nH^p(X, \mathcal{F})
\longrightarrow
H^p(X, \mathcal{F}/I^n\mathcal{F})
$$
는 극한의 동형사상
$$
H^p(X, \mathcal{F})^\wedge
\longrightarrow
\lim_n H^p(X, \mathcal{F}/I^n\mathcal{F})
$$
을 정의한다. 여기서 왼쪽은 $A$-가군 $H^p(X, \mathcal{F})$를
아이디얼 $I$에 관해 완비화한 것이다. Algebra, Section
\ref{algebra-section-completion}을 보라. 더 나아가 이는 실제로
극한 위상들에 관한 위상동형이다.
\end{theorem}

\begin{proof}
이는 Lemma \ref{coherent-lemma-ML-cohomology-powers-ideal}에서 다음과 같이
따른다. $M = H^p(X, \mathcal{F})$,
$M_n = H^p(X, \mathcal{F}/I^n\mathcal{F})$로 놓고,
$N_n = \Im(M \to M_n)$로 나타내자. Lemma
\ref{coherent-lemma-ML-cohomology-powers-ideal}의 (2)와 (3)에 의해
$(M_n)$은 Mittag-Leffler 계이고 $N_n \subset M_n$은 $M_k$의 상과
같으며, 이는 모든 $k \gg n$에 대해 성립한다. 따라서 위상 가군으로서
$\lim M_n = \lim N_n$이다(극한 위상을 취한다). 한편 $N_n$들은
가군 $M$의 몫들로 이루어진 역계를 이루므로 $\lim N_n$은
$M$을 핵들 $K_n = \Ker(M \to N_n)$이 주는 위상에 관해 완비화한
것이다. Lemma \ref{coherent-lemma-ML-cohomology-powers-ideal}의 (1)에 의해
$K_n \subset I^{n - c}M$이고, $N_n \subset M_n$은 $I^n$에 의해
소멸되므로 $I^n M \subset K_n$이다. 따라서 부분가군 $K_n$들을
$0$의 열린 근방의 기본계로 사용하여 정의한 위상은 $I$-진 위상과
같다. 그러므로 유도된 사상
$M^\wedge = \lim M/I^nM \to \lim N_n = \lim M_n$은 위상 가군의
동형사상이다\footnote{정확히 말하면, $M^\wedge$ 위의 극한 위상은
그 $I$-진 위상과 같다. 이는 Algebra, Lemma
\ref{algebra-lemma-hathat-finitely-generated}에서 따른다.
More on Algebra, Section \ref{more-algebra-section-topological-ring}도 보라.}.
\end{proof}

\begin{lemma}
\label{coherent-lemma-spell-out-theorem-formal-functions}
$A$를 환이라 하고 $I \subset A$를 아이디얼이라 하자. $A$가 뇌터이고
$I$에 관해 완비라고 가정하자. $f : X \to \Spec(A)$를 고유 사상이라
하자. 연접층 $\mathcal{F}$를 $X$ 위에서 택하자. 그러면 모든
$p \geq 0$에 대해
$$
H^p(X, \mathcal{F}) = \lim_n H^p(X, \mathcal{F}/I^n\mathcal{F})
$$
이다.
\end{lemma}

\begin{proof}
이는 완비 뇌터 밑환의 경우에 형식 함수 정리(Theorem
\ref{coherent-theorem-formal-functions})를 다시 쓴 것이다. 실제로 이 경우
$A$-가군 $H^p(X, \mathcal{F})$는 유한이고(Lemma
\ref{coherent-lemma-proper-over-affine-cohomology-finite}), 따라서
$I$-진 완비이다(Algebra, Lemma
\ref{algebra-lemma-completion-tensor}). 그러므로 왼쪽에서는 완비화가
필요하지 않음을 알 수 있다.
\end{proof}

\begin{lemma}
\label{coherent-lemma-formal-functions-stalk}
스킴의 사상 $f : X \to Y$와 준연접층 $\mathcal{F}$가 $X$ 위에
주어졌다고 하자. 다음을 가정하자.
\begin{enumerate}
\item $Y$는 국소 뇌터이다.
\item $f$는 고유이다.
\item $\mathcal{F}$는 연접이다.
\end{enumerate}
점 $y \in Y$를 택하자. 올
$$
\xymatrix{
X_n =
\Spec(\mathcal{O}_{Y, y}/\mathfrak m_y^n) \times_Y X
\ar[r]_-{i_n} \ar[d]_{f_n} &
X \ar[d]^f \\
\Spec(\mathcal{O}_{Y, y}/\mathfrak m_y^n) \ar[r]^-{c_n} & Y
}
$$
$X_1 = X_y$의 무한소 근방들을 생각하고
$\mathcal{F}_n = i_n^*\mathcal{F}$로 놓자. 그러면
$$
\left(R^pf_*\mathcal{F}\right)_y^\wedge
\cong
\lim_n H^p(X_n, \mathcal{F}_n)
$$
이 $\mathcal{O}_{Y, y}^\wedge$-가군의 동형으로 성립한다.
\end{lemma}

\begin{proof}
이는 형식 함수 정리 Theorem \ref{coherent-theorem-formal-functions}의 특수한
경우를 다시 쓴 것일 뿐이다. 자세히 써 보자.
$\mathcal{O}_{Y, y}$가 뇌터 국소환임을 주목하자. 표준 사상
$c : \Spec(\mathcal{O}_{Y, y}) \to Y$를 생각하자. Schemes, Equation
(\ref{schemes-equation-canonical-morphism})을 보라. 이 사상은 국소환들을
동일시하므로 평탄하다. 잠시
$f' : X' \to \Spec(\mathcal{O}_{Y, y})$로 $f$를 이 국소환으로
밑변환한 사상을 나타내자. Lemma
\ref{coherent-lemma-flat-base-change-cohomology}에 의해
$c^*R^pf_*\mathcal{F} = R^pf'_*\mathcal{F}'$임을 알 수 있다. 또한 올
$X_y$와 $X'_y$의 무한소 근방들은 동일시된다(확인은 생략한다. 힌트:
사상 $c_n$들은 $c$를 경유한다).

\medskip\noindent
따라서 $Y = \Spec(A)$가 뇌터 국소환 $A$의 스펙트럼이고, 이 환의
극대 아이디얼이 $\mathfrak m$이며, $y \in Y$는 닫힌점, 즉 $\mathfrak m$에
대응한다고 가정할 수 있다. 특히
$$
\left(R^pf_*\mathcal{F}\right)_y =
\Gamma(Y, R^pf_*\mathcal{F}) =
H^p(X, \mathcal{F}).
$$

\medskip\noindent
이 경우 사상 $c_n$들도 각각 닫힌 몰입이다. 따라서 그 밑변환인
$i_n$들도 닫힌 몰입이다.
$i_{n, *}\mathcal{F}_n = i_{n, *}i_n^*\mathcal{F}
= \mathcal{F}/\mathfrak m^n\mathcal{F}$임을 주목하자. $i_n$에 대한
Leray 스펙트럼 수열과 Lemma
\ref{coherent-lemma-finite-pushforward-coherent}에 의해
$$
H^p(X_n, \mathcal{F}_n) =
H^p(X, i_{n, *}\mathcal{F}_n) =
H^p(X, \mathcal{F}/\mathfrak m^n\mathcal{F})
$$
임을 알 수 있다. 따라서 보조정리의 서술에 있는 극한을 계산하는 데
형식 함수 정리를 실제로 적용할 수 있고, 이로써 끝난다.
\end{proof}

\noindent
다음 보조정리는 나중에 차원이 $ > 0$인 올로 일반화할 결과이다.

\begin{lemma}
\label{coherent-lemma-higher-direct-images-zero-finite-fibre}
$f : X \to Y$를 스킴의 사상이라 하고 $y \in Y$라 하자. 다음을
가정하자.
\begin{enumerate}
\item $Y$는 국소 뇌터이다.
\item $f$는 고유이다.
\item $f^{-1}(\{y\})$는 유한이다.
\end{enumerate}
그러면 임의의 연접층 $\mathcal{F}$를 $X$ 위에서 택했을 때
$(R^pf_*\mathcal{F})_y = 0$이고, 이는 모든 $p > 0$에 대해 성립한다.
\end{lemma}

\begin{proof}
올 $X_y$는 유한이고, Morphisms, Lemma
\ref{morphisms-lemma-finite-fibre}에 의해 유한 이산공간이다. 또한 각
무한소 근방 $X_n$의 바탕 위상공간도 같다. 따라서 Schemes, Lemma
\ref{schemes-lemma-scheme-finite-discrete-affine}에 의해 각 스킴
$X_n$은 아핀이다. 그러므로 $H^p(X_n, \mathcal{F}_n) = 0$이고, 이는
모든 $p > 0$에 대해 성립한다. 따라서 Lemma
\ref{coherent-lemma-formal-functions-stalk}에 의해
$(R^pf_*\mathcal{F})_y^\wedge = 0$임을 알 수 있다.
$R^pf_*\mathcal{F}$는 Proposition
\ref{coherent-proposition-proper-pushforward-coherent}에 의해 연접이고, 따라서
$R^pf_*\mathcal{F}_y$는 유한 $\mathcal{O}_{Y, y}$-가군이다. Nakayama
보조정리(Algebra, Lemma \ref{algebra-lemma-NAK})에 의해 국소환 위의
유한 가군의 완비화가 영이면 그 가군도 영이다. 따라서
$(R^pf_*\mathcal{F})_y = 0$이다.
\end{proof}

\begin{lemma}
\label{coherent-lemma-higher-direct-images-zero-above-dimension-fibre}
$f : X \to Y$를 스킴의 사상이라 하고 $y \in Y$라 하자. 다음을
가정하자.
\begin{enumerate}
\item $Y$는 국소 뇌터이다.
\item $f$는 고유이다.
\item $\dim(X_y) = d$이다.
\end{enumerate}
그러면 임의의 연접층 $\mathcal{F}$를 $X$ 위에서 택했을 때
$(R^pf_*\mathcal{F})_y = 0$이고, 이는 모든 $p > d$에 대해 성립한다.
\end{lemma}

\begin{proof}
올 $X_y$는 $\Spec(\kappa(y))$ 위 유한형이다. 따라서 Morphisms, Lemma
\ref{morphisms-lemma-finite-type-noetherian}에 의해 $X_y$는 뇌터
스킴이다. 그러므로 $X_y$의 바탕 위상공간은 뇌터이다. Properties,
Lemma \ref{properties-lemma-Noetherian-topology}를 보라. 또한 각 무한소
근방 $X_n$의 바탕 위상공간은 $X_y$의 바탕 위상공간과 같다. 따라서
Cohomology, Proposition
\ref{cohomology-proposition-vanishing-Noetherian}에 의해
$H^p(X_n, \mathcal{F}_n) = 0$이고, 이는 모든 $p > d$에 대해
성립한다. 그러므로 Lemma \ref{coherent-lemma-formal-functions-stalk}에 의해
$(R^pf_*\mathcal{F})_y^\wedge = 0$이고, 이는 $p > d$일 때 성립한다.
$R^pf_*\mathcal{F}$는 Proposition
\ref{coherent-proposition-proper-pushforward-coherent}에 의해 연접이고, 따라서
$R^pf_*\mathcal{F}_y$는 유한 $\mathcal{O}_{Y, y}$-가군이다. Nakayama
보조정리(Algebra, Lemma \ref{algebra-lemma-NAK})에 의해 국소환 위의
유한 가군의 완비화가 영이면 그 가군도 영이다. 따라서
$(R^pf_*\mathcal{F})_y = 0$이다.
\end{proof}

\section{형식 함수 정리의 응용}
\label{coherent-section-applications-formal-functions}


\noindent
필요에 따라 여기에 더 추가할 것이다. 지금은 뇌터인 경우의 유한 사상에
대한 다음 특징짓기가 필요하다.

\begin{lemma}
\label{coherent-lemma-characterize-finite}
(더 일반적인 판은 More on Morphisms, Lemma
\ref{more-morphisms-lemma-characterize-finite}를 보라.)
$f : X \to S$를 스킴의 사상이라 하자. $S$가 국소 뇌터라고 가정하자.
다음은 동치이다.
\begin{enumerate}
\item $f$는 유한이다.
\item $f$는 고유이고 유한 올들을 갖는다.
\end{enumerate}
\end{lemma}

\begin{proof}
Morphisms, Lemma \ref{morphisms-lemma-finite-proper}에 의해 유한 사상은
고유이다. Morphisms, Lemma
\ref{morphisms-lemma-finite-quasi-finite}에 의해 유한 사상은 준유한이다.
Morphisms, Lemma \ref{morphisms-lemma-quasi-finite}를 보면 준유한 사상은
유한 올들을 갖는다. 따라서 유한 사상은 고유이고 유한 올들을 갖는다.

\medskip\noindent
$f$가 고유이고 유한 올들을 갖는다고 가정하자. $f$가 유한임을 보이고자
한다. 실제로 $f$가 아핀임을 증명하면 충분하다. $f$가 아핀이면
Morphisms, Lemma \ref{morphisms-lemma-integral-universally-closed}에 의해
$f$가 정수적이고, 이어서 Morphisms, Lemma
\ref{morphisms-lemma-finite-integral}에 의해 $f$가 유한이기 때문이다.

\medskip\noindent
$f$가 아핀임을 보이기 위해 $S$가 아핀이라고 가정해도 되고, 목표는
$X$도 아핀임을 보이는 것이다. $f$가 고유이므로 $X$는 분리이고
준콤팩트이다. 따라서 Lemma
\ref{coherent-lemma-quasi-separated-h1-zero-covering}의 판정법을 사용하여 $X$가
아핀임을 증명할 수 있다. 이를 위해
$\mathcal{I} \subset \mathcal{O}_X$를 유한형 아이디얼 층이라 하자.
특히 $\mathcal{I}$는 $X$ 위의 연접층이다. Lemma
\ref{coherent-lemma-higher-direct-images-zero-finite-fibre}에 의해
$R^1f_*\mathcal{I}_s = 0$이고, 이는 모든 $s \in S$에 대해 성립한다.
다시 말해 $R^1f_*\mathcal{I} = 0$이다. 따라서 $f$에 대한 Leray
스펙트럼 수열에서
$H^1(X , \mathcal{I}) = H^1(S, f_*\mathcal{I})$임을 알 수 있다.
$S$는 아핀이고 $f_*\mathcal{I}$는 준연접이므로(Schemes, Lemma
\ref{schemes-lemma-push-forward-quasi-coherent}), Lemma
\ref{coherent-lemma-quasi-coherent-affine-cohomology-zero}에 의해 원하는 대로
$H^1(S, f_*\mathcal{I}) = 0$이다. 따라서 원하는 대로
$H^1(X, \mathcal{I}) = 0$이다.
\end{proof}

\noindent
그 결과로 다음 유용한 명제를 얻는다.

\begin{lemma}
\label{coherent-lemma-proper-finite-fibre-finite-in-neighbourhood}
\begin{slogan}
고유 사상은 유한 올의 근방에서 유한이다.
\end{slogan}
(더 일반적인 판은 More on Morphisms, Lemma
\ref{more-morphisms-lemma-proper-finite-fibre-finite-in-neighbourhood}를 보라.)
$f : X \to S$를 스킴의 사상이라 하고 $s \in S$라 하자. 다음을
가정하자.
\begin{enumerate}
\item $S$는 국소 뇌터이다.
\item $f$는 고유이다.
\item $f^{-1}(\{s\})$는 유한집합이다.
\end{enumerate}
그러면 열린 근방 $V \subset S$가 존재하여 $s$를 포함하고,
$f|_{f^{-1}(V)} : f^{-1}(V) \to V$는 유한이다.
\end{lemma}

\begin{proof}
Morphisms, Lemma \ref{morphisms-lemma-finite-fibre}에 의해 $f$는
$f^{-1}(\{s\})$의 모든 점에서 준유한이다. Morphisms, Lemma
\ref{morphisms-lemma-quasi-finite-points-open}에 의해 $f$가 준유한인
점들의 집합은 어떤 열린집합 $U \subset X$이다.
$Z = X \setminus U$로 놓자. 그러면 $s \not \in f(Z)$이다. $f$가
고유이므로 집합 $f(Z) \subset S$는 닫혀 있다. 열린 근방
$V \subset S$를 아무거나 택하여 $s$를 포함하고
$Z \cap V = \emptyset$이 되게 하자.
그러면 $f^{-1}(V) \to V$는 국소 준유한이고 고유이다. 따라서 이는
준유한이고(Morphisms, Lemma
\ref{morphisms-lemma-quasi-finite-locally-quasi-compact}), 유한 올들을
가지며(Morphisms, Lemma \ref{morphisms-lemma-quasi-finite}), 따라서
Lemma \ref{coherent-lemma-characterize-finite}에 의해 유한이다.
\end{proof}

\begin{lemma}
\label{coherent-lemma-ample-on-fibre}
$f : X \to Y$를 스킴의 고유 사상이라 하고 $Y$가 뇌터라고 하자.
$\mathcal{L}$을 가역 $\mathcal{O}_X$-가군이라 하고, $\mathcal{F}$를
연접 $\mathcal{O}_X$-가군이라 하자. 점 $y \in Y$가 주어져서
$\mathcal{L}_y$가 $X_y$ 위에서 풍부하다고 하자. 그러면 어떤
$d_0$가 존재하여 모든 $d \geq d_0$에 대해
$$
R^pf_*(\mathcal{F} \otimes_{\mathcal{O}_X} \mathcal{L}^{\otimes d})_y = 0
\text{ for }p > 0
$$
이고 사상
$$
f_*(\mathcal{F} \otimes_{\mathcal{O}_X} \mathcal{L}^{\otimes d})_y
\longrightarrow
H^0(X_y, \mathcal{F}_y \otimes_{\mathcal{O}_{X_y}} \mathcal{L}_y^{\otimes d})
$$
은 전사이다.
\end{lemma}

\begin{proof}
$\mathcal{O}_{Y, y}$가 뇌터 국소환임을 주목하자. 표준 사상
$c : \Spec(\mathcal{O}_{Y, y}) \to Y$를 생각하자. Schemes, Equation
(\ref{schemes-equation-canonical-morphism})을 보라. 이 사상은 국소환들을
동일시하므로 평탄하다. 잠시
$f' : X' \to \Spec(\mathcal{O}_{Y, y})$로 $f$를 이 국소환으로
밑변환한 사상을 나타내자. Lemma
\ref{coherent-lemma-flat-base-change-cohomology}에 의해
$c^*R^pf_*\mathcal{F} = R^pf'_*\mathcal{F}'$임을 알 수 있다. 또한 올
$X_y$와 $X'_y$는 동일시된다. 따라서 $Y = \Spec(A)$가 뇌터 국소환
$(A, \mathfrak m, \kappa)$의 스펙트럼이고 $y \in Y$가
$\mathfrak m$에 대응한다고 가정할 수 있다. 이 경우
$R^pf_*(\mathcal{F} \otimes_{\mathcal{O}_X} \mathcal{L}^{\otimes d})_y =
H^p(X, \mathcal{F} \otimes_{\mathcal{O}_X} \mathcal{L}^{\otimes d})$이고,
이는 모든 $p \geq 0$에 대해 성립한다.
$f_y : X_y \to \Spec(\kappa)$로 사영을 나타내자.

\medskip\noindent
$B = \text{Gr}_\mathfrak m(A) =
\bigoplus_{n \geq 0} \mathfrak m^n/\mathfrak m^{n + 1}$로 놓자.
$\mathcal{B} = f_y^*\widetilde{B}$라는 준연접 등급
$\mathcal{O}_{X_y}$-대수의 층을 생각하자. Section
\ref{coherent-section-theorem-formal-functions}의 표기를 사용하되, $I$를
$\mathfrak m$으로 바꿀 것이다. $X_y$는 $X$의 닫힌 부분스킴으로서
$\mathfrak m\mathcal{O}_X$에 의해 잘려 나오므로
$\mathfrak m^n\mathcal{F}/\mathfrak m^{n + 1}\mathcal{F}$를 연접
$\mathcal{O}_{X_y}$-가군으로 볼 수 있다. Lemma
\ref{coherent-lemma-i-star-equivalence}를 보라. 그러면
$\bigoplus_{n \geq 0} \mathfrak m^n\mathcal{F}/\mathfrak m^{n + 1}\mathcal{F}$는
유한형 준연접 등급 $\mathcal{B}$-가군이다. 그 이유는 이것이
$\mathcal{B}$ 위에서 차수 영의 성분으로 생성되고, 차수 영의 성분은
$\mathcal{F}_y = \mathcal{F}/\mathfrak m \mathcal{F}$이며 이는 연접
$\mathcal{O}_{X_y}$-가군이기 때문이다. 따라서 Lemma
\ref{coherent-lemma-graded-finiteness}의 (2)에 의해
$$
H^p(X_y, \mathfrak m^n \mathcal{F}/ \mathfrak m^{n + 1}\mathcal{F}
\otimes_{\mathcal{O}_{X_y}} \mathcal{L}_y^{\otimes d}) = 0
$$
이고, 이는 모든 $p > 0$, $d \geq d_0$, $n \geq 0$에 대해 성립한다.
Lemma \ref{coherent-lemma-relative-affine-cohomology}에 의해 이것은
$
H^p(X, \mathfrak m^n \mathcal{F}/ \mathfrak m^{n + 1}\mathcal{F}
\otimes_{\mathcal{O}_X} \mathcal{L}^{\otimes d}) = 0
$
가 모든 $p > 0$, $d \geq d_0$, $n \geq 0$에 대해 성립한다는 것과
같다.

\medskip\noindent
연접 $\mathcal{O}_X$-가군들의 짧은 완전열
$$
0 \to \mathfrak m^n\mathcal{F}/\mathfrak m^{n + 1} \mathcal{F}
\to \mathcal{F}/\mathfrak m^{n + 1} \mathcal{F}
\to \mathcal{F}/\mathfrak m^n \mathcal{F} \to 0
$$
을 생각하자. $\mathcal{L}^{\otimes d}$와의 텐서곱은 완전 함자이므로
다음 짧은 완전열들을 얻는다.
$$
0 \to
\mathfrak m^n\mathcal{F}/\mathfrak m^{n + 1} \mathcal{F}
\otimes_{\mathcal{O}_X} \mathcal{L}^{\otimes d}
\to \mathcal{F}/\mathfrak m^{n + 1} \mathcal{F}
\otimes_{\mathcal{O}_X} \mathcal{L}^{\otimes d}
\to \mathcal{F}/\mathfrak m^n \mathcal{F}
\otimes_{\mathcal{O}_X} \mathcal{L}^{\otimes d} \to 0
$$
긴 코호몰로지 완전열과 위의 소멸을 사용하면 귀납법으로 다음을 얻는다.
\begin{enumerate}
\item $H^p(X, \mathcal{F}/\mathfrak m^n \mathcal{F}
\otimes_{\mathcal{O}_X} \mathcal{L}^{\otimes d}) = 0$이고, 이는 모든
$p > 0$, $d \geq d_0$, $n \geq 0$에 대해 성립한다.
\item 사상 $H^0(X, \mathcal{F}/\mathfrak m^n \mathcal{F}
\otimes_{\mathcal{O}_X} \mathcal{L}^{\otimes d}) \to
H^0(X_y, \mathcal{F}_y \otimes_{\mathcal{O}_{X_y}} \mathcal{L}_y^{\otimes d})$
은 모든 $d \geq d_0$ 및 $n \geq 1$에 대해 전사이다.
\end{enumerate}
형식 함수 정리(Theorem \ref{coherent-theorem-formal-functions})에 의해
$\mathfrak m$-진 완비화된
$H^p(X, \mathcal{F} \otimes_{\mathcal{O}_X} \mathcal{L}^{\otimes d})$는
모든 $d \geq d_0$ 및 $p > 0$에 대해 영이다.
$H^p(X, \mathcal{F} \otimes_{\mathcal{O}_X} \mathcal{L}^{\otimes d})$는
Lemma \ref{coherent-lemma-proper-over-affine-cohomology-finite}에 의해 유한
$A$-가군이므로, Nakayama 보조정리(Algebra, Lemma
\ref{algebra-lemma-NAK})에 의해
$H^p(X, \mathcal{F} \otimes_{\mathcal{O}_X} \mathcal{L}^{\otimes d})$는
모든 $d \geq d_0$ 및 $p > 0$에 대해 영이다. $p = 0$인 경우에는
Lemma \ref{coherent-lemma-ML-cohomology-powers-ideal}의 (3)에서 사상
$H^0(X, \mathcal{F} \otimes_{\mathcal{O}_X} \mathcal{L}^{\otimes d}) \to
H^0(X_y, \mathcal{F}_y \otimes_{\mathcal{O}_{X_y}} \mathcal{L}_y^{\otimes d})$
이 전사임을 얻으며, 이것이 보조정리의 마지막 명제를 준다.
\end{proof}

\begin{lemma}
\label{coherent-lemma-ample-in-neighbourhood}
(더 일반적인 판은 More on Morphisms, Lemma
\ref{more-morphisms-lemma-ample-in-neighbourhood}를 보라.)
$f : X \to Y$를 스킴의 고유 사상이라 하고 $Y$가 뇌터라고 하자.
$\mathcal{L}$을 가역 $\mathcal{O}_X$-가군이라 하자. 점 $y \in Y$가
주어져서 $\mathcal{L}_y$가 $X_y$ 위에서 풍부하다고 하자. 그러면
열린 근방 $V \subset Y$가 존재하여 $y$를 포함하고,
$\mathcal{L}|_{f^{-1}(V)}$는 $f^{-1}(V)/V$ 위에서 풍부하다.
\end{lemma}

\begin{proof}
Lemma \ref{coherent-lemma-ample-on-fibre}에서 얻은 $d_0$를 택하되
$\mathcal{F} = \mathcal{O}_X$인 경우에 적용하자. $d \geq d_0$를 택하여
어떤 $r \geq 0$와 단면들
$s_{y, 0}, \ldots, s_{y, r} \in H^0(X_y, \mathcal{L}_y^{\otimes d})$를
찾을 수 있게 하고, 이 단면들이
닫힌 몰입
$$
\varphi_y =
\varphi_{\mathcal{L}_y^{\otimes d}, (s_{y, 0}, \ldots, s_{y, r})} :
X_y \to \mathbf{P}^r_{\kappa(y)}.
$$
을 정의하게 하자. Morphisms, Lemma
\ref{morphisms-lemma-finite-type-over-affine-ample-very-ample}에 의해 이는
가능하다. 다만 $\varphi_y$가 닫힌 몰입임을 알기 위해 Morphisms,
Lemma \ref{morphisms-lemma-image-proper-scheme-closed}도 사용하고,
가역층과 단면들로 사영공간으로 가는 사상을 기술하기 위해 Constructions,
Section \ref{constructions-section-projective-space}도 사용한다.
$d_0$의 선택에 의해 $Y$를 $y$의 열린 근방으로 바꾼 뒤,
$s_0, \ldots, s_r \in H^0(X, \mathcal{L}^{\otimes d})$를 택하여
$s_{y, 0}, \ldots, s_{y, r}$로 사상되게 할 수 있다.
$X_{s_i} \subset X$를 $s_i$가 $\mathcal{L}^{\otimes d}$의 생성원이 되는
열린 부분집합이라 하자. $s_{y, i}$들이 $\mathcal{L}_y^{\otimes d}$를
생성하므로 $X_y \subset U = \bigcup X_{s_i}$임을 알 수 있다.
$X \to Y$가 닫힌 사상이므로 $y \in V \subset Y$인 열린 근방이
존재하여 $f^{-1}(V) \subset U$가 된다. $Y$를 $V$로 바꾼 뒤에는
$s_i$들이 $\mathcal{L}^{\otimes d}$를 생성한다고 가정할 수 있다.
따라서 사상
$$
\varphi = \varphi_{\mathcal{L}^{\otimes d}, (s_0, \ldots, s_r)} :
X \longrightarrow \mathbf{P}^r_Y
$$
을 얻고, $\mathcal{L}^{\otimes d} \cong \varphi^*\mathcal{O}_{\mathbf{P}^r_Y}(1)$이며
이를 $y$로 밑변환하면
$\varphi_y$를 얻는다.

\medskip\noindent
약간의 요령으로 증명을 마치겠다. ``올바른'' 증명은 $\varphi$가 닫힌
몰입임을 $y$의 열린 근방으로 밑변환한 뒤 직접 보이는 방식으로 진행된다.
실제로 Lemma \ref{coherent-lemma-proper-finite-fibre-finite-in-neighbourhood}에 의해
$\varphi$는 올 $\mathbf{P}^r_{\kappa(y)}$의 열린 근방 위에서 유한이다.
여기서 이 올은 $\mathbf{P}^r_Y \to Y$의 $y$ 위의 올이다.
$\mathbf{P}^r_Y \to Y$가 닫힌 사상임을 사용하여 $Y$를 줄인 뒤에는
$\varphi$가 유한이라고 가정할 수 있다. 그러면
$\mathcal{L}^{\otimes d} \cong \varphi^*\mathcal{O}_{\mathbf{P}^r_Y}(1)$은
아주 일반적인 Morphisms, Lemma
\ref{morphisms-lemma-pullback-ample-tensor-relatively-ample}에 의해
풍부하다.
\end{proof}

\section{코호몰로지와 밑변환, III}
\label{coherent-section-cohomology-and-base-change-perfect}

\noindent
이 절에서는 매우 일반적인 현상의 가장 간단한 경우를 증명한다. 그
일반적인 현상은 Derived Categories of Schemes, Section
\ref{perfect-section-cohomology-and-base-change-perfect}에서 논의한다.
다음 보조정리를 대수의 언어로 옮긴 내용은 Remark
\ref{coherent-remark-explain-perfect-direct-image}를 보라.

\begin{lemma}
\label{coherent-lemma-perfect-direct-image}
$A$를 뇌터 환이라 하고 $S = \Spec(A)$로 놓자. $f : X \to S$를 스킴의
고유 사상이라 하자. $\mathcal{F}$를 연접 $\mathcal{O}_X$-가군이라 하고
$S$ 위에서 평탄이라고 하자. 그러면
\begin{enumerate}
\item $R\Gamma(X, \mathcal{F})$는 $D(A)$의 완전 대상이다.
\item 임의의 환 사상 $A \to A'$에 대해 밑변환 사상
$$
R\Gamma(X, \mathcal{F}) \otimes_A^{\mathbf{L}} A'
\longrightarrow
R\Gamma(X_{A'}, \mathcal{F}_{A'})
$$
은 동형사상이다.
\end{enumerate}
\end{lemma}

\begin{proof}
유한 아핀 열린 덮개 $X = \bigcup_{i = 1, \ldots, n} U_i$를 택하자.
Lemmas \ref{coherent-lemma-separated-case-relative-cech}와
\ref{coherent-lemma-base-change-complex}에 의해 \v{C}ech 복합체
$K^\bullet = {\check C}^\bullet(\mathcal{U}, \mathcal{F})$는
$$
K^\bullet \otimes_A A' = R\Gamma(X_{A'}, \mathcal{F}_{A'})
$$
을 모든 환 사상 $A \to A'$에 대해 만족한다.
$K_{alt}^\bullet = {\check C}_{alt}^\bullet(\mathcal{U}, \mathcal{F})$를
교대 \v{C}ech 복합체라 하자. Cohomology, Lemma
\ref{cohomology-lemma-alternating-usual}에 의해 호모토피 동치
$K_{alt}^\bullet \to K^\bullet$가 $A$-가군들 사이에 존재한다. 특히
$$
K_{alt}^\bullet \otimes_A A' = R\Gamma(X_{A'}, \mathcal{F}_{A'})
$$
도 성립한다. $\mathcal{F}$는 $A$ 위에서 평탄이므로 각
$K_{alt}^n$도 $A$ 위에서 평탄이다(Morphisms, Lemma
\ref{morphisms-lemma-flat-module-characterize}를 보라). 또한
$K_{alt}^\bullet$은 위로 유계이므로(이것이 교대 \v{C}ech 복합체로 바꾼
이유이다), 유도 텐서곱의 정의에 의해
$K_{alt}^\bullet \otimes_A A' = K_{alt}^\bullet \otimes_A^{\mathbf{L}} A'$이다
(More on Algebra, Section
\ref{more-algebra-section-derived-tensor-product}을 보라). Lemma
\ref{coherent-lemma-proper-over-affine-cohomology-finite}에 의해 코호몰로지 군
$H^i(K_{alt}^\bullet)$은 유한 $A$-가군이다. $K_{alt}^\bullet$이
유계이므로 $K_{alt}^\bullet$은 유사연접이다. More on Algebra, Lemma
\ref{more-algebra-lemma-Noetherian-pseudo-coherent}를 보라. 임의의
$A$-가군 $M$이 주어졌을 때 $A' = A \oplus M$으로 놓되 $M$이
제곱영 아이디얼이 되게 하자. 즉
$(a, m) \cdot (a', m') = (aa', am' + a'm)$이다. 위의 결과에 의해
$K_{alt}^\bullet \otimes_A^\mathbf{L} A'$는 차수 $0, \ldots, n$에서만
코호몰로지를 갖는다. 따라서
$K_{alt}^\bullet \otimes_A^\mathbf{L} M$도 차수 $0, \ldots, n$에서만
코호몰로지를 갖는다. 그러므로 $K_{alt}^\bullet$은 유한 Tor 차원을
갖는다. More on Algebra, Definition
\ref{more-algebra-definition-tor-amplitude}을 보라. 이제 More on Algebra,
Lemma \ref{more-algebra-lemma-perfect}를 적용하면 끝난다.
\end{proof}

\begin{remark}
\label{coherent-remark-explain-perfect-direct-image}
Lemma \ref{coherent-lemma-perfect-direct-image}의 한 결과로, 유한 사영
$A$-가군들의 유한 복합체 $M^\bullet$이 존재하여
$$
H^i(X_{A'}, \mathcal{F}_{A'}) = H^i(M^\bullet \otimes_A A')
$$
이 $A'$에 대해 함자적으로 성립한다. $\mathcal{F}$가 $A$ 위에서
평탄이라는 조건은 본질적이다. \cite{Hartshorne}을 보라.
\end{remark}

\section{연접 형식 가군}
\label{coherent-section-coherent-formal}

\noindent
아직 형식 스킴 이론을 쓸 수 없으므로, ``뇌터 진 형식 스킴 위의 연접
가군''이라는 개념을 대신할 표현을 조금 마련한다.

\medskip\noindent
$X$를 뇌터 스킴이라 하고 $\mathcal{I} \subset \mathcal{O}_X$를 준연접
아이디얼 층이라 하자. 다음을 만족하는 역계 $(\mathcal{F}_n)$을 생각할
것이다. 각 항은 연접 $\mathcal{O}_X$-가군이다.
\begin{enumerate}
\item $\mathcal{F}_n$은 $\mathcal{I}^n$에 의해 소멸된다.
\item 전이 사상들은 동형사상
$\mathcal{F}_{n + 1}/\mathcal{I}^n\mathcal{F}_{n + 1} \to \mathcal{F}_n$을
유도한다.
\end{enumerate}
이러한 역계들의 사상은 보통대로 정의한다. 이 역계들의 범주를
$\textit{Coh}(X, \mathcal{I})$로 나타내자. 이제 이 범주의 대상에 관한
여러 보조정리를 증명할 것이다. 실제로 뒤따르는 보조정리 대부분은
아핀인 경우의 범주에 대한 다음 기술에서 곧바로 나온다.

\begin{lemma}
\label{coherent-lemma-inverse-systems-affine}
$X = \Spec(A)$가 뇌터 환의 스펙트럼이고 $\mathcal{I}$가 아이디얼
$I \subset A$에 결부된 준연접 아이디얼 층이면,
$\textit{Coh}(X, \mathcal{I})$는 유한 $A^\wedge$-가군들의 범주와
동치이다. 여기서 $A^\wedge$는 $A$를 $I$에 관해 완비화한 것이다.
\end{lemma}

\begin{proof}
$\text{Mod}^{fg}_{A, I}$를 다음을 만족하는 역계 $(M_n)$의 범주라 하자.
각 항은 유한 $A$-가군이다. (1) $M_n$은 $I^n$에 의해 소멸되고, (2)
$M_{n + 1}/I^nM_{n + 1} = M_n$이다. $X$ 위의 연접층과 유한
$A$-가군 사이의 대응(Lemma \ref{coherent-lemma-coherent-Noetherian})에 의해
$\text{Mod}^{fg}_{A, I}$가 유한 $A^\wedge$-가군들의 범주와 동치임을
보이면 충분하다. 이를 위해 대상 $(M_n)$이 $\text{Mod}^{fg}_{A, I}$에
주어졌을 때 가군
$$
M = \lim M_n
$$
이 유한 $A^\wedge$-가군이고 $M/I^nM = M_n$임을 증명하면 충분하다.
전이 사상들은 전사이므로 $M \to M_1$이 전사임을 알 수 있다.
$x_1, \ldots, x_t \in M$을 택하여 그 상들이 $M_1$을 생성하게 하자.
그러면 계의 사상 $(A/I^n)^{\oplus t} \to M_n$이 유도된다. Nakayama
보조정리(Algebra, Lemma \ref{algebra-lemma-NAK})에 의해 이 사상들은
전사이다. $K_n \subset (A/I^n)^{\oplus t}$를 그 핵이라 하자.
성질 (2)에 의해 $K_{n + 1} \to K_n$은 전사이고, 특히 계 $(K_n)$은
Mittag-Leffler 조건을 만족한다. Homology, Lemma
\ref{homology-lemma-Mittag-Leffler}에 의해 완전열
$0 \to K \to (A^\wedge)^{\oplus t} \to M \to 0$을 얻으며
$K = \lim K_n$이다. 따라서 $M$은 유한 $A^\wedge$-가군이다.
$K \to K_n$이 전사이므로 원하는 대로
$$
M/I^nM = \Coker(K \to (A/I^n)^{\oplus t}) =
(A/I^n)^{\oplus t}/K_n = M_n
$$
이다.
\end{proof}

\begin{lemma}
\label{coherent-lemma-inverse-systems-abelian}
$X$를 뇌터 스킴이라 하고 $\mathcal{I} \subset \mathcal{O}_X$를 준연접
아이디얼 층이라 하자.
\begin{enumerate}
\item 범주 $\textit{Coh}(X, \mathcal{I})$는 아벨 범주이다.
\item 열린집합 $U \subset X$에 대해 제한 함자
$\textit{Coh}(X, \mathcal{I}) \to \textit{Coh}(U, \mathcal{I}|_U)$는
완전하다.
\item $\textit{Coh}(X, \mathcal{I})$에서의 완전성은 $X$의 열린 덮개의
원소들로 제한하여 확인할 수 있다.
\end{enumerate}
\end{lemma}

\begin{proof}
$\alpha =(\alpha_n) : (\mathcal{F}_n) \to (\mathcal{G}_n)$를
$\textit{Coh}(X, \mathcal{I})$의 사상이라 하자. $\alpha$의 여핵은
역계 $(\Coker(\alpha_n))$이다(세부사항은 생략한다). 핵을 기술하기 위해
$$
\mathcal{K}'_{l, m} = \Im(\Ker(\alpha_l) \to \mathcal{F}_m)
$$
으로 $l \geq m$에 대해 놓자. 다음을 주장한다.
\begin{enumerate}
\item[(a)] 역계 $(\mathcal{K}'_{l, m})_{l \geq m}$은 궁극적으로
상수이고, 그 값을 $\mathcal{K}'_m$이라 하자.
\item[(b)] 계 $(\mathcal{K}'_m/\mathcal{I}^n\mathcal{K}'_m)_{m \geq n}$은
궁극적으로 상수이고, 그 값을 $\mathcal{K}_n$이라 하자.
\item[(c)] 계 $(\mathcal{K}_n)$은 $\textit{Coh}(X, \mathcal{I})$의
대상을 이룬다.
\item[(d)] 이 대상은 $\alpha$의 핵이다.
\end{enumerate}
(a), (b), (c)를 보이기 위해 아핀 국소적으로 작업해도 된다. 즉
$X = \Spec(A)$이고 $\mathcal{I}$가 아이디얼 $I \subset A$에 대응한다고
하자. Lemma \ref{coherent-lemma-inverse-systems-affine}에 의해 $\alpha$는 사상
$f : M \to N$에 대응하며, 이는 유한 $A^\wedge$-가군들의 사상이다.
$K = \Ker(f)$로 나타내자. $A^\wedge$가 뇌터 환임을 주목하자(Algebra,
Lemma \ref{algebra-lemma-completion-Noetherian-Noetherian}). 정수
$c \geq 0$를 택하여 $K \cap I^n M \subset I^{n - c}K$가
$n \geq c$에 대해 성립하게 하고(Algebra, Lemma
\ref{algebra-lemma-Artin-Rees}), 사상 $f$와 아이디얼
$I^\wedge = IA^\wedge$에 대해 Algebra, Lemma
\ref{algebra-lemma-map-AR}도 만족하게 하자. 그러면
$\mathcal{K}'_{l, m}$은 다음 $A$-가군에 대응한다.
$$
K'_{l, m} = \frac{a^{-1}(I^lN) + I^mM}{I^mM} =
\frac{K + I^{l - c}f^{-1}(I^cN) + I^mM}{I^mM} =
\frac{K + I^mM}{I^mM}
$$
여기서 마지막 등식은 $l \geq m + c$이면 성립한다. 따라서
$\mathcal{K}'_m$은 $A$-가군 $K/K \cap I^mM$에 대응하고,
$\mathcal{K}'_m/\mathcal{I}^n\mathcal{K}'_m$은
$$
\frac{K}{K \cap I^mM + I^nK} = \frac{K}{I^nK}
$$
에 대응한다. 이는 $m \geq n + c$이면 위에서 택한 $c$에 의해
성립한다. 따라서 $\mathcal{K}_n$은 $K/I^nK$에 대응한다.

\medskip\noindent
(d)를 증명하자. 위의 아핀 기술에서 합성
$(\mathcal{K}_n) \to (\mathcal{F}_n) \to (\mathcal{G}_n)$이 영임은
분명하다. $\beta : (\mathcal{H}_n) \to (\mathcal{F}_n)$를
$\alpha \circ \beta = 0$인 사상이라 하자. 그러면
$\mathcal{H}_l \to \mathcal{F}_l$의 상은 $\Ker(\alpha_l)$에 포함된다.
$\mathcal{H}_m = \mathcal{H}_l/\mathcal{I}^m\mathcal{H}_l$이 $l \geq m$에
대해 성립하므로 사상들의
계 $\mathcal{H}_m \to \mathcal{K}'_{l, m}$을 얻는다. 따라서 사상
$\mathcal{H}_m \to \mathcal{K}_m'$을 얻는다.
$\mathcal{H}_n = \mathcal{H}_m/\mathcal{I}^n\mathcal{H}_m$이므로 사상들의
계 $\mathcal{H}_n \to \mathcal{K}'_m/\mathcal{I}^n\mathcal{K}'_m$을 얻고,
따라서 원하는 사상 $\mathcal{H}_n \to \mathcal{K}_n$을 얻는다.

\medskip\noindent
(1)의 증명을 마치려면 $\Coim = \Im$임을 $\textit{Coh}(X, \mathcal{I})$에서
아직 보여야 한다. 위에서 아핀인 경우 핵과 여핵을
취하는 것이 $\textit{Coh}(X, \mathcal{I})$를 가군의 범주로 기술하는
것과 가환함을 보았다. 가군의 범주에서는 $\Im = \Coim$이 성립하므로
이것은 $\Coim = \Im$을 $\textit{Coh}(X, \mathcal{I})$에서 준다. 보조정리의
(2)와 (3)은 우리가 구성한 핵과 여핵에서 즉시 따른다.
\end{proof}

\begin{lemma}
\label{coherent-lemma-inverse-systems-surjective}
$X$를 뇌터 스킴이라 하고 $\mathcal{I} \subset \mathcal{O}_X$를 준연접
아이디얼 층이라 하자. 사상
$(\mathcal{F}_n) \to (\mathcal{G}_n)$이
$\textit{Coh}(X, \mathcal{I})$에서 전사일 필요충분조건은
$\mathcal{F}_1 \to \mathcal{G}_1$이 전사인 것이다.
\end{lemma}

\begin{proof}
생략한다. 힌트: 아핀 열린집합들 위에서 살펴보고, Lemma
\ref{coherent-lemma-inverse-systems-affine}와 Algebra, Lemma
\ref{algebra-lemma-NAK}를 사용하라.
\end{proof}

\noindent
$X$를 뇌터 스킴이라 하고 $\mathcal{I} \subset \mathcal{O}_X$를 준연접
아이디얼 층이라 하자. 다음 함자가 있다.
\begin{equation}
\label{coherent-equation-completion-functor}
\textit{Coh}(\mathcal{O}_X) \longrightarrow \textit{Coh}(X, \mathcal{I}), \quad
\mathcal{F} \longmapsto \mathcal{F}^\wedge
\end{equation}
이 함자는 연접 $\mathcal{O}_X$-가군 $\mathcal{F}$에 대상
$\mathcal{F}^\wedge = (\mathcal{F}/\mathcal{I}^n\mathcal{F})$를 결부시키며,
이 대상은 $\textit{Coh}(X, \mathcal{I})$에 속한다.

\begin{lemma}
\label{coherent-lemma-exact}
함자 (\ref{coherent-equation-completion-functor})는 완전하다.
\end{lemma}

\begin{proof}
$X$ 위에서 국소적으로 확인하면 충분하다. 따라서 $X$가 아핀이라고,
즉 Lemma \ref{coherent-lemma-inverse-systems-affine}의 상황이라고 가정해도 된다.
이 함자는 함자 $\text{Mod}^{fg}_A \to \text{Mod}^{fg}_{A^\wedge}$이다.
이는 유한 $A$-가군 $M$에 완비화 $M^\wedge$를 결부시킨다. 따라서 결과는
Algebra, Lemma \ref{algebra-lemma-completion-flat}에서 따른다.
\end{proof}

\begin{lemma}
\label{coherent-lemma-completion-internal-hom}
$X$를 뇌터 스킴이라 하고 $\mathcal{I} \subset \mathcal{O}_X$를 준연접
아이디얼 층이라 하자. $\mathcal{F}$, $\mathcal{G}$를 연접
$\mathcal{O}_X$-가군이라 하자.
$\mathcal{H} = \SheafHom_{\mathcal{O}_X}(\mathcal{G}, \mathcal{F})$로 놓자.
그러면
$$
\lim H^0(X, \mathcal{H}/\mathcal{I}^n\mathcal{H}) =
\Mor_{\textit{Coh}(X, \mathcal{I})}
(\mathcal{G}^\wedge, \mathcal{F}^\wedge).
$$
\end{lemma}

\begin{proof}
이를 증명하기 위해 $X$ 위에서 아핀 국소적으로 작업해도 된다. 따라서
$X = \Spec(A)$라고 가정하고 $\mathcal{F}$, $\mathcal{G}$가 유한
$A$-가군 $M$과 $N$으로 주어진다고 하자. 그러면 $\mathcal{H}$는 유한
$A$-가군 $H = \Hom_A(M, N)$에 대응한다. 보조정리의 명제는 Lemma
\ref{coherent-lemma-inverse-systems-affine}의 동치를 통해
$$
H^\wedge = \Hom_{A^\wedge}(M^\wedge, N^\wedge)
$$
이라는 명제가 된다. Algebra, Lemma
\ref{algebra-lemma-completion-flat}를 세 번 사용하면
$$
H^\wedge = \Hom_A(M, N) \otimes_A A^\wedge =
\Hom_{A^\wedge}(M \otimes_A A^\wedge, N \otimes_A A^\wedge) =
\Hom_{A^\wedge}(M^\wedge, N^\wedge)
$$
이다. 여기서 둘째 등식에는 $A^\wedge$가 $A$ 위에서 평탄임을 사용했다
(More on Algebra, Lemma
\ref{more-algebra-lemma-pseudo-coherence-and-base-change-ext}를 보라).
이로써 보조정리가 증명되었다.
\end{proof}

\noindent
$X$를 뇌터 스킴이라 하자. $\mathcal{I} \subset \mathcal{O}_X$를 준연접
아이디얼 층이라 하자. 대상 $(\mathcal{F}_n)$이
$\textit{Coh}(X, \mathcal{I})$에 속한다고 하자. 이 대상이
{\it $\mathcal{I}$-멱 비틀림} 또는 {\it $\mathcal{I}$의 한 멱에 의해
소멸된다}고 함은 어떤 $c \geq 1$이 존재하여
$\mathcal{F}_n = \mathcal{F}_c$가 모든 $n \geq c$에 대해 성립한다는
뜻이다. 이 경우 $(\mathcal{F}_n)$이 {\it $\mathcal{I}^c$에 의해
소멸된다}고 말한다. $X = \Spec(A)$가 아핀이면 Lemma
\ref{coherent-lemma-inverse-systems-affine}의 동치를 통해 이러한 대상들은 유한
$A$-가군 가운데 $I$의 한 멱 또는 $I^c$에 의해 소멸되는 것들과 정확히
대응한다.

\begin{lemma}
\label{coherent-lemma-existence-easy}
$X$를 뇌터 스킴이라 하고 $\mathcal{I} \subset \mathcal{O}_X$를 준연접
아이디얼 층이라 하자. $\mathcal{G}$를 연접 $\mathcal{O}_X$-가군이라
하자. $(\mathcal{F}_n)$을 $\textit{Coh}(X, \mathcal{I})$의 대상이라 하자.
\begin{enumerate}
\item $\alpha : (\mathcal{F}_n) \to \mathcal{G}^\wedge$가 핵과 여핵이
$\mathcal{I}$의 한 멱에 의해 소멸되는 사상이면, 다음을 만족하는
삼중항 $(\mathcal{F}, a, \beta)$가 유일한 동형사상까지 유일하게
존재한다.
\begin{enumerate}
\item $\mathcal{F}$는 연접 $\mathcal{O}_X$-가군이다.
\item $a : \mathcal{F} \to \mathcal{G}$는 $\mathcal{O}_X$-가군 사상이고,
그 핵과 여핵은 $\mathcal{I}$의 한 멱에 의해 소멸된다.
\item $\beta : (\mathcal{F}_n) \to \mathcal{F}^\wedge$는 동형사상이다.
\item $\alpha = a^\wedge \circ \beta$이다.
\end{enumerate}
\item $\alpha : \mathcal{G}^\wedge \to (\mathcal{F}_n)$가 핵과 여핵이
$\mathcal{I}$의 한 멱에 의해 소멸되는 사상이면, 다음을 만족하는
삼중항 $(\mathcal{F}, a, \beta)$가 유일한 동형사상까지 유일하게
존재한다.
\begin{enumerate}
\item $\mathcal{F}$는 연접 $\mathcal{O}_X$-가군이다.
\item $a : \mathcal{G} \to \mathcal{F}$는 $\mathcal{O}_X$-가군 사상이고,
그 핵과 여핵은 $\mathcal{I}$의 한 멱에 의해 소멸된다.
\item $\beta : \mathcal{F}^\wedge \to (\mathcal{F}_n)$는 동형사상이다.
\item $\alpha = \beta \circ a^\wedge$이다.
\end{enumerate}
\end{enumerate}
\end{lemma}

\begin{proof}
(1)의 증명. 유일성 때문에 $(\mathcal{F}, a, \beta)$를 $X$ 위에서
Zariski 국소적으로 구성하면 충분하다. 따라서 $X = \Spec(A)$이고
$\mathcal{I}$가 아이디얼 $I \subset A$에 대응한다고 가정해도 된다.
이 상황에는 Lemma \ref{coherent-lemma-inverse-systems-affine}가 적용된다.
$M'$을 유한 $A^\wedge$-가군이라 하되 $(\mathcal{F}_n)$에 대응하는 것으로
택하자. $N$을 유한 $A$-가군이라 하되 $\mathcal{G}$에 대응하는 것으로
택하자. 그러면
$\alpha$는 사상
$$
\varphi : M' \longrightarrow N^\wedge
$$
에 대응하며, 그 핵과 여핵은 $I^t$에 의해 소멸된다. 여기서 $t$는 어떤
정수이다.
$N^\wedge = N \otimes_A A^\wedge$임을 상기하자(Algebra, Lemma
\ref{algebra-lemma-completion-tensor}). More on Algebra, Lemma
\ref{more-algebra-lemma-application-formal-glueing}에 의해 $A$-가군 사상
$\psi : M \to N$이 존재하고 그 핵과 여핵은 $I$-멱 비틀림이다. 또한
동형사상 $M \otimes_A A^\wedge = M'$이 존재하여 $\varphi$와 양립한다.
$N$과 $M'$은 유한
가군이므로 $M$도 유한 $A$-가군이다. More on Algebra, Remark
\ref{more-algebra-remark-formal-glueing-algebras}를 보라. 따라서
$M \otimes_A A^\wedge = M^\wedge$이다. 이렇게 얻은 삼중항
$(M, N \to M, M^\wedge \to M')$이 유일한 동형사상까지 유일하다는
확인은 생략한다.

\medskip\noindent
(2)의 증명도 완전히 같으므로 생략한다.
\end{proof}

\begin{lemma}
\label{coherent-lemma-torsion-hom-ext}
$X$를 뇌터 스킴이라 하고 $\mathcal{I} \subset \mathcal{O}_X$를 준연접
아이디얼 층이라 하자. $\textit{Coh}(X, \mathcal{I})$의 대상 가운데
$\mathcal{I}$의 한 멱에 의해 소멸되는 것은 모두
(\ref{coherent-equation-completion-functor})의 본질적 상에 속한다. 더 나아가
$\mathcal{F}$, $\mathcal{G}$가 $\textit{Coh}(\mathcal{O}_X)$에 속하고
$\mathcal{F}$ 또는 $\mathcal{G}$가 $\mathcal{I}$의 한 멱에 의해
소멸되면, 사상들
$$
\xymatrix{
\Hom_X(\mathcal{F}, \mathcal{G}) \ar[d] &
\Ext_X(\mathcal{F}, \mathcal{G}) \ar[d] \\
\Hom_{\textit{Coh}(X, \mathcal{I})}(\mathcal{F}^\wedge, \mathcal{G}^\wedge) &
\Ext_{\textit{Coh}(X, \mathcal{I})}(\mathcal{F}^\wedge, \mathcal{G}^\wedge)
}
$$
은 동형사상이다.
\end{lemma}

\begin{proof}
$(\mathcal{F}_n)$이 $\textit{Coh}(X, \mathcal{I})$의 대상이고
$\mathcal{I}^c$에 의해 소멸된다고 하자. 여기서 $c \geq 1$이다. 그러면
$\mathcal{F}_n \to \mathcal{F}_c$는 $n \geq c$에 대해 동형사상이다.
따라서 $\mathcal{F} = \mathcal{F}_c$로 놓으면
$\mathcal{F}^\wedge \cong (\mathcal{F}_n)$임을 알 수 있다. 이로써 첫째
주장이 증명되었다.

\medskip\noindent
$\mathcal{F}$, $\mathcal{G}$를 $\textit{Coh}(\mathcal{O}_X)$의 대상이라
하고, $\mathcal{F}$ 또는 $\mathcal{G}$가 $\mathcal{I}^c$에 의해
소멸된다고 하자. 여기서 $c \geq 1$이다.
$\mathcal{H} = \SheafHom_{\mathcal{O}_X}(\mathcal{G}, \mathcal{F})$는
연접 $\mathcal{O}_X$-가군이고 $\mathcal{I}^c$에 의해 소멸된다. 따라서
$$
\Hom_X(\mathcal{G}, \mathcal{F}) =
H^0(X, \mathcal{H}) =
\lim H^0(X, \mathcal{H}/\mathcal{I}^n\mathcal{H}) =
\Mor_{\textit{Coh}(X, \mathcal{I})}
(\mathcal{G}^\wedge, \mathcal{F}^\wedge).
$$
임을 알 수 있다. Lemma \ref{coherent-lemma-completion-internal-hom}를 보라.
이로써 준동형사상에 관한 명제가 증명되었다.

\medskip\noindent
표기 $\Ext$는 Homology, Section \ref{homology-section-extensions}에서
정의한 확장을 뜻한다. $\Ext$ 위의 사상이 단사임은 $\Hom$ 위의 사상이
전단사임에서 즉시 따른다. 전사성을 위해 $\mathcal{F}$가 $I$의 한 멱에
의해 소멸된다고 가정하자. 그러면 Lemma
\ref{coherent-lemma-existence-easy}의 (1)에 의해 확장
$$
0 \to \mathcal{G}^\wedge \to (\mathcal{E}_n) \to \mathcal{F}^\wedge \to 0
$$
이 $\textit{Coh}(U, I\mathcal{O}_U)$ 안에서 주어졌을 때, 사상
$\mathcal{G}^\wedge \to (\mathcal{E}_n)$은
$\mathcal{G} \to \mathcal{E}^\wedge$와 동형이다. 여기서 어떤 사상
$\mathcal{G} \to \mathcal{E}$가 $\textit{Coh}(\mathcal{O}_U)$에 속한다.
다른 경우도
마찬가지이다.
\end{proof}

\begin{lemma}
\label{coherent-lemma-finite-over-rees-algebra}
$X$를 뇌터 스킴이라 하고 $\mathcal{I} \subset \mathcal{O}_X$를 준연접
아이디얼 층이라 하자. $(\mathcal{F}_n)$이
$\textit{Coh}(X, \mathcal{I})$의 대상이면
$\bigoplus \Ker(\mathcal{F}_{n + 1} \to \mathcal{F}_n)$은 유한형 등급
준연접 $\bigoplus \mathcal{I}^n/\mathcal{I}^{n + 1}$-가군이다.
\end{lemma}

\begin{proof}
문제는 $X$ 위에서 국소적이므로 $X$가 아핀이라고, 즉 Lemma
\ref{coherent-lemma-inverse-systems-affine}의 상황이라고 가정해도 된다. 이 경우
$(\mathcal{F}_n)$이 유한 $A^\wedge$-가군 $M$에 대응하면,
$\bigoplus \Ker(\mathcal{F}_{n + 1} \to \mathcal{F}_n)$은
$\bigoplus I^nM/I^{n + 1}M$에 대응하고, 이는 분명
$\bigoplus I^n/I^{n + 1}$ 위의 유한 가군이다.
\end{proof}

\begin{lemma}
\label{coherent-lemma-inverse-systems-pullback}
$f : X \to Y$를 뇌터 스킴들의 사상이라 하자.
$\mathcal{J} \subset \mathcal{O}_Y$를 준연접 아이디얼 층이라 하고
$\mathcal{I} = f^{-1}\mathcal{J} \mathcal{O}_X$로 놓자. 그러면 오른쪽
완전 함자
$$
f^* : \textit{Coh}(Y, \mathcal{J}) \longrightarrow \textit{Coh}(X, \mathcal{I})
$$
가 존재하고, 이는 $(\mathcal{G}_n)$을 $(f^*\mathcal{G}_n)$으로 보낸다.
$f$가 평탄이면 $f^*$는 완전 함자이다.
\end{lemma}

\begin{proof}
$f^* : \textit{Coh}(\mathcal{O}_Y) \to \textit{Coh}(\mathcal{O}_X)$는
오른쪽 완전이므로
$$
f^*\mathcal{G}_n =
f^*(\mathcal{G}_{n + 1}/\mathcal{I}^n\mathcal{G}_{n + 1}) =
f^*\mathcal{G}_{n + 1}/f^{-1}\mathcal{I}^nf^*\mathcal{G}_{n + 1} =
f^*\mathcal{G}_{n + 1}/\mathcal{J}^nf^*\mathcal{G}_{n + 1}
$$
이다. 따라서 계의 당김은 계이다. Lemma
\ref{coherent-lemma-inverse-systems-abelian}의 증명에서 여핵을 구성한 방식에
따르면 $f^* : \textit{Coh}(Y, \mathcal{J}) \to \textit{Coh}(X, \mathcal{I})$는
언제나 오른쪽 완전이다. $f$가 평탄이면
$f^* : \textit{Coh}(\mathcal{O}_Y) \to \textit{Coh}(\mathcal{O}_X)$는 완전
함자이다. Lemma
\ref{coherent-lemma-inverse-systems-abelian}의 증명에서 핵을 구성한 방식에 따르면
이 경우 $f^* : \textit{Coh}(Y, \mathcal{J}) \to \textit{Coh}(X, \mathcal{I})$는
핵도 핵으로 보낸다.
\end{proof}

\begin{lemma}
\label{coherent-lemma-inverse-systems-pullback-equivalence}
$f : X' \to X$를 뇌터 스킴들의 사상이라 하자. $Z \subset X$를 닫힌
부분스킴이라 하고 $Z' = f^{-1}Z$로 스킴론적 역상을 나타내자.
$\mathcal{I} \subset \mathcal{O}_X$와
$\mathcal{I}' \subset \mathcal{O}_{X'}$을 이에 대응하는 준연접 아이디얼
층들이라 하자. $f$가 평탄이고 유도된 사상 $Z' \to Z$가
동형사상이면, 당김 함자
$f^* : \textit{Coh}(X, \mathcal{I}) \to \textit{Coh}(X', \mathcal{I}')$
(Lemma \ref{coherent-lemma-inverse-systems-pullback})는 동치이다.
\end{lemma}

\begin{proof}
$X$와 $X'$이 아핀이면 More on Algebra, Lemma
\ref{more-algebra-lemma-neighbourhood-equivalence}에서 즉시 따른다.
일반적인 경우를 증명하기 위해 $Z_n \subset X$와
$Z'_n \subset X'$을 $n$번째 무한소 근방이라 하자. 이들은 각각
$Z$와 $Z'$의 무한소 근방이다.
유도된 사상 $Z_n \to Z'_n$은 바탕 위상공간 사이의 위상동형이다.
한편 $z' \in Z'$가 $z \in Z$로 사상되면 환 사상
$\mathcal{O}_{X, z} \to \mathcal{O}_{X', z'}$은 평탄이고 동형사상
$\mathcal{O}_{X, z}/\mathcal{I}_z \to \mathcal{O}_{X', z'}/\mathcal{I}'_{z'}$을
유도한다. 따라서 동형사상
$\mathcal{O}_{X, z}/\mathcal{I}_z^n \to
\mathcal{O}_{X', z'}/(\mathcal{I}'_{z'})^n$을
모든 $n \geq 1$에 대해 유도한다. 예를 들어 More
on Algebra, Lemma \ref{more-algebra-lemma-neighbourhood-isomorphism}을
보라. 그러므로 $Z'_n \to Z_n$은 스킴의 동형사상이다. 따라서
$f^*$는 연접 $\mathcal{O}_X$-가군 가운데 $\mathcal{I}^n$에 의해 소멸되는
것들의 범주와 연접 $\mathcal{O}_{X'}$-가군 가운데 $(\mathcal{I}')^n$에
의해 소멸되는 것들의 범주 사이의 동치를 유도한다. Lemma
\ref{coherent-lemma-i-star-equivalence}를 보라. 이것은 분명 보조정리를 함의한다.
\end{proof}

\begin{lemma}
\label{coherent-lemma-inverse-systems-ideals-equivalence}
$X$를 뇌터 스킴이라 하자.
$\mathcal{I}, \mathcal{J} \subset \mathcal{O}_X$를 준연접 아이디얼
층들이라 하자. $V(\mathcal{I}) = V(\mathcal{J})$가 $X$의 같은 닫힌
부분집합이면 $\textit{Coh}(X, \mathcal{I})$와
$\textit{Coh}(X, \mathcal{J})$는 동치이다.
\end{lemma}

\begin{proof}
먼저 $X = \Spec(A)$가 아핀이라고 가정하자. $I, J \subset A$를
$\mathcal{I}, \mathcal{J}$에 대응하는 아이디얼들이라 하자. 그러면
$V(I) = V(J)$에 의해 $I^c \subset J$ 및 $J^d \subset I$이다. 여기서
어떤 $c, d \geq 1$을 택할 수 있다. 이는 Zariski 위상의 초등적인
성질에서 따른다(Algebra, Section \ref{algebra-section-spectrum-ring}과
Lemma \ref{algebra-lemma-Noetherian-power}를 보라). 따라서 $I$-진 완비화와
$J$-진 완비화는 $A$에 대해 일치한다. Algebra, Lemma
\ref{algebra-lemma-change-ideal-completion}을 보라. 그러므로 이 경우
동치는 Lemma \ref{coherent-lemma-inverse-systems-affine}에서 따른다.

\medskip\noindent
일반적으로 위에서 말한 것과 $X$가 준콤팩트라는 사실을 사용하면
$c, d \geq 1$을 택하여 $\mathcal{I}^c \subset \mathcal{J}$ 및
$\mathcal{J}^d \subset \mathcal{I}$이 되게 할 수 있다. 이제
대상 $(\mathcal{F}_n)$이 $\textit{Coh}(X, \mathcal{I})$에 주어졌을 때
역계
$$
(\mathcal{F}_{cn}/\mathcal{J}^n\mathcal{F}_{cn})
$$
가 $\textit{Coh}(X, \mathcal{J})$에 속한다고 주장한다. 이는 아핀
덮개의 원소들 위에서 확인할 수 있다. 세부사항은 생략한다. 같은 방식으로
$\textit{Coh}(X, \mathcal{I})$의 대상을 구성할 수 있다. 이때
$\textit{Coh}(X, \mathcal{J})$의 대상에서 출발한다. 이 구성들이
서로 준역인 함자들을 정의한다는 확인은 생략한다.
\end{proof}

\section{그로텐디크 존재 정리 I}
\label{coherent-section-existence}

\noindent
이 절에서는 사영인 경우의 그로텐디크 존재 정리를 논한다. Section
\ref{coherent-section-coherent-formal}에서 전개한 연접 형식 가군의 개념을 사용할
것이다. 형식 스킴에 익숙한 독자에게는 \cite{EGA}에 있는 정리의 진술과
증명을 읽기를 권한다.

\begin{lemma}
\label{coherent-lemma-fully-faithful}
$A$를 아이디얼 $I$에 관해 완비인 뇌터 환이라 하자.
$f : X \to \Spec(A)$를 고유 사상이라 하자.
$\mathcal{I} = I\mathcal{O}_X$로 놓자.
그러면 함자 (\ref{coherent-equation-completion-functor})는 충실충만이다.
\end{lemma}

\begin{proof}
$\mathcal{F}$, $\mathcal{G}$를 연접 $\mathcal{O}_X$-가군이라 하자.
그러면 $\mathcal{H} = \SheafHom_{\mathcal{O}_X}(\mathcal{G}, \mathcal{F})$는
연접 $\mathcal{O}_X$-가군이다. Modules, Lemma
\ref{modules-lemma-internal-hom-locally-kernel-direct-sum}를 보라.
Lemma \ref{coherent-lemma-completion-internal-hom}에 의해 사상
$$
\lim_n H^0(X, \mathcal{H}/\mathcal{I}^n\mathcal{H})
\to
\Mor_{\textit{Coh}(X, \mathcal{I})}
(\mathcal{G}^\wedge, \mathcal{F}^\wedge)
$$
은 전단사이다. 따라서 함자 (\ref{coherent-equation-completion-functor})가
충실충만이라는 사실은 연접층 $\mathcal{H}$에 대한 형식 함수 정리
(Lemma \ref{coherent-lemma-spell-out-theorem-formal-functions})에서 따른다.
\end{proof}

\begin{lemma}
\label{coherent-lemma-vanishing-projective}
$A$를 뇌터 환이라 하고 $I \subset A$를 아이디얼이라 하자.
$f : X \to \Spec(A)$를 고유 사상이라 하고
$\mathcal{L}$을 $f$-풍부한 가역층이라 하자.
$\mathcal{I} = I\mathcal{O}_X$로 놓자. $(\mathcal{F}_n)$을
$\textit{Coh}(X, \mathcal{I})$의 대상이라 하자. 그러면 정수 $d_0$가
존재하여
$$
H^1(X, \Ker(\mathcal{F}_{n + 1} \to \mathcal{F}_n)
\otimes \mathcal{L}^{\otimes d} )
= 0
$$
이 모든 $n \geq 0$ 및 모든 $d \geq d_0$에 대해 성립한다.
\end{lemma}

\begin{proof}
$B = \bigoplus I^n/I^{n + 1}$로 놓고
$\mathcal{B} = \bigoplus \mathcal{I}^n/\mathcal{I}^{n + 1} = f^*\widetilde{B}$로
놓자. Lemma \ref{coherent-lemma-finite-over-rees-algebra}에 의해 등급 준연접
$\mathcal{B}$-가군
$\mathcal{G} = \bigoplus \Ker(\mathcal{F}_{n + 1} \to \mathcal{F}_n)$은
유한형이다. 따라서 보조정리는 Lemma \ref{coherent-lemma-graded-finiteness}의
(2)에서 따른다.
\end{proof}

\begin{lemma}
\label{coherent-lemma-existence-projective}
$A$를 아이디얼 $I$에 관해 완비인 뇌터 환이라 하자.
$f : X \to \Spec(A)$를 사영 사상이라 하자.
$\mathcal{I} = I\mathcal{O}_X$로 놓자.
그러면 함자 (\ref{coherent-equation-completion-functor})는 동치이다.
\end{lemma}

\begin{proof}
함자 (\ref{coherent-equation-completion-functor})가 충실충만임은 Lemma
\ref{coherent-lemma-fully-faithful}에서 이미 보였다. 따라서
이 함자가 본질적 전사임을 보이면 충분하다.

\medskip\noindent
먼저 모든 대상 $(\mathcal{F}_n)$이 $\textit{Coh}(X, \mathcal{I})$에 속하면
함자 (\ref{coherent-equation-completion-functor})의 상에 속하는 어떤 대상의
몫임을 보이자. $\mathcal{L}$을 $f$-풍부한 가역층으로서 $X$ 위에 주어진
것이라 하자.
Lemma \ref{coherent-lemma-vanishing-projective}에서와 같이 $d_0$를 택하자.
$d \geq d_0$를 택하여
$\mathcal{F}_1 \otimes \mathcal{L}^{\otimes d}$가 어떤 절단들
$s_{1, 1}, \ldots, s_{t, 1}$에 의해 대역적으로 생성되게 하자. 계의 전이
사상들
$$
H^0(X, \mathcal{F}_{n + 1} \otimes \mathcal{L}^{\otimes d})
\longrightarrow
H^0(X, \mathcal{F}_n \otimes \mathcal{L}^{\otimes d})
$$
은 $H^1$의 소멸에 의해 전사이므로, $s_{1, 1}, \ldots, s_{t, 1}$을
양립하는 대역 절단들의 계 $s_{1, n}, \ldots, s_{t, n}$으로 올릴 수 있다.
이들은 $\mathcal{F}_n \otimes \mathcal{L}^{\otimes d}$의 절단들이다. 이
절단들은 사상들의
양립하는 계
$$
(s_{1, n}, \ldots, s_{t, n}) :
(\mathcal{L}^{\otimes -d})^{\oplus t} \longrightarrow \mathcal{F}_n
$$
를 정한다. Lemma \ref{coherent-lemma-inverse-systems-surjective}를 사용하면 원하는
전사 사상
$$
\left((\mathcal{L}^{\otimes -d})^{\oplus t}\right)^\wedge
\longrightarrow
(\mathcal{F}_n)
$$
을 얻는다.

\medskip\noindent
앞 문단의 결과와 $\textit{Coh}(X, \mathcal{I})$가 아벨 범주라는 사실
(Lemma \ref{coherent-lemma-inverse-systems-abelian})에 의해,
$\textit{Coh}(X, \mathcal{I})$의 모든 대상은
$\textit{Coh}(\mathcal{O}_X)$에서 오는 두 대상 사이의 사상의 여핵이다.
함자 (\ref{coherent-equation-completion-functor})는 Lemmas
\ref{coherent-lemma-fully-faithful} 및 \ref{coherent-lemma-exact}에 의해 충실충만이고
완전하므로 결론을 얻는다.
\end{proof}

\section{그로텐디크 존재 정리 II}
\label{coherent-section-existence-proper}

\noindent
이 절에서는 고유인 경우의 그로텐디크 존재 정리를 논한다. 정리의 진술과
증명을 제시하기 전에 Section \ref{coherent-section-coherent-formal}에서 도입한 연접
형식 가군들의 범주 $\textit{Coh}(X, \mathcal{I})$에 관한 이론을 조금 더
전개해야 한다.

\begin{remark}
\label{coherent-remark-inverse-systems-kernel-cokernel-annihilated-by}
$X$를 뇌터 스킴이라 하자.
$\mathcal{I}, \mathcal{K} \subset \mathcal{O}_X$를 준연접 아이디얼 층들이라
하자. $\alpha : (\mathcal{F}_n) \to (\mathcal{G}_n)$를
$\textit{Coh}(X, \mathcal{I})$의 사상이라 하자. 아핀 열린부분집합
$\Spec(A) = U \subset X$가 주어지고 $\mathcal{I}|_U, \mathcal{K}|_U$가
아이디얼들 $I, K \subset A$에 대응한다고 하자. 사상을
$\alpha_U : M \to N$으로 나타내자. 이는 유한 $A^\wedge$-가군들의
사상으로서 $\alpha|_U$에 대응한다. Lemma
\ref{coherent-lemma-inverse-systems-affine}를 보라. 다음 조건들이 동치라고 주장한다.
\begin{enumerate}
\item 어떤 정수 $t \geq 1$이 존재하여 $\Ker(\alpha_n)$과
$\Coker(\alpha_n)$이 $\mathcal{K}^t$에 의해 소멸된다. 이는 모든
$n \geq 1$에 대해 성립한다.
\item 위와 같은 임의의 아핀 열린부분집합 $\Spec(A) = U \subset X$에
대하여, 가군들 $\Ker(\alpha_U)$와 $\Coker(\alpha_U)$은 $K^t$에 의해
소멸된다. 여기서 어떤 정수 $t \geq 1$을 택할 수 있다.
\item 유한 아핀 열린 덮개 $X = \bigcup U_i$가 존재하여
$\alpha_{U_i}$에 대해 (2)의 결론이 성립한다.
\end{enumerate}
이 동치 조건들이 성립하면 $\alpha$를
{\it 핵과 여핵이 $\mathcal{K}$의 한 멱에 의해 소멸되는 사상}이라고 한다.
동치성을 보이기 위해 다음 가환대수 사실을 사용한다. 완전열
$$
0 \to T \to M \to N \to Q \to 0
$$
이 $A$-가군들로 이루어지고 $T$와 $Q$는 $K^t$에 의해 소멸된다고 하자.
여기서 어떤 아이디얼 $K \subset A$가 주어져 있다. 그러면 모든
$f, g \in K^t$에 대해 표준 사상 $"fg": N \to M$이 존재하여 합성
$M \to N \to M$은 $fg$를 곱하는 사상과 같다. 실제로 $y \in N$에
대하여 $x \in M$을 택하되 그 상이 $fy$가 되게 하자. 이 상은 $N$에서
본 것이다. 그리고
$"fg"(y) = gx$로 놓으면 된다. 따라서 $\Ker(M/JM \to N/JN)$과
$\Coker(M/JM \to N/JN)$이 $K^{2t}$에 의해 소멸됨은 분명하다. 이는
임의의 아이디얼 $J \subset A$에 대해 성립한다.

\medskip\noindent
이 가환대수 사실을 $\alpha_{U_i}$와 $J = I^n$에 적용하면 (3)이 (1)을
함의함을 알 수 있다. 거꾸로 (1)이 성립하고 $M \to N$이
$\alpha_U$와 같다고 하자. 그러면 어떤 $t \geq 1$이 존재하여
$\Ker(M/I^nM \to N/I^nN)$과 $\Coker(M/I^nM \to N/I^nN)$이 $K^t$에 의해
소멸된다. 이는 모든 $n$에 대해 성립한다. 사상들
$"fg" : N/I^nN \to M/I^nM$을 얻으며, 이 사상들은 극한에서 사상
$N \to M$을 유도한다. 실제로 $N$과 $M$은 $I$-진 완비이다. 합성
$N \to M \to N$은 $fg$를 곱하는 사상이므로 $fg$가 $T$와 $Q$를
소멸시킨다는 결론을 얻는다. 다시 말해 원하는 대로 $T$와 $Q$는
$K^{2t}$에 의해 소멸된다.
\end{remark}

\begin{lemma}
\label{coherent-lemma-existence-tricky}
$X$를 뇌터 스킴이라 하자.
$\mathcal{I}, \mathcal{K} \subset \mathcal{O}_X$를 준연접 아이디얼 층들이라
하자. $X_e \subset X$를 $\mathcal{K}^e$에 의해 정의되는 닫힌 부분스킴이라
하자. $\mathcal{I}_e = \mathcal{I}\mathcal{O}_{X_e}$로 놓자.
$(\mathcal{F}_n)$을 $\textit{Coh}(X, \mathcal{I})$의 대상이라 하자. 다음을
가정하자.
\begin{enumerate}
\item 함자
$\textit{Coh}(\mathcal{O}_{X_e}) \to \textit{Coh}(X_e, \mathcal{I}_e)$는
모든 $e \geq 1$에 대해 동치이다.
\item 연접층 $\mathcal{H}$가 $X$ 위에 주어지고, 사상
$\alpha : (\mathcal{F}_n) \to \mathcal{H}^\wedge$가 존재하며, 그 핵과
여핵은 $\mathcal{K}$의 한 멱에 의해 소멸된다.
\end{enumerate}
그러면 $(\mathcal{F}_n)$은 함자 (\ref{coherent-equation-completion-functor})의 본질적
상에 속한다.
\end{lemma}

\begin{proof}
이 증명에서는 다음 사실을 더 언급하지 않고 사용하겠다. 닫힌 몰입
$i : Z \to X$에 대해 함자 $i_*$는 $Z$ 위의 연접 가군들의 범주와
$X$ 위의 연접 가군 가운데 $Z$의 아이디얼 층에 의해 소멸되는 것들의
범주 사이의 동치를 준다. Lemma \ref{coherent-lemma-i-star-equivalence}를 보라.
특히 $\textit{Coh}(\mathcal{O}_{X_e})$를 연접 $\mathcal{O}_X$-가군 가운데
$\mathcal{K}^e$에 의해 소멸되는 것들의 범주와 동일시할 수 있고,
$\textit{Coh}(X_e, \mathcal{I}_e)$를 $\textit{Coh}(X, \mathcal{I})$에서
$\mathcal{K}^e$에 의해 소멸되는 대상들로 이루어진 충만한 부분범주와
동일시할 수 있다. 더욱이 (1)에 따르면 이 두 범주는 완비화 함자
(\ref{coherent-equation-completion-functor}) 아래에서 동치이다.

\medskip\noindent
이 동치를 적용하면 연접 $\mathcal{O}_X$-가군 $\mathcal{G}_e$를 얻으며,
이는 $\mathcal{K}^e$에 의해 소멸된다. 이 가군은 계
$(\mathcal{F}_n/\mathcal{K}^e\mathcal{F}_n)$에 대응하며, 이 계는
$\textit{Coh}(X, \mathcal{I})$에 속한다. 사상들
$\mathcal{F}_n/\mathcal{K}^{e + 1}\mathcal{F}_n \to
\mathcal{F}_n/\mathcal{K}^e\mathcal{F}_n$은 표준 사상들
$\mathcal{G}_{e + 1} \to \mathcal{G}_e$에 대응하며, 이 사상들은
동형사상
$\mathcal{G}_{e + 1}/\mathcal{K}^e\mathcal{G}_{e + 1} \to \mathcal{G}_e$를
유도한다. 따라서 $(\mathcal{G}_e)$는
$\textit{Coh}(X, \mathcal{K})$의 대상이다. 사상 $\alpha$는 사상들의 계
$$
\mathcal{F}_n/\mathcal{K}^e\mathcal{F}_n
\longrightarrow
\mathcal{H}/(\mathcal{I}^n + \mathcal{K}^e)\mathcal{H}
$$
를 유도하고, 따라서 사상들
$\mathcal{G}_e \to \mathcal{H}/\mathcal{K}^e\mathcal{H}$을 유도한다(범주들의
동치를 다시 사용한다). 가정 (2)에 의해 존재하는 정수 $t \geq 1$을
택하여 $\mathcal{K}^t$가 모든 사상
$\mathcal{F}_n \to \mathcal{H}/\mathcal{I}^n\mathcal{H}$의 핵과 여핵을
소멸시키게 하자. 그러면 $\mathcal{K}^{2t}$는 사상
$\mathcal{F}_n/\mathcal{K}^e\mathcal{F}_n \to
\mathcal{H}/(\mathcal{I}^n + \mathcal{K}^e)\mathcal{H}$의 핵과 여핵을
소멸시킨다. Remark
\ref{coherent-remark-inverse-systems-kernel-cokernel-annihilated-by}를 보라. 따라서
사상들
$$
\mathcal{G}_e
\longrightarrow
\mathcal{H}/\mathcal{K}^e\mathcal{H},
$$
의 핵과 여핵은 $\mathcal{K}^{4t}$에 의해 소멸된다는 결론을 얻는다.
Remark \ref{coherent-remark-inverse-systems-kernel-cokernel-annihilated-by}를 보라.
Lemma \ref{coherent-lemma-existence-easy}를 적용하면 연접 $\mathcal{O}_X$-가군
$\mathcal{F}$, 사상 $a : \mathcal{F} \to \mathcal{H}$ 및 동형사상
$\beta : (\mathcal{G}_e) \to (\mathcal{F}/\mathcal{K}^e\mathcal{F})$을
$\textit{Coh}(X, \mathcal{K})$ 안에서 얻는다. 이제 거꾸로 추적하자.
주어진 $n$에 대해 삼중항
$(\mathcal{F}/\mathcal{I}^n\mathcal{F}, a \bmod \mathcal{I}^n, \beta
\bmod \mathcal{I}^n)$은 사상 $\alpha_n \bmod \mathcal{K}^e :
(\mathcal{F}_n/\mathcal{K}^e\mathcal{F}_n) \to
(\mathcal{H}/(\mathcal{I}^n + \mathcal{K}^e)\mathcal{H})$에 대해 보조정리에
나오는 삼중항이다. 이 사상은 $\textit{Coh}(X, \mathcal{K})$에 속한다.
따라서 Lemma \ref{coherent-lemma-existence-easy}의 유일성은
눈앞의 모든 사상과 양립하는 표준 동형사상
$\mathcal{F}/\mathcal{I}^n\mathcal{F} \to \mathcal{F}_n$을 준다. 이로써
보조정리의 증명이 끝난다.
\end{proof}

\begin{lemma}
\label{coherent-lemma-inverse-systems-push-pull}
$Y$를 뇌터 스킴이라 하자.
$\mathcal{J}, \mathcal{K} \subset \mathcal{O}_Y$를 준연접 아이디얼 층들이라
하자. $f : X \to Y$를 고유 사상이라 하되
$V = Y \setminus V(\mathcal{K})$ 위에서 동형사상이라고 하자.
$\mathcal{I} = f^{-1}\mathcal{J} \mathcal{O}_X$로 놓자.
$(\mathcal{G}_n)$을 $\textit{Coh}(Y, \mathcal{J})$의 대상이라 하고,
$\mathcal{F}$를 연접 $\mathcal{O}_X$-가군이라 하며,
$\beta : (f^*\mathcal{G}_n)  \to \mathcal{F}^\wedge$를
$\textit{Coh}(X, \mathcal{I})$ 안의 동형사상이라 하자. 그러면 사상
$$
\alpha :
(\mathcal{G}_n)
\longrightarrow
(f_*\mathcal{F})^\wedge
$$
가 $\textit{Coh}(Y, \mathcal{J})$ 안에 존재하고, 그 핵과 여핵은
$\mathcal{K}$의 한 멱에 의해 소멸된다.
\end{lemma}

\begin{proof}
$f$가 고유 사상이므로 $f_*\mathcal{F}$는 연접 $\mathcal{O}_Y$-가군이다
(Proposition \ref{coherent-proposition-proper-pushforward-coherent}). 따라서
보조정리의 진술은 의미가 있다. 합성들을 생각하자.
$$
\gamma_n : \mathcal{G}_n \to
f_*f^*\mathcal{G}_n \to
f_*(\mathcal{F}/\mathcal{I}^n\mathcal{F}).
$$
여기서 첫째 사상은 수반 사상이고 둘째 사상은 $f_*\beta_n$이다.
보조정리의 사상 $\alpha$가 유일하게 존재하여 합성
$$
\mathcal{G}_n \xrightarrow{\alpha_n}
f_*\mathcal{F}/\mathcal{J}^nf_*\mathcal{F} \to
f_*(\mathcal{F}/\mathcal{I}^n\mathcal{F})
$$
이 $\gamma_n$과 모든 $n$에 대해 같도록 할 수 있다고 주장한다. 유일성
때문에 $Y = \Spec(B)$가 아핀이라고
가정해도 된다. $J \subset B$를 $\mathcal{J}$에 대응하는 아이디얼이라
하자. 다음과 같이 놓자.
$$
M_n = H^0(X, \mathcal{F}/\mathcal{I}^n\mathcal{F})
\quad\text{and}\quad
M = H^0(X, \mathcal{F})
$$
Lemma \ref{coherent-lemma-ML-cohomology-powers-ideal} 및 Theorem
\ref{coherent-theorem-formal-functions}에 의해 가군들 $M_n$의 역극한은 완비화
$M^\wedge = \lim M/J^nM$과 같다. $N_n = H^0(Y, \mathcal{G}_n)$ 및
$N = \lim N_n$으로 놓자. Lemma \ref{coherent-lemma-inverse-systems-affine}의
범주 동치를 통해 유한 $B^\wedge$-가군들 $N$과 $M^\wedge$은 각각
$(\mathcal{G}_n)$과 $f_*\mathcal{F}^\wedge$에 대응한다. 따라서
$\alpha$는 $\textit{Coh}(Y, \mathcal{J})$의 사상이어야 하며, 이는
준동형사상
$$
\lim \gamma_n : N = \lim_n N_n \longrightarrow \lim M_n = M^\wedge
$$
에 대응한다. 이 준동형사상은 유한 $B^\wedge$-가군들의 사상이다.

\medskip\noindent
이제 $\alpha$의 핵과 여핵이 $\mathcal{K}$의 한 멱에 의해 소멸됨을 보여야
한다. $Y' = \Spec(B^\wedge)$ 및 $X' = Y' \times_Y X$로 놓자.
$\mathcal{K}'$, $\mathcal{J}'$, $\mathcal{G}'_n$ 및 $\mathcal{I}'$,
$\mathcal{F}'$을 각각 $\mathcal{K}$, $\mathcal{J}$, $\mathcal{G}_n$ 및
$\mathcal{I}$, $\mathcal{F}$의 $Y'$와 $X'$로의 당김이라 하자. 사영 사상
$f' : X' \to Y'$은 $f$의 $Y' \to Y$에 의한 밑변환이다. 사상
$Y' \to Y$는 스킴들의 평탄 사상임을 주목하자. 실제로
$B \to B^\wedge$는 Algebra, Lemma \ref{algebra-lemma-completion-flat}에
의해 평탄이다. 따라서
$f'_*\mathcal{F}'$, 각각 $f'_*(f')^*\mathcal{G}_n'$은 Lemma
\ref{coherent-lemma-flat-base-change-cohomology}에 의해 $f_*\mathcal{F}$, 각각
$f_*f^*\mathcal{G}_n$의 $Y'$로의 당김이다. 우리 구성의 유일성에 따라
$\alpha$의 $Y'$로의 당김은 대응 사상 $\alpha'$이다. 이 사상은 $Y'$ 위의
상황에서 구성한 것이다. 더욱이 $\alpha$의 핵과 여핵이 $\mathcal{K}^t$에 의해
소멸되는지를 확인하려면 $\alpha'$의 핵과 여핵이
$(\mathcal{K}')^t$에 의해 소멸되는지만 확인하면 충분하다. 실제로 이를
보려면 사상들 $\alpha_n$과 $\alpha'_n$의 핵과 여핵에 대해 이 사실을
확인해야 하고(Remark
\ref{coherent-remark-inverse-systems-kernel-cokernel-annihilated-by}를 보라), 환 사상
$B \to B^\wedge$은 $J^n$에 의해 소멸되는 가군들의 범주와
$(J')^n$에 의해 소멸되는 가군들의 범주 사이의 동치를 유도한다. More on
Algebra, Lemma \ref{more-algebra-lemma-neighbourhood-equivalence}를 보라.
따라서 $B$가 $J$에 관해 완비라고 가정해도 된다.

\medskip\noindent
$Y = \Spec(B)$가 아핀이고, $\mathcal{J}$가 아이디얼 $J \subset B$에
대응하며, $B$가 $J$에 관해 완비라고 가정하자. 이 경우
$(\mathcal{G}_n)$은 함자
$\textit{Coh}(\mathcal{O}_Y) \to \textit{Coh}(Y, \mathcal{J})$의 본질적
상에 속한다. $\mathcal{G}$를 연접 $\mathcal{O}_Y$-가군이라 하되
$(\mathcal{G}_n) = \mathcal{G}^\wedge$을 만족하는 것으로 택하자.
$f^*(\mathcal{G}^\wedge) = (f^*\mathcal{G})^\wedge$임을 주목하자. 따라서
Lemma \ref{coherent-lemma-fully-faithful}에 의해 $\beta$는 동형사상
$b : f^*\mathcal{G} \to \mathcal{F}$에서 오며, $\alpha$는 다음 사상에
완비화 함자를 적용한 것이다.
$$
\mathcal{G} \to f_*f^*\mathcal{G} \cong f_*\mathcal{F}
$$
따라서 수반 사상 $c : \mathcal{G} \to f_*f^*\mathcal{G}$의 핵과 여핵이
$\mathcal{K}$의 한 멱에 의해 소멸됨을 확인하면 된다. 그런데 제한
$f|_{f^{-1}(V)} : f^{-1}(V) \to V$는 동형사상이므로 $c|_V$도
동형사상이다. 따라서 연접층들 $\Ker(c)$와 $\Coker(c)$의 지지는
$V(\mathcal{K})$에 포함되고, 그러므로 이들은 원하는 대로
$\mathcal{K}$의 한 멱에 의해 소멸된다(Lemma
\ref{coherent-lemma-power-ideal-kills-sheaf}).
\end{proof}

\noindent
다음 명제는 실제로 가장 자주 쓰이는 그로텐디크 존재 정리의 형태이다.

\begin{proposition}
\label{coherent-proposition-existence-proper}
$A$를 아이디얼 $I$에 관해 완비인 뇌터 환이라 하자.
$f : X \to \Spec(A)$를 스킴들의 고유 사상이라 하자.
$\mathcal{I} = I\mathcal{O}_X$로 놓자.
그러면 함자 (\ref{coherent-equation-completion-functor})는 동치이다.
\end{proposition}

\begin{proof}
함자 (\ref{coherent-equation-completion-functor})가 충실충만임은 Lemma
\ref{coherent-lemma-fully-faithful}에서 이미 보았다. 따라서 이 함자가 본질적
전사임을 보이면 충분하다.

\medskip\noindent
모음 $\Xi$를 다음과 같은 준연접 아이디얼 층
$\mathcal{K} \subset \mathcal{O}_X$들로 이루어진 것으로 정의하자. 모든
대상 $(\mathcal{F}_n)$ 가운데 $\mathcal{K}$에 의해 소멸되는 것은 본질적
상에 속한다.
$(0)$이 $\Xi$에 속함을 보이고자 한다. 그렇지 않다면 $X$가 뇌터이므로
극대 준연접 아이디얼 층 $\mathcal{K}$가 존재하되 $\Xi$에 속하지 않는다.
Lemma \ref{coherent-lemma-acc-coherent}를 보라. $X$를 그 안의 닫힌 부분스킴으로
대체하자. 이 부분스킴은 $X$에서 $\mathcal{K}$에 대응한다. 그러면 모든
영이 아닌 $\mathcal{K}$가
$\Xi$에 속한다고 가정해도 된다. 여기서는 $\mathcal{K}$에 의해 소멸되는
연접 가군들과 $\mathcal{K}$에 대응하는 닫힌 부분스킴 위의 연접
가군들의 대응을 사용한다. Lemma \ref{coherent-lemma-i-star-equivalence}를 보라.
$(\mathcal{F}_n)$을 $\textit{Coh}(X, \mathcal{I})$의 대상이라 하자. 이
대상이 함자 (\ref{coherent-equation-completion-functor})의 본질적 상에 속함을 보여
명제의 증명을 끝내겠다.

\medskip\noindent
Chow의 보조정리(Lemma \ref{coherent-lemma-chow-Noetherian})를 적용하여 고유 전사
사상 $f : X' \to X$를 택하자. 이 사상은 조밀한 열린부분집합
$U \subset X$ 위에서 동형사상이고 $X'$은 $A$ 위에서 사영이다.
$\mathcal{K}$를 축약된 여집합 $X \setminus U$를 정의하는 준연접 아이디얼
층이라 하자. 그로텐디크 존재 정리의 사영인 경우(Lemma
\ref{coherent-lemma-existence-projective})에 의해 연접 가군 $\mathcal{F}'$이
$X'$ 위에 존재하여
$(\mathcal{F}')^\wedge \cong (f^*\mathcal{F}_n)$이다. Proposition
\ref{coherent-proposition-proper-pushforward-coherent}에 의해 $\mathcal{O}_X$-가군
$\mathcal{H} = f_*\mathcal{F}'$은 연접이고, Lemma
\ref{coherent-lemma-inverse-systems-push-pull}에 의해 사상
$(\mathcal{F}_n) \to \mathcal{H}^\wedge$가
$\textit{Coh}(X, \mathcal{I})$ 안에 존재한다. 그 핵과 여핵은
$\mathcal{K}$의 한 멱에 의해 소멸된다. 모든 멱 $\mathcal{K}^e$는
$\Xi$에 속하므로 함자 (\ref{coherent-equation-completion-functor})는 닫힌
부분스킴들 $X_e = V(\mathcal{K}^e)$에 대해 동치이다. Lemma
\ref{coherent-lemma-existence-tricky}에 의해 결론을 얻는다.
\end{proof}

\section{밑 위에서의 고유성}
\label{coherent-section-proper-over-base}

\noindent
이 짧은 절에서는 밑 위에서 고유인 닫힌 부분집합과, 지지가 밑 위에서
고유인 유한형 준연접 가군의 몇 가지 유용한 성질을 지적한다.

\begin{lemma}
\label{coherent-lemma-closed-proper-over-base}
$f : X \to S$를 국소 유한형인 스킴 사상이라 하자.
$Z \subset X$를 닫힌 부분집합이라 하자. 다음 조건들은 동치이다.
\begin{enumerate}
\item 사상 $Z \to S$는 고유이다. 여기서 $Z$에는 유도된 축약 닫힌
부분스킴 구조를 준다(Schemes, Definition
\ref{schemes-definition-reduced-induced-scheme}).
\item $Z$ 위의 어떤 닫힌 부분스킴 구조에 대해 사상 $Z \to S$는
고유이다.
\item $Z$ 위의 임의의 닫힌 부분스킴 구조에 대해 사상 $Z \to S$는
고유이다.
\end{enumerate}
\end{lemma}

\begin{proof}
(3) $\Rightarrow$ (1) 및 (1) $\Rightarrow$ (2)는 즉시 따른다. 따라서
(2)가 (3)을 함의함을 보이면 충분하다. 독자에게 이 사실의 증명을 직접
찾아보기를 권한다. $Z'$과 $Z''$을 $Z$ 위의 닫힌 부분스킴 구조들이라
하고 $Z' \to S$가 고유라고 하자. $Z'' \to S$가 고유임을 보여야 한다.
$Z''' = Z' \cup Z''$을 스킴론적 합집합이라 하자. Morphisms, Definition
\ref{morphisms-definition-scheme-theoretic-intersection-union}을 보라. 그러면
$Z'''$은 $Z$ 위의 또 다른 닫힌 부분스킴 구조이다. 예를 들어 Morphisms,
Lemma \ref{morphisms-lemma-scheme-theoretic-union}에 있는 스킴론적 합집합의
기술에서 이 사실이 따른다. $Z'' \to Z'''$이 닫힌 몰입이므로
$Z''' \to S$가 고유임을 증명하면 충분하다(Morphisms, Lemmas
\ref{morphisms-lemma-closed-immersion-proper} 및
\ref{morphisms-lemma-composition-proper}를 보라). 사상
$Z' \to Z'''$은 전단사 닫힌 몰입이고, 특히 전사이며 보편적으로 닫혀
있다. 따라서 $Z' \to S$가 분리라는 사실로부터 $Z''' \to S$가 분리임이
따른다. Morphisms, Lemma
\ref{morphisms-lemma-image-universally-closed-separated}를 보라. 또한
$Z''' \to S$는 국소 유한형이다. 실제로 $X \to S$가 국소 유한형이다
(Morphisms, Lemmas \ref{morphisms-lemma-immersion-locally-finite-type} 및
\ref{morphisms-lemma-composition-finite-type}). $Z' \to S$가 준콤팩트이고
$Z' \to Z'''$이 위상동형이므로 $Z''' \to S$도 준콤팩트이다. 마지막으로
$Z' \to S$가 보편적으로 닫혀 있으므로 Morphisms, Lemma
\ref{morphisms-lemma-image-proper-is-proper}에 의해 $Z''' \to S$도
그러하다. 이로써 증명이 끝난다.
\end{proof}

\begin{definition}
\label{coherent-definition-proper-over-base}
$f : X \to S$를 국소 유한형인 스킴 사상이라 하자.
$Z \subset X$를 닫힌 부분집합이라 하자. Lemma
\ref{coherent-lemma-closed-proper-over-base}의 동치 조건들이 성립하면
{\it $Z$는 $S$ 위에서 고유이다}라고 한다.
\end{definition}

\noindent
위 정의에 사용한 보조정리는 사상 $f : X \to S$가 국소 유한형이 아니면
거짓이다. 따라서 $f$가 국소 유한형이 아닐 때는 이 용어를 사용하지
말기를 독자에게 강하게 권한다.

\begin{lemma}
\label{coherent-lemma-closed-closed-proper-over-base}
$f : X \to S$를 국소 유한형인 스킴 사상이라 하자.
$Y \subset Z \subset X$를 닫힌 부분집합들이라 하자. $Z$가 $S$ 위에서
고유이면 $Y$도 그러하다.
\end{lemma}

\begin{proof}
생략한다.
\end{proof}

\begin{lemma}
\label{coherent-lemma-base-change-closed-proper-over-base}
스킴들의 데카르트 도식
$$
\xymatrix{
X' \ar[d]_{f'} \ar[r]_{g'} & X \ar[d]^f \\
S' \ar[r]^g & S
}
$$
을 생각하자. 여기서 $f$는 국소 유한형이다. $Z$가 $X$의 닫힌 부분집합이고
$S$ 위에서 고유이면, $(g')^{-1}(Z)$는 $X'$의 닫힌 부분집합이고 $S'$ 위에서
고유이다.
\end{lemma}

\begin{proof}
Morphisms, Lemma \ref{morphisms-lemma-base-change-finite-type}에 의해
$f'$이 국소 유한형이므로 진술은 의미가 있다. $Z$에 유도된 축약 닫힌
부분스킴 구조를 주자. 스킴론적 역상
$Z' = (g')^{-1}(Z)$으로 나타내자(Schemes, Definition
\ref{schemes-definition-inverse-image-closed-subscheme}). 그러면
$Z' = X' \times_X Z = (S' \times_S X) \times_X Z = S' \times_S Z$는
$S'$ 위에서 고유이다. 실제로 이는 $Z$를 $S$ 위에서 밑변환한 것이다
(Morphisms, Lemma
\ref{morphisms-lemma-base-change-proper}).
\end{proof}

\begin{lemma}
\label{coherent-lemma-functoriality-closed-proper-over-base}
$S$를 스킴이라 하자. $f : X \to Y$를 스킴들의 사상이라 하고, 두 스킴은
$S$ 위에서 국소 유한형이라고 하자.
\begin{enumerate}
\item $Y$가 $S$ 위에서 분리이고 $Z \subset X$가 $S$ 위에서 고유인 닫힌
부분집합이면, $f(Z)$는 $Y$의 닫힌 부분집합이고 $S$ 위에서 고유이다.
\item $f$가 보편적으로 닫혀 있고 $Z \subset X$가 $S$ 위에서 고유인
닫힌 부분집합이면, $f(Z)$는 $Y$의 닫힌 부분집합이고 $S$ 위에서 고유이다.
\item $f$가 고유이고 $Z \subset Y$가 $S$ 위에서 고유인 닫힌
부분집합이면, $f^{-1}(Z)$는 $X$의 닫힌 부분집합이고 $S$ 위에서 고유이다.
\end{enumerate}
\end{lemma}

\begin{proof}
(1)의 증명. $Y$가 $S$ 위에서 분리이고 $Z \subset X$가 $S$ 위에서
고유인 닫힌 부분집합이라고 하자. $Z$에 유도된 축약 닫힌 부분스킴
구조를 주고 $Z \to Y$에 $S$ 위에서 Morphisms, Lemma
\ref{morphisms-lemma-scheme-theoretic-image-is-proper}를 적용하면 된다.

\medskip\noindent
(2)의 증명. $f$가 보편적으로 닫혀 있고 $Z \subset X$가 $S$ 위에서
고유인 닫힌 부분집합이라고 하자. $Z$와 $Z' = f(Z)$에 유도된 축약
닫힌 부분스킴 구조를 주자. 유도된 사상 $Z \to Z'$을 얻는다. 스킴론적
역상 $Z'' = f^{-1}(Z')$으로 나타내자(Schemes, Definition
\ref{schemes-definition-inverse-image-closed-subscheme}). 그러면
$Z'' \to Z'$은 $f$의 밑변환이므로 보편적으로 닫혀 있다(Morphisms,
Lemma \ref{morphisms-lemma-base-change-proper}). 따라서 $Z \to Z'$은 닫힌
몰입 $Z \to Z''$과 $Z'' \to Z'$의 합성이므로 보편적으로 닫혀 있다
(Morphisms, Lemmas \ref{morphisms-lemma-closed-immersion-proper} 및
\ref{morphisms-lemma-composition-proper}). Morphisms, Lemma
\ref{morphisms-lemma-image-universally-closed-separated}에 의해
$Z' \to S$가 분리라는 결론을 얻는다. $Z \to S$가 준콤팩트이고
$Z \to Z'$이 전사이므로 $Z' \to S$는 준콤팩트이다. 사상
$Z' \to S$는 $Z' \to Y$와 $Y \to S$의 합성이다. 따라서
$Z' \to S$는 국소 유한형이다
(Morphisms, Lemmas \ref{morphisms-lemma-immersion-locally-finite-type} 및
\ref{morphisms-lemma-composition-finite-type}). 마지막으로 $Z \to S$가
보편적으로 닫혀 있으므로 Morphisms, Lemma
\ref{morphisms-lemma-image-proper-is-proper}에 의해 $Z' \to S$도 그러하다.
이로써 증명이 끝난다.

\medskip\noindent
(3)의 증명. $f$가 고유이고 $Z \subset Y$가 $S$ 위에서 고유인 닫힌
부분집합이라고 하자. $Z$에 유도된 축약 닫힌 부분스킴 구조를 주자.
스킴론적 역상 $Z' = f^{-1}(Z)$으로 나타내자(Schemes, Definition
\ref{schemes-definition-inverse-image-closed-subscheme}). 그러면
$Z' \to Z$은 $f$의 밑변환이므로 고유이다(Morphisms, Lemma
\ref{morphisms-lemma-base-change-proper}). 따라서 $Z' \to S$는
$Z' \to Z$와 $Z \to S$의 합성이므로 고유이다(Morphisms, Lemma
\ref{morphisms-lemma-composition-proper}). 이로써 증명이 끝난다.
\end{proof}

\begin{lemma}
\label{coherent-lemma-union-closed-proper-over-base}
$f : X \to S$를 국소 유한형인 스킴 사상이라 하자.
$Z_i \subset X$, $i = 1, \ldots, n$을 닫힌 부분집합들이라 하자.
$Z_i$, $i = 1, \ldots, n$이 $S$ 위에서 고유이면
$Z_1 \cup \ldots \cup Z_n$도 그러하다.
\end{lemma}

\begin{proof}
$Z_i$에 각각 유도된 축약 닫힌 부분스킴 구조를 주자. 사상
$$
Z_1 \amalg \ldots \amalg Z_n \longrightarrow X
$$
은 Morphisms, Lemmas \ref{morphisms-lemma-closed-immersion-finite} 및
\ref{morphisms-lemma-finite-union-finite}에 의해 유한이다. 유한 사상은
보편적으로 닫혀 있고(Morphisms, Lemma
\ref{morphisms-lemma-finite-proper}),
$Z_1 \amalg \ldots \amalg Z_n$은 $S$ 위에서 고유이다. 따라서 Lemma
\ref{coherent-lemma-functoriality-closed-proper-over-base}의 (2)에 의해 그 상
$Z_1 \cup \ldots \cup Z_n$은 $S$ 위에서 고유이다.
\end{proof}

\noindent
$f : X \to S$를 국소 유한형인 스킴 사상이라 하자. $\mathcal{F}$를
유한형 준연접 $\mathcal{O}_X$-가군이라 하자. 그러면 지지
$\text{Supp}(\mathcal{F})$, 즉 $\mathcal{F}$의 지지는 $X$의 닫힌
부분집합이다. Morphisms, Lemma
\ref{morphisms-lemma-support-finite-type}을 보라. 따라서 ``$\mathcal{F}$의
지지는 $S$ 위에서 고유이다''라고 말하는 것이 의미가 있다.

\begin{lemma}
\label{coherent-lemma-module-support-proper-over-base}
$f : X \to S$를 국소 유한형인 스킴 사상이라 하자. $\mathcal{F}$를
유한형 준연접 $\mathcal{O}_X$-가군이라 하자. 다음 조건들은 동치이다.
\begin{enumerate}
\item $\mathcal{F}$의 지지는 $S$ 위에서 고유이다.
\item $\mathcal{F}$의 스킴론적 지지는 $S$ 위에서 고유이다
(Morphisms, Definition \ref{morphisms-definition-scheme-theoretic-support}).
\item 닫힌 부분스킴 $Z \subset X$와 유한형 준연접
$\mathcal{O}_Z$-가군 $\mathcal{G}$가 존재하여 (a) $Z \to S$가
고유이고, (b) $(Z \to X)_*\mathcal{G} = \mathcal{F}$이다.
\end{enumerate}
\end{lemma}

\begin{proof}
지지 $\text{Supp}(\mathcal{F})$, 즉 $\mathcal{F}$의 지지는 $X$의 닫힌
부분집합이다. Morphisms, Lemma \ref{morphisms-lemma-support-finite-type}을
보라. 따라서 Definition \ref{coherent-definition-proper-over-base}를 적용할 수
있다. $\mathcal{F}$의 스킴론적 지지는 그 바탕 닫힌 부분집합이
$\text{Supp}(\mathcal{F})$인 닫힌 부분스킴이므로 Definition
\ref{coherent-definition-proper-over-base}에 의해 (1)과 (2)는 동치이다. (2)가
(3)을 함의함은 분명하다. 거꾸로 (3)이 성립하면
$\text{Supp}(\mathcal{F}) \subset Z$이고, 이는 $X$의 닫힌 부분집합들의
포함관계이다. 따라서 예를 들어 Lemma
\ref{coherent-lemma-closed-closed-proper-over-base}에 의해
$\text{Supp}(\mathcal{F})$는 $S$ 위에서 고유이다.
\end{proof}

\begin{lemma}
\label{coherent-lemma-base-change-module-support-proper-over-base}
스킴들의 데카르트 도식
$$
\xymatrix{
X' \ar[d]_{f'} \ar[r]_{g'} & X \ar[d]^f \\
S' \ar[r]^g & S
}
$$
을 생각하자. 여기서 $f$는 국소 유한형이다. $\mathcal{F}$를 유한형
준연접 $\mathcal{O}_X$-가군이라 하자. $\mathcal{F}$의 지지가 $S$ 위에서
고유이면 $(g')^*\mathcal{F}$의 지지는 $S'$ 위에서 고유이다.
\end{lemma}

\begin{proof}
Modules, Lemma \ref{modules-lemma-pullback-finite-type}에 의해
$(g')*\mathcal{F}$가 유한형이므로 진술은 의미가 있다. Morphisms, Lemma
\ref{morphisms-lemma-support-finite-type}에 의해
$\text{Supp}((g')^*\mathcal{F}) = (g')^{-1}(\text{Supp}(\mathcal{F}))$이다.
따라서 보조정리는 Lemma
\ref{coherent-lemma-base-change-closed-proper-over-base}에서 따른다.
\end{proof}

\begin{lemma}
\label{coherent-lemma-cat-module-support-proper-over-base}
$f : X \to S$를 국소 유한형인 스킴 사상이라 하자.
$\mathcal{F}$, $\mathcal{G}$를 유한형 준연접 $\mathcal{O}_X$-가군이라 하자.
\begin{enumerate}
\item $\mathcal{F}$와 $\mathcal{G}$의 지지가 $S$ 위에서 고유이면 다음
가군들의 지지도 그러하다. $\mathcal{F} \oplus \mathcal{G}$,
$\mathcal{G}$의 $\mathcal{F}$에 의한 임의의 확장, $\Im(u)$와
$\Coker(u)$가 여기에 해당한다. 여기서 후자의 두 가군은
$\mathcal{O}_X$-가군 사상 $u : \mathcal{F} \to \mathcal{G}$가 주어졌을
때의 것이다.
또한 $\mathcal{F}$ 또는 $\mathcal{G}$의 임의의
준연접 몫이 여기에 해당한다.
\item $S$가 국소 뇌터이면, 연접 $\mathcal{O}_X$-가군 가운데 지지가
$S$ 위에서 고유인 것들의 범주는 연접 $\mathcal{O}_X$-가군들의 아벨
범주의 세르 부분범주이다(Homology, Definition
\ref{homology-definition-serre-subcategory}).
\end{enumerate}
\end{lemma}

\begin{proof}
(1)의 증명. $Z$, $Z'$을 각각 $\mathcal{F}$와 $\mathcal{G}$의 지지라
하자. 그러면 (1)에서 언급한 모든 층의 지지는 $Z \cup Z'$에 포함된다.
따라서 이 층들이 유한형이고 준연접임을 확인한다면 주장 자체는 Lemmas
\ref{coherent-lemma-closed-closed-proper-over-base} 및
\ref{coherent-lemma-union-closed-proper-over-base}에서 분명히 따른다. 준연접성에
대해서는 Schemes, Section \ref{schemes-section-quasi-coherent}을 보라.
``유한형''에 대해서는 Modules, Section
\ref{modules-section-finite-type}을 살펴보기를 권한다.

\medskip\noindent
(2)의 증명은 (1)의 증명과 같다. Morphisms, Lemma
\ref{morphisms-lemma-finite-type-noetherian}에 의해 $X$가 국소 뇌터이고,
Section \ref{coherent-section-coherent-sheaves}에서 연접 가군의 범주를 기술했으므로
주장들은 의미가 있음을 주목하자.
\end{proof}

\begin{lemma}
\label{coherent-lemma-support-proper-over-base-pushforward}
$S$를 국소 뇌터 스킴이라 하자. $f : X \to S$를 국소 유한형인 스킴
사상이라 하자. $\mathcal{F}$를 연접 $\mathcal{O}_X$-가군이라 하고 그
지지는 $S$ 위에서 고유라고 하자. 그러면 $R^pf_*\mathcal{F}$는 연접
$\mathcal{O}_S$-가군이며, 이는 모든 $p \geq 0$에 대해 성립한다.
\end{lemma}

\begin{proof}
Lemma \ref{coherent-lemma-module-support-proper-over-base}에 의해 닫힌 몰입
$i : Z \to X$와 유한형 준연접 $\mathcal{O}_Z$-가군 $\mathcal{G}$가
존재하여 (a) $g = f \circ i : Z \to S$가 고유이고, (b)
$i_*\mathcal{G} = \mathcal{F}$이다. Proposition
\ref{coherent-proposition-proper-pushforward-coherent}에 의해
$R^pg_*\mathcal{G}$는 $S$ 위에서 연접이다. 한편
$R^qi_*\mathcal{G} = 0$이며 이는 $q > 0$에 대해 성립한다(Lemma
\ref{coherent-lemma-finite-pushforward-coherent}).
Cohomology, Lemma \ref{cohomology-lemma-relative-Leray}에 의해
$R^pf_*\mathcal{F} = R^pg_*\mathcal{G}$를 얻으며, 이것으로 증명이 끝난다.
\end{proof}

\begin{lemma}
\label{coherent-lemma-systems-with-proper-support}
$S$를 뇌터 스킴이라 하자. $f : X \to S$를 유한형 사상이라 하자.
$\mathcal{I} \subset \mathcal{O}_X$를 준연접 아이디얼 층이라 하자. 다음은
$\textit{Coh}(X, \mathcal{I})$의 세르 부분범주들이다.
\begin{enumerate}
\item $\textit{Coh}(X, \mathcal{I})$의 충만한 부분범주로서, 대상
$(\mathcal{F}_n)$ 가운데 $\mathcal{F}_1$의 지지가 $S$ 위에서 고유인
것들로 이루어진 범주.
\item $\textit{Coh}(X, \mathcal{I})$의 충만한 부분범주로서, 대상
$(\mathcal{F}_n)$ 가운데 닫힌 부분스킴 $Z \subset X$가 존재하고 이
부분스킴이 $S$ 위에서 고유이며 모든 항에 대해
$\mathcal{I}_Z \mathcal{F}_n = 0$인 것들로 이루어진 범주. 여기서 모든
$n \geq 1$을 고려한다.
\end{enumerate}
\end{lemma}

\begin{proof}
Homology, Lemma \ref{homology-lemma-characterize-serre-subcategory}의 판정법을
사용하겠다. 또한
$0 \to (\mathcal{G}_n) \to (\mathcal{F}_n) \to (\mathcal{H}_n) \to 0$이
$\textit{Coh}(X, \mathcal{I})$의 짧은 완전열이면 다음 사실들을 사용하겠다.
(a) $\mathcal{G}_n \to \mathcal{F}_n \to \mathcal{H}_n \to 0$은 모든
$n \geq 1$에 대해 완전하고, (b) $\mathcal{G}_n$은
$\Ker(\mathcal{F}_m \to \mathcal{H}_m)$의 몫이다. 여기서 어떤
$m \geq n$을 택할 수 있다. Lemma
\ref{coherent-lemma-inverse-systems-abelian}의 증명을 보라.

\medskip\noindent
(1)의 증명. $(\mathcal{F}_n)$을 $\textit{Coh}(X, \mathcal{I})$의 대상이라
하자. 그러면 $\text{Supp}(\mathcal{F}_n) = \text{Supp}(\mathcal{F}_1)$이고
이는 모든 $n \geq 1$에 대해 성립한다. 따라서 위의 (a)와 (b)에 의해
임의의 짧은 완전열
$0 \to (\mathcal{G}_n) \to (\mathcal{F}_n) \to (\mathcal{H}_n) \to 0$이
$\textit{Coh}(X, \mathcal{I})$에 주어지면
$\text{Supp}(\mathcal{G}_1) \cup \text{Supp}(\mathcal{H}_1) =
\text{Supp}(\mathcal{F}_1)$이다. 이것으로 (1)에서 정의한 범주가
$\textit{Coh}(X, \mathcal{I})$의 세르 부분범주임이 증명된다.

\medskip\noindent
(2)의 증명. 여기서도 같은 방식으로 논증한다.
$0 \to (\mathcal{G}_n) \to (\mathcal{F}_n) \to (\mathcal{H}_n) \to 0$을
$\textit{Coh}(X, \mathcal{I})$의 짧은 완전열이라 하자. $Z \subset X$가
닫힌 부분스킴이고 $\mathcal{I}_Z$가 $\mathcal{F}_n$을 모든 $n$에 대해
소멸시키면, 위의 (a)와 (b)에 의해 $\mathcal{I}_Z$는 $\mathcal{G}_n$과
$\mathcal{H}_n$도 모든 $n$에 대해 소멸시킨다. 따라서 $Z \to S$가
고유이면 (2)에서 정의한 범주는 $\textit{Coh}(X, \mathcal{I})$ 안에서
부분대상과 몫대상을 취하는 것에 대해 닫혀 있다는 결론을 얻는다.
마지막으로 $Z \subset X$와 $Y \subset X$가 $S$ 위에서 고유인 닫힌
부분스킴들이고,
$\mathcal{I}_Z \mathcal{G}_n = 0$ 및
$\mathcal{I}_Y \mathcal{H}_n = 0$이라고 하자. 이는 모든 $n \geq 1$에
대해 성립한다고 가정한다. 그러면 위의 (a)에 의해
$\mathcal{I}_{Z \cup Y} = \mathcal{I}_Z \cdot \mathcal{I}_Y$는
$\mathcal{F}_n$을 모든 $n$에 대해 소멸시킨다. Lemma
\ref{coherent-lemma-union-closed-proper-over-base}에 의해 $Z \cup Y \to S$는
고유이다. 여기서는 Definition \ref{coherent-definition-proper-over-base}에 따라
합집합에 임의의 스킴 구조를 택할 수 있다는 사실도 사용했다. 이로써
증명이 끝난다.
\end{proof}

\section{그로텐디크 존재 정리 III}
\label{coherent-section-existence-proper-support}

\noindent
그로텐디크 존재 정리의 일반형을 진술하기 위해 표기를 조금 더 도입한다.
$A$를 아이디얼 $I$에 관해 완비인 뇌터 환이라 하자.
$f : X \to \Spec(A)$를 분리 유한형 스킴 사상이라 하자.
$\mathcal{I} = I\mathcal{O}_X$로 놓자. 이 상황에서
$$
\textit{Coh}_{\text{support proper over } A}(\mathcal{O}_X)
$$
를 $\textit{Coh}(\mathcal{O}_X)$의 충만한 부분범주로 정의하되, 그 대상은
연접 $\mathcal{O}_X$-가군들로 한다. 그 지지는 $\Spec(A)$ 위에서 고유해야
한다.
이는 $\textit{Coh}(\mathcal{O}_X)$의 세르 부분범주이다. Lemma
\ref{coherent-lemma-cat-module-support-proper-over-base}를 보라. 마찬가지로
$$
\textit{Coh}_{\text{support proper over } A}(X, \mathcal{I})
$$
를 $\textit{Coh}(X, \mathcal{I})$의 충만한 부분범주로 정의하되, 대상
$(\mathcal{F}_n)$ 가운데 $\mathcal{F}_1$의 지지가 $\Spec(A)$ 위에서
고유인 것들로 한다. Lemma \ref{coherent-lemma-systems-with-proper-support}의
(1)에 의해 이는 $\textit{Coh}(X, \mathcal{I})$의 세르 부분범주이다.
몫가군의 지지는 원래 가군의 지지에 포함되므로
(\ref{coherent-equation-completion-functor})는 함자
\begin{equation}
\label{coherent-equation-completion-functor-proper-over-A}
\textit{Coh}_{\text{support proper over }A}(\mathcal{O}_X)
\longrightarrow
\textit{Coh}_{\text{support proper over }A}(X, \mathcal{I})
\end{equation}
를 유도한다. 이제 이 절의 주정리를 진술할 준비가 되었다.

\begin{theorem}[그로텐디크 존재 정리]
\label{coherent-theorem-grothendieck-existence}
\begin{reference}
\cite[III Theorem 5.1.5]{EGA}
\end{reference}
$A$를 아이디얼 $I$에 관해 완비인 뇌터 환이라 하자. $X$를 $A$ 위의
분리 유한형 스킴이라 하자. 그러면 함자
(\ref{coherent-equation-completion-functor-proper-over-A})
$$
\textit{Coh}_{\text{support proper over }A}(\mathcal{O}_X)
\longrightarrow
\textit{Coh}_{\text{support proper over }A}(X, \mathcal{I})
$$
는 동치이다.
\end{theorem}

\begin{proof}
Lemma \ref{coherent-lemma-i-star-equivalence}의 범주 동치를 더 언급하지 않고
사용하겠다. 이 증명에서 닫힌 부분스킴 $Z \subset X$가 $A$ 위에서
고유할 때, $X$ 위의 연접 가군이 ``$Z$ 위에 지지된다''고
함은 그것이 $Z$의 아이디얼 층에 의해 소멸되는 것, 또는 동치로 $Z$ 위의
연접 가군의 직상인 것을 뜻한다. Proposition
\ref{coherent-proposition-existence-proper}에 의해 결과가 $Z$ 위에 지지되는 연접
가군들과 연접 가군들의 계 사이의 함자에 대해서는 참임을 안다. 따라서
$\textit{Coh}_{\text{support proper over }A}(\mathcal{O}_X)$의 모든 대상과
$\textit{Coh}_{\text{support proper over }A}(X, \mathcal{I})$의 모든 대상이
어떤 닫힌 부분스킴 $Z \subset X$ 위에 지지되고 이 부분스킴이 $A$ 위에서
고유임을 보이면
충분하다. $\textit{Coh}_{\text{support proper over }A}(\mathcal{O}_X)$의
대상에 대해서는 정의상 이 사실이 성립한다.
$\textit{Coh}_{\text{support proper over }A}(X, \mathcal{I})$의 대상에
대해서는 Proposition \ref{coherent-proposition-existence-proper}의 증명 방법을
사용하여 이 명제를 증명하겠다. 먼저 그 증명을 읽기를 독자에게 권한다.

\medskip\noindent
모음 $\Xi$를 다음과 같은 준연접 아이디얼 층
$\mathcal{K} \subset \mathcal{O}_X$들로 이루어진 것으로 정의하자. 명제는
모든 대상 $(\mathcal{F}_n)$에 대해 성립해야 하는데, 이 대상은
$\textit{Coh}_{\text{support proper over }A}(X, \mathcal{I})$에 속하고
$\mathcal{K}$에 의해 소멸되어야 한다. $(0)$이 $\Xi$에 속함을 보이고자
한다. 그렇지 않다면 $X$가 뇌터이므로 극대 준연접 아이디얼 층
$\mathcal{K}$가 존재하되 $\Xi$에는 속하지 않는다. Lemma
\ref{coherent-lemma-acc-coherent}를 보라. $X$를 그 안의 닫힌 부분스킴으로
대체하자. 이 부분스킴은 $X$에서 $\mathcal{K}$에 대응한다. 그러면 모든
영이 아닌 $\mathcal{K}$가 $\Xi$에 속한다고 가정해도 된다.
$(\mathcal{F}_n)$을
$\textit{Coh}_{\text{support proper over }A}(X, \mathcal{I})$의 대상이라
하자. 이 대상이 닫힌 부분스킴 $Z \subset X$ 위에 지지되고 이 부분스킴이
$A$ 위에서 고유임을 보여 정리의 증명을 끝내겠다.

\medskip\noindent
Chow의 보조정리(Lemma \ref{coherent-lemma-chow-Noetherian})를 적용하여 고유 전사
사상 $f : Y \to X$를 택하자. 이 사상은 조밀한 열린부분집합
$U \subset X$ 위에서 동형사상이고, $Y$는 $A$ 위에서 H-준사영이다.
열린 몰입 $j : Y \to Y'$을 택하되 $Y'$이 $A$ 위에서 사영이 되게 하자.
Morphisms, Lemma
\ref{morphisms-lemma-H-quasi-projective-open-H-projective}를 보라. 다음을
주목하자.
$$
\text{Supp}(f^*\mathcal{F}_n) = f^{-1}\text{Supp}(\mathcal{F}_n) =
f^{-1}\text{Supp}(\mathcal{F}_1)
$$
첫째 등식은 Morphisms, Lemma
\ref{morphisms-lemma-support-finite-type}에서 따른다. 가정과 Lemma
\ref{coherent-lemma-functoriality-closed-proper-over-base}의 (3)에 의해
$f^{-1}\text{Supp}(\mathcal{F}_1)$은 $A$ 위에서 고유이다. 따라서
$f^{-1}\text{Supp}(\mathcal{F}_1)$의 $j$에 의한 상은 Lemma
\ref{coherent-lemma-functoriality-closed-proper-over-base}의 (1)에 의해 $Y'$에서
닫혀 있다. 그러므로 $\mathcal{F}'_n = j_*f^*\mathcal{F}_n$은 Lemma
\ref{coherent-lemma-pushforward-coherent-on-open}에 의해 $Y'$ 위에서 연접이다.
따라서 $(\mathcal{F}_n')$은
$\textit{Coh}(Y', I\mathcal{O}_{Y'})$의 대상이다. 그로텐디크 존재 정리의
사영인 경우(Lemma \ref{coherent-lemma-existence-projective})에 의해 연접
$\mathcal{O}_{Y'}$-가군 $\mathcal{F}'$과 동형사상
$(\mathcal{F}')^\wedge \cong (\mathcal{F}'_n)$이 존재하며, 이 동형사상은
$\textit{Coh}(Y', I\mathcal{O}_{Y'})$ 안에 있다.
$\mathcal{F}'/I\mathcal{F}' = \mathcal{F}'_1$이므로 다음을 얻는다.
$$
\text{Supp}(\mathcal{F}') \cap V(I\mathcal{O}_{Y'}) =
\text{Supp}(\mathcal{F}'_1) = j(f^{-1}\text{Supp}(\mathcal{F}_1))
$$
구조 사상 $p' : Y' \to \Spec(A)$는 고유이므로
$p'(\text{Supp}(\mathcal{F}') \setminus j(Y))$는 $\Spec(A)$에서 닫혀 있다.
$\Spec(A)$의 영이 아닌 닫힌 부분집합은 $V(I)$의 점을 하나 포함한다. 실제로
Algebra, Lemma \ref{algebra-lemma-radical-completion}에 의해 $I$는 $A$의
Jacobson 근기에 포함된다. 위 등식은
$\text{Supp}(\mathcal{F}') \cap (p')^{-1}V(I) \subset j(Y)$임을 보여주므로
$\text{Supp}(\mathcal{F}') \subset j(Y)$라는 결론을 얻는다. 따라서
$\mathcal{F}'|_Y = j^*\mathcal{F}'$은 닫힌 부분스킴 $Z'$ 위에 지지된다.
이 부분스킴은 $Y$에 속하고 $A$ 위에서 고유이다. 또한
$(\mathcal{F}'|_Y)^\wedge = (f^*\mathcal{F}_n)$이다.

\medskip\noindent
$\mathcal{K}$를 축약된 여집합 $X \setminus U$를 정의하는 준연접 아이디얼
층이라 하자. Proposition \ref{coherent-proposition-proper-pushforward-coherent}에
의해 $\mathcal{O}_X$-가군 $\mathcal{H} = f_*(\mathcal{F}'|_Y)$은 연접이고,
Lemma \ref{coherent-lemma-inverse-systems-push-pull}에 의해 사상
$\alpha : (\mathcal{F}_n) \to \mathcal{H}^\wedge$가 존재한다. 이 사상은
$\textit{Coh}(X, \mathcal{I})$에 속하며, 그 핵과 여핵은 멱
$\mathcal{K}^t$, 즉 $\mathcal{K}$의 한 멱에 의해 소멸된다. 완전열
$$
0 \to \Ker(\alpha) \to (\mathcal{F}_n) \to
\mathcal{H}^\wedge \to \Coker(\alpha) \to 0
$$
을 $\textit{Coh}(X, \mathcal{I})$ 안에서 얻는다. $Z_0 \subset X$가
$\mathcal{H}$의 스킴론적 지지이면 집합론적으로 $Z_0 \subset f(Z')$임은
분명하다. 따라서 Lemma \ref{coherent-lemma-closed-closed-proper-over-base} 및
Lemma \ref{coherent-lemma-functoriality-closed-proper-over-base}의 (2)에 의해 $Z_0$는
$A$ 위에서 고유이다. 그러므로 $\mathcal{H}^\wedge$은 Lemma
\ref{coherent-lemma-systems-with-proper-support}의 (2)에서 정의한 부분범주에 속하고,
더구나
$\textit{Coh}_{\text{support proper over }A}(X, \mathcal{I})$에도 속한다.
Lemma \ref{coherent-lemma-systems-with-proper-support}의 (1)에 의해
$\Ker(\alpha)$와 $\Coker(\alpha)$도
$\textit{Coh}_{\text{support proper over }A}(X, \mathcal{I})$에 속한다는
결론을 얻는다. 귀납가정, 더 정확히는 $\mathcal{K}^t$가 $\Xi$에 속한다는
사실에 의해 $\Ker(\alpha)$와 $\Coker(\alpha)$는 Lemma
\ref{coherent-lemma-systems-with-proper-support}의 (2)에서 정의한 부분범주에
속한다. 이 부분범주는 그 보조정리에 의해 세르 부분범주이므로
$(\mathcal{F}_n)$도 그러하다는 결론을 얻는다. 이것이 보이고자 한
것이다.
\end{proof}

\begin{remark}[그로텐디크 존재 정리를 풀어 쓴 형태]
\label{coherent-remark-reformulate-existence-theorem}
$A$를 아이디얼 $I$에 관해 완비인 뇌터 환이라 하자.
$S = \Spec(A)$ 및 $S_n = \Spec(A/I^n)$으로 쓰자. $X \to S$를 분리
유한형 사상이라 하자. $n \geq 1$에 대해 $X_n = X \times_S S_n$으로
놓자. 도식은 다음과 같다.
$$
\xymatrix{
X_1 \ar[r]_{i_1} \ar[d] & X_2 \ar[r]_{i_2} \ar[d] & X_3 \ar[r] \ar[d] &
\ldots & X \ar[d] \\
S_1 \ar[r] & S_2 \ar[r] & S_3 \ar[r] & \ldots & S
}
$$
이 상황에서 다음을 만족하는 계 $(\mathcal{F}_n, \varphi_n)$을 생각한다.
\begin{enumerate}
\item $\mathcal{F}_n$은 연접 $\mathcal{O}_{X_n}$-가군이다.
\item $\varphi_n : i_n^*\mathcal{F}_{n + 1} \to \mathcal{F}_n$은
동형사상이다.
\item $\text{Supp}(\mathcal{F}_1)$은 $S_1$ 위에서 고유이다.
\end{enumerate}
Theorem \ref{coherent-theorem-grothendieck-existence}은 완비화 함자
$$
\begin{matrix}
\text{coherent }\mathcal{O}_X\text{-modules }\mathcal{F} \\
\text{with support proper over }A
\end{matrix}
\quad
\longrightarrow
\quad
\begin{matrix}
\text{systems }(\mathcal{F}_n) \\
\text{as above}
\end{matrix}
$$
가 범주의 동치임을 말한다. 특별히 $X$가 $A$ 위에서 고유이면 지지에
관한 조건들을 생략할 수 있다.
\end{remark}

\section{그로텐디크 대수화 정리}
\label{coherent-section-algebraization}

\noindent
첫 결과는 그로텐디크 존재 정리를 닫힌 부분스킴과 유한 사상의 말로
옮긴 것이다.

\begin{lemma}
\label{coherent-lemma-algebraize-formal-closed-subscheme}
$A$를 아이디얼 $I$에 관해 완비인 뇌터 환이라 하자.
$S = \Spec(A)$ 및 $S_n = \Spec(A/I^n)$으로 쓰자. $X \to S$를 분리
유한형 사상이라 하자. $n \geq 1$에 대해 $X_n = X \times_S S_n$으로
놓자. 데카르트 정사각형들로 이루어진 가환 도식
$$
\xymatrix{
Z_1 \ar[r] \ar[d] & Z_2 \ar[r] \ar[d] & Z_3 \ar[r] \ar[d] & \ldots \\
X_1 \ar[r]^{i_1} & X_2 \ar[r]^{i_2} & X_3 \ar[r] & \ldots
}
$$
이 주어졌다고 하자. 다음을 가정하자.
\begin{enumerate}
\item $Z_1 \to X_1$은 닫힌 몰입이다.
\item $Z_1 \to S_1$은 고유이다.
\end{enumerate}
그러면 스킴들의 닫힌 몰입 $Z \to X$가 존재하여
$Z_n = Z \times_S S_n$이다. 더 나아가 $Z$는 $S$ 위에서 고유이다.
\end{lemma}

\begin{proof}
수직 사상들을 $j_n : Z_n \to X_n$으로 쓰자. 진술에 나오는 정사각형들이
데카르트이므로 $j_n$을 $X_1$로 밑변환한 것은 $j_1$이다. 따라서
Morphisms, Lemma \ref{morphisms-lemma-check-closed-infinitesimally}에 의해
$j_n$은 닫힌 몰입이다. $\mathcal{F}_n = j_{n, *}\mathcal{O}_{Z_n}$으로
놓자. 그러면 $j_n^\sharp$은 전사
$\mathcal{O}_{X_n} \to \mathcal{F}_n$이다. 정사각형들이 데카르트라는
사실을 다시 사용하면 $\mathcal{F}_{n + 1}$을 $X_n$으로 당긴 것이
$\mathcal{F}_n$임을 알 수 있다. 따라서 Remark
\ref{coherent-remark-reformulate-existence-theorem}에서 다시 서술한 그로텐디크
존재 정리에 의해 사상 $\mathcal{O}_X \to \mathcal{F}$가 존재하며, 이는
연접 $\mathcal{O}_X$-가군들의 사상이다. 이를 $X_n$에 제한하면
$\mathcal{O}_{X_n} \to \mathcal{F}_n$을 되찾는다. 더 나아가
$\mathcal{F}$의 지지는 $S$ 위에서 고유이다. 완비화 함자는 완전하므로
(Lemma \ref{coherent-lemma-exact}), 여핵 $\mathcal{Q}$의 완비화는 영이다. 이 여핵은
$\mathcal{O}_X \to \mathcal{F}$의 여핵이다.
$\mathcal{F}$의 지지가 $S$ 위에서
고유이고 $\mathcal{Q}$의 지지도 그러하므로, 예를 들어 그로텐디크 존재
정리에 의해 함자 (\ref{coherent-equation-completion-functor-proper-over-A})가
동치라는 사실로부터 $\mathcal{Q} = 0$이 따른다. 따라서
$\mathcal{F} = \mathcal{O}_X/\mathcal{J}$이다. 여기서 $\mathcal{J}$는
어떤 준연접 아이디얼 층이다.
$Z = V(\mathcal{J})$로 놓으면 증명이 끝난다.
\end{proof}

\noindent
다음 보조정리에서는 사실 $Y_1 \to X_1$이 유한이라고만 가정해도
충분하다. 이 가정으로부터 $Y_n \to X_n$도 유한임이 따르기 때문이다
(More on Morphisms, Lemma
\ref{more-morphisms-lemma-thicken-property-morphisms-cartesian}를 보라).

\begin{lemma}
\label{coherent-lemma-algebraize-formal-scheme-finite-over-proper}
$A$를 아이디얼 $I$에 관해 완비인 뇌터 환이라 하자.
$S = \Spec(A)$ 및 $S_n = \Spec(A/I^n)$으로 쓰자. $X \to S$를 분리
유한형 사상이라 하자. $n \geq 1$에 대해 $X_n = X \times_S S_n$으로
놓자. 데카르트 정사각형들로 이루어진 가환 도식
$$
\xymatrix{
Y_1 \ar[r] \ar[d] & Y_2 \ar[r] \ar[d] & Y_3 \ar[r] \ar[d] & \ldots \\
X_1 \ar[r]^{i_1} & X_2 \ar[r]^{i_2} & X_3 \ar[r] & \ldots
}
$$
이 주어졌다고 하자. 다음을 가정하자.
\begin{enumerate}
\item $Y_n \to X_n$은 유한 사상이다.
\item $Y_1 \to S_1$은 고유이다.
\end{enumerate}
그러면 스킴들의 유한 사상 $Y \to X$가 존재하여
$Y_n = Y \times_S S_n$이다. 더 나아가 $Y$는 $S$ 위에서 고유이다.
\end{lemma}

\begin{proof}
수직 사상들을 $f_n : Y_n \to X_n$으로 쓰자.
$\mathcal{F}_n = f_{n, *}\mathcal{O}_{Y_n}$으로 놓자. 이는 연접
$\mathcal{O}_{X_n}$-가군이다. 실제로 $f_n$은 유한이다(Lemma
\ref{coherent-lemma-finite-pushforward-coherent}). 정사각형들이 데카르트라는
사실을 사용하면 $\mathcal{F}_{n + 1}$을 $X_n$으로 당긴 것이
$\mathcal{F}_n$임을 알 수 있다. 따라서 Remark
\ref{coherent-remark-reformulate-existence-theorem}에서 다시 서술한 그로텐디크
존재 정리에 의해 연접 $\mathcal{O}_X$-가군 $\mathcal{F}$가 존재하며,
이를 $X_n$에 제한하면 $\mathcal{F}_n$을 되찾는다. 더 나아가
$\mathcal{F}$의 지지는 $S$ 위에서 고유이다. 완비화 함자는 충실충만이므로
(Theorem \ref{coherent-theorem-grothendieck-existence}), 곱셈 사상들
$\mathcal{F}_n \otimes_{\mathcal{O}_{X_n}} \mathcal{F}_n \to
\mathcal{F}_n$은 서로 맞물려 $\mathcal{F}$ 위의 대수 구조를 준다.
$Y = \underline{\Spec}_X(\mathcal{F})$로 놓으면 증명이 끝난다.
\end{proof}

\begin{lemma}
\label{coherent-lemma-algebraize-morphism}
$A$를 아이디얼 $I$에 관해 완비인 뇌터 환이라 하자.
$S = \Spec(A)$ 및 $S_n = \Spec(A/I^n)$으로 쓰자. $X$, $Y$를 $S$ 위의
스킴들이라 하자. $n \geq 1$에 대해 $X_n = X \times_S S_n$ 및
$Y_n = Y \times_S S_n$으로 놓자. 서로 양립하는 가환 도식들의 계
$$
\xymatrix{
& & X_{n + 1} \ar[rd] \ar[rr]_{g_{n + 1}} & & Y_{n + 1} \ar[ld] \\
X_n \ar[rru] \ar[rd] \ar[rr]_{g_n} & & Y_n \ar[rru] \ar[ld] & S_{n + 1} \\
& S_n \ar[rru]
}
$$
이 주어졌다고 하자. 다음을 가정하자.
\begin{enumerate}
\item $X \to S$는 고유이다.
\item $Y \to S$는 분리 유한형이다.
\end{enumerate}
그러면 스킴 사상 $g : X \to Y$가 유일하게 존재하며 이는 $S$ 위의
사상이다. 또한
$g_n$은 $g$를 $S_n$으로 밑변환한 것이다.
\end{lemma}

\begin{proof}
사상들 $(1, g_n) : X_n \to X_n \times_S Y_n$은 닫힌 몰입이다. 실제로
$Y_n \to S_n$이 분리이다(Schemes, Lemma
\ref{schemes-lemma-section-immersion}). 따라서 Lemma
\ref{coherent-lemma-algebraize-formal-closed-subscheme}에 의해 닫힌 부분스킴
$Z \subset X \times_S Y$가 존재하며 이는 $S$ 위에서 고유이다. 또한
이를 $S_n$으로 밑변환하면 $X_n \subset X_n \times_S Y_n$을 되찾는다.
첫째 사영 $p : Z \to X$는 고유 사상이다. 실제로 $Z$는 $S$ 위에서
고유이다. Morphisms, Lemma
\ref{morphisms-lemma-image-proper-scheme-closed}를 보라. 이 사상을
$S_n$으로 밑변환하면 모든 $n$에 대해 동형사상이 된다. 특히 Lemma
\ref{coherent-lemma-proper-finite-fibre-finite-in-neighbourhood}에 의해
$p : Z \to X$는 $X_0$의 어떤 열린 근방 위에서 유한이다. $X$가 $S$
위에서 고유이므로 이 열린 근방은 $X$ 전체이고, 따라서
$p : Z \to X$가 유한이라는 결론을 얻는다. Proposition
\ref{coherent-proposition-existence-proper}의 동치를 적용하면
$p_*\mathcal{O}_Z = \mathcal{O}_X$임을 알 수 있다. 실제로 이 사실은 법
$I^n$으로 모든 $n$에 대해 참이다. 따라서 $p$는 동형사상이고, 사상
$g$를 합성 $X \cong Z \to Y$로 얻는다. 유일성의
증명은 생략한다.
\end{proof}

\noindent
``추상적'' 대수화 정리를 증명하려면 풍부한 가역층이 존재한다고 가정해야
한다. 그러한 가정 없이는 결과가 거짓이기 때문이다.

\begin{theorem}[그로텐디크 대수화 정리]
\label{coherent-theorem-algebraization}
$A$를 아이디얼 $I$에 관해 완비인 뇌터 환이라 하자.
$S = \Spec(A)$ 및 $S_n = \Spec(A/I^n)$으로 놓자. 데카르트
정사각형들로 이루어진 가환 도식
$$
\xymatrix{
X_1 \ar[r]_{i_1} \ar[d] & X_2 \ar[r]_{i_2} \ar[d] & X_3 \ar[r] \ar[d] &
\ldots \\
S_1 \ar[r] & S_2 \ar[r] & S_3 \ar[r] & \ldots
}
$$
을 생각하자. $(\mathcal{L}_n, \varphi_n)$이 주어졌다고 하자. 여기서 각
$\mathcal{L}_n$은 $X_n$ 위의 가역층이고,
$\varphi_n : i_n^*\mathcal{L}_{n + 1} \to \mathcal{L}_n$은
동형사상이다. 다음을 가정하자.
\begin{enumerate}
\item $X_1 \to S_1$은 고유이다.
\item $\mathcal{L}_1$은 $X_1$ 위에서 풍부하다.
\end{enumerate}
그러면 스킴들의 고유 사상 $X \to S$, 풍부한 가역
$\mathcal{O}_X$-가군 $\mathcal{L}$ 및 동형사상들
$X_n \cong X \times_S S_n$과
$\mathcal{L}_n \cong \mathcal{L}|_{X_n}$이 존재하며, 이들은 사상들
$i_n$ 및 $\varphi_n$과 양립한다.
\end{theorem}

\begin{proof}
도식의 정사각형들이 데카르트이고 사상들 $S_n \to S_{n + 1}$이 닫힌
몰입이므로 사상들 $i_n$도 닫힌 몰입이다. 특히 $X_m$을 $X_n$의 닫힌
부분스킴으로 보아도 되며, 이는 $m < n$일 때 적용된다. 실제로 $X_m$은
준연접 아이디얼
층 $I^m\mathcal{O}_{X_n}$에 의해 정의되는 닫힌 부분스킴이다. 더욱이
스킴들 $X_1, X_2, X_3, \ldots$의 바탕 위상공간은 모두 동일시되므로,
층들 $\mathcal{O}_{X_n}$이 같은 바탕 위상공간 위에 놓여 있다고 볼 수
있고 실제로 그렇게 보겠다. 연접 $\mathcal{O}_{X_n}$-가군들도 마찬가지다.
다음과 같이 놓자.
$$
\mathcal{F}_n =
\Ker(\mathcal{O}_{X_{n + 1}} \to \mathcal{O}_{X_n})
$$
그러면 짧은 완전열들
$$
0 \to \mathcal{F}_n \to \mathcal{O}_{X_{n + 1}} \to \mathcal{O}_{X_n} \to 0
$$
을 얻는다. 위에서 본 바에 따라
$\mathcal{F}_n = I^n\mathcal{O}_{X_{n + 1}}$이다. 따라서
$\mathcal{F}_n$은 $X_{n + 1}$ 위의 연접층이고 $I$에 의해 소멸된다.
그러므로 이를 연접 $\mathcal{O}_{X_1}$-가군으로 보아도 되며 실제로
그렇게 보겠다. $m > n$일 때 층
$$
I^n\mathcal{O}_{X_m}/I^{n + 1}\mathcal{O}_{X_m}
$$
은 $\mathcal{F}_n$으로 동형사상으로 사상됨을 주목하자. 여기서 쓰는
사상은 $\mathcal{O}_{X_m} \to \mathcal{O}_{X_{n + 1}}$이다. 따라서
$n_1, n_2 \geq 0$이 주어지면 $m > n_1 + n_2$를 만족하는 정수를 택하여
곱셈 사상
$$
I^{n_1}\mathcal{O}_{X_m} \times I^{n_2}\mathcal{O}_{X_m}
\longrightarrow
I^{n_1 + n_2}\mathcal{O}_{X_m} \to \mathcal{F}_{n_1 + n_2}
$$
을 생각할 수 있다. 이는 $\mathcal{O}_{X_1}$-쌍선형 사상
$$
\mathcal{F}_{n_1} \times \mathcal{F}_{n_2} \longrightarrow
\mathcal{F}_{n_1 + n_2}
$$
을 유도하고, 이 사상은 다시 $\mathcal{O}_{X_1}$-대수
$\mathcal{F} = \bigoplus_{n \geq 0} \mathcal{F}_n$ 위에 등급 대수 구조를
정의한다.

\medskip\noindent
$B = \bigoplus I^n/I^{n + 1}$로 놓자. 이는 유한 생성 등급
$A/I$-대수이다. $\mathcal{B} = (X_1 \to S_1)^*\widetilde{B}$로 놓자. 위
논의는 등급 $\mathcal{O}_{X_1}$-대수들의 표준 전사
$$
\mathcal{B} \longrightarrow \mathcal{F}
$$
를 준다. 특히 $\mathcal{F}$는 유한형 준연접 등급 $\mathcal{B}$-가군이다.
Lemma \ref{coherent-lemma-graded-finiteness}에 의해 정수 $d_0$가 존재하여
$H^1(X_1, \mathcal{F} \otimes \mathcal{L}^{\otimes d}) = 0$이고, 이는 모든
$d \geq d_0$에 대해 성립한다. $d \geq d_0$를 택하여 절단들
$s_{0, 1}, \ldots, s_{N, 1} \in \Gamma(X_1, \mathcal{L}_1^{\otimes d})$이
존재하고 몰입
$$
\psi_1 : X_1 \to \mathbf{P}^N_{S_1}
$$
을 $S_1$ 위에서 유도하게 하자. Morphisms, Lemma
\ref{morphisms-lemma-finite-type-over-affine-ample-very-ample}를 보라.
$X_1$은 $S_1$ 위에서 고유이므로 $\psi_1$은 닫힌 몰입이다. Morphisms,
Lemma \ref{morphisms-lemma-image-proper-scheme-closed} 및 Schemes, Lemma
\ref{schemes-lemma-immersion-when-closed}를 보라. $\psi_1$을
$X_n$에서 $\mathbf{P}^N$으로 가는 닫힌 몰입들의 양립하는 계로
``올리겠다''.

\medskip\noindent
증명의 첫 문단에 나오는 짧은 완전열들에
$\mathcal{L}_{n + 1}^{\otimes d}$를 텐서하면 짧은 완전열들
$$
0 \to \mathcal{F}_n \otimes \mathcal{L}_{n + 1}^{\otimes d} \to
\mathcal{L}_{n + 1}^{\otimes d} \to \mathcal{L}_{n + 1}^{\otimes d} \to 0
$$
을 얻는다. 동형사상들 $\varphi_n$을 사용하면 동형사상들
$\mathcal{L}_{n + 1} \otimes \mathcal{O}_{X_l} = \mathcal{L}_l$을 얻으며,
이는 $l \leq n$에 대해 성립한다.
따라서 위 열은 다음과 같이 된다.
$$
0 \to \mathcal{F}_n \otimes \mathcal{L}_1^{\otimes d} \to
\mathcal{L}_{n + 1}^{\otimes d} \to \mathcal{L}_n^{\otimes d} \to 0
$$
$H^1(X, \mathcal{F}_n \otimes \mathcal{L}_1^{\otimes d})$의 소멸에 의해
$s_{0, 1}, \ldots, s_{N, 1} \in \Gamma(X_1, \mathcal{L}_1^{\otimes d})$을
절단들
$s_{0, n}, \ldots, s_{N, n} \in \Gamma(X_n, \mathcal{L}_n^{\otimes d})$로
귀납적으로 올릴 수 있다. 따라서 가환 도식
$$
\xymatrix{
X_1 \ar[r]_{i_1} \ar[d]_{\psi_1} &
X_2 \ar[r]_{i_2} \ar[d]_{\psi_2} &
X_3 \ar[r] \ar[d]_{\psi_3} &
\ldots \\
\mathbf{P}^N_{S_1} \ar[r] &
\mathbf{P}^N_{S_2} \ar[r] &
\mathbf{P}^N_{S_3} \ar[r] & \ldots
}
$$
을 얻는다. 여기서 Constructions, Section
\ref{constructions-section-projective-space}의 표기로
$\psi_n = \varphi_{(\mathcal{L}_n, (s_{0, n}, \ldots, s_{N, n}))}$이다.
정리의 진술에 나오는 정사각형들이 데카르트이므로 위 도식의
정사각형들도 데카르트이다. Lemma
\ref{coherent-lemma-algebraize-formal-closed-subscheme}를 적용하면 결론을 얻는다.
\end{proof}
